\chapter{Combining the bases and applying the classical test}
\label{chap:qpbt-combining}
% Paper subsection label sec:combining kept for cross-reference fidelity; the
% companion label sec:apply-ldt is attached to the corresponding section below.
% Source: references/qpbt-paper/14_analysis_of_the_pauli_basis_test.tex,
% lines 680-1414 (subsections "Combining the X and Z measurements" and
% "Applying the classical low-degree test").
\label{sec:combining}

Throughout this chapter we work in the setting of
Chapter~\ref{chap:qpbt-observables}: a projective strategy for the Pauli basis
test game succeeds with probability at least $1-\eps$, and all operators act
relative to the expanded state $\ket{\hat\psi}$ on registers
$\mathsf{A}\,\mathsf{A}'\,\mathsf{A}''\,\mathsf{B}\,\mathsf{B}'\,\mathsf{B}''$
(Definition~\ref{def:expanded-state}), where each of the four ancilla
registers $\mathsf{A}'$, $\mathsf{A}''$, $\mathsf{B}'$, and $\mathsf{B}''$
--- two appended per player, the pair $\mathsf{A}'\mathsf{A}''$
(respectively $\mathsf{B}'\mathsf{B}''$) holding $M$ EPR pairs of local
dimension $q$ --- is a register of $M = 2^m$ qudits with Hilbert space
isomorphic to $(\C^q)^{\ot M}$.
We write $q = 2^t$ (Definition~\ref{def:admissible-size}) and
$\tr:\F_q\to\F_2$ for the subfield trace. The points measurements
$\hat{M}^{(\mathrm{Point},W),u}_a$ for $W \in \{X,Z\}$
(Definition~\ref{def:expanded-point-measurement}), their binary refinements
\[
	\hat{M}^{(\mathrm{Point},W),u,r}_b
	= \sum_{a \,:\, \tr(ar) = b} \hat{M}^{(\mathrm{Point},W),u}_a\,,
	\qquad r \in \F_q,\ b \in \F_2
\]
(Definition~\ref{def:expanded-point-trace-projection}), and the lines
measurements $\hat{M}^{(\mathrm{Line},W),\ell}_{f}$
(Definition~\ref{def:expanded-line-measurement}) are those constructed in
Chapter~\ref{chap:qpbt-observables}, and $\delta_{\mathrm{acom}}(\eps)$ and
$\delta_{\mathrm{Line}}(\eps)$ denote the error functions of
Lemma~\ref{lem:qld-comm-cons} and Lemma~\ref{lem:qld-comm-line-cons},
respectively.

A subscript such as $(\,\cdot\,)_{\mathsf{A}\mathsf{A}'}$ indicates the pair of
registers on which an operator acts; the two bipartitions in use are
$\mathsf{A}\mathsf{A}'$ versus $\mathsf{B}\mathsf{A}''$ and
$\mathsf{B}\mathsf{B}'$ versus $\mathsf{A}\mathsf{B}''$. The state-dependent
distances $\approx_\delta$ (Definition~\ref{def:povm-distance}) and
$\simeq_\delta$ (Definition~\ref{def:consistency}) are taken with respect to
$\ket{\hat\psi}$ and averaged over the stated question distribution, and
``all symmetric equivalents'' of an approximation is meant in the sense of
Definition~\ref{def:symmetric-equivalents}. Between scalars, $\alpha
\approx_\delta \beta$ means $|\alpha - \beta| \le O(\delta)$. For sums of
measurement operators over outcomes satisfying a predicate we use the bracket
notation of Definition~\ref{def:bracket}; for instance
$T^{\ell_X,\ell_Z}_{[\mathrm{eval}_x(\cdot)=a,\ \mathrm{eval}_z(\cdot)=b]}
= \sum_{f_X,f_Z \,:\, f_X(x)=a,\ f_Z(z)=b} T^{\ell_X,\ell_Z}_{f_X,f_Z}$.
For a line $\ell$ (Definition~\ref{def:line}) and $k \in \N$ we write
$\deg_k(\ell)$ for the univariate polynomial representatives of degree at most
$k$ in the line parameter, each given by its $k+1$ coefficients, with the
parameterization and evaluation convention of
Definition~\ref{def:ideg-deg-polynomials}. Distinct representatives remain
distinct outcomes even when they induce the same function on $\F_q$; for
example, $0$ and $t^q-t$ are distinct when $k\geq q$. The source describes
these polynomials as functions defined on the line, while the coefficient-list
answer format retains their representatives
(Remark~\ref{rem:deg-line-representatives}). In particular, admissibility
does not require $md<q$. We write
$\mathrm{ideg}_{d,m}(\F_q) = \polyfunc{m}{q}{d}$
(Definition~\ref{def:polyfunc}) for the $m$-variate polynomial representatives
over $\F_q$ of individual degree at most $d$. Elements of
$\mathrm{ideg}_{d,m}(\F_q)$ enter the statements and proofs of this chapter
only through the representatives themselves and their evaluations; evaluation
identifies the representatives with the induced polynomial functions
(Definition~\ref{def:polynomials-degree}), bijectively precisely when
$d \le q-1$, a bound that admissibility of $(q,m,d)$ does not imply and that
is not assumed here. The restriction of an element of
$\mathrm{ideg}_{d,m}(\F_q)$ to a line has total degree at most $md$.

\section{Combined point measurements}

The measurements $\hat{M}^{(\mathrm{Point},X),x}$ and
$\hat{M}^{(\mathrm{Point},Z),z}$ commute only approximately. The first goal of
this chapter is to construct a single projective measurement that
simultaneously measures both, first at points and then along lines. The
construction of the combined point measurement proceeds in three steps: the
binary refinements $\hat{M}^{(\mathrm{Point},X),x,r}$ and
$\hat{M}^{(\mathrm{Point},Z),z,s}$ are combined into a sandwich POVM, the
resulting family is shown to satisfy an approximate linearity in the index
pair $(r,s)$, and the following result of Natarajan and
Vidick~\cite{Natarajan2017Linearity} converts approximate linearity into exact
linearity, from which a projective measurement with outcomes in $\F_q^2$ is
assembled by Fourier averaging.

% Statement as in Theorem 10 of the cited source \cite{Natarajan2017Linearity}
% (arXiv:1610.03574): hypothesis as in its display (7), conclusion as in its
% display (8) with the normalization of the distance bound corrected, see
% rem:linearity-distance-normalization. The cited theorem imposes no upper bound
% on delta, asserts the exact linearity for all pairs (u, u'), and guarantees
% the closeness of L^u to O^u on average over uniformly random u only, in mean
% squared state-dependent distance; it is stated for an arbitrary density
% operator rho on a single Hilbert space, with the index group written {0,1}^n
% under bitwise XOR (here F_2^t, additively) and Obs(H) denoting the Hermitian
% unitaries, i.e. the binary observables of def:povm-conventions. In the
% closeness conclusion an operator on H is compared with an operator on the
% extension via its ampliation S ot Id. The cited theorem prints the closeness
% bound as delta for the squared operator distance; what its argument
% establishes is delta for the squared binary-measurement distance, i.e.
% 2 delta for the squared operator distance, and the printed operator bound is
% false (docs/paper-gaps/qpbt_linearity-distance-normalization.tex). The
% statement below carries the corrected operator bound together with the
% equivalent measurement form. The mirror's quotation
% (references/qpbt-paper/14_analysis_of_the_pauli_basis_test.tex:713-725)
% inserts the restriction 0 <= delta <= 1, places the closeness conclusion,
% together with the exact linearity, under the quantifier "for all u, u' in
% F_2^t", and appends "L^u = Id when u = 0"; none of these is part of the
% cited theorem, and none is reproduced here. See rem:linearity-import.
\begin{theorem}[Quantum linearity test~\cite{Natarajan2017Linearity}]
	\label{thm:linearity}
	\lean{MIPStarRE.QPBT.IsBinaryObservable, MIPStarRE.QPBT.stateDepDistSq}
	\lean{MIPStarRE.QPBT.exists_exactly_linear_observables}
	\lean{MIPStarRE.QPBT.exists_exactly_linear_observables_binaryObservableDistSq}
	\leanok
	\uses{def:povm-conventions, lem:binary-measurement-distance}
	Let $t$ be a positive integer and $\delta \ge 0$. Let $\rho$ be a
	density operator on a Hilbert space $\hilb$ and let
	$\{\mathcal{O}^u\}$ be a set of binary observables indexed by
	$u \in \F_2^t$ acting on $\hilb$, such that on average over $u, u'$
	chosen uniformly at random from $\F_2^t$,
	\[
		\E_{u,u'} \Re \Tr\big(\mathcal{O}^u \, \mathcal{O}^{u'} \,
		\mathcal{O}^{u+u'} \, \rho\big) \ge 1 - \delta\,.
	\]
	Then there exist a Hilbert space $\hilb'$, a pure state
	$\ket{\mathrm{anc}} \in \hilb'$, and binary observables
	$\{\mathcal{L}^u\}$ indexed by $u \in \F_2^t$ acting on
	$\hilb \ot \hilb'$ such that, writing
	$\rho' = \rho \ot \ket{\mathrm{anc}}\bra{\mathrm{anc}}$, the exact
	linearity
	\[
		\mathcal{L}^u \, \mathcal{L}^{u'} = \mathcal{L}^{u+u'}
	\]
	holds for all $u, u' \in \F_2^t$, and
	\[
		\E_{u} \, d_{\rho'}\big(\mathcal{L}^u,\ \mathcal{O}^u \ot
		\Id_{\hilb'}\big)^2 \le 2\delta\,,
	\]
	where the expectation is over $u$ chosen uniformly at random from
	$\F_2^t$ and
	$d_{\sigma}(S,T)
	= \big(\Tr\big((S-T)^\dagger (S-T)\, \sigma\big)\big)^{1/2}$
	denotes the state-dependent distance with respect to a density
	operator $\sigma$. Equivalently, in the distance $d^{\mathrm{bin}}_\sigma$
	between the binary projective measurements associated with the
	observables (Lemma~\ref{lem:binary-measurement-distance}),
	\[
		\E_{u} \, d^{\mathrm{bin}}_{\rho'}\big(\mathcal{L}^u,\
		\mathcal{O}^u \ot \Id_{\hilb'}\big)^2 \le \delta\,.
	\]
	The cited source prints the first bound with $\delta$ in place of
	$2\delta$; that normalization is false, and the bounds stated here are
	the ones its argument establishes
	(Remark~\ref{rem:linearity-distance-normalization}).
\end{theorem}
\begin{proof}
	\leanok
	\uses{lem:binary-measurement-distance, lem:blr-defect-identity,
		lem:boolean-representation-stability}
	This is Theorem~10 of Natarajan and Vidick~\cite{Natarajan2017Linearity},
	whose Fourier--Naimark argument is followed; \cite{Ji2020MIPStar} invokes
	the result without proof. Write $\chi_v(u) = (-1)^{v \cdot u}$ and
	$\widehat{\mathcal{O}}^v = \E_u \, \chi_v(u) \, \mathcal{O}^u$ as in
	Lemma~\ref{lem:boolean-fourier-transform}.

	\emph{From the correlation to the multiplicative defect.} For
	$u, u' \in \F_2^t$, expanding
	$(\mathcal{O}^u \mathcal{O}^{u'} - \mathcal{O}^{u+u'})^\dagger
	(\mathcal{O}^u \mathcal{O}^{u'} - \mathcal{O}^{u+u'})$ with
	$(\mathcal{O}^w)^2 = \Id$ and $\Tr \rho = 1$, as in the proof of
	Lemma~\ref{lem:binary-measurement-distance}, gives
	\[
		d_\rho\big(\mathcal{O}^u \mathcal{O}^{u'}, \mathcal{O}^{u+u'}\big)^2
		= 2 - 2 \Re \Tr\big(\mathcal{O}^{u'} \mathcal{O}^u
		\mathcal{O}^{u'+u} \rho\big)\,,
	\]
	the two cross terms being adjoint to each other. Exchanging the two
	coordinates of the uniform pair $(u,u')$, the hypothesis yields
	$\E_{u,u'} d_\rho(\mathcal{O}^u \mathcal{O}^{u'},
	\mathcal{O}^{u+u'})^2 \le 2\delta$. No commutation among the
	observables is used.

	\emph{The Fourier-square measurement and its Naimark rounding.} By
	Lemma~\ref{lem:boolean-fourier-transform} the operators
	$B^v = (\widehat{\mathcal{O}}^v)^2$ are positive semidefinite and sum to
	$\Id$, so they form a POVM indexed by $\F_2^t$.
	Lemma~\ref{lem:naimark-helper}, applied with an ancillary space $\hilb'$
	of dimension $2^t + 1$ and with $\ket{\mathrm{anc}}$ the basis vector of
	its extra direction, produces orthogonal projectors on $\hilb \ot \hilb'$
	compressing along $\ket{\mathrm{anc}}$ to the $B^v$; the residual
	projector of the construction is folded into the outcome $v = 0$, which
	gives a projective measurement $\{C^v\}$ resolving the identity on
	$\hilb \ot \hilb'$ whose compression is still $\{B^v\}$, because the
	$B^v$ themselves sum to $\Id$. Set
	$\mathcal{L}^u = \sum_v \chi_v(u) \, C^v$. Each $\mathcal{L}^u$ is
	Hermitian, and the orthogonality of the $C^v$ together with
	$\chi_v(u) \chi_v(u') = \chi_v(u+u')$ gives
	$\mathcal{L}^u \mathcal{L}^{u'} = \mathcal{L}^{u+u'}$ for all $u, u'$;
	in particular $(\mathcal{L}^u)^2 = \mathcal{L}^0 = \sum_v C^v = \Id$,
	so the $\mathcal{L}^u$ are binary observables.

	\emph{The distance/defect identity.} Since $\mathcal{L}^u$ and
	$\mathcal{O}^u \ot \Id$ are binary observables and $\rho'$ is a density
	operator, Lemma~\ref{lem:binary-measurement-distance} gives
	$d_{\rho'}(\mathcal{L}^u, \mathcal{O}^u \ot \Id)^2
	= 2 - 2 \Re \Tr(\mathcal{L}^u (\mathcal{O}^u \ot \Id) \rho')$.
	Compressing along $\ket{\mathrm{anc}}$,
	$\Re \Tr(\mathcal{L}^u (\mathcal{O}^u \ot \Id) \rho')
	= \Re \Tr(\mathcal{O}^u \sum_v \chi_v(u) B^v \rho)$; averaging over $u$
	turns $\E_u \chi_v(u) \mathcal{O}^u$ into $\widehat{\mathcal{O}}^v$, and
	expanding the three Fourier coefficients and summing the characters
	(Lemma~\ref{lem:boolean-fourier-transform}) gives
	\[
		\E_u \Re \Tr\big(\mathcal{L}^u (\mathcal{O}^u \ot \Id) \rho'\big)
		= \sum_v \Re \Tr\big((\widehat{\mathcal{O}}^v)^3 \rho\big)
		= \E_{u,u'} \Re \Tr\big(\mathcal{O}^u \mathcal{O}^{u'}
		\mathcal{O}^{u+u'} \rho\big)\,.
	\]
	Hence
	$\E_u d_{\rho'}(\mathcal{L}^u, \mathcal{O}^u \ot \Id)^2
	= 2 - 2 \E_{u,u'} \Re \Tr(\mathcal{O}^u \mathcal{O}^{u'}
	\mathcal{O}^{u+u'} \rho)
	= \E_{u,u'} d_\rho(\mathcal{O}^u \mathcal{O}^{u'},
	\mathcal{O}^{u+u'})^2 \le 2\delta$ by the first step. The measurement
	form follows by halving, through the normalization identity of
	Lemma~\ref{lem:binary-measurement-distance}.

	The obligations attached to the theorem at its point of use --- the
	correspondence between the correlation hypothesis above and the relation
	$\approx_\delta$ of Definition~\ref{def:povm-distance}, the
	instantiation of $\rho$ as a reduced state of $\ket{\hat\psi}$, the
	absorption of the ancilla register into padding of the strategy state,
	and the identity $\mathcal{L}^0 = \Id$ --- are enumerated in
	Remark~\ref{rem:linearity-import} and addressed in the proof of
	Lemma~\ref{lem:qld-4-10}.
\end{proof}

\begin{remark}[Obligations of the linearity-theorem import]
	\label{rem:linearity-import}
	Theorem~\ref{thm:linearity} is stated in the form established by the
	argument of \cite{Natarajan2017Linearity}, with the normalization of its
	distance bound corrected as described in
	Remark~\ref{rem:linearity-distance-normalization}. Three features are
	load-bearing for the proof of Lemma~\ref{lem:qld-4-10}: the error is
	preserved without polynomial loss ($\delta$ in the hypothesis becomes
	$2\delta$ in the conclusion for the squared operator distance, and
	$\delta$ for the squared measurement distance, a constant factor absorbed
	by the relation $\approx_\delta$); the linearity of the
	$\{\mathcal{L}^u\}$ is exact and holds for all pairs $(u,u')$, not merely
	on average; and the hypothesis carries no upper bound on $\delta$, so the
	theorem may be applied separately for each pair $(x,z)$ with hypothesis
	error the conditional linearity defect given $(x,z)$, a quantity for
	which no a priori bound is available. The closeness of $\mathcal{L}^u$ to
	$\mathcal{O}^u$ is guaranteed on average over a uniformly random $u$
	only; this averaged form suffices because the $\mathcal{L}^u$ enter
	the proof of Lemma~\ref{lem:qld-4-10} only through averages over $u$.
	The quotation of the theorem in \cite{Ji2020MIPStar} differs from the
	cited statement in three respects: it inserts the restriction
	$0 \le \delta \le 1$, which would block the fiberwise application
	just described; it places the closeness conclusion, together with the
	exact linearity, under the quantifier ``for all $u, u' \in \F_2^t$'',
	a strengthening of the theorem proved in
	\cite{Natarajan2017Linearity}; and it appends the identity
	$\mathcal{L}^u = \Id$ for $u = 0$. None of the three is
	reproduced here: the restriction is dropped, the pointwise closeness
	is not used, and the appended identity is instead derived from the
	conclusion, since exact linearity at $u = u' = 0$ gives
	$(\mathcal{L}^0)^2 = \mathcal{L}^0$, and an idempotent unitary equals
	the identity. The factor of two of
	Remark~\ref{rem:linearity-distance-normalization} is not one of these
	differences. It is a defect of the cited statement itself: the mean
	squared operator bound printed in
	\cite[Theorem~10]{Natarajan2017Linearity} is $\delta$, whereas the
	calculation behind \cite[equation~(3)]{Natarajan2017Linearity} bounds
	the mean squared distance between the associated binary projective
	measurements, which is half the squared operator distance; the
	corrected bound $2\delta$ is stated in
	Theorem~\ref{thm:linearity}. The quotation neither reproduces nor
	asserts the erroneous constant, since it states its conclusion as
	$\mathcal{L}^u \approx_\delta \mathcal{O}^u$, which by
	Definition~\ref{def:povm-distance} bounds the associated expectation
	only by $O(\delta)$ and so absorbs a constant factor. These three
	discrepancies, together with the ancillary-padding obligation recorded
	below, are analyzed in
	\cite{gap:qpbt_linearity-theorem-quotation}, while the normalization defect
	in \cite[Theorem~10]{Natarajan2017Linearity} is analyzed in
	\cite{gap:qpbt_linearity-distance-normalization}. The source's
	linearity-based argument also calls for two conversions: the theorem is
	instantiated with $\rho$ the reduced density operator of
	$\ket{\hat\psi}$ on the registers carrying the observables, so that
	traces against $\rho$ are expectation values in $\ket{\hat\psi}$; and
	the correlation hypothesis and the squared-distance conclusion are
	converted to and from the relation $\approx_\delta$ of
	Definition~\ref{def:povm-distance} by expanding the square, using
	unitarity of the observables, and relabelling the averaged index
	pair, at the cost of the constant factors absorbed by that
	definition. This source route is not the formal proof of
	Lemma~\ref{lem:qld-4-10}. The optional ancilla can be chosen uniformly
	in $(x,z)$ from its fixed index length (Lemma~\ref{lem:linearity-common-ancilla}).
	The separate absorption of this ancilla into zero-state padding of the
	fixed strategy state remains open: the source assumes that padding without
	a construction, and the direct proof of Lemma~\ref{lem:qld-4-10} does not
	use it.
\end{remark}

\begin{lemma}[Formalization support for Theorem~\ref{thm:linearity}: common ancilla]
	\label{lem:linearity-common-ancilla}
	\lean{MIPStarRE.QPBT.exists_exactly_linear_observables_commonAncilla}
	\leanok
	\uses{def:povm-conventions}
	This formalization-only auxiliary is not a named statement of
	\cite{Natarajan2017Linearity}; it records that the optional ancilla of
	Theorem~\ref{thm:linearity} depends only on the index length. For every
	positive integer $t$ there are a finite-dimensional complex Hilbert space
	$\hilb'_t$ and a unit vector $\ket{\mathrm{anc}_t} \in \hilb'_t$ such that,
	for every finite-dimensional complex Hilbert space $\hilb$, every density
	operator $\rho$ on $\hilb$, every $\delta \ge 0$, and every family of binary
	observables $\{\mathcal{O}^u\}_{u \in \F_2^t}$ on $\hilb$ satisfying
	the averaged correlation hypothesis of Theorem~\ref{thm:linearity}, there
	are binary observables $\{\mathcal{L}^u\}$ on $\hilb \ot \hilb'_t$ with
	$\mathcal{L}^u\mathcal{L}^{u'}=\mathcal{L}^{u+u'}$ for all $u,u'$ and
	$\E_u d_{\rho'}(\mathcal{L}^u,\mathcal{O}^u\ot\Id)^2 \le 2\delta$,
	where $\rho'=\rho\ot\ket{\mathrm{anc}_t}\bra{\mathrm{anc}_t}$.
\end{lemma}
\begin{proof}
	\leanok
	\uses{lem:naimark-rounding, lem:boolean-representation-stability,
		lem:blr-defect-identity}
	The Naimark space has an orthonormal basis indexed by
	$\F_2^t\sqcup\{\bot\}$ and the ancilla is the basis vector at $\bot$,
	independently of $\hilb$, $\rho$, $\delta$, and the family. Apply the
	Fourier--Naimark rounding construction to each family, followed by the
	averaged BLR estimate.
\end{proof}

\begin{lemma}[Linearity with unrestricted numerical parameters]
	\label{lem:linearity-unrestricted-parameters}
	\lean{MIPStarRE.QPBT.exists_exactly_linear_observables_of_correlation}
	\leanok
	\uses{lem:blr-defect-identity, lem:boolean-representation-stability}
	The corrected conclusion of Theorem~\ref{thm:linearity} also holds for
	every $t\in\mathbb N$ and every real $\delta$, assuming only that $\rho$
	is a density operator, the $\mathcal O^u$ are binary observables, and
	$\E_{u,v}\Re\Tr(\mathcal O^u\mathcal O^v\mathcal O^{u+v}\rho)
	\geq1-\delta$. Thus there are a finite ancillary space, a unit vector
	$\mathrm{anc}$ in it, and exactly linear binary observables $\mathcal L^u$
	on the extension with
	$\E_u d_{\rho\otimes|\mathrm{anc}\rangle\langle\mathrm{anc}|}
	(\mathcal L^u,\mathcal O^u\otimes\Id)^2\leq2\delta$.
	This records the full numerical domain of the provider theorem; no
	positivity of $t$ or separate nonnegativity premise on $\delta$ is needed.
\end{lemma}
\begin{proof}
	\leanok
	\uses{lem:naimark-rounding, lem:blr-defect-identity,
		lem:boolean-representation-stability}
	The rounded family and its distance/defect identity are defined also
	when $t=0$. Apply that identity and the correlation bound directly.
\end{proof}

\begin{remark}[Normalization of the distance in the linearity conclusion]
	\label{rem:linearity-distance-normalization}
	As printed in \cite{Natarajan2017Linearity}, the conclusion of the quantum
	linearity theorem bounds the mean squared state-dependent operator distance
	$\E_u \, d_{\rho'}(\mathcal{L}^u, \mathcal{O}^u \ot \Id_{\hilb'})^2$ by
	$\delta$. That exact constant is the one printed in
	\cite[Theorem~10]{Natarajan2017Linearity}
	(\texttt{references/nv-paper/fullpaper.tex:1079--1085}), once the distance
	there is read as the state-dependent operator distance. The quotation in
	\cite{Ji2020MIPStar} does not assert it: there the conclusion reads
	$\mathcal{L}^u \approx_\delta \mathcal{O}^u$, and
	Definition~\ref{def:povm-distance} bounds the associated expectation only
	by $O(\delta)$, so that quotation absorbs a constant factor rather than
	claiming an exact one. The calculation behind the printed bound in
	\cite[Section~3.1]{Natarajan2017Linearity} computes the squared distance
	between the binary projective measurements associated with the
	observables, which is half the squared operator distance, and then
	identifies the two quantities; the identification loses that factor. What
	the Fourier argument proves is the bound $\delta$ for the mean squared
	measurement distance, equivalently $2\delta$ for the mean squared operator
	distance, and the printed operator bound $\delta$ is refuted by a
	four-point counterexample
	\cite{gap:qpbt_linearity-distance-normalization}. The normalization
	identities are recorded in Lemma~\ref{lem:binary-measurement-distance},
	and Theorem~\ref{thm:linearity} is stated and proved with the corrected
	bounds: $2\delta$ for the mean squared state-dependent operator distance
	and $\delta$ for the mean squared measurement distance. The same
	conclusion, carrying the mean multiplicative defect as its hypothesis
	parameter, is recorded as
	Lemma~\ref{lem:boolean-representation-stability}, from which the theorem
	is obtained by the substitution $\eta = 2\delta$. The universal
	coefficient $2$ is optimal: the scalar quadratic families
	$f_k(x,y)=(-1)^{x\cdot y}$ on $\mathbb F_2^{2k}$ have correlation
	$4^{-k}$ and best possible error $2(1-2^{-k})$, even with arbitrary
	ancillas. The ratio to $1-4^{-k}$ tends to $2$. This sharpness proof is
	mathematical, not a claim of a formalized asymptotic theorem. The
	Fourier--Naimark construction need not attain the best error for an
	individual family. The normalization register remains pending independent
	review and adoption, separately from the open padding obligation.
\end{remark}

\begin{definition}[Unasserted printed linearity propositions]
	\label{def:linearity-printed-claim}
	\lean{MIPStarRE.QPBT.PrintedLinearityClaim}
	\lean{MIPStarRE.QPBT.PrintedLinearityLiteralDistanceClaim}
	\leanok
	Define the first proposition to be the printed universal assertion: for every
	$t\in\mathbb N$, real $\delta$, finite-dimensional space, density operator
	$\rho$, and binary family with correlation at least $1-\delta$, there
	exist an arbitrary finite pure ancillary extension and exactly linear
	binary observables whose average squared operator error is at most
	$\delta$. Define the second proposition by retaining also the duplicated
	density in the literal source definition of the distance, using
	$d_{\rho'^2}$ in our notation. Neither proposition is asserted or supplied to
	a construction. Both are false by the following example.
\end{definition}

\begin{lemma}[Four-point obstruction to the printed bound]
	\label{lem:linearity-four-point-obstruction}
	\lean{MIPStarRE.QPBT.fourPointObservable}
	\lean{MIPStarRE.QPBT.fourPointObservable_isBinaryObservable}
	\lean{MIPStarRE.QPBT.blrCorrelation_fourPointObservable}
	\lean{MIPStarRE.QPBT.one_le_avg_fourPointObservable_distance}
	\lean{MIPStarRE.QPBT.avg_fourPointObservable_distance_one}
	\lean{MIPStarRE.QPBT.not_printedLinearityClaim}
	\lean{MIPStarRE.QPBT.not_printedLinearityLiteralDistanceClaim}
	\leanok
	\uses{lem:binary-measurement-distance, def:linearity-printed-claim}
	Let $A(1,1)=-\Id$ and $A(a)=\Id$ at the other three points of
	$\mathbb F_2^2$. In every density operator its correlation is $1/4$.
	Every exactly linear binary family $L$ on the same space satisfies
	$\E_a d_\rho(L(a),A(a))^2\geq1$, with equality for $L(a)=\Id$.
	For the scalar pure input, the same bound holds on every ancillary
	extension, refuting both printed propositions at $\delta=3/4$.
\end{lemma}
\begin{proof}
	\leanok
	\uses{lem:binary-measurement-distance, lem:boolean-representation-stability}
	Put $X=L(1,0)$ and $Y=L(0,1)$. Exact linearity gives commuting
	involutions with $L(0)=\Id$ and $L(1,1)=XY$. The identity
	$((\Id-X)(\Id-Y))^\dagger((\Id-X)(\Id-Y))=4(\Id-X-Y+XY)$
	shows that their correlation with $A$ is at most $1/2$ in every density
	operator. The distance is therefore at least $1$. The scalar pure state
	and every pure ancilla have idempotent density, so duplicating that
	density does not change the obstruction.
\end{proof}

\begin{lemma}[Formalization support for the Boolean characters]
	\label{lem:boolean-characters}
	\lean{MIPStarRE.QPBT.booleanCharacter_add_right}
	\lean{MIPStarRE.QPBT.booleanCharacter_add_left}
	\lean{MIPStarRE.QPBT.booleanCharacter_comm}
	\lean{MIPStarRE.QPBT.booleanCharacter_zero_right}
	\lean{MIPStarRE.QPBT.booleanCharacter_zero_left}
	\lean{MIPStarRE.QPBT.booleanCharacter_mul_self}
	\lean{MIPStarRE.QPBT.star_booleanCharacter}
	\lean{MIPStarRE.QPBT.sum_booleanCharacter}
	\lean{MIPStarRE.QPBT.sum_booleanCharacter_swap}
	\lean{MIPStarRE.QPBT.sum_booleanCharacter_mul}
	\leanok
	This is not a named statement of \cite{Natarajan2017Linearity}; it collects
	the character identities on which the Fourier transform of
	Lemma~\ref{lem:boolean-fourier-transform} rests. Write
	$\chi_u(a) = (-1)^{u \cdot a}$ for $u, a \in \F_2^t$. Then, for all
	$u, v, a, b \in \F_2^t$,
	\[
		\chi_u(a + b) = \chi_u(a) \, \chi_u(b) ,
		\quad
		\chi_{u+v}(a) = \chi_u(a) \, \chi_v(a) ,
		\quad
		\chi_u(a) = \chi_a(u) ,
	\]
	\[
		\chi_u(0) = \chi_0(a) = 1 ,
		\quad
		\chi_u(a)^2 = 1 ,
		\quad
		\overline{\chi_u(a)} = \chi_u(a) ,
	\]
	and the character sums are
	\[
		\sum_a \chi_u(a) =
		\begin{cases}
			2^t & \text{if } u = 0, \\
			0 & \text{otherwise,}
		\end{cases}
		\qquad
		\sum_u \chi_u(a) =
		\begin{cases}
			2^t & \text{if } a = 0, \\
			0 & \text{otherwise,}
		\end{cases}
	\]
	\[
		\sum_u \chi_u(a) \, \chi_u(b) =
		\begin{cases}
			2^t & \text{if } a = b, \\
			0 & \text{otherwise.}
		\end{cases}
	\]
\end{lemma}
\begin{proof}
	\leanok
	\uses{prop:fourier-fact-vector}
	The multiplicativity, symmetry and normalization identities hold because
	$\chi$ is the additive character pairing of $\F_2^t$, and
	$\chi_u(a)^2 = \chi_u(a + a) = \chi_u(0) = 1$ because $a + a = 0$ in
	characteristic two; realness follows from $\chi_u(a)^2 = 1$. The first
	character sum is Lemma~\ref{prop:fourier-fact-vector} for $q = 2$, rescaled
	from the average to the sum, and the second follows from it by the symmetry
	$\chi_u(a) = \chi_a(u)$. The third follows from the second together with
	$\chi_u(a) \, \chi_u(b) = \chi_u(a + b)$ and the equivalence of
	$a + b = 0$ with $a = b$.
\end{proof}

\begin{lemma}[Formalization support for the Boolean Fourier transform]
	\label{lem:boolean-fourier-transform}
	\lean{MIPStarRE.QPBT.booleanCharacter, MIPStarRE.QPBT.operatorFourier}
	\lean{MIPStarRE.QPBT.operatorFourier_isHermitian}
	\lean{MIPStarRE.QPBT.operatorFourier_inversion}
	\lean{MIPStarRE.QPBT.sum_operatorFourier_mul_operatorFourier}
	\lean{MIPStarRE.QPBT.sum_operatorFourier_sq}
	\lean{MIPStarRE.QPBT.sum_operatorFourier_sq_eq_one}
	\leanok
	\uses{def:povm-conventions}
	This is not a named statement of \cite{Natarajan2017Linearity}; it isolates
	the finite Fourier identities used in the proof given there for
	Theorem~\ref{thm:linearity}. Write $\chi_u(a) = (-1)^{u \cdot a}$ for
	$u, a \in \F_2^t$ and, for a family $\{O(a)\}_{a \in \F_2^t}$ of operators,
	set $\widehat{O}^u = \E_a \, \chi_u(a) \, O(a)$. Each $\widehat{O}^u$ is
	Hermitian when every $O(a)$ is, and for two such families $A$ and $B$,
	\[
		O(a) = \sum_u \chi_u(a) \, \widehat{O}^u ,
		\qquad
		\sum_u \widehat{A}^u \, \widehat{B}^u = \E_a \, A(a) B(a) ;
	\]
	taking $B = A$ gives $\sum_u (\widehat{O}^u)^2 = \E_a \, O(a)^2$ for a
	single family $O$. If every $O(a)$ is a binary observable, then
	$\sum_u (\widehat{O}^u)^2 = \Id$.
\end{lemma}
\begin{proof}
	\leanok
	\uses{lem:boolean-characters}
	Orthogonality of the characters of $\F_2^t$ gives the inversion formula;
	the Parseval identity follows by substituting it into the product and
	exchanging the two sums; the last identity uses $O(a)^2 = \Id$.
\end{proof}

\begin{lemma}[Formalization support for the binary-measurement distance]
	\label{lem:binary-measurement-distance}
	\lean{MIPStarRE.QPBT.binaryObservableEffect}
	\lean{MIPStarRE.QPBT.binaryObservableDistSq}
	\lean{MIPStarRE.QPBT.binaryObservableDistSq_eq_stateDepDistSq_div_two}
	\lean{MIPStarRE.QPBT.stateDepDistSq_eq_two_sub_two_mul_correlation}
	\lean{MIPStarRE.QPBT.binaryObservableDistSq_eq_one_sub_correlation}
	\leanok
	\uses{def:povm-conventions}
	This is not a named statement of \cite{Natarajan2017Linearity}; it records
	the normalization discussed in
	Remark~\ref{rem:linearity-distance-normalization}. For operators $S$, $T$
	and $\rho$ on $\hilb$, write
	$d_\rho(S,T)^2 = \Re \Tr\big((S-T)^\dagger (S-T) \, \rho\big)$ and
	$P_S^\pm = (\Id \pm S)/2$ for the affine operators that are effects of a
	binary projective measurement when $S$ is a binary observable, and set
	$d^{\mathrm{bin}}_\rho(S,T)^2
	= \sum_{s \in \{+,-\}} d_\rho(P_S^s, P_T^s)^2$. Then, with no hypothesis
	on $S$, $T$ or $\rho$,
	\[
		d^{\mathrm{bin}}_\rho(S,T)^2 = \tfrac{1}{2} \, d_\rho(S,T)^2 .
	\]
	If moreover $S$ and $T$ are binary observables and $\rho$ is a density
	operator, then
	\[
		d_\rho(S,T)^2 = 2 - 2 \Re \Tr(S T \rho) ,
		\qquad
		d^{\mathrm{bin}}_\rho(S,T)^2 = 1 - \Re \Tr(S T \rho) .
	\]
\end{lemma}
\begin{proof}
	\leanok
	Expand $P_S^s - P_T^s = \pm (S - T)/2$ in each of the two outcomes and
	sum; this uses no hypothesis on $S$, $T$ or $\rho$. The two correlation
	identities use $S^2 = T^2 = \Id$, $\Tr \rho = 1$, and the Hermiticity of
	$S$, $T$ and $\rho$.
\end{proof}

\begin{lemma}[Formalization support for the operator-valued BLR defect]
	\label{lem:blr-defect-identity}
	\lean{MIPStarRE.QPBT.blrCorrelation}
	\lean{MIPStarRE.QPBT.multiplicativeDefect}
	\lean{MIPStarRE.QPBT.multiplicativeDefect_eq_two_sub_two_mul_correlation}
	\lean{MIPStarRE.QPBT.avg_multiplicativeDefect_eq_two_sub_two_mul_correlation}
	\lean{MIPStarRE.QPBT.avg_multiplicativeDefect_le_two_mul_error}
	\leanok
	\uses{def:povm-conventions}
	This is not a named statement of \cite{Natarajan2017Linearity}; it records
	the step of the proof given there for Theorem~\ref{thm:linearity} that
	converts the correlation hypothesis into approximate linearity of the
	family, namely display~(7) of that proof
	(\texttt{references/nv-paper/fullpaper.tex:1078--1080}) together with the
	expansion at \texttt{references/nv-paper/fullpaper.tex:900--912}. For a
	family $\{O(a)\}_{a \in \F_2^t}$ of operators on $\hilb$ and an operator
	$\rho$ on $\hilb$, with no hypothesis on either, set
	\[
		C_\rho(O) = \E_{u,u'} \, \Re \Tr\big(O(u) \, O(u') \,
		O(u+u') \, \rho\big) ,
		\qquad
		D_\rho(u,v) = d_\rho\big(O(u) \, O(v), \, O(u+v)\big)^2 ,
	\]
	the two-query correlation of the linearity test and the multiplicative
	defect at $(u,v)$, where
	$d_\rho(S,T)^2 = \Re \Tr\big((S-T)^\dagger (S-T) \, \rho\big)$ as in
	Lemma~\ref{lem:binary-measurement-distance} and every expectation is over
	uniform arguments in $\F_2^t$. If every $O(a)$ is a binary observable and
	$\rho$ is a density operator, then for all $u, v \in \F_2^t$
	\[
		D_\rho(u,v) = 2 - 2 \Re \Tr\big(O(v) \, O(u) \, O(v+u) \, \rho\big) ,
	\]
	hence
	\[
		\E_{u,v} \, D_\rho(u,v) = 2 - 2 \, C_\rho(O) ,
	\]
	and if in addition $C_\rho(O) \ge 1 - \delta$, then
	$\E_{u,v} \, D_\rho(u,v) \le 2\delta$. The constant $2$ is the
	normalization discussed in
	Remark~\ref{rem:linearity-distance-normalization}.
\end{lemma}
\begin{proof}
	\leanok
	Expanding
	$\big(O(u) O(v) - O(u+v)\big)^\dagger \big(O(u) O(v) - O(u+v)\big)$ and
	using $O(a)^2 = \Id$ leaves
	$2 \, \Id - O(v) O(u) O(u+v) - O(u+v) O(u) O(v)$; the two triple products
	are adjoint to one another, so their traces against $\rho$ are complex
	conjugates and have equal real parts, and $\Tr \rho = 1$ then gives the
	pointwise identity. Averaging it and exchanging the two uniform
	coordinates gives the identity for the mean defect, and the bound follows
	from the correlation hypothesis.
\end{proof}

\begin{lemma}[Formalization support for the Fourier-square measurement]
	\label{lem:fourier-square-measurement}
	\lean{MIPStarRE.QPBT.sum_operatorFourier_cube}
	\lean{MIPStarRE.QPBT.operatorFourier_cube_trace_eq_correlation}
	\lean{MIPStarRE.QPBT.operatorFourier_mul_self_posSemidef}
	\lean{MIPStarRE.QPBT.fourierSquareMeasurement}
	\lean{MIPStarRE.QPBT.fourierSquareMeasurement_effect}
	\lean{MIPStarRE.QPBT.avg_overlap_fourierSquare_eq_sum_cube}
	\lean{MIPStarRE.QPBT.one_sub_error_le_avg_overlap_fourierSquare}
	\leanok
	\uses{def:povm-conventions, lem:boolean-fourier-transform,
		lem:blr-defect-identity}
	This is not a named statement of \cite{Natarajan2017Linearity}; it isolates
	the measurement to which Naimark's theorem is applied in the proof given
	there for Theorem~\ref{thm:linearity}
	(\texttt{references/nv-paper/fullpaper.tex:1096--1112}). Let
	$\{O(a)\}_{a \in \F_2^t}$ be a family of operators on $\hilb$, with Fourier
	coefficients $\widehat{O}^u$ and characters $\chi_u$ as in
	Lemma~\ref{lem:boolean-fourier-transform}, let $\rho$ be an operator on
	$\hilb$, and write $B^u = (\widehat{O}^u)^2$. With no hypothesis on the
	family or on $\rho$,
	\[
		\sum_u (\widehat{O}^u)^3 = \E_{a,b} \, O(a) \, O(b) \, O(a+b) ,
		\qquad
		\sum_u \Re \Tr\big((\widehat{O}^u)^3 \, \rho\big) = C_\rho(O) ,
	\]
	where $C_\rho(O)$ is the two-query correlation of
	Lemma~\ref{lem:blr-defect-identity}, and
	\[
		\E_a \, \Re \Tr\Big(O(a) \,
		\Big(\sum_u \chi_u(a) \, B^u\Big) \, \rho\Big)
		= \sum_u \Re \Tr\big((\widehat{O}^u)^3 \, \rho\big) .
	\]
	If every $O(a)$ is Hermitian, then every $B^u$ is positive semidefinite;
	if every $O(a)$ is a binary observable, then the $B^u$ are the effects of a
	POVM with outcomes in $\F_2^t$, the Fourier-square measurement. Finally,
	if $C_\rho(O) \ge 1 - \delta$, then
	\[
		1 - \delta \le \E_a \, \Re \Tr\Big(O(a) \,
		\Big(\sum_u \chi_u(a) \, B^u\Big) \, \rho\Big) .
	\]
\end{lemma}
\begin{proof}
	\leanok
	\uses{lem:blr-defect-identity, lem:boolean-characters,
		lem:boolean-fourier-transform}
	Expanding one Fourier coefficient at a time through the inversion formula
	of Lemma~\ref{lem:boolean-fourier-transform} and collapsing the inner
	character sum with Lemma~\ref{lem:boolean-characters} gives the cubic
	identity; multiplying it by $\rho$, taking the trace and passing to real
	parts gives the correlation identity, and the same expansion applied to
	$\sum_u \widehat{O}^u \, B^u$ gives the overlap identity. Positivity of
	$B^u$ holds because $\widehat{O}^u$ is Hermitian, so that
	$B^u = (\widehat{O}^u)^\dagger \widehat{O}^u$, and completeness is the
	identity $\sum_u (\widehat{O}^u)^2 = \Id$ of
	Lemma~\ref{lem:boolean-fourier-transform}. The final inequality
	substitutes the correlation hypothesis into the overlap identity.
\end{proof}

\begin{lemma}[Formalization support for the Naimark rounding of the
	Fourier-square measurement]
	\label{lem:naimark-rounding}
	\lean{MIPStarRE.QPBT.naimarkAncilla}
	\lean{MIPStarRE.QPBT.norm_naimarkAncilla}
	\lean{MIPStarRE.QPBT.ancProj_naimarkAncilla}
	\lean{MIPStarRE.QPBT.fourierNaimarkData}
	\lean{MIPStarRE.QPBT.fourierNaimarkData_source}
	\lean{MIPStarRE.QPBT.roundedFourierPVM}
	\lean{MIPStarRE.QPBT.roundedFourierPVM_isProj}
	\lean{MIPStarRE.QPBT.roundedFourierPVM_sum_eq_one}
	\lean{MIPStarRE.QPBT.roundedFourierPVM_mul_eq_zero_of_ne}
	\lean{MIPStarRE.QPBT.roundedFourierMeasurement}
	\lean{MIPStarRE.QPBT.roundedFourierMeasurement_effect}
	\lean{MIPStarRE.QPBT.roundedFourierMeasurement_isProjective}
	\lean{MIPStarRE.QPBT.trace_heteroKron_mul_roundedFourierPVM}
	\lean{MIPStarRE.QPBT.roundedObservable}
	\lean{MIPStarRE.QPBT.roundedObservable_mul}
	\lean{MIPStarRE.QPBT.roundedObservable_zero}
	\lean{MIPStarRE.QPBT.roundedObservable_isHermitian}
	\lean{MIPStarRE.QPBT.roundedObservable_isBinaryObservable}
	\lean{MIPStarRE.QPBT.trace_heteroKron_mul_roundedObservable}
	\leanok
	\uses{def:povm-conventions, lem:boolean-fourier-transform,
		lem:fourier-square-measurement, lem:naimark-helper}
	This is not a named statement of \cite{Natarajan2017Linearity}; it isolates
	the dilation step of the proof given there for
	Theorem~\ref{thm:linearity}, namely the passage from the Fourier-square
	measurement to the exactly linear observables
	(\texttt{references/nv-paper/fullpaper.tex:1099--1110}). Let
	$\{O(a)\}_{a \in \F_2^t}$ be binary observables on a finite-dimensional
	Hilbert space $\hilb$, with characters $\chi_u$ and Fourier coefficients
	$\widehat{O}^u$ as in Lemma~\ref{lem:boolean-fourier-transform}, and let
	$\{B^u = (\widehat{O}^u)^2\}_{u \in \F_2^t}$ be the Fourier-square
	measurement of Lemma~\ref{lem:fourier-square-measurement}. Write
	$\hilb' = \C^{\F_2^t \sqcup \{\bot\}}$ for the auxiliary space of
	Lemma~\ref{lem:naimark-helper} and
	$\ket{\mathrm{anc}} = \ket{\bot} \in \hilb'$ for its initial vector, so
	that $\|\ket{\mathrm{anc}}\| = 1$ and
	$\ket{\mathrm{anc}}\bra{\mathrm{anc}}$ is the auxiliary projector of that
	lemma. Applying Lemma~\ref{lem:naimark-helper} to $\{B^u\}$ yields
	projections $\{\hat{P}_a\}_{a \in \F_2^t \sqcup \{\bot\}}$ on
	$\hilb \ot \hilb'$ with $\sum_a \hat{P}_a \le \Id$ --- a projective
	submeasurement, not a measurement --- and such that, for every operator
	$X$ on $\hilb$ and every $u \in \F_2^t$,
	\[
		\Tr\big((X \ot \ket{\mathrm{anc}}\bra{\mathrm{anc}}) \, \hat{P}_u\big)
		= \Tr\big(X \, B^u\big) .
	\]
	Define the rounded family $\{C^u\}_{u \in \F_2^t}$ by enlarging the zero
	outcome to the complement of the nonzero ones,
	\[
		C^u =
		\begin{cases}
			\Id - \sum_{v \ne 0} \hat{P}_v & \text{if } u = 0, \\
			\hat{P}_u & \text{otherwise,}
		\end{cases}
	\]
	and set $\mathcal{L}^a = \sum_u \chi_u(a) \, C^u$ for $a \in \F_2^t$. Then
	every $C^u$ is an orthogonal projection, $\sum_u C^u = \Id$, and
	$C^u C^v = 0$ for $u \ne v$; hence $\{C^u\}$ is a projective measurement
	with outcomes in $\F_2^t$ on $\hilb \ot \hilb'$, whose compression along
	$\ket{\mathrm{anc}}$ is $\{B^u\}$: for every operator $X$ on $\hilb$ and
	every $u \in \F_2^t$,
	\[
		\Tr\big((X \ot \ket{\mathrm{anc}}\bra{\mathrm{anc}}) \, C^u\big)
		= \Tr\big(X \, B^u\big) .
	\]
	Moreover $\mathcal{L}^0 = \Id$, each $\mathcal{L}^a$ is a binary
	observable, the linearity
	\[
		\mathcal{L}^a \, \mathcal{L}^b = \mathcal{L}^{a+b}
	\]
	holds exactly for all $a, b \in \F_2^t$, and for every operator $X$ on
	$\hilb$ and every $a \in \F_2^t$,
	\[
		\Tr\big((X \ot \ket{\mathrm{anc}}\bra{\mathrm{anc}}) \,
		\mathcal{L}^a\big)
		= \Tr\Big(X \sum_u \chi_u(a) \, B^u\Big) .
	\]
\end{lemma}
\begin{proof}
	\leanok
	\uses{lem:boolean-characters, lem:boolean-fourier-transform,
		lem:fourier-square-measurement, lem:naimark-helper}
	Projections summing to at most the identity are mutually orthogonal, so
	$\hat{P}_a \hat{P}_b = 0$ for $a \ne b$. Consequently
	$Q = \sum_{v \ne 0} \hat{P}_v$ absorbs each of its summands on either
	side, $Q \hat{P}_v = \hat{P}_v Q = \hat{P}_v$ for $v \ne 0$, whence $Q$ is
	itself a projection and $C^0 = \Id - Q$ is a projection orthogonal to
	every $\hat{P}_v$ with $v \ne 0$. Splitting off the zero outcome in
	$\sum_u C^u$ leaves $(\Id - Q) + Q = \Id$.

	For $u \ne 0$ the compression identity is the expectation-preservation
	identity supplied with the dilation of Lemma~\ref{lem:naimark-helper},
	rewritten with unnormalized traces; the normalizing factors of the two
	spaces cancel, the enlarged space having dimension
	$(\lvert \F_2^t \rvert + 1)$ times that of $\hilb$. For $u = 0$,
	linearity of the trace reduces the claim to
	$\Tr(X \ot \ket{\mathrm{anc}}\bra{\mathrm{anc}}) = \Tr(X)$ together with
	the identity $B^0 = \Id - \sum_{v \ne 0} B^v$, which is the completeness
	$\sum_u B^u = \Id$ of Lemma~\ref{lem:boolean-fourier-transform}.

	Expanding $\mathcal{L}^a \mathcal{L}^b$ and using that $C^u C^v$ vanishes
	for $u \ne v$ and equals $C^u$ for $u = v$ leaves the diagonal terms
	$\sum_u \chi_u(a) \chi_u(b) \, C^u$, and
	$\chi_u(a) \chi_u(b) = \chi_u(a+b)$ by
	Lemma~\ref{lem:boolean-characters}. Taking $a = 0$ in the definition of
	$\mathcal{L}^a$ and using $\chi_u(0) = 1$ gives
	$\mathcal{L}^0 = \sum_u C^u = \Id$. Hermiticity of $\mathcal{L}^a$ follows
	from the realness of the characters and the self-adjointness of the
	$C^u$, and $(\mathcal{L}^a)^2 = \mathcal{L}^{a+a} = \mathcal{L}^0 = \Id$
	because $a + a = 0$ in characteristic two, so each $\mathcal{L}^a$ is a
	binary observable. The last identity is the compression identity for the
	$C^u$ summed against the characters.
\end{proof}

\begin{lemma}[Formalization support for Boolean representation stability]
	\label{lem:boolean-representation-stability}
	\lean{MIPStarRE.QPBT.ancProj_posSemidef}
	\lean{MIPStarRE.QPBT.trace_ancProj}
	\lean{MIPStarRE.QPBT.heteroKron_ancProj_posSemidef}
	\lean{MIPStarRE.QPBT.trace_heteroKron_ancProj}
	\lean{MIPStarRE.QPBT.isBinaryObservable_heteroKron_one}
	\lean{MIPStarRE.QPBT.re_trace_roundedObservable_mul_heteroKron}
	\lean{MIPStarRE.QPBT.stateDepDistSq_roundedObservable}
	\lean{MIPStarRE.QPBT.avg_stateDepDistSq_roundedObservable_eq_avg_multiplicativeDefect}
	\lean{MIPStarRE.QPBT.exists_exact_boolean_representation}
	\leanok
	\uses{def:povm-conventions, lem:binary-measurement-distance,
		lem:blr-defect-identity, lem:naimark-rounding}
	This is not a named statement of \cite{Natarajan2017Linearity}; it records
	the conclusion that the Fourier--Naimark argument given there for
	Theorem~\ref{thm:linearity}
	(\texttt{references/nv-paper/fullpaper.tex:1096--1112}) establishes, with
	the correlation hypothesis replaced by the mean multiplicative defect of
	Lemma~\ref{lem:blr-defect-identity}. Let $t$ be a positive integer, let
	$\eta \ge 0$, let $\rho$ be a density operator on $\hilb$, and let
	$\{\mathcal{O}^u\}_{u \in \F_2^t}$ be binary observables on $\hilb$ with
	\[
		\E_{u,u'} \, d_\rho\big(\mathcal{O}^u \, \mathcal{O}^{u'},\
		\mathcal{O}^{u+u'}\big)^2 \le \eta .
	\]
	Then there are a finite-dimensional space $\hilb'$, a unit vector
	$\ket{\mathrm{anc}} \in \hilb'$ and binary observables
	$\{\mathcal{L}^u\}_{u \in \F_2^t}$ on $\hilb \ot \hilb'$ such that
	$\mathcal{L}^u \mathcal{L}^{u'} = \mathcal{L}^{u+u'}$ for all
	$u, u' \in \F_2^t$ and, writing
	$\rho' = \rho \ot \ket{\mathrm{anc}}\bra{\mathrm{anc}}$, which is again a
	density operator,
	\[
		\E_u \, d_{\rho'}\big(\mathcal{L}^u,\ \mathcal{O}^u \ot
		\Id_{\hilb'}\big)^2 \le \eta .
	\]
	The bound is the specialization of the identity
	\[
		\E_u \, d_{\rho'}\big(\mathcal{L}^u,\ \mathcal{O}^u \ot
		\Id_{\hilb'}\big)^2
		= \E_{u,u'} \, d_\rho\big(\mathcal{O}^u \, \mathcal{O}^{u'},\
		\mathcal{O}^{u+u'}\big)^2 ,
	\]
	which holds for the family constructed below and whose two sides both
	equal $2 - 2 \, C_\rho(\mathcal{O})$ for the two-query correlation
	$C_\rho(\mathcal{O})$ of Lemma~\ref{lem:blr-defect-identity}.

	This is the mathematical content of Theorem~\ref{thm:linearity}: the
	ancillary space and state are produced explicitly, and by
	Lemma~\ref{lem:blr-defect-identity} a correlation hypothesis
	$C_\rho(\mathcal{O}) \ge 1 - \delta$ is exactly the hypothesis above with
	$\eta = 2\delta$. The conclusion obtained from such a hypothesis is
	therefore the bound $2\delta$, which is the bound displayed in
	Theorem~\ref{thm:linearity}; the constant $2$ separating it from the bound
	$\delta$ printed in the source is the normalization recorded in
	Remark~\ref{rem:linearity-distance-normalization}. The present statement is
	the defect-parameterized form of that theorem, and the substitution
	$\eta = 2\delta$ is the whole of the proof of
	Theorem~\ref{thm:linearity} from it.
\end{lemma}
\begin{proof}
	\leanok
	\uses{lem:binary-measurement-distance, lem:blr-defect-identity,
		lem:fourier-square-measurement, lem:naimark-rounding}
	Let $\hilb'$ have orthonormal basis indexed by $\F_2^t$ together with one
	extra element $\bot$, and take $\ket{\mathrm{anc}} = \ket{\bot}$, the
	initial state of the auxiliary register in the one-measurement Naimark
	construction of Lemma~\ref{lem:naimark-helper}. Apply that construction to
	the Fourier-square POVM $\{B^u\}$ of
	Lemma~\ref{lem:fourier-square-measurement}. It returns
	projectors indexed by $\F_2^t \cup \{\bot\}$; the residual projector at
	$\bot$ is folded into the outcome $u = 0$, so that the resulting family
	$\{C^u\}_{u \in \F_2^t}$ is still a family of orthogonal projectors summing
	to $\Id_{\hilb \ot \hilb'}$ and still compresses to $B^u$ along
	$\ket{\mathrm{anc}}$, the latter because $\sum_u B^u = \Id_{\hilb}$. Set
	$\mathcal{L}^a = \sum_u \chi_u(a) \, C^u$; orthogonality of the $C^u$ and
	multiplicativity of the characters give
	$\mathcal{L}^a \mathcal{L}^b = \mathcal{L}^{a+b}$ and
	$(\mathcal{L}^a)^2 = \Id$, and each $\mathcal{L}^a$ is Hermitian because
	the characters are real. Since $\rho'$ is a density operator and both
	$\mathcal{L}^a$ and $\mathcal{O}^a \ot \Id_{\hilb'}$ are binary
	observables, Lemma~\ref{lem:binary-measurement-distance} gives
	\[
		d_{\rho'}\big(\mathcal{L}^a, \mathcal{O}^a \ot \Id_{\hilb'}\big)^2
		= 2 - 2 \, \Re \Tr\big(\mathcal{L}^a \,
		(\mathcal{O}^a \ot \Id_{\hilb'}) \, \rho'\big) ,
	\]
	and taking the adjoint of the triple product, absorbing
	$\mathcal{O}^a \ot \Id_{\hilb'}$ into the weight and compressing
	$\mathcal{L}^a$ along $\ket{\mathrm{anc}}$ rewrites the trace as
	$\Re \Tr\big(\mathcal{O}^a \, (\sum_u \chi_u(a) B^u) \, \rho\big)$.
	Averaging over $a$ and applying the overlap and cubic identities of
	Lemma~\ref{lem:fourier-square-measurement} turns the average into
	$2 - 2 \, C_\rho(\mathcal{O})$, which is the mean multiplicative defect by
	Lemma~\ref{lem:blr-defect-identity}. No commutation among the
	$\mathcal{O}^u$ is used. The hypothesis then bounds the average by $\eta$.
\end{proof}

\begin{lemma}[Field-valued approximate commutation of the expanded point
	measurements]
	\label{lem:expanded-point-field-commutation}
	\lean{MIPStarRE.QPBT.ProjectiveSetting.expPoint_comm}
	\lean{MIPStarRE.QPBT.sum_norm_sum_smul_sq_of_orthogonal}
	\lean{MIPStarRE.QPBT.sum_fixedCharacter_mul}
	\leanok
	\uses{lem:qld-comm-cons, def:expanded-state,
		def:expanded-point-measurement, def:povm-distance,
		def:symmetric-equivalents}
	There is a constant $C \geq 1$ such that the following holds for every
	admissible parameter set, every $\eps$, every projective strategy of value
	at least $1-\eps$, and each of the four register placements of
	Definition~\ref{def:symmetric-equivalents}. Writing $(\cdot)_p$ for the
	chosen placement, on average over uniformly random $x, z \in \F_q^m$,
	\[
		\E_{x,z} \sum_{a, b \in \F_q}
		\bigl\lVert
		\bigl(\hat{M}^{(\mathrm{Point},X),x}_{a}\,
		\hat{M}^{(\mathrm{Point},Z),z}_{b}
		- \hat{M}^{(\mathrm{Point},Z),z}_{b}\,
		\hat{M}^{(\mathrm{Point},X),x}_{a}\bigr)_{p}
		\ket{\hat{\psi}}
		\bigr\rVert^2
		\;\leq\; C\sqrt{\eps}\,.
	\]
	The outcomes $a, b$ range over $\F_q$, whereas the commutation conclusion
	of item~2 of Lemma~\ref{lem:qld-comm-cons} is stated for the binary
	refinements in $\F_2$. The error is the same $\sqrt{\eps}$: the transfer
	to field-valued outcomes loses no factor depending on the field size.
	This is the estimate actually used by the formal proof of
	Lemma~\ref{lem:qld-4-10}; paper
	\texttt{14\_analysis\_of\_the\_pauli\_basis\_test.tex:466--505}.

	The formal proof isolates two finite-sum ingredients, both
	formalization-only auxiliaries and neither a named statement of
	\cite{Ji2020MIPStar}. The first is Parseval's identity in the form used
	here: if $c_{j,i}$ are complex coefficients whose columns are orthogonal
	with one common squared norm $N$, that is
	$\sum_j \overline{c_{j,i}}\, c_{j,i'} = N \, [i = i']$, then
	$\sum_j \lVert \sum_i c_{j,i} v_i \rVert^2 = N \sum_i \lVert v_i \rVert^2$
	for any vectors $v_i$ of a complex inner product space. The second is the
	one-dimensional case of the cancellation identity of
	Lemma~\ref{lem:cancellation}: for the fixed trace character $\chi$ of
	$\F_q$ and every $c \in \F_q$,
	$\sum_{a \in \F_q} \chi(a c) = q \, [c = 0]$.
\end{lemma}
\begin{proof}
	\leanok
	\uses{lem:qld-comm-cons, lem:twisted-commutation, lem:cancellation,
		def:expanded-observables, def:expanded-point-measurement}
	Each field-valued effect is the character average of the expanded point
	observables over the dual scalar, so the commutator of the two effects is
	the corresponding average of the commutators of the observables. The
	characters of $\F_q$ are orthogonal, so, for each fixed $x,z$ and placement
	$p$, Parseval's identity identifies the unnormalized sum over the outcomes
	$a,b$ of the squared effect-commutator norms with the uniform average over
	the dual scalars $r,s$ of the squared observable-commutator norms.
	Averaging this identity over uniformly random $x,z$ gives the
	twisted-commutator average, with no factor depending on $q$
	(Lemma~\ref{lem:twisted-commutation}) already bounded by $C\sqrt{\eps}$ in
	the proof of item~2 of Lemma~\ref{lem:qld-comm-cons}.
\end{proof}

\begin{lemma}[Formalization support for combined-point distances]
	\label{lem:combined-point-distance-support}
	\lean{MIPStarRE.QPBT.opFamilyDistSq_uniform_le_four}
	\lean{MIPStarRE.QPBT.sum_mul_conjTranspose_mul_self_eq_one}
	\lean{MIPStarRE.QPBT.sum_mul_conjTranspose_mul_self_eq_one'}
	\lean{MIPStarRE.QPBT.ProjectiveSetting.chain_bounds}
	\lean{MIPStarRE.QPBT.error_terms_le_rpow_eighth}
	\leanok
	\uses{def:povm-distance, def:povm-conventions,
		def:expanded-point-measurement, def:expanded-state,
		def:symmetric-equivalents}
	These formalization-only auxiliary estimates support the direct proof
	of Lemma~\ref{lem:qld-4-10}; they are not separately named source results.
	For a nonempty finite question set $\Omega$, finite outcome set $I$,
	operator families $A,B$ on a finite-dimensional Hilbert space, and a
	unit vector $\psi$, write
	$D_\psi(A,B)=\E_{\omega\in\Omega}\sum_{a\in I}
	\|(A^\omega_a-B^\omega_a)\psi\|^2$.
	\begin{enumerate}
		\item If $\sum_a(A^\omega_a)^\dagger A^\omega_a\leq\Id$ and
		$\sum_a(B^\omega_a)^\dagger B^\omega_a\leq\Id$ for every
		$\omega$, then $D_\psi(A,B)\leq4$.
		\item For any two finite projective measurements $\{A_a\}$ and
		$\{B_b\}$ on a common finite-dimensional space,
		\[
			\sum_{a,b}(A_aB_b)^\dagger(A_aB_b)=\Id,
			\qquad
			\sum_{a,b}(B_bA_a)^\dagger(B_bA_a)=\Id.
		\]
		No commutation between the measurements is required.
		\item In the chapter's setting, let $Q^{x,z}$ be any family of
		POVMs on each player's expanded local space, and let $Q_i$ denote
		its effects at placement $i$. In the following bounds, $D$ denotes
		$D_{\hat\psi}$ with uniform question $(x,z)$ and outcome $(a,b)$;
		write $X_iZ_i$ for the family $X_{i,a}Z_{i,b}$, and similarly for
		$Z_iX_i$. Let $R_i$ be the $Z$--$X$--$Z$ sandwich POVM formed from
		the expanded point measurements at placement $i$. Suppose real
		numbers $\eta,C_1,C_2,C_3$ satisfy, for every placement $i$ and
		every placement $j$ opposite to $i$,
		\begin{align*}
			D(Q_i,R_i)&\leq\eta,& D(R_i,Z_iX_i)&\leq C_1\sqrt\eps,\\
			D(X_iZ_i,Z_iX_i)&\leq C_2\sqrt\eps,&
			D(Z_iX_i,X_jZ_j)&\leq C_3\eps.
		\end{align*}
		Then, for every such directed pair $i,j$,
		\begin{align*}
			D(Q_i,X_jZ_j)&\leq2\eta+4C_1\sqrt\eps+4C_3\eps,\\
			D(Q_i,Z_jX_j)&\leq4\eta+8C_1\sqrt\eps+8C_3\eps+2C_2\sqrt\eps,\\
			D(Q_i,Q_j)&\leq12\eta+20C_1\sqrt\eps+16C_3\eps+4C_2\sqrt\eps.
		\end{align*}
		\item If $0\leq\eps\leq1$ and $C_0\geq1$, then
		\[
			\eps\leq\eps^{1/8},\qquad \sqrt\eps\leq\eps^{1/8},\qquad
			\bigl(C_0(\eps+\sqrt\eps)\bigr)^{1/4}
			\leq2C_0\eps^{1/8}.
		\]
	\end{enumerate}
\end{lemma}
\begin{proof}
	\leanok
	\uses{fact:triangle}
	For the first bound, use
	$\|(A_a-B_a)\psi\|^2\leq2\|A_a\psi\|^2+2\|B_a\psi\|^2$,
	sum over outcomes, and average. For the second, projectivity gives
	$(A_aB_b)^\dagger(A_aB_b)=B_bA_aB_b$; sum first over $a$ and then
	over $b$. Interchanging the measurements proves the other identity.
	For the third, repeatedly apply the squared triangle inequality
	$D(A,C)\leq2D(A,B)+2D(B,C)$ along the indicated comparisons, using
	symmetry to compare $Z_jX_j$ with $Q_j$.
	For the last, the case $\eps=0$ is immediate. For $0<\eps\leq1$,
	monotonicity in the exponent gives the first two inequalities and
	$\eps+\sqrt\eps\leq2\sqrt\eps$. Taking fourth roots and using
	$(2C_0)^{1/4}\leq2C_0$ proves the final bound.
\end{proof}

\begin{lemma}[Formalization support: the crosswise EPR exchange]
	\label{lem:qld-epr-cross-swap}
	\lean{MIPStarRE.QPBT.ProjectiveSetting.eprCrossSwap}
	\lean{MIPStarRE.QPBT.ProjectiveSetting.reindexState_eprCrossSwap_psiHat}
	\lean{MIPStarRE.QPBT.ProjectiveSetting.norm_place_AB''_eq_norm_place_AA'}
	\lean{MIPStarRE.QPBT.ProjectiveSetting.norm_place_BB'_eq_norm_place_BA''}
	\leanok
	\uses{def:expanded-state, def:symmetric-equivalents, def:EPR}
	This formalization-only auxiliary is not a named statement of
	\cite{Ji2020MIPStar}; it transfers norms between register placements
	without assuming that the strategy is symmetric.
	Let $U$ be the permutation of the six registers of
	Definition~\ref{def:symmetric-equivalents} that exchanges
	$\mathsf{A}' \leftrightarrow \mathsf{B}''$ and
	$\mathsf{A}'' \leftrightarrow \mathsf{B}'$ and fixes the two strategy
	registers $\mathsf{A}$ and $\mathsf{B}$. Then $U$ is an involution, it
	fixes the expanded state $\ket{\hat\psi}$ of
	Definition~\ref{def:expanded-state}, and, writing $(\cdot)_p$ for a
	placement,
	\[
		\bigl\lVert (P)_{\mathsf{A}\mathsf{B}''}\ket{\hat\psi}\bigr\rVert
		= \bigl\lVert (P)_{\mathsf{A}\mathsf{A}'}\ket{\hat\psi}\bigr\rVert,
		\qquad
		\bigl\lVert (P')_{\mathsf{B}\mathsf{B}'}\ket{\hat\psi}\bigr\rVert
		= \bigl\lVert (P')_{\mathsf{B}\mathsf{A}''}\ket{\hat\psi}\bigr\rVert,
	\]
	for every operator $P$ on Alice's expanded local space and every operator
	$P'$ on Bob's. The only input is that the EPR vector of
	Definition~\ref{def:EPR} takes the same value on $(x,y)$ and on $(y,x)$;
	no symmetry of the strategy is used. This exchange is a separate
	involution from the four-cycle on the ancillas that transports the
	symmetric equivalents of Definition~\ref{def:symmetric-equivalents}.
\end{lemma}
\begin{proof}
	\leanok
	\uses{def:expanded-state, def:EPR}
	The permutation is its own inverse by inspection of the coordinate tuple.
	On the expanded state it multiplies the same strategy amplitude by the two
	EPR amplitudes with their registers exchanged, so the symmetry of the EPR
	vector returns the original state. Reindexing an operator by $U$ carries
	$(P)_{\mathsf{A}\mathsf{B}''}$ to $(P)_{\mathsf{A}\mathsf{A}'}$ and
	$(P')_{\mathsf{B}\mathsf{B}'}$ to $(P')_{\mathsf{B}\mathsf{A}''}$
	entrywise, and reindexing a state by a permutation of its coordinates
	preserves norms.
\end{proof}

\begin{lemma}[Formalization support: the sandwich point POVM]
	\label{lem:qld-sandwich-povm}
	\lean{MIPStarRE.QPBT.ProjectiveSetting.sandwichPoint}
	\lean{MIPStarRE.QPBT.ProjectiveSetting.sandwichPoint_ordered_dist_le}
	\lean{MIPStarRE.QPBT.ProjectiveSetting.ordered_cross_dist_le}
	\leanok
	\uses{def:expanded-point-measurement, def:expanded-state,
		def:povm-distance, def:symmetric-equivalents,
		lem:expanded-point-field-commutation, lem:qld-comm-cons,
		lem:qld-placement-completeness}
	This formalization-only auxiliary is not a named statement of
	\cite{Ji2020MIPStar}; it is the first step of the proof of Lemma~4.10 of
	\cite{Ji2020MIPStar}
	(\texttt{references/qpbt-paper/14\_analysis\_of\_the\_pauli\_basis\_test.tex},
	lines~731--770), carried out with field-valued outcomes rather than with
	the binary refinements of the source proof; the replacement is recorded in
	the paper-gap note
	\texttt{docs/paper-gaps/qpbt\_combined-points-field-valued.tex}.
	Write $\hat{M}^{(\mathrm{Point},X),x}_{a}$ and
	$\hat{M}^{(\mathrm{Point},Z),z}_{b}$ for the expanded point measurements
	of Definition~\ref{def:expanded-point-measurement} on one player's
	expanded local space, and set, for $x, z \in \F_q^m$,
	\[
		R^{x,z}_{a,b}
		= \hat{M}^{(\mathrm{Point},Z),z}_{b}\,
		\hat{M}^{(\mathrm{Point},X),x}_{a}\,
		\hat{M}^{(\mathrm{Point},Z),z}_{b},
		\qquad a, b \in \F_q .
	\]
	\begin{enumerate}
		\item $\{R^{x,z}_{a,b}\}_{a,b \in \F_q}$ is a POVM: every effect is
		positive semidefinite and the effects sum to the identity. Only
		projectivity of the two point measurements is used; no commutation
		between them is assumed.
		\item There is a constant $C \ge 1$ such that, for every admissible
		parameter set, every $\eps$, every projective strategy of value at
		least $1-\eps$, and every one of the four placements $p$ of
		Definition~\ref{def:symmetric-equivalents},
		\[
			\E_{x,z} \sum_{a,b \in \F_q}
			\bigl\lVert \bigl(R^{x,z}_{a,b}
			- \hat{M}^{(\mathrm{Point},Z),z}_{b}
			\hat{M}^{(\mathrm{Point},X),x}_{a}\bigr)_{p}
			\ket{\hat\psi} \bigr\rVert^2
			\;\le\; C \sqrt{\eps}\,.
		\]
		\item There is a constant $C \ge 1$ such that, with the same
		quantification and for every directed pair of opposite placements
		$p_1, p_2$,
		\[
			\E_{x,z} \sum_{a,b \in \F_q}
			\bigl\lVert \bigl(
			(\hat{M}^{(\mathrm{Point},Z),z}_{b}
			\hat{M}^{(\mathrm{Point},X),x}_{a})_{p_1}
			- (\hat{M}^{(\mathrm{Point},X),x}_{a}
			\hat{M}^{(\mathrm{Point},Z),z}_{b})_{p_2}
			\bigr) \ket{\hat\psi} \bigr\rVert^2
			\;\le\; C \eps\,.
		\]
	\end{enumerate}
	Both averages are over a uniformly random point pair
	$(x,z) \in \F_q^m \times \F_q^m$.
\end{lemma}
\begin{proof}
	\leanok
	\uses{lem:expanded-point-field-commutation, lem:qld-comm-cons,
		lem:qld-placement-completeness, fact:triangle, def:povm-distance}
	Each effect equals
	$(\hat{M}^{X,x}_{a}\hat{M}^{Z,z}_{b})^\dagger
	(\hat{M}^{X,x}_{a}\hat{M}^{Z,z}_{b})$ and is therefore positive
	semidefinite; summing first over $a$ and then over $b$ and using
	idempotence of the $Z$-effects gives the identity.
	For item~2, projectivity gives
	$R^{x,z}_{a,b} - \hat{M}^{Z,z}_{b}\hat{M}^{X,x}_{a}
	= \hat{M}^{Z,z}_{b}\,[\hat{M}^{X,x}_{a},\hat{M}^{Z,z}_{b}]$; a placement
	is multiplicative and carries a projector to a contraction
	(Lemma~\ref{lem:qld-placement-completeness}), so each summand is at most
	the corresponding placed commutator norm, whose average is
	$C\sqrt\eps$ by Lemma~\ref{lem:expanded-point-field-commutation}.
	For item~3, pass through the intermediate family
	$(\hat{M}^{Z,z}_{b})_{p_1}(\hat{M}^{X,x}_{a})_{p_2}$, transporting one
	factor at a time: each transport costs the self-consistency of that single
	point measurement, which is item~1 of Lemma~\ref{lem:qld-comm-cons}, with a
	placed contraction as the remaining factor. The squared triangle
	inequality (Lemma~\ref{fact:triangle}) combines the two estimates and
	accounts for the constant.
\end{proof}

\begin{lemma}[Formalization support: consistency of the sandwich point POVM]
	\label{lem:qld-sandwich-consistency-defect}
	\lean{MIPStarRE.QPBT.sandwich_overlap_identity}
	\lean{MIPStarRE.QPBT.sandwich_defect_pointwise_le}
	\lean{MIPStarRE.QPBT.ProjectiveSetting.sandwichDefectBound}
	\lean{MIPStarRE.QPBT.ProjectiveSetting.sandwich_offDiagonal_le_sandwichDefectBound}
	\lean{MIPStarRE.QPBT.ProjectiveSetting.avg_sandwichDefectBound_le}
	\leanok
	\uses{lem:qld-sandwich-povm, lem:qld-placement-completeness,
		def:consistency, def:expanded-state, def:expanded-point-measurement,
		lem:qld-comm-cons, lem:expanded-point-field-commutation}
	This formalization-only auxiliary is not a named statement of
	\cite{Ji2020MIPStar}; it replaces the Cauchy--Schwarz chain in the first
	step of the proof of Lemma~4.10 of \cite{Ji2020MIPStar}
	(\texttt{references/qpbt-paper/14\_analysis\_of\_the\_pauli\_basis\_test.tex},
	lines~743--790) by an exact identity and a bound derived from it.

	\emph{Abstract form.} Let $\ket\psi$ be a vector of a finite-dimensional
	complex Hilbert space and let $\{X^1_a\}$, $\{X^2_a\}$, $\{Z^1_b\}$,
	$\{Z^2_b\}$ be projective measurements on it with finite outcome sets,
	such that $X^1_a Z^2_b = Z^2_b X^1_a$ and $Z^1_b X^2_a = X^2_a Z^1_b$ for
	all $a, b$, and $Z^1_b Z^2_b = Z^2_b Z^1_b$ for every $b$. Put
	$R^i_{a,b} = Z^i_b X^i_a Z^i_b$ and $W_b = Z^1_b Z^2_b$. Then
	\[
		\sum_{a,b} \operatorname{Re}\!\left(\bra\psi R^1_{a,b} R^2_{a,b} \ket\psi\right)
		= \lVert\psi\rVert^2
		- \tfrac12 \sum_b \lVert (Z^1_b - Z^2_b)\ket\psi \rVert^2
		- \tfrac12 \sum_{a,b} \lVert (X^1_a - X^2_a) W_b \ket\psi \rVert^2 ,
	\]
	and consequently
	\[
		\lVert\psi\rVert^2
		- \sum_{a,b} \operatorname{Re}\!\left(\bra\psi R^1_{a,b} R^2_{a,b} \ket\psi\right)
		\;\le\; \tfrac12 \sum_b \lVert (Z^1_b - Z^2_b)\ket\psi \rVert^2
		+ \tfrac32 \Bigl( \sum_a \lVert (X^1_a - X^2_a)\ket\psi \rVert^2
		+ K_1 + K_2 \Bigr) ,
	\]
	where $K_i = \sum_{a,b} \lVert [X^i_a, Z^i_b] \ket\psi \rVert^2$. No
	commutation between $X^i$ and $Z^i$ with the same superscript is assumed.

	\emph{In the chapter's setting.} Fix a directed pair of opposite
	placements $p_1, p_2$ and let $c_{x,z}$ denote the right-hand side of the
	second display evaluated at $\ket{\hat\psi}$ with the two placements of the
	expanded point measurements at the point pair $(x,z)$. Then the
	off-diagonal mass
	$\sum_{(a,b) \ne (a',b')} \bra{\hat\psi} (R^{x,z}_{a,b})_{p_1}
	(R^{x,z}_{a',b'})_{p_2} \ket{\hat\psi}$ of the two placed sandwich POVMs
	of Lemma~\ref{lem:qld-sandwich-povm} is at most $c_{x,z}$, and on average
	over a uniformly random point pair $\E_{x,z} c_{x,z} \le C(\eps+\sqrt\eps)$
	for a constant $C \ge 1$ independent of all parameters.
\end{lemma}
\begin{proof}
	\leanok
	\uses{lem:qld-sandwich-povm, lem:qld-comm-cons,
		lem:expanded-point-field-commutation, lem:distance-point-defect,
		lem:qld-placement-completeness, def:consistency}
	The commutation hypotheses and idempotence collapse each summand to
	$\bra{W_b \psi} X^1_a X^2_a \ket{W_b \psi}$. The identity
	$\sum_c \operatorname{Re}\!\left(\bra{w} P_c Q_c \ket{w}\right) = \lVert w \rVert^2
	- \tfrac12 \sum_c \lVert (P_c - Q_c) w \rVert^2$ for two projective
	measurements, applied once with $w = W_b \psi$ and once with $w = \psi$,
	gives the first display. For the second, write
	$(X^1_a - X^2_a) W_b = W_b (X^1_a - X^2_a)
	+ Z^2_b [X^1_a, Z^1_b] - Z^1_b [X^2_a, Z^2_b]$ and use that the $Z^i_b$
	are contractions whose squares sum to the identity; splitting the three
	terms costs the factor $3/2$.
	In the chapter's setting the abstract hypotheses hold because operators on
	opposite placements commute and each placement is multiplicative
	(Lemma~\ref{lem:qld-placement-completeness}); the off-diagonal mass is the
	consistency defect of Definition~\ref{def:consistency}, which
	Lemma~\ref{lem:distance-point-defect} rewrites in the form just bounded.
	The average of $c_{x,z}$ is controlled by item~1 of
	Lemma~\ref{lem:qld-comm-cons} for the two distance terms and by
	Lemma~\ref{lem:expanded-point-field-commutation} for the two commutator
	terms.
\end{proof}

\begin{lemma}[Formalization support: orthonormalizing the sandwich point POVM]
	\label{lem:qld-sandwich-orthonormalization}
	\lean{MIPStarRE.QPBT.avgOver_rpow_quarter_le}
	\lean{MIPStarRE.QPBT.baAaBipartition}
	\lean{MIPStarRE.QPBT.ProjectiveSetting.exists_projective_close_sandwich}
	\lean{MIPStarRE.QPBT.ProjectiveSetting.exists_projective_close_sandwich_alice}
	\lean{MIPStarRE.QPBT.ProjectiveSetting.exists_projective_close_sandwich_bob}
	\leanok
	\uses{lem:qld-sandwich-consistency-defect, lem:ortho,
		lem:ortho-explicit-constant, def:qld-opposite-bipartition,
		def:povm-distance, def:symmetric-equivalents, def:expanded-state}
	This formalization-only auxiliary is not a named statement of
	\cite{Ji2020MIPStar}; it applies Lemma~\ref{lem:ortho} once for every
	point pair, as in the proof of Lemma~4.10 of \cite{Ji2020MIPStar}
	(\texttt{references/qpbt-paper/14\_analysis\_of\_the\_pauli\_basis\_test.tex},
	lines~782--785).

	First, a scalar fact: if the weights of a distribution $\mathcal D$ sum to
	at most one and $b \ge 0$, then
	$\E_{\mathcal D}\, b^{1/4} \le \bigl(\E_{\mathcal D}\, b\bigr)^{1/4}$.

	Second, let $p_1, p_2$ be a directed pair of opposite placements, let
	$R^{x,z}$ be the sandwich POVM of
	Lemma~\ref{lem:qld-sandwich-povm} and let $c_{x,z}$ be the pointwise bound
	of Lemma~\ref{lem:qld-sandwich-consistency-defect} for that pair. Given a
	regrouping of the six registers that exhibits the registers of $p_1$ as
	the first tensor factor and those of $p_2$ as the second, with the
	remaining EPR pair as an unused third factor, there is a family
	$\{\hat{Q}^{x,z}_{a,b}\}_{a,b \in \F_q}$ of projective measurements on the
	expanded local space of the side of $p_1$, one for each point pair, with
	\[
		\E_{x,z} \sum_{a,b \in \F_q}
		\bigl\lVert \bigl(\hat{Q}^{x,z}_{a,b} - R^{x,z}_{a,b}\bigr)_{p_1}
		\ket{\hat\psi} \bigr\rVert^2
		\;\le\; 220 \, \bigl(\E_{x,z} c_{x,z}\bigr)^{1/4} .
	\]
	Such a regrouping exists for the directed pair
	$(\mathsf{A}\mathsf{A}', \mathsf{B}\mathsf{A}'')$, giving Alice's family,
	and for the reversed pair
	$(\mathsf{B}\mathsf{A}'', \mathsf{A}\mathsf{A}')$, which regroups the six
	registers as
	$\mathsf{B}\mathsf{A}'' \mid \mathsf{A}\mathsf{A}'(\mathsf{B}'\mathsf{B}'')$
	and gives Bob's family.
\end{lemma}
\begin{proof}
	\leanok
	\uses{lem:qld-sandwich-consistency-defect, lem:ortho-explicit-constant,
		def:qld-opposite-bipartition}
	On the nonnegative reals $t^{1/4} = \sqrt{\sqrt t}$, so the scalar fact is
	Jensen's inequality for the square root applied twice.
	For the second part, read the expanded state in the given regrouping: it
	stays a unit vector, operators at $p_1$ act on the first tensor factor and
	operators at $p_2$ on the second, and reindexing preserves the squared
	distances. For each fixed $(x,z)$ the consistency defect of the sandwich
	POVM at $p_1$ against its copy at $p_2$ is at most $c_{x,z}$ by
	Lemma~\ref{lem:qld-sandwich-consistency-defect}, so
	Lemma~\ref{lem:ortho-explicit-constant} supplies a projective measurement
	on the first factor at squared distance at most $220\, c_{x,z}^{1/4}$.
	Choosing one for every point pair and averaging, the scalar fact turns
	$\E_{x,z} 220\, c_{x,z}^{1/4}$ into $220 (\E_{x,z} c_{x,z})^{1/4}$.
\end{proof}

\begin{lemma}[Combined point measurements]
	\label{lem:qld-4-10}
	\lean{MIPStarRE.QPBT.CombinedPointsWitness}
	\lean{MIPStarRE.QPBT.exists_combinedPointsWitness}
	\leanok
	\uses{def:expanded-state, def:expanded-point-measurement,
		def:povm-distance, def:symmetric-equivalents}
	There exists a function $\delta_Q(\eps) = \poly(\eps)$ such that the
	following holds. For every $x, z \in \F_q^m$ and
	$\hilb \in \{\hilb_{\mathsf{A}}, \hilb_{\mathsf{B}}\}$ there are
	projective measurements $\{\hat{Q}^{x,z}_{a,b}\}_{a,b \in \F_q}$ acting
	on $\hilb \ot (\C^q)^{\ot M}$ such that:
	\begin{enumerate}
		\item (\textbf{Self-consistency}) On average over uniformly random
		$x, z \in \F_q^m$,
		\begin{equation}
			\label{eq:qld-q-self-cons}
			(\hat{Q}^{x,z}_{a,b})_{\mathsf{A}\mathsf{A}'}
			\approx_{\delta_Q}
			(\hat{Q}^{x,z}_{a,b})_{\mathsf{B}\mathsf{A}''}\,.
		\end{equation}
		\item (\textbf{Consistency with $\hat{M}$}) On average over uniformly
		random $x, z \in \F_q^m$,
		\begin{gather}
			\label{eq:qld-q-cons-m-hat-xz}
			(\hat{Q}^{x,z}_{a,b})_{\mathsf{A}\mathsf{A}'}
			\approx_{\delta_Q}
			(\hat{M}^{(\mathrm{Point},X),x}_{a} \,
			\hat{M}^{(\mathrm{Point},Z),z}_{b})_{\mathsf{B}\mathsf{A}''}\,,
			\\
			\label{eq:qld-q-cons-m-hat-zx}
			(\hat{Q}^{x,z}_{a,b})_{\mathsf{A}\mathsf{A}'}
			\approx_{\delta_Q}
			(\hat{M}^{(\mathrm{Point},Z),z}_{b} \,
			\hat{M}^{(\mathrm{Point},X),x}_{a})_{\mathsf{B}\mathsf{A}''}\,.
		\end{gather}
	\end{enumerate}
	Furthermore, all symmetric equivalents of these approximations also hold.
	The operators $\hat{M}^{(\mathrm{Point},X),x}_{a}$ and
	$\hat{M}^{(\mathrm{Point},Z),z}_{b}$ do not commute in general, so the
	products in \eqref{eq:qld-q-cons-m-hat-xz} and
	\eqref{eq:qld-q-cons-m-hat-zx} differ and the two relations are separate
	conclusions.
\end{lemma}
\begin{proof}
	\leanok
	\uses{lem:qld-comm-cons, lem:expanded-point-field-commutation,
		lem:ortho, fact:triangle, lem:combined-point-distance-support,
		lem:qld-epr-cross-swap, lem:qld-sandwich-povm,
		lem:qld-sandwich-consistency-defect,
		lem:qld-sandwich-orthonormalization,
		def:expanded-point-measurement, def:symmetric-equivalents,
		def:povm-distance}
	The formal proof works with field-valued outcomes directly. On each
	expanded local space form the complete sandwich POVM
	$R^{x,z}_{a,b}=M^{Z,z}_b M^{X,x}_a M^{Z,z}_b$.
	Field-valued approximate commutation bounds its distance from the ordered
	point product. Point self-consistency and the same commutation estimate
	bound its off-diagonal overlap across opposite placements.
	Apply orthonormalization to obtain a projective measurement for each
	point pair, on the same expanded local space, close to this POVM.
	The squared-distance triangle inequality gives self-consistency of the
	projective families and their distances from both ordered products.
	The placement comparisons give all symmetric equivalents.
	The estimates are absorbed into $\delta_Q(\eps)=K\eps^{1/8}$ for one
	universal constant $K$.

	This proves the stated measurement construction without using the source's
	intermediate binary refinements and quantum linearity construction.
	The zero-state padding needed by that source proof remains the separate
	import issue in Remark~\ref{rem:linearity-import} and
	\cite{gap:qpbt_linearity-theorem-quotation}.
\end{proof}

\begin{lemma}[Formalization support for the $X$ marginal of the combined point measurement]
	\label{lem:qld-combined-point-x-marginal-distance}
	\lean{MIPStarRE.QPBT.CombinedPointsWitness.marginal_X_distance_le}
	\leanok
	\uses{lem:qld-4-10, def:povm-distance, def:expanded-point-measurement,
		def:expanded-state, def:symmetric-equivalents}
	This is not a named statement of \cite{Ji2020MIPStar}; it isolates
	Equation~\eqref{eq:qld-qxz-close-to-point}. Fix an admissible parameter
	tuple, a projective setting, a real number $\delta_Q$, and any combined point
	measurements satisfying the conclusions of Lemma~\ref{lem:qld-4-10} with
	error $\delta_Q$. For every directed opposite pair of placements $p_1,p_2$,
	on average over independent uniform $x,z\in\F_q^m$,
	\[
		\E_{x,z}\sum_a \left\|
		\left(\sum_b (\hat Q^{x,z}_{a,b})_{p_1}
			- (\hat M^{(\mathrm{Point},X),x}_a)_{p_2}\right)
		\ket{\hat\psi}\right\|^2 \leq 4\delta_Q.
		\label{eq:qld-qxz-close-to-point}
	\]
\end{lemma}
\begin{proof}
	\leanok
	\uses{lem:qld-4-10}
	For fixed $x,z,a$, contract the projective refinement over the $Z$ outcome.
	The ordered $Z$-then-$X$ consistency relation bounds the resulting sum of
	squared distances by $\delta_Q$; the squared triangle estimate in the
	contraction contributes the factor four. Sum over $a$ and average over
	$x,z$.
\end{proof}

\begin{lemma}[Formalization support for the $Z$ marginal of the combined point measurement]
	\label{lem:qld-combined-point-z-marginal-distance}
	\lean{MIPStarRE.QPBT.CombinedPointsWitness.marginal_Z_distance_le}
	\leanok
	\uses{lem:qld-4-10, def:povm-distance, def:expanded-point-measurement,
		def:expanded-state, def:symmetric-equivalents}
	This is not a named statement of \cite{Ji2020MIPStar}; it isolates the
	symmetric Equation~\eqref{eq:qld-qxz-close-to-point-2}. Fix an admissible
	parameter tuple, a projective setting, a real number $\delta_Q$, and any
	combined point measurements satisfying the conclusions of
	Lemma~\ref{lem:qld-4-10} with error $\delta_Q$. For every directed opposite
	pair of placements $p_1,p_2$,
	\[
		\E_{x,z}\sum_b \left\|
		\left(\sum_a (\hat Q^{x,z}_{a,b})_{p_1}
			- (\hat M^{(\mathrm{Point},Z),z}_b)_{p_2}\right)
		\ket{\hat\psi}\right\|^2 \leq 4\delta_Q,
		\label{eq:qld-qxz-close-to-point-2}
	\]
	where $x,z\in\F_q^m$ are independent and uniform.
\end{lemma}
\begin{proof}
	\leanok
	\uses{lem:qld-4-10}
	Contract the projective refinement over the $X$ outcome and apply the ordered
	$X$-then-$Z$ consistency relation. The same squared triangle estimate gives
	the factor four, after summing over $b$ and averaging over $x,z$.
\end{proof}

\begin{lemma}[Formalization support: unrestricted error in the point witness domain]
	\label{lem:combined-points-unrestricted-error}
	\lean{MIPStarRE.QPBT.CombinedPointsWitness.exists_error_of_projective}
	\leanok
	\uses{lem:qld-4-10, def:povm-distance, def:expanded-state}
	This formalization-only auxiliary is not a named statement of the source
	article, and it does not certify Lemma~\ref{lem:qld-4-10}. Let $S$ be a
	projective setting and let $\{Q^{x,z}\}$ be any family of projective
	measurements with outcomes in $\F_q \times \F_q$, one for each player side
	and each pair $(x,z)$ of supplied points. Then there is a real number
	$\delta$ such that $\{Q^{x,z}\}$ satisfies, at error $\delta$, all three
	placement comparisons required of a combined-points witness. The witness
	domain by itself therefore carries no quantitative content: $\delta$ is
	obtained as an upper bound of a finite range of real numbers, and no
	dependence of $\delta$ on the strategy error is asserted. This is the
	obstruction to the unrestricted supplied-point domain recorded in
	\cite{gap:qpbt_combined-lines-error-term}, for the supplied-point statement
	at \texttt{14\_analysis\_of\_the\_pauli\_basis\_test.tex:993-1034}.
\end{lemma}
\begin{proof}
	\leanok
	\uses{def:povm-distance}
	There are finitely many ordered pairs of placements, and for each such pair
	the three comparison quantities are real numbers. The pointwise maximum of
	the three therefore has a finite range, which is bounded above; any upper
	bound serves as $\delta$, and each of the three required inequalities
	follows from the corresponding bound on the maximum.
\end{proof}

\begin{lemma}[Formalization support: linear answers rarely agree with a quadratic]
	\label{lem:linear-quadratic-agreement}
	\lean{MIPStarRE.QPBT.linear_quadratic_agreement_le}
	\leanok
	\uses{lem:coefficient-polynomial-agreement, def:ld-combining-parameters}
	This formalization-only auxiliary is not a named statement of the source
	article. Fix directly indexed parameters with field size $q$, let
	$f = (f_0, f_1, f_2)$ be a coefficient answer of degree at most two whose
	quadratic coefficient vanishes, so that $f$ evaluates to $f_0 + f_1 t$, and
	let $\alpha \neq 0$ be a scalar. Then
	\[
		\Pr_{t}\bigl[\,f_0 + f_1 t = \alpha t^2\,\bigr] \leq \frac{2}{q}\;,
	\]
	the probability being over a uniformly random field element $t$. No
	relation between the degree and $q$ is assumed, so the bound is not
	truncated at one.
\end{lemma}
\begin{proof}
	\leanok
	\uses{lem:coefficient-polynomial-agreement}
	The coefficient answer of $\alpha T^2$ has quadratic coefficient
	$\alpha \neq 0$, so it differs from $f$. The averaged collision estimate for
	two distinct coefficient answers of degree at most two bounds the agreement
	probability by $2/q$; that estimate is the uniform-parameter form of the
	root count of Lemma~\ref{lem:coefficient-polynomial-agreement}.
\end{proof}

\section{Combined line measurements}

\begin{theorem}[Resampling a line parameter]
	\label{thm:line-parameter-resampling}
	\lean{MIPStarRE.QPBT.avgOver_uniform_lineRepMap_resample_parameter}
	\leanok
	\uses{def:line-representative}
	This auxiliary identity isolates the conditional sampling used in the
	proof of Lemma~\ref{lem:qld-xz-lines}; it is not a named statement of
	\cite{Ji2020MIPStar}. Let $K$ be a finite field, $m \geq 0$, and
	$v \in K^m$, and write $R_v = L^{\mathrm{Ln}}_v$ for the canonical line
	representative map of Definition~\ref{def:line-representative}.
	For every $f : K^m \times K^m \to \mathbb{R}$,
	\[
		\mathbb{E}_{u \in K^m} f(R_v(u),u)
		= \mathbb{E}_{u \in K^m}\mathbb{E}_{t \in K}
		f(R_v(u),R_v(u)+tv),
	\]
	where all averages are uniform. The identity includes $v=0$.
\end{theorem}
\begin{proof}
	\leanok
	\uses{def:line-representative, lem:line-representative-idempotent}
	The point $R_v(u)+tv$ is uniform on $K^m$: each point has exactly
	$|K|$ preimages under $(u,t) \mapsto R_v(u)+tv$. Moreover,
	$R_v(R_v(u)+tv)=R_v(u)$. Applying these two identities to $f$ gives
	the result. When $v=0$, $R_v$ is the identity and $t$ is free.
\end{proof}

\begin{corollary}[Resampling a point from the line-point distribution]
	\label{cor:line-point-parameter-resampling}
	\lean{MIPStarRE.QPBT.avgOver_aLinePointDist_resample_parameter}
	\lean{MIPStarRE.QPBT.avgOver_dLinePointDist_resample_parameter}
	\lean{MIPStarRE.QPBT.avgOver_linePointDist_resample_parameter}
	\leanok
	\uses{def:ld-game, def:line-point-dist}
	This auxiliary identity is not a named statement of
	\cite{Ji2020MIPStar}; it isolates the line-point resampling used in the
	proof of Lemma~\ref{lem:qld-xz-lines}. For low-degree parameters
	$(q,m,d,k)$ as in Definition~\ref{def:ld-game}, let $\widetilde D$ be any
	of the seed-bearing laws $\widetilde D_{\mathrm{ALine}}$,
	$\widetilde D_{\mathrm{DLine}}$, or their equal mixture
	$\widetilde D_{\mathrm{Line}}$ from
	Definition~\ref{def:line-point-dist}. Write $b_\ell$ and $v_\ell$ for the
	canonical base and direction of a seed-bearing line. For every real-valued
	function $f$ on seed-bearing line-point pairs,
	\[
		\E_{(\ell,u)\sim \widetilde D} f(\ell,u)
		= \E_{(\ell,u)\sim \widetilde D}\E_{t\in\F_q}
		f(\ell,b_\ell+t v_\ell),
	\]
	where $t$ is uniform and independent. This includes diagonal lines of
	direction zero, for which the resampled point is always $b_\ell$.
	For each choice of $\widetilde D$, let $D$ be the corresponding source law
	and $\rho$ the corresponding map among
	$\rho_{\mathrm{ALine}}$, $\rho_{\mathrm{DLine}}$, and
	$\rho_{\mathrm{Line}}$ from Definition~\ref{def:line-point-dist}. Since
	$\rho$ preserves the canonical base and direction, substituting
	$f=g\circ\rho$ gives
	\[
		\E_{(\bar\ell,u)\sim D} g(\bar\ell,u)
		=\E_{(\bar\ell,u)\sim D}\E_{t\in\F_q}
		g(\bar\ell,b_{\bar\ell}+t v_{\bar\ell}),
	\]
	for every real-valued function $g$ on source line-point pairs. Thus the
	linked declarations prove the refined identity, including functions that
	depend on the retained seed, and the source identity is its push-forward.
\end{corollary}
\begin{proof}
	\leanok
	\uses{thm:line-parameter-resampling, def:line-point-dist,
		lem:uniform-map-equivalence, lem:line-representative-idempotent}
	In each component, condition on the scalar seed and direction data.
	The remaining point is uniform in $\F_q^m$, so
	Theorem~\ref{thm:line-parameter-resampling} applies with that fixed
	direction. Averaging over these data proves the axis and diagonal
	identities. Taking their equal mixture proves the mixed identity. Composing
	the test function with the forgetful map gives the source push-forward
	identity.
\end{proof}

\begin{lemma}[Formalization support for coefficient polynomial agreement]
	\label{lem:coefficient-polynomial-agreement}
	\lean{MIPStarRE.QPBT.evalCoefficient_collision_card_le}
	\lean{MIPStarRE.QPBT.linePolynomialOfCoefficients_eq_ofFn,
		MIPStarRE.QPBT.linePolynomialOfCoefficients_eval,
		MIPStarRE.QPBT.linePolynomialOfCoefficients_injective,
		MIPStarRE.QPBT.linePolynomialOfCoefficients_natDegree_le}
	\leanok
	Let $K$ be a finite integral domain, let $c\in\N$, and let
	$f,g\in K^{c+1}$ be distinct coefficient vectors. Their polynomials
	$p_f(T)=\sum_{i=0}^{c}f_iT^i$ and $p_g(T)=\sum_{i=0}^{c}g_iT^i$
	agree at most $c$ elements of $K$.
	The coefficient representation is injective, preserves evaluation,
	and has natural degree at most $c$.
	This is a formalization-only auxiliary for the Schwartz--Zippel step
	in the proof of Lemma~\ref{lem:qld-xz-lines}
	(\texttt{14\_analysis\_of\_the\_pauli\_basis\_test.tex:955}).
	There is no restriction $c<|K|$.
\end{lemma}
\begin{proof}
	\leanok
	The coefficient representation is injective, preserves evaluation,
	and gives polynomials of degree at most $c$. If the polynomials agreed
	at more than $c$ elements, the univariate root bound applied to their
	difference would force them to be equal, contradicting $f\ne g$.
\end{proof}

\begin{lemma}[Evaluation at the parameter of a nondegenerate affine line]
	\label{lem:affine-line-evaluation}
	\lean{MIPStarRE.QPBT.evalOpt_affine_parameter_of_direction_ne_zero}
	\leanok
	\uses{def:ld-line-description, def:ideg-deg-polynomials, def:line}
	Let $\ell$ be a line description with base point $u_0$ and direction
	$v \neq 0$, let $c \in \N$, let $f \in \deg_{c}(\ell)$ be given by its
	coefficient list, and let $t \in \F_q$. Then the completed evaluation of
	$f$ at the point $u_0 + t v$ is defined and equals the value of the
	coefficient polynomial of $f$ at $t$.
	This is the geometric step of the collision argument in the proof of
	Lemma~\ref{lem:qld-xz-lines}
	(\texttt{14\_analysis\_of\_the\_pauli\_basis\_test.tex:950--955}): on a
	line with nonzero direction a point determines its affine parameter, so
	completed evaluation is unambiguous there.
\end{lemma}
\begin{proof}
	\leanok
	The point $u_0 + t v$ lies on $\ell$ with parameter $t$. If
	$u_0 + t v = u_0 + t' v$ for some $t'$, then $(t - t')v = 0$, and
	$v \neq 0$ forces $t = t'$. The parameter witnessing membership is
	therefore unique, so completed evaluation is defined at $u_0 + t v$ and
	returns the value of the coefficient polynomial at $t$.
\end{proof}

\begin{lemma}[Formalization support: collision mass on a fixed nondegenerate line]
	\label{lem:nondegenerate-fiber-line-collision}
	\lean{MIPStarRE.QPBT.linePointDist_nondegenerate_fiber_collision_le}
	\leanok
	\uses{def:line-point-dist, def:ideg-deg-polynomials}
	Let $\ell_0$ be a line of $\F_q^m$ with direction $v_{\ell_0} \ne 0$ and
	let $f \ne g$ be polynomials of degree at most $c$. Then
	\[
		\E_{(\ell,x)}\bigl[\bone{\ell = \ell_0}\,\bone{f(x) = g(x)}\bigr]
		\le \frac{c}{q}\,
			\E_{(\ell,x)}\bigl[\bone{\ell = \ell_0}\bigr],
	\]
	the expectation being over the line-point distribution
	(Definition~\ref{def:line-point-dist}). Both sides are unnormalized
	masses; when the fiber mass on the right is positive, dividing by it gives
	the conditional collision probability at most $c/q$ required by
	Lemma~\ref{lem:pasting}. The nonzero-direction hypothesis on $\ell_0$ is
	part of the assertion. This is a formalization-only auxiliary for the
	Schwartz--Zippel step in the proof of Lemma~\ref{lem:qld-xz-lines}
	(\texttt{14\_analysis\_of\_the\_pauli\_basis\_test.tex:950--955}).
\end{lemma}
\begin{proof}
	\leanok
	\uses{lem:affine-line-evaluation, lem:coefficient-polynomial-agreement,
		cor:line-point-parameter-resampling}
	Specialize the nonnegative weighted line-collision estimate to
	$w(\ell)=\bone{\ell=\ell_0}$. Since $v_{\ell_0}\ne0$, the
	nondegenerate-direction indicator is redundant on the support of this
	weight, giving the stated fixed-fiber inequality.
\end{proof}

\begin{definition}[Formalization support: retained nondegenerate mass]
	\label{def:nondegenerate-line-pasting-mass}
	\lean{MIPStarRE.QPBT.nondegenerateLinePastingMass}
	\leanok
	\uses{def:line-point-dist}
	Fix an arbitrary low-degree parameter tuple $L=(q,m,d,k)$ as in
	Definition~\ref{def:ld-game}. Write $D_{\mathsf{Line}}^L$ for the source
	line-point law at $L$ and $\widetilde D_{\mathsf{Line}}^L$ for its
	seed-bearing refinement in Definition~\ref{def:line-point-dist}. In the
	refined law, $\ell$ is a line description retaining its kind, canonical base
	point and seed $s$: an axis description records $(u_0,s)$ and a diagonal
	description records $(u_0,s,v)$, with the base and prefix conditions of that
	definition. Set $\operatorname{dir}(\ell)=e_{\chi(s)}$ in the axis case and
	$\operatorname{dir}(\ell)=v$ in the diagonal case; the latter is the stored,
	already projected direction. Mapping the axis description to
	$(u_0,e_{\chi(s)})$ and leaving the diagonal description unchanged sends
	$\widetilde D_{\mathsf{Line}}^L$ to $D_{\mathsf{Line}}^L$.
	For two independent samples
	$(\ell_X,x),(\ell_Z,z)\sim\widetilde D_{\mathsf{Line}}^L$, define the retained mass
	\[
		\mu_{\mathrm{nd}}(L)
		:= \sum_{\substack{(\ell_X,x),(\ell_Z,z)\\
			\operatorname{dir}(\ell_X)\ne 0}}
		\widetilde D_{\mathsf{Line}}^L(\ell_X,x)
		\widetilde D_{\mathsf{Line}}^L(\ell_Z,z).
	\]
	Thus $\mu_{\mathrm{nd}}(L)$ is the probability that the first sampled line
	has nonzero direction; the second line-point sample is unrestricted. This is
	a formalization-only normalization factor for the conditioning step in the
	proof of Lemma~\ref{lem:qld-xz-lines}. It does not replace the independent,
	unconditioned product law in that lemma's statement.
\end{definition}

\begin{theorem}[Formalization support: retained-mass bounds]
	\label{thm:nondegenerate-line-pasting-mass-bounds}
	\lean{MIPStarRE.QPBT.nondegenerateLinePastingMass_bounds}
	\leanok
	\uses{def:nondegenerate-line-pasting-mass}
	For every low-degree parameter tuple $L$,
	\[
		\frac12\leq \mu_{\mathrm{nd}}(L)\leq 1.
	\]
\end{theorem}
\begin{proof}
	\leanok
	\uses{def:nondegenerate-line-pasting-mass, def:line-point-dist,
		lem:axis-line-zero-direction-mass}
	The axis-parallel component has weight one half and always has nonzero
	direction, while the diagonal component contributes nonnegative mass. This
	gives the lower bound. The upper bound follows because the retained samples
	form an event in the product of two probability distributions.
\end{proof}

\begin{theorem}[Formalization support: positivity of the retained mass]
	\label{thm:nondegenerate-line-pasting-mass-positive}
	\lean{MIPStarRE.QPBT.prod_linePointDist_nondegenerate_mass_pos}
	\leanok
	\uses{def:nondegenerate-line-pasting-mass}
	For every low-degree parameter tuple $L$, the event
	$\operatorname{dir}(\ell_X)\ne 0$ has positive mass under
	$\widetilde D_{\mathsf{Line}}^L\times\widetilde D_{\mathsf{Line}}^L$; equivalently,
	$0<\mu_{\mathrm{nd}}(L)$.
\end{theorem}
\begin{proof}
	\leanok
	\uses{thm:nondegenerate-line-pasting-mass-bounds}
	The lower bound $\mu_{\mathrm{nd}}(L)\geq\frac12$ is strictly positive.
\end{proof}

\begin{definition}[Formalization support: conditioned pasting distribution]
	\label{def:nondegenerate-line-pasting-dist}
	\lean{MIPStarRE.QPBT.nondegenerateLinePastingDist}
	\leanok
	\uses{def:nondegenerate-line-pasting-mass,
		thm:nondegenerate-line-pasting-mass-positive}
	Condition $\widetilde D_{\mathsf{Line}}^L\times\widetilde D_{\mathsf{Line}}^L$ on
	$\operatorname{dir}(\ell_X)\ne 0$, using the positive retained mass from
	Theorem~\ref{thm:nondegenerate-line-pasting-mass-positive}, and push the
	result forward along
	\[
		\big((\ell_X,x),(\ell_Z,z)\big)
		\longmapsto
		\big(((\ell_X,\ell_Z),(\ell_Z,z)),(\ell_X,x)\big).
	\]
	Denote the resulting distribution by $D_{\mathrm{nd}}^L$; every line
	description in its carrier still retains its seed. Thus the paired lines
	are ordered $X,Z$, followed by the full $Z$ sample and then the full
	$X$ sample, and only the first, $X$-line direction is restricted. This
	formalization-only law supports the pasting construction. The source
	conclusion of Lemma~\ref{lem:qld-xz-lines} remains an average over two
	independent, unconditioned samples from $D_{\mathsf{Line}}^L$.
	Lemma~\ref{lem:qld-line-conditioning-restoration} transfers the internal
	estimate back to the unconditioned product of the refined laws, as used in
	Lemma~\ref{lem:qld-paired-line-restored-defect}. The resulting comparison is
	on the completed answer alphabet $(\F_q\cup\{\bot\})^2$: line evaluation
	returns $\bot$ when no value at the sampled point is determined, and the
	joint point measurement assigns zero effect to outcomes containing $\bot$.
	The source comparison instead has answers $(a,b)\in\F_q^2$.
	Relating these answer conventions remains necessary to identify the
	completed-answer estimate with Lemma~\ref{lem:qld-xz-lines}; restoring the
	unconditioned law does not by itself make that identification.
\end{definition}

\begin{theorem}[Formalization support: normalization of the conditioned law]
	\label{thm:nondegenerate-line-pasting-dist-probability}
	\lean{MIPStarRE.QPBT.nondegenerateLinePastingDist_isProbability}
	\leanok
	\uses{def:nondegenerate-line-pasting-dist}
	For every low-degree parameter tuple $L$, $D_{\mathrm{nd}}^L$ is a
	probability distribution.
\end{theorem}
\begin{proof}
	\leanok
	\uses{def:nondegenerate-line-pasting-dist,
		thm:nondegenerate-line-pasting-mass-positive}
	Normalization by the positive retained mass gives a probability
	distribution on the restricted event, and push-forward preserves total mass.
\end{proof}

% Unmarked, tracked in issue #667; no permanent exemption. The concrete
% unconditioned line measurement and its full product-law consistency are proved.
% CombinedLinesWitness compares completed evaluations; the field-valued source
% comparison, including the zero-direction convention, is not certified here.
\begin{lemma}[Combined line measurements]
	\label{lem:qld-xz-lines}
	\lean{MIPStarRE.QPBT.CombinedLinesWitness}
	\lean{MIPStarRE.QPBT.exists_combinedLinesWitness}
	\notready
	\uses{lem:qld-4-10, def:line, def:line-point-dist, def:consistency,
		def:bracket, def:expanded-state, def:symmetric-equivalents,
		def:ideg-deg-polynomials}
	There exists a function $\delta_P(\eps,m,d,q) = \poly(\eps, md/q)$, the
	two-variable polynomial error read in the corrected sum sense of
	\cite{gap:qpbt_pasting-product-error}, that is, a bound
	$C\bigl(\eps^{r} + (md/q)^{s}\bigr)$ with $C \ge 1$ and $r, s > 0$ on the
	closed nonnegative quadrant rather than the literal product form of the
	shorthand, such that the following holds. For
	$\hilb \in \{\hilb_{\mathsf{A}}, \hilb_{\mathsf{B}}\}$ and every pair of
	lines $\ell_X$ and $\ell_Z$ in $\F_q^m$ there exists a POVM
	$\{T^{\ell_X,\ell_Z}_{f_X,f_Z}\}$ acting on $\hilb \ot (\C^q)^{\ot M}$,
	with answers consisting of pairs of polynomials
	$(f_X, f_Z) \in \deg_{md}(\ell_X) \times \deg_{md}(\ell_Z)$, such that on
	average over pairs $(\ell_X, x)$ and $(\ell_Z, z)$ independently sampled
	from the line-point distribution
	(Definition~\ref{def:line-point-dist}),
	\[
		\big(T^{\ell_X,\ell_Z}
		_{[\mathrm{eval}_x(\cdot)=a,\ \mathrm{eval}_z(\cdot)=b]}
		\big)_{\mathsf{A}\mathsf{A}'}
		\simeq_{\delta_P}
		\big(\hat{Q}^{x,z}_{a,b}\big)_{\mathsf{B}\mathsf{A}''}\,,
	\]
	as well as all symmetric equivalents of this approximation. Moreover, if
	$\ell_W$ is axis-parallel, then $f_W \in \deg_{d}(\ell_W)$.
\end{lemma}
\begin{proof}
	\uses{def:expanded-line-measurement, lem:qld-4-10, fact:add-a-proj,
		lem:cool-closeness-fact, lem:qld-comm-line-cons, fact:agreement,
		fact:data-processing, lem:pasting, lem:schwartz-zippel-total-degree,
		def:symmetric-equivalents, def:concrete-paired-line-measurement,
		lem:concrete-paired-line-degree-support}
	Define, for all lines $\ell_X, \ell_Z$ and
	$(f_X, f_Z) \in \deg_{md}(\ell_X) \times \deg_{md}(\ell_Z)$, the sandwich
	operator
	$T^{\ell_X,\ell_Z}_{f_X,f_Z}
	= \hat{M}^{(\mathrm{Line},X),\ell_X}_{f_X}
	\hat{M}^{(\mathrm{Line},Z),\ell_Z}_{f_Z}
	\hat{M}^{(\mathrm{Line},X),\ell_X}_{f_X}$,
	which forms a valid POVM; since the lines measurements come from a valid
	strategy for the low-degree test, on an axis-parallel line $\ell_W$ the
	outcome $f_W$ lies in $\deg_d(\ell_W)$, and this property is inherited by
	$T$. As intermediate consistency relations, write
	$\hat{Q}^{x,z}_{a,\bullet} = \sum_b \hat{Q}^{x,z}_{a,b}$ and
	$\hat{Q}^{x,z}_{\bullet,b} = \sum_a \hat{Q}^{x,z}_{a,b}$ for the two
	marginals of the combined point measurement of
	Lemma~\ref{lem:qld-4-10}. Applying
	Lemma~\ref{fact:add-a-proj} to \eqref{eq:qld-q-cons-m-hat-zx} and then
	Lemma~\ref{lem:cool-closeness-fact} gives, on average over uniform
	$x, z$,
	$(\hat{Q}^{x,z}_{a,\bullet})_{\mathsf{A}\mathsf{A}'}
	\approx_\delta
	\sum_b (\hat{Q}^{x,z}_{a,b})_{\mathsf{A}\mathsf{A}'}
	(\hat{M}^{(\mathrm{Point},Z),z}_{b})_{\mathsf{B}\mathsf{A}''}$,
	and a direct norm computation from this display and
	\eqref{eq:qld-q-cons-m-hat-zx} yields the marginal consistencies
	$(\hat{Q}^{x,z}_{a,\bullet})_{\mathsf{A}\mathsf{A}'} \approx_{\delta_Q}
	(\hat{M}^{(\mathrm{Point},X),x}_{a})_{\mathsf{B}\mathsf{A}''}$ and
	$(\hat{Q}^{x,z}_{\bullet,b})_{\mathsf{A}\mathsf{A}'} \approx_{\delta_Q}
	(\hat{M}^{(\mathrm{Point},Z),z}_{b})_{\mathsf{B}\mathsf{A}''}$.
	From items 1 and 3 of Lemma~\ref{lem:qld-comm-line-cons}, converting
	between $\approx$ and $\simeq$ with Lemma~\ref{fact:agreement} and
	summing outcomes with Lemma~\ref{fact:data-processing}, the lines
	measurements are consistent with the points measurements, and combining
	with the marginal consistencies gives, on average over the line-point
	distribution,
	$(\hat{Q}^{x,z}_{a,\bullet})_{\mathsf{A}\mathsf{A}'} \simeq_{\delta_Q}
	(\hat{M}^{(\mathrm{Line},X),\ell_X}
	_{[\mathrm{eval}_x(\cdot)=a]})_{\mathsf{B}\mathsf{A}''}$
	together with the analogous relation for $Z$. The conclusion now follows
	from the pasting lemma (Lemma~\ref{lem:pasting}) applied with
	$\hat{Q}^{x,z}_{a,b}$ in the role of the combined measurement,
	$\hat{M}^{(\mathrm{Line},X),\ell_X}$ and
	$\hat{M}^{(\mathrm{Line},Z),\ell_Z}$ in the roles of the two factor
	measurements, $T^{\ell_X,\ell_Z}$ in the role of the pasted measurement,
	$\delta := \delta_Q$, and $\eta := md/q$, the last justified by the
	Schwartz--Zippel bound (Lemma~\ref{lem:schwartz-zippel-total-degree})
	applied to the polynomials $f_X, f_Z$ of total degree at most $md$; the
	two hypotheses of the pasting lemma are supplied by the preceding display
	and by \eqref{eq:qld-q-self-cons}, and the distribution over
	$(\ell_X, \ell_Z)$ is that of two independent uniformly random lines. The
	register-exchanged comparison \eqref{eq:pasting-1-sym} is supplied by the
	symmetric equivalents of that preceding display, which
	Lemma~\ref{lem:qld-4-10} establishes for all of its approximations.
	This gives the stated relation with $\delta_P$ equal to the pasting
	error, a $\poly(\eps, md/q)$ function, and the identical argument with
	the registers modified gives all symmetric equivalents.
\end{proof}

	\begin{definition}[Retained mass of the nondegenerate-line restriction]
		\label{def:qld-nondegenerate-pasting-mass}
		\lean{MIPStarRE.QPBT.nondegenerateLinePastingMass}
		\leanok
		\uses{def:line-point-dist}
		This quantity is not named in the source article; it is the normalization
		constant of the first-line conditioning used in the proof of
		Lemma~\ref{lem:qld-xz-lines}. Draw two independent samples from the
		line-point distribution of Definition~\ref{def:line-point-dist}. The
		\emph{retained mass} is the probability that the line of the first sample
		has nonzero direction.
	\end{definition}

	\begin{lemma}[The retained mass lies between one half and one]
		\label{lem:qld-nondegenerate-pasting-mass-bounds}
		\lean{MIPStarRE.QPBT.nondegenerateLinePastingMass_bounds}
		\leanok
		\uses{def:qld-nondegenerate-pasting-mass, def:line-point-dist}
		The retained mass of Definition~\ref{def:qld-nondegenerate-pasting-mass} is
		at least $1/2$ and at most $1$.
	\end{lemma}

	\begin{lemma}[The first-line nondegeneracy event has positive mass]
		\label{lem:qld-nondegenerate-pasting-mass-pos}
		\lean{MIPStarRE.QPBT.prod_linePointDist_nondegenerate_mass_pos}
		\leanok
		\uses{def:qld-nondegenerate-pasting-mass, def:line-point-dist}
		Under two independent samples from the line-point distribution of
		Definition~\ref{def:line-point-dist}, the event that the line of the first
		sample has nonzero direction has positive probability.
	\end{lemma}
	\begin{proof}
		\leanok
		\uses{lem:qld-nondegenerate-pasting-mass-bounds}
		That probability is the retained mass of
		Definition~\ref{def:qld-nondegenerate-pasting-mass}, which is at least
		$1/2$ by Lemma~\ref{lem:qld-nondegenerate-pasting-mass-bounds}.
	\end{proof}

	\begin{definition}[Nondegenerate-line pasting distribution]
		\label{def:qld-nondegenerate-pasting-dist}
		\lean{MIPStarRE.QPBT.nondegenerateLinePastingDist}
		\leanok
		\uses{def:line-point-dist, lem:qld-nondegenerate-pasting-mass-pos}
		This law is not named in the source article; it is the proof-only question
		law of the first-line conditioning in the proof of
		Lemma~\ref{lem:qld-xz-lines}. Condition two independent samples from the
		line-point distribution of Definition~\ref{def:line-point-dist} on the
		event that the line of the first sample has nonzero direction, which is
		legitimate by Lemma~\ref{lem:qld-nondegenerate-pasting-mass-pos}, and
		present the outcome as the pair of the two lines, followed by the second
		line-point sample and then the first.
	\end{definition}

	\begin{lemma}[The nondegenerate-line pasting law is normalized]
		\label{lem:qld-nondegenerate-pasting-dist-prob}
		\lean{MIPStarRE.QPBT.nondegenerateLinePastingDist_isProbability}
		\leanok
		\uses{def:qld-nondegenerate-pasting-dist}
		The law of Definition~\ref{def:qld-nondegenerate-pasting-dist} has total
		mass one.
	\end{lemma}
	\begin{proof}
		\leanok
		\uses{lem:qld-nondegenerate-pasting-mass-pos}
		Conditioning on an event of positive mass yields a normalized law by
		Lemma~\ref{lem:qld-nondegenerate-pasting-mass-pos}, and the relabeling of
		the sample preserves the total mass.
	\end{proof}

	\begin{lemma}[Formalization support for conditioning a nonnegative finite average]
		\label{lem:qld-restricted-average-le-div-mass}
		\lean{MIPStarRE.QPBT.avgOver_restrict_le_div_mass}
		\leanok
		Let $\mu$ be a finite-support nonnegative weighting on a finite set $\Omega$,
		let $G\subseteq\Omega$, and put
		$\rho=\sum_{\omega\in\operatorname{supp}(\mu)\cap G}\mu(\omega)$.
		If $\rho>0$ and $f\colon\Omega\to\R$ is nonnegative, then the normalized
		restriction of $\mu$ to $G$ satisfies
		\[
			\operatorname{Avg}_{\mu|G}(f)
			\leq \frac{\operatorname{Avg}_{\mu}(f)}{\rho}.
		\]
		No assumption that $\mu$ has total mass one is required.
	\end{lemma}
	\begin{proof}
		\leanok
		Multiplying the left-hand side by the positive number $\rho$ gives the
		weighted sum over $G$. This is at most the full weighted sum because every
		weight and every value of $f$ is nonnegative.
	\end{proof}

	\begin{lemma}[Formalization support for positivity of the consistency integrand]
		\label{lem:qld-consistency-integrand-nonnegative}
		\lean{MIPStarRE.QPBT.consistencyDefect_integrand_nonneg}
		\leanok
		\uses{def:consistency, def:expanded-state,
			def:distance-state-quadratic-form, def:symmetric-equivalents}
		Fix a projective setting, a directed opposite pair of placements $p_1,p_2$,
		and two complete measurements $M^1,M^2$ with the same finite outcome set,
		acting on the corresponding expanded local spaces. Then
		\[
			0\leq \sum_{a\ne b}
			q_{\hat\psi}\!\left((M^1_a)_{p_1}(M^2_b)_{p_2}\right).
		\]
	\end{lemma}
	\begin{proof}
		\leanok
		\uses{lem:distance-qform-nonnegative}
		For each fixed $a$, sum the positive effects $M^2_b$ over $b\ne a$.
		Opposite placements commute, so its product with the positive effect
		$M^1_a$ is positive. Sum over $a$ and apply positivity of the state
		quadratic form.
	\end{proof}

	\begin{lemma}[Conditioning inflates the consistency defect by the inverse retained mass]
		\label{lem:qld-conditioned-pasting-defect}
		\lean{MIPStarRE.QPBT.consistencyDefect_nondegenerateLinePastingDist_le}
		\leanok
		\uses{def:qld-nondegenerate-pasting-dist, def:qld-nondegenerate-pasting-mass,
			def:consistency, def:line-point-dist}
		This estimate is not a named statement of the source article; it quantifies
		the cost of the first-line conditioning used in the proof of
		Lemma~\ref{lem:qld-xz-lines}. Fix admissible parameters, an error
		parameter, and two \emph{opposite} placements of the two players. Let two
		families of measurements on the corresponding expanded local spaces be
		indexed by a pair of line-point samples. Then the consistency defect of
		Definition~\ref{def:consistency} measured under the conditioned law of
		Definition~\ref{def:qld-nondegenerate-pasting-dist} is at most the
		consistency defect measured under two independent line-point samples,
		divided by the retained mass of
		Definition~\ref{def:qld-nondegenerate-pasting-mass}.
	\end{lemma}
	\begin{proof}
		\leanok
		\uses{lem:qld-restricted-average-le-div-mass,
			lem:qld-consistency-integrand-nonnegative}
		The relabeling built into
		Definition~\ref{def:qld-nondegenerate-pasting-dist} only restores the
		original order of the two samples, so both sides average the same
		integrand. It is nonnegative at every sample by
		Lemma~\ref{lem:qld-consistency-integrand-nonnegative}. Apply
		Lemma~\ref{lem:qld-restricted-average-le-div-mass} to the first-line
		nondegeneracy event and its positive retained mass.
	\end{proof}

\begin{definition}[The concrete paired-line POVM]
	\label{def:concrete-paired-line-measurement}
	\lean{MIPStarRE.QPBT.ProjectiveSetting.combinedLineMeasurement}
	\lean{MIPStarRE.QPBT.ProjectiveSetting.combinedLineMeasurement_effect}
	\leanok
	\uses{def:expanded-line-measurement, def:povm-conventions,
		def:ideg-deg-polynomials}
	For any projective setting, either player, and any two source lines,
	define the POVM on the expanded local space, with each outcome a pair
	of coefficient representatives as in
	Definition~\ref{def:ideg-deg-polynomials}, by
	\[
		T^{\ell_X,\ell_Z}_{f_X,f_Z}
		=\hat M^{(\mathrm{Line},X),\ell_X}_{f_X}
		 \hat M^{(\mathrm{Line},Z),\ell_Z}_{f_Z}
		 \hat M^{(\mathrm{Line},X),\ell_X}_{f_X},
		\qquad (f_X,f_Z)\in\deg_{md}(\ell_X)\times\deg_{md}(\ell_Z).
	\]
	Each effect is positive. Summing first over $f_Z$ and then over $f_X$
	gives the identity by completeness and projectivity of the line
	measurements. This is the construction in the proof of
	Lemma~\ref{lem:qld-xz-lines} in \cite{Ji2020MIPStar}.
\end{definition}

\begin{lemma}[Orientation and axis-degree support of the concrete paired-line POVM]
	\label{lem:concrete-paired-line-degree-support}
	\lean{MIPStarRE.QPBT.ProjectiveSetting.combinedLineMeasurement_effect_eq_pastedMeasurement}
	\lean{MIPStarRE.QPBT.ProjectiveSetting.combinedLineMeasurement_axis_degree_X}
	\lean{MIPStarRE.QPBT.ProjectiveSetting.combinedLineMeasurement_axis_degree_Z}
	\leanok
	\uses{def:concrete-paired-line-measurement, def:expanded-line-measurement}
	For either player and any pair of line outcomes, the one-sided pasted
	effect $J$ of \eqref{eq:pasting-2a} with
	$G_1=\hat M^{(\mathrm{Line},Z),\ell_Z}$ and
	$G_2=\hat M^{(\mathrm{Line},X),\ell_X}$ has the exchanged outcome order
	$T^{\ell_X,\ell_Z}_{f_X,f_Z}=J_{f_Z,f_X}$.
	For $W\in\{X,Z\}$, if $\ell_W$ is axis-parallel
	and $f_W$ does not have degree at most $d$, then
	$T^{\ell_X,\ell_Z}_{f_X,f_Z}=0$.
\end{lemma}
\begin{proof}
	\leanok
	\uses{def:concrete-paired-line-measurement, def:expanded-line-measurement}
	The pasting identity follows by expanding the two sandwich products with
	$G_1=Z$ and $G_2=X$; it asserts no point-consistency estimate.
	The corresponding expanded line effect vanishes outside degree $d$.
	It is a factor of the defining sandwich, so the sandwich vanishes.
\end{proof}

\begin{lemma}[Formalization support: point marginals and completed answers]
	\label{lem:qld-line-point-marginal-identities}
	\lean{MIPStarRE.QPBT.avgOver_linePointDist_point}
	\lean{MIPStarRE.QPBT.avgOver_prod_linePointDist_points}
	\lean{MIPStarRE.QPBT.CombinedPointsWitness.postprocess_fst_effect}
	\lean{MIPStarRE.QPBT.CombinedPointsWitness.postprocess_snd_effect}
	\lean{MIPStarRE.QPBT.ProjectiveSetting.opFamilyDistSq_postprocess_some}
	\lean{MIPStarRE.QPBT.CombinedPointsWitness.completed_fst_effect}
	\lean{MIPStarRE.QPBT.CombinedPointsWitness.completed_snd_effect}
	\leanok
	\uses{def:line-point-dist, def:povm-distance, lem:qld-4-10}
	The point marginal of the line-point distribution is uniform, and two
	independent line-point samples have independent uniform point marginals.
	For a joint point measurement, postprocessing by either coordinate gives
	the corresponding marginal sum. Adjoining an undefined answer with zero
	effect preserves the state-dependent squared distance between two
	measurements. Moreover, taking either marginal commutes with adjoining this
	answer, both before and after completing a joint point measurement.
\end{lemma}
\begin{proof}
	\leanok
	\uses{def:line-point-dist, def:povm-distance, lem:qld-4-10}
	The marginal identities follow by expanding the line-point mixture and the
	postprocessing sums. The additional answer contributes the zero operator,
	so it contributes neither to a marginal nor to the squared distance.
\end{proof}

\begin{lemma}[Formalization support: point self-consistency defect]
	\label{lem:qld-point-self-consistency-defect}
	\lean{MIPStarRE.QPBT.CombinedPointsWitness.self_consistency_defect_le}
	\lean{MIPStarRE.QPBT.CombinedPointsWitness.self_consistency_linePoint_defect_le}
	\leanok
	\uses{lem:qld-4-10, def:consistency, def:povm-distance,
		def:line-point-dist, lem:qld-line-point-marginal-identities}
	Let $\{\hat Q^{x,z}_{a,b}\}$ be the joint point measurements of
	Lemma~\ref{lem:qld-4-10}, placed on two opposite placements. Their
	self-consistency estimate, stated there as a state-dependent squared
	distance, also holds in the consistency-defect convention of
	Definition~\ref{def:consistency}, with the same bound $\delta_Q$, on
	average over uniformly random $x,z$. The same bound holds when the pair
	$(x,z)$ is instead read off as the point coordinates of two independent
	samples from the line-point distribution.
\end{lemma}
\begin{proof}
	\leanok
	\uses{lem:qld-4-10, def:consistency, def:povm-distance,
		def:line-point-dist, lem:qld-line-point-marginal-identities}
	The two placed families are projective, so the consistency defect is bounded
	by the state-dependent squared distance, which is the hypothesis. The point
	marginal of the line-point distribution is uniform and the two samples are
	independent, so the average defining the defect is unchanged by reading
	$(x,z)$ off the two line-point samples.
\end{proof}

\begin{lemma}[Formalization support: completing the point answers]
	\label{lem:qld-completed-point-answers}
	\lean{MIPStarRE.QPBT.ProjectiveSetting.consistencyDefect_postprocess_some_pair}
	\leanok
	\uses{def:consistency, def:povm-conventions, def:expanded-state}
	Let $\iota \colon \F_q \to \F_q \cup \{\bot\}$ be the inclusion and write
	$\iota \times \iota$ for the coordinatewise inclusion of
	$\F_q \times \F_q$ into $(\F_q \cup \{\bot\})^2$. For two families
	$\{A^s_{a,b}\}$, $\{B^s_{a,b}\}$ of measurements with outcomes in
	$\F_q \times \F_q$, indexed by a sample $s$ and placed on two directed
	registers, the postprocessings of $A$ and $B$ along $\iota \times \iota$
	have the same consistency defect as $A$ and $B$ themselves.
	This node and the two that follow it are auxiliary leaf support recorded
	for tracking issue 485. No node of this blueprint uses them: the pasting
	route to Lemma~\ref{lem:combined-line-measurement-consistency} runs
	through Lemma~\ref{lem:qld-point-line-marginal-bounds} and
	Lemma~\ref{lem:qld-line-conditioning-restoration} instead, and
	no consumer is claimed for them here.
\end{lemma}
\begin{proof}
	\leanok
	\uses{def:consistency, def:povm-conventions}
	The map $\iota \times \iota$ is injective, so each of its fibers over a
	pair in its image is a single outcome and each fiber over a pair with an
	absent coordinate is empty; the corresponding effects are therefore either
	the original effects or zero, and the double outcome sum defining the
	defect is unchanged term by term.
\end{proof}

\begin{lemma}[Formalization support: completed point self-consistency]
	\label{lem:qld-completed-point-self-consistency}
	\lean{MIPStarRE.QPBT.CombinedPointsWitness.self_consistency_completed_linePoint_defect_le}
	\lean{MIPStarRE.QPBT.CombinedPointsWitness.self_consistency_completed_orderedZX_linePoint_defect_le}
	\leanok
	\uses{lem:qld-4-10, lem:qld-completed-point-answers, def:line-point-dist,
		def:consistency, def:expanded-state,
		lem:qld-point-self-consistency-defect}
	Let $(\ell_X, x)$ and $(\ell_Z, z)$ be drawn independently from the
	line-point distribution of Definition~\ref{def:line-point-dist}, and let
	$\hat{Q}^{x,z}$ be the joint point measurement of
	Lemma~\ref{lem:qld-4-10} with both answers completed as in
	Lemma~\ref{lem:qld-completed-point-answers}. On every pair of directed
	opposite placements, the consistency defect of
	$\big(\hat{Q}^{x,z}_{\iota a,\,\iota b}\big)$ against its copy on the
	opposite placement, averaged over the product law of $(\ell_X,x)$ and
	$(\ell_Z,z)$, is at most $\delta_Q$, and the same bound holds for the
	exchanged answer pair $(\iota b, \iota a)$, ordered $Z$ then $X$.
	This is a formalization-only consequence
	of~\eqref{eq:qld-q-self-cons}, not that equation: its question law is the
	product line-point law rather than the uniform law on pairs of points, and
	its outcome alphabet is the completed one. The exchanged ordering is the
	one that would be required if the X-Z-X sandwich were read as a pasting
	instance with inner family the $Z$-line family and outer family the
	$X$-line family. It completes nothing that the source leaves open.
\end{lemma}
\begin{proof}
	\leanok
	\uses{lem:qld-4-10, lem:qld-completed-point-answers, def:line-point-dist,
		lem:qld-point-self-consistency-defect}
	Completion leaves the defect unchanged
	(Lemma~\ref{lem:qld-completed-point-answers}), so the claim reduces to the
	line-point form already recorded in
	Lemma~\ref{lem:qld-point-self-consistency-defect}, which
	bounds~\eqref{eq:qld-q-self-cons} transported to the product law by
	$\delta_Q$. Exchanging the two completed coordinates is a bijective
	relabeling of the outcome set, under which the defect is invariant.
\end{proof}

\begin{lemma}[Formalization support: conditioned completed point
	self-consistency]
	\label{lem:qld-conditioned-completed-point-self-consistency}
	\lean{MIPStarRE.QPBT.CombinedPointsWitness.self_consistency_conditioned_completed_defect_le}
	\lean{MIPStarRE.QPBT.CombinedPointsWitness.self_consistency_conditioned_completed_orderedZX_defect_le}
	\leanok
	\uses{lem:qld-completed-point-self-consistency, def:line-point-dist,
		def:consistency, def:expanded-state}
	Let $\mu$ be the law of two independent line-point samples
	$(\ell_X, x)$, $(\ell_Z, z)$ conditioned on the event that the direction
	of $\ell_X$ is nonzero, and let $\rho > 0$ be the probability of that
	event under the unconditioned product law. With the completed joint point
	measurement of Lemma~\ref{lem:qld-completed-point-self-consistency}, the
	consistency defect on every pair of directed opposite placements, averaged
	over $\mu$, is at most $\delta_Q / \rho$, in both orderings of the
	completed answer pair.
	This is a formalization-only restriction of
	Lemma~\ref{lem:qld-completed-point-self-consistency} to the question law
	on which the pasting argument is run; it is neither~\eqref{eq:qld-q-self-cons}
	nor the hypothesis~\eqref{eq:pasting-2} as printed, both of which average
	over the unconditioned law with error $\delta_Q$.
\end{lemma}
\begin{proof}
	\leanok
	\uses{lem:qld-completed-point-self-consistency}
	The summand defining the defect at a fixed question is nonnegative, so
	discarding the questions outside an event of probability $\rho$ and
	renormalizing increases the average by a factor of at most $1/\rho$.
\end{proof}

\begin{lemma}[Formalization support: point-to-line marginal estimates]
	\label{lem:qld-point-line-marginal-bounds}
	\lean{MIPStarRE.QPBT.CombinedPointsWitness.marginal_X_linePoint_distance_le}
	\lean{MIPStarRE.QPBT.CombinedPointsWitness.marginal_Z_linePoint_distance_le}
	\lean{MIPStarRE.QPBT.CombinedPointsWitness.marginal_X_option_linePoint_distance_le}
	\lean{MIPStarRE.QPBT.CombinedPointsWitness.marginal_Z_option_linePoint_distance_le}
	\lean{MIPStarRE.QPBT.exists_expLine_point_same_placement_distance_le}
	\lean{MIPStarRE.QPBT.exists_combinedPoints_line_marginal_distance_le}
	\lean{MIPStarRE.QPBT.exists_combinedPoints_line_marginal_defect_le}
	\leanok
	\uses{lem:qld-line-point-marginal-identities, lem:qld-4-10,
		lem:qld-comm-line-cons, fact:triangle, def:consistency}
	Let the joint point witness have error $\delta_Q$. On every directed
	opposite placement, each of its $X$ and $Z$ marginals has squared distance
	at most $4\delta_Q$ from the corresponding point measurement, both on the
	original and on the completed answer alphabet. There is a universal
	$C_0\geq1$ such that an evaluated line measurement and the corresponding
	completed point measurement on the same placement have squared distance at
	most $C_0(\eps+\sqrt\eps)$.

	Consequently, for universal constants $C_1,C_2\geq1$, which need not be
	equal, both completed joint-point marginals have squared distance at most
	\[
		8\delta_Q+C_1(\eps+\sqrt\eps)
	\]
	from the corresponding evaluated line measurements, and the analogous two
	consistency defects are at most
	$8\delta_Q+C_2(\eps+\sqrt\eps)$.
\end{lemma}
\begin{proof}
	\leanok
	\uses{lem:qld-line-point-marginal-identities, lem:qld-4-10,
		lem:qld-comm-line-cons, fact:triangle, def:consistency}
	Rewrite the line-point averages using their uniform point marginals. Apply
	the point-witness estimates, then insert the evaluated line measurement and
	use the squared-distance triangle inequality. The consistency estimates
	follow from the corresponding distance bounds.
\end{proof}

\begin{lemma}[Formalization support: discarded zero directions]
	\label{lem:qld-zero-direction-mass}
	\lean{MIPStarRE.QPBT.dLinePointDist_zero_direction_mass_le}
	\lean{MIPStarRE.QPBT.linePointDist_zero_direction_mass_le}
	\lean{MIPStarRE.QPBT.prod_linePointDist_zero_X_direction_mass_le}
	\leanok
	\uses{def:line-point-dist}
	Under the diagonal line-point law, the probability of a zero direction is
	at most $1/q$. Under the equal mixture of axis-parallel and diagonal laws it
	is at most $1/(2q)$. Hence, for two independent line-point samples, removing
	the event that the $X$-line has zero direction discards mass at most
	$1/(2q)$.
\end{lemma}
\begin{proof}
	\leanok
	\uses{def:line-point-dist}
	A zero diagonal direction forces the final unprojected direction coordinate
	to vanish, and that coordinate is uniform. Axis-parallel directions are
	nonzero. Averaging the two cases and then the unused independent sample gives
	the stated bounds.
\end{proof}

\begin{lemma}[Formalization support: retained nonzero-direction mass]
	\label{lem:qld-nondegenerate-line-mass}
	\lean{MIPStarRE.QPBT.linePointDist_nondegenerate_mass_ge}
	\leanok
	\uses{def:line-point-dist, lem:qld-zero-direction-mass}
	For every admissible low-degree parameter tuple, the unchanged line-point
	law assigns probability at least $3/4$ to lines with nonzero direction.
	This is an auxiliary bound for the conditioning in the proof of
	Lemma~\ref{lem:qld-xz-lines}, not an additional source assertion.
\end{lemma}
\begin{proof}
	\leanok
	\uses{lem:qld-zero-direction-mass}
	The retained probability is one minus the zero-direction probability.
	Lemma~\ref{lem:qld-zero-direction-mass} bounds the latter by $1/(2q)$,
	which is at most $1/4$ because a field has at least two elements.
\end{proof}

\begin{lemma}[Formalization support: conditioning and restoring line defects]
	\label{lem:qld-line-conditioning-restoration}
	\lean{MIPStarRE.QPBT.consistencyDefect_nondegenerateLinePastingDist_le}
	\lean{MIPStarRE.QPBT.exists_combinedPoints_conditioned_line_marginal_defect_le}
	\lean{MIPStarRE.QPBT.consistencyDefect_le_nondegenerateLinePastingDist_add_mass}
	\lean{MIPStarRE.QPBT.consistencyDefect_le_nondegenerateLinePastingDist_add_discarded_mass}
	\lean{MIPStarRE.QPBT.ProjectiveSetting.combinedLineMeasurement_consistency_le_conditioned}
	\leanok
	\uses{def:nondegenerate-line-pasting-mass, def:nondegenerate-line-pasting-dist,
		def:consistency, def:concrete-paired-line-measurement}
	Write $\mu=\mu_{\mathrm{nd}}(L)$ for the retained mass and use
	$D_{\mathrm{nd}}^L$ for the conditioned law. For arbitrary complete
	measurements on directed opposite placements, conditioning on nonzero $X$ direction
	increases the consistency defect by at most a factor $1/\mu$, while restoring
	the original law gives the exact bound
	$\operatorname{defect}_{\mathrm{all}}\leq
	\mu\operatorname{defect}_{\mathrm{cond}}+
	\Pr_{(\ell_X,x),(\ell_Z,z)}[\operatorname{dir}(\ell_X)=0]$.
	Bounding that discarded probability gives
	\[
		\operatorname{defect}_{\mathrm{all}}
		\leq \mu\operatorname{defect}_{\mathrm{cond}}+\frac1{2q}.
	\]
	For the completed joint point marginals and the evaluated $Z$- and $X$-line
	measurements in the order required by the $X$-outer sandwich, there is a
	universal $C\geq1$ for which both conditioned defects are at most
	\[
		\frac{8\delta_Q+C(\eps+\sqrt\eps)}{\mu}.
	\]
	The same restoration inequality applies to the concrete paired-line
	measurement and the completed joint points.
\end{lemma}
\begin{proof}
	\leanok
	\uses{lem:qld-point-line-marginal-bounds, lem:qld-zero-direction-mass,
		def:nondegenerate-line-pasting-mass, def:nondegenerate-line-pasting-dist,
		thm:nondegenerate-line-pasting-mass-positive,
		def:consistency, def:concrete-paired-line-measurement}
	Restrict the finite expectation to the nondegenerate event and divide by its
	positive mass. The integrand is nonnegative and at most one on opposite
	placements, so restoring the discarded event costs at most its probability.
\end{proof}

\begin{theorem}[Formalization support: positivity on two register pairs]
	\label{thm:qld-specialized-placement-positivity}
	\lean{MIPStarRE.QPBT.ProjectiveSetting.place_AA'_mul_place_BA''_nonneg}
	\lean{MIPStarRE.QPBT.ProjectiveSetting.place_AB''_mul_place_BB'_nonneg}
	\leanok
	\uses{def:expanded-state, def:symmetric-equivalents}
	For every projective setting, let $X$ and $Y$ be positive operators on
	Alice's and Bob's expanded local spaces, respectively. Then
	\[
		X_{AA'}Y_{BA''}\geq 0,
		\qquad X_{AB''}Y_{BB'}\geq 0.
	\]
	These are formalization-only consequences of the register placements,
	not additional hypotheses on the source consistency statements.
\end{theorem}
\begin{proof}
	\leanok
	\uses{def:expanded-state, def:symmetric-equivalents}
	Placement preserves positivity. On each displayed pair the two operators
	act on disjoint registers and therefore commute. The product of commuting
	positive operators is positive.
\end{proof}

\begin{definition}[The opposite-placement bipartition]
	\label{def:qld-opposite-bipartition}
	\lean{MIPStarRE.QPBT.bbAbBipartition}
	\leanok
	\uses{def:expanded-state}
	For the directed placement with the $BB'$ registers first and the $AB''$
	registers second, regroup the six registers as
	\[
		BB'\mid AB''(A'A'').
	\]
	Concretely, the coordinate tuple
	$((A,(A',A'')),(B,(B',B'')))$ is sent to
	$((B,B'),((A,B''),(A',A'')))$. The final factor is retained as an unused
	tensor factor.
\end{definition}

\begin{lemma}[Formalization support: opposite-placement transport]
	\label{lem:qld-opposite-bipartition-transport}
	\lean{MIPStarRE.QPBT.bbAbBipartition_apply}
	\lean{MIPStarRE.QPBT.ProjectiveSetting.exists_opposite_bipartition}
	\lean{MIPStarRE.QPBT.ProjectiveSetting.consistencyDefect_bipartition}
	\lean{MIPStarRE.QPBT.evaluated_pastedMeasurement_heteroKron_one}
	\lean{MIPStarRE.QPBT.postprocess_pair_swap_effect}
	\lean{MIPStarRE.QPBT.ProjectiveSetting.combinedLineMeasurement_evaluated_eq_pasted}
	\lean{MIPStarRE.QPBT.consistencyDefect_opposite_symm}
	\leanok
	\uses{def:qld-opposite-bipartition, def:concrete-paired-line-measurement,
		def:consistency, def:expanded-state}
	Every directed pair of opposite placements admits a bipartition of the six
	registers in which operators on the first placement act as $A\otimes I$ and
	operators on the second act as $I\otimes B\otimes I$; the remaining EPR
	pair is unused. Reindexing through this bipartition preserves consistency
	defects. Tensoring the pasted measurement with the identity on the unused
	factor commutes with outcome evaluation.

	For any paired postprocessing, swapping the two output coordinates swaps the
	corresponding effect. Evaluating the concrete $X$-$Z$-$X$ paired-line
	measurement gives exactly the evaluated sandwich used by one-sided pasting.
	Finally, the consistency defect of two measurements on opposite placements
	is unchanged when their order is reversed.
\end{lemma}
\begin{proof}
	\leanok
	\uses{def:qld-opposite-bipartition, def:concrete-paired-line-measurement,
		def:consistency, def:expanded-state}
	Check the four directed placement pairs by their coordinate equivalences.
	The remaining assertions follow by expanding reindexing, tensor products,
	postprocessing sums, the sandwich definition, and the consistency defect.
\end{proof}

\begin{lemma}[Formalization support: the restored paired-line defect]
	\label{lem:qld-paired-line-restored-defect}
	\lean{MIPStarRE.QPBT.exists_combinedLine_conditioned_defect_le}
	\lean{MIPStarRE.QPBT.exists_combinedLine_restored_defect_le}
	\leanok
	\uses{lem:qld-line-conditioning-restoration,
		lem:qld-opposite-bipartition-transport, lem:pasting,
		def:concrete-paired-line-measurement}
	There are a universal constant $C\geq1$ and a polynomially small function
	$p$ of two variables such that, on every directed opposite placement, the
	conditioned consistency defect between the completed joint points and the
	evaluated concrete paired-line measurement is at most
	\[
		p\left(\frac{md}{q},
		\frac{8\delta_Q+C(\eps+\sqrt\eps)}{\mu}\right).
	\]
	Here $\mu$ is the nonzero-$X$-direction mass. On the original product
	line-point law, the corresponding defect is at most
	\[
		\mu p\left(\frac{md}{q},
		\frac{8\delta_Q+C(\eps+\sqrt\eps)}{\mu}\right)+\frac1{2q}.
	\]
\end{lemma}
\begin{proof}
	\leanok
	\uses{lem:qld-line-conditioning-restoration,
		lem:qld-opposite-bipartition-transport, lem:pasting,
		def:nondegenerate-line-pasting-dist,
		thm:nondegenerate-line-pasting-mass-positive,
		thm:nondegenerate-line-pasting-dist-probability,
		def:concrete-paired-line-measurement}
	Transport the two measurement families to a bipartite tensor product and
	apply heterogeneous one-sided pasting with the two conditioned marginal
	bounds. Positivity of the retained mass makes the error parameter
	nonnegative, and the conditioned law is normalized by
	Theorem~\ref{thm:nondegenerate-line-pasting-dist-probability}.
	Identify its sandwich with the concrete paired-line measurement,
	then restore the discarded event and its mass factor.
\end{proof}

\begin{lemma}[Consistency of the constructed paired-line measurement]
	\label{lem:combined-line-measurement-consistency}
	\lean{MIPStarRE.QPBT.combined_line_measurement_consistency}
	\leanok
	\uses{lem:qld-4-10, lem:qld-comm-line-cons, lem:pasting,
		def:concrete-paired-line-measurement, def:consistency,
		thm:qld-conditioned-polynomial-bound}
	This is a formalization-only construction theorem. Fix a polynomially small
	function $\delta_Q(\eps)$. There is a function
	$\delta_P(\eps,md/q)=\poly(\eps,md/q)$, in the corrected sum sense,
	such that for every projective strategy passing with error $\eps$ and
	every joint point family satisfying the conclusions of
	Lemma~\ref{lem:qld-4-10} with error $\delta_Q(\eps)$, the concrete POVM
	\[
		T^{\ell_X,\ell_Z}_{f_X,f_Z}
		=\hat M^{(\mathrm{Line},X),\ell_X}_{f_X}
		 \hat M^{(\mathrm{Line},Z),\ell_Z}_{f_Z}
		 \hat M^{(\mathrm{Line},X),\ell_X}_{f_X}
	\]
	has evaluated consistency defect at most $\delta_P(\eps,md/q)$ with that
	joint point family, averaged over two independent line-point samples,
	for every directed opposite placement.

	The two measurements are compared on the completed answer alphabet
	$(\F_q\cup\{\bot\})^2$. An outcome $(f_X,f_Z)$ of $T^{\ell_X,\ell_Z}$ is
	postprocessed to the pair of evaluations of $f_X$ and $f_Z$ at the two
	sampled points, with $\bot$ recorded in a coordinate whose polynomial
	admits no evaluation at its point. The joint point family is extended to
	the same alphabet by the inclusion $\F_q\subseteq\F_q\cup\{\bot\}$ in each
	coordinate, so it assigns the zero operator to every answer having $\bot$
	in a coordinate. The quantity bounded is therefore the full off-diagonal
	defect on $(\F_q\cup\{\bot\})^2$: pairs in which a line evaluation is
	undefined while the point family answers in $\F_q$ contribute to it, and
	it is not the summation over answers in $\F_q$ alone.

	The line-error function is chosen after $\delta_Q$ and its polynomial
	bound. This isolates the consistency part of the preceding construction
	while retaining its defining measurement; see
	\cite{gap:qpbt_subline-claims-line-marginal}.
\end{lemma}
\begin{proof}
	\leanok
	\uses{lem:qld-paired-line-restored-defect,
		thm:nondegenerate-line-pasting-mass-bounds,
		lem:qld-opposite-bipartition-transport,
		lem:consistency-defect-integrand-le-one,
		thm:qld-conditioned-polynomial-bound}
	Condition on the first line having nonzero direction. The two completed
	point-to-line marginal comparisons and heterogeneous pasting bound the
	actual sandwich measurement on that law. Restore the discarded probability,
	at most $1/(2q)$, and multiply the conditioned error by its retained mass.
	The bounds $1/2\leq\mu_{\mathrm{nd}}(L)\leq1$ allow polynomial absorption;
	together with the independent unit defect bound, this gives the result.
\end{proof}

\section{Applying the classical low-degree test}
\label{sec:apply-ldt}

The combined point and line measurements of the previous two sections are next
merged into a single measurement returning a pair of global polynomials that
encode the $X$-basis and the $Z$-basis information simultaneously. This is
done by encoding the pair as one polynomial in $2m+2$ variables and invoking
the quantum soundness of the classical low individual degree test
(Theorem~\ref{lem:ld-soundness}). The classical game, its question
distribution, and the line-point distribution are defined at dimension $2m+2$
only when $2m+2$ divides $q$ (Definition~\ref{def:ld-game}), a condition that
admissibility of $(q,m,d)$ does not supply. The seed-bearing source route is
blocked at the instantiation step for admissible parameters with $m \ge 2$.
The global-pair conclusion also has the directly indexed proof below. The
obstruction is recorded in Remark~\ref{rem:qld-4-7-divisibility} and
analyzed, together with candidate repairs, in
\cite{gap:qpbt_ld-dimension-divisibility}.

% Unmarked, tracked in issues #667 and #695. Projection inclusion now supplies
% affine compatibility and the coefficient formula. The coefficient carrier
% still does not identify all its answers with functions on singleton lines.
\begin{definition}[The combining map]
	\label{def:combine-map}
	\lean{MIPStarRE.QPBT.combinePoly, MIPStarRE.QPBT.combinePoly_eval}
	\lean{MIPStarRE.QPBT.combinePoly_mem_polyFunc}
	\lean{MIPStarRE.QPBT.combineLinePoly}
	\lean{MIPStarRE.QPBT.combineLinePolynomial_natDegree_le}
	\lean{MIPStarRE.QPBT.IsCombineLineCompatible, MIPStarRE.QPBT.combineLinePoly_spec}
	\notready
	\uses{def:polyfunc, def:line, def:ideg-deg-polynomials}
	For $f, g \in \mathrm{ideg}_{d,m}(\F_q)$, the polynomial
	$\mathrm{combine}_{f,g} \in \mathrm{ideg}_{d,2m+2}(\F_q)$ is defined by
	\[
		\mathrm{combine}_{f,g}(x, z, \alpha, \beta)
		= \alpha f(x) + \beta g(z)\,,
		\qquad x \in \F_q^m,\ z \in \F_q^m,\ \alpha \in \F_q,\
		\beta \in \F_q\,.
	\]
	The combining map is also defined for polynomials restricted to lines.
	Let $\ell$ be a line in $\F_q^{2m+2}$ with intercept
	$u = (u_X, u_Z, u_\alpha, u_\beta)$ and slope
	$v = (v_X, v_Z, v_\alpha, v_\beta)$, and let $\ell_X$ and $\ell_Z$ be
	lines in $\F_q^m$ such that for every point
	$(x, z, \alpha, \beta) \in \ell$ it holds that $x \in \ell_X$ and
	$z \in \ell_Z$. Then for any two polynomials
	$f \in \deg_{md}(\ell_X)$ and $g \in \deg_{md}(\ell_Z)$, the polynomial
	$\mathrm{combine}_{f,g} \in \deg_{md+1}(\ell)$ is defined by
	\begin{equation}
		\label{eq:combine-lines}
		\mathrm{combine}_{f,g}(u + t \cdot v)
		= (u_\alpha + t \cdot v_\alpha)\, f(u_X + t \cdot v_X)
		+ (u_\beta + t \cdot v_\beta)\, g(u_Z + t \cdot v_Z)\,,
		\qquad t \in \F_q\,.
	\end{equation}
	The degree bound holds because $f$ and $g$, restricted to lines, are
	polynomials of degree at most $md$ in the line parameter, and each is
	multiplied by an affine factor of degree at most one.
\end{definition}

\begin{lemma}[Formalization support for the combining map: affine parameters]
	\label{lem:combine-map-affine-parameters}
	\lean{MIPStarRE.QPBT.exists_affine_parameters_of_linePoints_subset}
	\lean{MIPStarRE.QPBT.exists_isCombineLineCompatible_of_projection_mem}
	\leanok
	\uses{def:line, def:combine-map}
	This auxiliary isolates the geometric consequence of the projection
	hypothesis in Definition~\ref{def:combine-map}. Over any field $K$, if
	$\ell(u,v)\subseteq\ell(u',v')$ in $K^m$, there are $a,b\in K$ such that
	$u=u'+av'$ and $v=bv'$. Consequently, under the projection inclusions
	of that definition, there are $a_X,b_X,a_Z,b_Z\in K$ such that
	\[
		(u+tv)_W=u_W'+(a_W+b_Wt)v_W',
		\qquad W\in\{X,Z\},\quad t\in K,
	\]
	where $\ell_W=\ell(u_W',v_W')$. No direction is required to be nonzero.
\end{lemma}
\begin{proof}
	\leanok
	\uses{def:line}
	Choose parameters $a,s$ for the points $u$ and $u+v$ on the containing
	line. Subtracting $u=u'+av'$ from $u+v=u'+sv'$ gives
	$v=(s-a)v'$. Take $b=s-a$. Apply this argument to each projected line.
\end{proof}

\begin{lemma}[Formalization support for the combining map: parameter polynomials]
	\label{lem:combine-map-parameter-polynomial}
	\lean{MIPStarRE.QPBT.exists_combineLinePoly_of_projection_mem}
	\leanok
	\uses{def:line, def:combine-map, def:ideg-deg-polynomials}
	This auxiliary concerns coefficient polynomials in a line parameter.
	Let $K$ be a field, let $c\geq0$, and assume the projection inclusions
	of Definition~\ref{def:combine-map}. There are compatible affine
	parameters $a_X,b_X,a_Z,b_Z$, independent of the coefficient lists,
	such that for every $f,g\in K[T]$ of degree at most $c$, the polynomial
	\[
		H(T)=(u_\alpha+v_\alpha T)f(a_X+b_XT)
		      +(u_\beta+v_\beta T)g(a_Z+b_ZT)
	\]
	has degree at most $c+1$. Its coefficient list, through degree $c+1$,
	evaluates at every $t\in K$ to
	\[
		H(t)=(u+tv)_\alpha f(a_X+b_Xt)
		     +(u+tv)_\beta g(a_Z+b_Zt).
	\]
	This assertion includes constant projections and zero directions. It
	does not identify every parameter polynomial with a function on the
	geometric point set of its line.
\end{lemma}
\begin{proof}
	\leanok
	\uses{lem:combine-map-affine-parameters, def:combine-map}
	Construct the affine parameters by the preceding lemma. Affine
	substitution preserves the degree bound $c$; multiplication by an affine
	factor raises it by at most one. The coefficient evaluation formula
	follows by evaluation of the two compositions and their sum.
\end{proof}

\begin{lemma}[Formalization support for the combining map: singleton evaluation]
	\label{lem:combine-map-singleton-evaluation}
	\lean{MIPStarRE.QPBT.evaluatesTo_zero_direction_iff}
	\lean{MIPStarRE.QPBT.not_evaluatesTo_parameter_on_zero_direction}
	\leanok
	\uses{def:ideg-deg-polynomials, def:ld-line-description}
	For the coefficient interpretation of
	Definition~\ref{def:ideg-deg-polynomials}, let $\ell=\ell(u_0,0)$ be a
	line description and let $f$ be a coefficient polynomial of degree at
	most $c$. It evaluates to $a$ at $u_0$ if and only if
	$f(t)=a$ for every $t\in\F_q$. If $c\geq1$, the answer $f(T)=T$
	therefore has no value at $u_0$.
\end{lemma}
\begin{proof}
	\leanok
	\uses{def:ideg-deg-polynomials}
	Every parameter represents $u_0$ when the direction is zero. The
	definition of evaluation thus requires the same value at all
	parameters. For $f(T)=T$, the parameters zero and one give distinct
	values.
\end{proof}

\begin{remark}[The function domain of the combining map]
	\label{rem:combine-map-function-domain}
	Equation~\eqref{eq:combine-lines} is well defined for polynomial
	functions on the geometric lines. The coefficient interpretation in
	Remark~\ref{rem:deg-line-representatives} also admits answers that do
	not define such functions, as the preceding lemma shows. Choosing an
	affine parameter for a singleton does not resolve this discrepancy:
	for $f(T)=T$, two constant parameter choices give different values.
	There is also a distinction between formal representatives and their
	functions even on a nondegenerate line: over $\F_q$, the distinct
	polynomials $0$ and $T^q-T$ induce the same function. Thus the condition
	in the preceding lemma is constancy as a function, rather than formal
	degree zero. The geometric parameter construction does not by itself
	establish the identification of these carriers; see
	\cite{gap:qpbt_subline-claims-line-marginal}.
\end{remark}

\begin{definition}[Formalization support for Definition~\ref{def:combine-map}: the separated locus]
	\label{def:combine-map-separated-locus}
	\lean{MIPStarRE.QPBT.combinedBlockVar}
	\lean{MIPStarRE.QPBT.IsSeparatedCombined}
	\lean{MIPStarRE.QPBT.recoverCombinedPoly}
	\lean{MIPStarRE.QPBT.SeparatedCombinedPoly}
	\leanok
	\uses{def:combine-map, def:polyfunc, def:ideg-deg-polynomials}
	This formalization-only auxiliary is not a named statement of the paper; it
	isolates a step used in the Lean proof of Lemma~\ref{lem:qld-4-7}, where the
	set of combined polynomials is written $G$ and is then reindexed by pairs
	$(f, g)$. Separation itself carries no degree bound and no restriction on
	the coefficients: call a polynomial $p$ in the $2m+2$ coordinates, over an
	arbitrary commutative semiring, \emph{separated} when, viewed as a
	polynomial in $\alpha$ and $\beta$ over the point variables, every
	coefficient of $p$ vanishes except those of $\alpha$ and of $\beta$, the
	coefficient of $\alpha$ involves only the $x$ variables, and the
	coefficient of $\beta$ involves only the $z$ variables. The recovery map is
	likewise defined for every such $p$, with no degree parameter: setting
	$(\alpha, \beta) = (1, 0)$ and $z = 0$ returns the first component, and
	setting $(\alpha, \beta) = (0, 1)$ and $x = 0$ returns the second. The
	\emph{bounded separated locus} is the intersection of this unrestricted
	condition with the individual-degree class, that is, the set of separated
	$p \in \mathrm{ideg}_{d,2m+2}(\F_q)$; only this second object depends on
	$d$. The condition is stated by coefficient vanishing and variable
	containment, not by assuming a pair $(f, g)$ with
	$p = \mathrm{combine}_{f,g}$.
\end{definition}

\begin{lemma}[Formalization support for Definition~\ref{def:combine-map}: the separated image]
	\label{lem:combine-map-separated-image}
	\lean{MIPStarRE.QPBT.combinedBlockVar_injective}
	\lean{MIPStarRE.QPBT.combinePoly_eq_combinePolyTuple}
	\lean{MIPStarRE.QPBT.recoverCombinedPoly_combinePoly}
	\lean{MIPStarRE.QPBT.isSeparatedCombined_combinePoly}
	\lean{MIPStarRE.QPBT.rename_recoverCombinedPoly}
	\lean{MIPStarRE.QPBT.combinePoly_recoverCombinedPoly}
	\lean{MIPStarRE.QPBT.isSeparatedCombined_iff_exists_combinePoly}
	\lean{MIPStarRE.QPBT.recoverCombinedPoly_mem_polyFunc}
	\lean{MIPStarRE.QPBT.combinePoly_mem_polyFunc_iff}
	\lean{MIPStarRE.QPBT.existsUnique_bounded_combinePoly}
	\lean{MIPStarRE.QPBT.combinedPolynomialImageEquiv}
	\leanok
	\uses{def:combine-map, def:combine-map-separated-locus, def:polyfunc,
		def:ideg-deg-polynomials, lem:qld-combined-pair-completion}
	This formalization-only auxiliary is not a named statement of the paper; it
	records the algebraic identification of the set $G$ of the proof of
	Lemma~\ref{lem:qld-4-7} with pairs of bounded-degree polynomials. Each of
	the two point-block inclusions is injective. The map
	$(f, g) \mapsto \mathrm{combine}_{f,g}$ is injective, which is recorded in
	Lemma~\ref{lem:qld-combined-pair-completion}. Its image is exactly
	the separated locus of
	Definition~\ref{def:combine-map-separated-locus}, and the two specialization
	maps invert it there. For every $d$, including $d = 0$, the components
	recovered from a separated $p \in \mathrm{ideg}_{d,2m+2}(\F_q)$ again lie in
	$\mathrm{ideg}_{d,m}(\F_q)$, and they form the only pair of components of
	$p$ in $\mathrm{ideg}_{d,m}(\F_q)$; neither statement needs $d \ge 1$. That
	hypothesis enters only in the forward degree bound: for $d \ge 1$, a
	combined polynomial lies in $\mathrm{ideg}_{d,2m+2}(\F_q)$ if and only if
	both components lie in $\mathrm{ideg}_{d,m}(\F_q)$, and together with the
	preceding statements this gives a bijection between pairs of
	bounded-degree polynomials and the bounded separated locus. The
	statements hold over an arbitrary commutative semiring of coefficients. No
	bound on the mass the strategy places on the separated image is asserted
	here; that estimate belongs to the proof of Lemma~\ref{lem:qld-4-7}.
\end{lemma}
\begin{proof}
	\leanok
	\uses{def:combine-map, def:combine-map-separated-locus,
		def:ld-combining-map, lem:ld-combining-split,
		lem:ld-combining-exact-linearity}
	Write $\mathrm{combine}_{f,g}$ as the tuple combining map applied to $f$ and
	to $g$ after renaming their variables into the $x$ block and the $z$ block
	respectively; the two block inclusions are injective. The component
	splitting of Lemma~\ref{lem:ld-combining-split} then recovers each component
	from the coefficient of the corresponding combining variable, which gives
	both round-trip identities. If $p$ is separated, the
	coefficient characterization of
	Lemma~\ref{lem:ld-combining-exact-linearity} writes $p$ as a combined tuple,
	and its two coefficients are supported in the respective blocks, so each is
	the renaming of a polynomial in $m$ variables and recombining the two
	returns $p$; this argument uses no bound on $d$. Individual degrees are
	unchanged by renaming along an injection, so each component inherits the
	bound of $p$ from the bound on the coefficients of $p$, again for every $d$;
	conversely the degree bound of the tuple combining map of
	Definition~\ref{def:ld-combining-map} supplies the bound for
	$\mathrm{combine}_{f,g}$ when $d \ge 1$, which is the hypothesis that also
	bounds the two fresh variables $\alpha$ and $\beta$.
\end{proof}

\begin{lemma}[Formalization support for the algebra of register placements]
	\label{lem:qld-placement-completeness}
	\lean{MIPStarRE.QPBT.place_sum}
	\lean{MIPStarRE.QPBT.sum_placed_measurement_eq_one}
	\lean{MIPStarRE.QPBT.sum_placed_measurement_products_eq_one}
	\lean{MIPStarRE.QPBT.ProjectiveSetting.place_mul}
	\lean{MIPStarRE.QPBT.ProjectiveSetting.place_comm}
	\lean{MIPStarRE.QPBT.ProjectiveSetting.place_conjTranspose_mul_self_le_one}
	\leanok
	\uses{def:symmetric-equivalents, def:povm-conventions}
	This auxiliary lemma isolates the finite-sum identities used in
	Lemma~\ref{lem:qld-4-12}; it is not a named statement of
	\cite{Ji2020MIPStar}. Fix any register placement $p$ from
	Definition~\ref{def:symmetric-equivalents}, and write $(A)_p$ for the
	operator acting as $A$ on those registers and as the identity elsewhere.
	For a finite operator family $\{A_a\}_{a\in I}$ on the corresponding
	expanded local space,
	\[
		\Big(\sum_{a\in I}A_a\Big)_p=\sum_{a\in I}(A_a)_p.
	\]
	For any two complete measurements $\{M_a\}_{a\in I}$ and
	$\{N_b\}_{b\in J}$ on that space, with finite outcome sets $I,J$,
	\[
		\sum_{a\in I}(M_a)_p=\Id,
		\qquad
		\sum_{(a,b)\in I\times J}(M_aN_b)_p=\Id.
	\]
	Neither projectivity nor commutation between the measurements is required.

	The same placement is moreover multiplicative,
	$(AB)_p = (A)_p (B)_p$ for all operators $A, B$ on the expanded local
	space of the side of $p$, and operators on opposite placements commute:
	if $p_1$ and $p_2$ are opposite, then
	$(A)_{p_1} (B)_{p_2} = (B)_{p_2} (A)_{p_1}$ for every $A$ on the side of
	$p_1$ and every $B$ on the side of $p_2$. If moreover $E$ is a projector
	on that space, then $(E)_p$ is a contraction,
	$(E)_p^\dagger (E)_p \le \Id$.
\end{lemma}
\begin{proof}
	\leanok
	\uses{def:symmetric-equivalents, def:povm-conventions}
	Each placement preserves finite sums and identity by its matrix-entry
	formula. The first completeness identity follows by placing
	$\sum_a M_a=\Id$. For the second, distributivity and completeness give
	\[
		\sum_{a,b}M_aN_b
		=\Big(\sum_a M_a\Big)\Big(\sum_b N_b\Big)=\Id;
	\]
	apply preservation of sums and identity once more.
	Multiplicativity and opposite-placement commutation are read off from the
	expression of each placement as a reindexed tensor factor: the two
	operands of a product sit in the same factor, and two opposite placements
	sit in different factors of one bipartition. For the contraction bound,
	multiplicativity and compatibility with the adjoint turn
	$(E)_p^\dagger (E)_p$ into $(E^\dagger E)_p = (E)_p$ for a projector
	$E$, and $E \le \Id$ is carried to $(E)_p \le \Id$ because a placement
	is positive and preserves the identity.
\end{proof}

\begin{lemma}[Formalization support for averaged affine coarse-graining]
	\label{lem:qld-affine-distance-contraction}
	\lean{MIPStarRE.QPBT.opFamilyDistSq_uniform_affine_postprocess_le}
	\leanok
	\uses{def:povm-distance}
	This auxiliary lemma isolates the squared-distance calculation used in
	Lemma~\ref{lem:qld-4-12}; it is not a named statement of the source paper.
	Let $U$ be a nonempty finite set, let $K$ be a finite field of order $q$,
	and let $\ket{\psi}$ be a vector in a finite-dimensional complex Hilbert
	space. For each $u\in U$, let $\{A^u_p\}_{p\in K^2}$ and
	$\{B^u_p\}_{p\in K^2}$ be operator families satisfying
	$\sum_p A^u_p=\sum_p B^u_p$. For $\alpha,\beta\in K$, put
	$h_{\alpha,\beta}(a,b)=\alpha a+\beta b$ and define
	\[
		A^{u,\alpha,\beta}_c=\sum_{p:h_{\alpha,\beta}(p)=c}A^u_p,
		\qquad
		B^{u,\alpha,\beta}_c=\sum_{p:h_{\alpha,\beta}(p)=c}B^u_p.
	\]
	Then, with all expectations uniform and independent,
	\[
		\E_{u,\alpha,\beta}\sum_{c\in K}
		\big\|(A^{u,\alpha,\beta}_c-B^{u,\alpha,\beta}_c)\ket{\psi}\big\|^2
		\le
		\E_u\sum_{p\in K^2}\big\|(A^u_p-B^u_p)\ket{\psi}\big\|^2.
	\]
\end{lemma}
\begin{proof}
	\leanok
	Fix $u$ and put $v_p=(A^u_p-B^u_p)\ket{\psi}$. Equality of the
	operator totals gives $\sum_p v_p=0$. For distinct $p,p'\in K^2$,
	the equation $h_{\alpha,\beta}(p)=h_{\alpha,\beta}(p')$ is one
	nonzero homogeneous linear equation in $\alpha,\beta$, so it has $q$
	solutions among the $q^2$ coefficient pairs. Thus its probability is
	$1/q$, whereas a pair collides with itself with probability one.
	Expanding the squared norms and averaging gives
	\[
		\E_{\alpha,\beta}\sum_c
		\Big\|\sum_{p:h_{\alpha,\beta}(p)=c}v_p\Big\|^2
		=\Big(1-\frac1q\Big)\sum_p\|v_p\|^2
		 +\frac1q\Big\|\sum_p v_p\Big\|^2
		=\Big(1-\frac1q\Big)\sum_p\|v_p\|^2
		\le\sum_p\|v_p\|^2.
	\]
	Averaging over $u$ proves the assertion. No positivity or commutativity
	of the operator families is required.
\end{proof}

\begin{lemma}[Combined point measurement on the extended coordinates]
	\label{lem:qld-4-12}
	\lean{MIPStarRE.QPBT.exists_extendedQ}
	\leanok
	\uses{lem:qld-4-10, def:expanded-state, def:expanded-point-measurement,
		def:povm-distance, def:symmetric-equivalents}
	For all $x, z \in \F_q^m$, $\alpha, \beta \in \F_q$, and
	$\hilb \in \{\hilb_{\mathsf{A}}, \hilb_{\mathsf{B}}\}$ there exists a
	projective measurement
	$\{\hat{Q}^{x,z,\alpha,\beta}_{c}\}_{c \in \F_q}$ acting on
	$\hilb \ot (\C^q)^{\ot M}$ such that the following holds, where
	$\delta_Q$ is the function of Lemma~\ref{lem:qld-4-10}:
	\begin{enumerate}
		\item (\textbf{Self-consistency}) On average over uniformly random
		$(x, z, \alpha, \beta) \in \F_q^{2m+2}$,
		\begin{equation}
			\label{eq:qld-4-12-self-cons}
			(\hat{Q}^{x,z,\alpha,\beta}_{c})_{\mathsf{A}\mathsf{A}'}
			\approx_{\delta_Q}
			(\hat{Q}^{x,z,\alpha,\beta}_{c})_{\mathsf{B}\mathsf{A}''}\,.
		\end{equation}
		\item (\textbf{Consistency with $\hat{M}$}) On average over uniformly
		random $(x, z, \alpha, \beta) \in \F_q^{2m+2}$,
		\begin{gather}
			\label{eq:qld-4-12-cons-m-hat}
			(\hat{Q}^{x,z,\alpha,\beta}_{c})_{\mathsf{A}\mathsf{A}'}
			\approx_{\delta_Q}
			\Big( \sum_{\substack{a,b:\\ \alpha a + \beta b = c}}
			\hat{M}^{(\mathrm{Point},X),x}_{a} \,
			\hat{M}^{(\mathrm{Point},Z),z}_{b}
			\Big)_{\mathsf{B}\mathsf{A}''}\,, \\
			\label{eq:qld-4-12-cons-m-hat2}
			(\hat{Q}^{x,z,\alpha,\beta}_{c})_{\mathsf{A}\mathsf{A}'}
			\approx_{\delta_Q}
			\Big( \sum_{\substack{a,b:\\ \alpha a + \beta b = c}}
			\hat{M}^{(\mathrm{Point},Z),z}_{b} \,
			\hat{M}^{(\mathrm{Point},X),x}_{a}
			\Big)_{\mathsf{B}\mathsf{A}''}\,.
		\end{gather}
	\end{enumerate}
	Furthermore, all symmetric equivalents of these approximations also hold.
\end{lemma}
\begin{proof}
	\lean{MIPStarRE.QPBT.CombinedPointsWitness.extendedQ}
	\lean{MIPStarRE.QPBT.extendedQ_spec}
	\leanok
	\uses{lem:qld-4-10, fact:data-processing,
		lem:qld-placement-completeness, lem:qld-affine-distance-contraction}
	Define
	$\hat{Q}^{x,z,\alpha,\beta}_{c}
	= \sum_{a,b \,:\, \alpha a + \beta b = c} \hat{Q}^{x,z}_{a,b}$,
	a coarse-graining of the projective measurement of
	Lemma~\ref{lem:qld-4-10}, hence itself projective. A uniformly random
	point of $\F_q^{2m+2}$ has uniformly random coordinates $(x,z)$, so the
	self-consistency and consistency properties follow from those of
	$\{\hat{Q}^{x,z}_{a,b}\}$ established in Lemma~\ref{lem:qld-4-10}
	together with Lemma~\ref{fact:data-processing}, which shows that
	coarse-graining the outcomes preserves the $\delta_Q$ bounds.
\end{proof}

\begin{lemma}[Combined line measurement on the extended coordinates]
	\label{lem:qld-4-13}
	\lean{MIPStarRE.QPBT.ExtendedLinesWitness}
	\lean{MIPStarRE.QPBT.PrintedExtendedLinesWitnessClaim}
	\notready
	\uses{lem:qld-4-12, def:line, def:line-point-dist, def:consistency,
		def:bracket, def:expanded-state, def:ideg-deg-polynomials}
	% Statement as in the source (references/qpbt-paper/
	% 14_analysis_of_the_pauli_basis_test.tex, lem:qld-4-13), including the
	% error form poly(m^2 eps, md/q); neither of the source's two proof
	% routes delivers that form, and the route retained below establishes
	% m * poly(eps, md/q), which is the form the sequel uses. See
	% rem:qld-4-13-source-defects and the analysis in
	% docs/paper-gaps/qpbt_combined-lines-error-term.tex; the open printed-error
	% form is tracked by GitHub issue 598. Since 2026-09-19 the printed claim is
	% carried in Lean by MIPStarRE.QPBT.PrintedExtendedLinesWitnessClaim, a
	% definition of type Prop that states the printed conclusion without
	% asserting it; it replaced a theorem of the same statement whose proof was
	% an open sorry. This node therefore still carries no formalization mark,
	% and the established content of the node is lem:qld-4-13-established.
	% The line-point
	% distribution over F_q^{2m+2} presupposes that 2m+2 divides q; see
	% rem:qld-4-7-divisibility.
	There exists a function
	$\delta_{\mathrm{combine}}(\eps,m,d,q) = \poly(m^2\eps, md/q)$, the
	two-variable polynomial error read in the corrected sum sense of
	\cite{gap:qpbt_pasting-product-error}, that is, a bound
	$C\bigl((m^2\eps)^{r} + (md/q)^{s}\bigr)$ with $C \ge 1$ and $r, s > 0$ on
	the closed nonnegative quadrant rather than the literal product form of the
	shorthand, such that the following holds. For
	every line $\ell$ in $\F_q^{2m+2}$ and
	$\hilb \in \{\hilb_{\mathsf{A}}, \hilb_{\mathsf{B}}\}$ there exists a
	POVM $\{\hat{Q}^{\ell}_{f}\}$ acting on $\hilb \ot (\C^q)^{\ot M}$ with
	outcomes $f \in \deg_{md+1}(\ell)$ such that on average over a uniformly
	random pair $(\ell, (x,z,\alpha,\beta))$ drawn from the line-point
	distribution over $\F_q^{2m+2}$,
	\begin{gather}
		\label{eq:qld-4-13}
		\big(\hat{Q}^{\ell}
		_{[\mathrm{eval}_{(x,z,\alpha,\beta)}(\cdot)=a]}
		\big)_{\mathsf{A}\mathsf{A}'}
		\simeq_{\delta_{\mathrm{combine}}}
		\big(\hat{Q}^{x,z,\alpha,\beta}_{a}\big)_{\mathsf{B}\mathsf{A}''}\,,
		\\
		\big(\hat{Q}^{\ell}
		_{[\mathrm{eval}_{(x,z,\alpha,\beta)}(\cdot)=a]}
		\big)_{\mathsf{B}\mathsf{B}'}
		\simeq_{\delta_{\mathrm{combine}}}
		\big(\hat{Q}^{x,z,\alpha,\beta}_{a}\big)_{\mathsf{A}\mathsf{B}''}\,,
		\notag
	\end{gather}
	where the answer summation is over $a \in \F_q$. Moreover, if $\ell$ is
	axis-parallel, then the outcome $f \in \deg_{d}(\ell)$.
\end{lemma}

Before proving the lemma we record definitions and estimates concerning
distributions over lines and points. Recall from
Definition~\ref{def:line-point-dist} that the line-point distribution over
$\F_q^m$ is an even mixture of source laws $D_{\mathrm{ALine}}$ and
$D_{\mathrm{DLine}}$, and that their seed-bearing refinements are denoted by
$\widetilde D_{\mathrm{ALine}}$, $\widetilde D_{\mathrm{DLine}}$, and
$\widetilde D_{\mathrm{Line}}$. In the auxiliary statements below, tilded
symbols denote these refined laws and untilded symbols denote the source laws.
Coordinate indices below are zero-based: at ambient dimension $r$, the index
$i\in\{0,\dots,r-1\}$ denotes coordinate $i+1$ in the one-based notation of
\cite{Ji2020MIPStar}.

\begin{definition}[Restricted line distributions]
	\label{def:ith-restricted-line}
	\lean{MIPStarRE.QPBT.restrictedALineDist, MIPStarRE.QPBT.restrictedDLineDist}
	\lean{MIPStarRE.QPBT.restrictedALineDist_isProbability}
	\lean{MIPStarRE.QPBT.restrictedDLineDist_isProbability}
	\leanok
	\uses{def:line-point-dist, def:ld-question-distribution}
	Use coordinate indices $i\in\{0,\dots,m-1\}$ throughout this definition,
	the following mixture estimates, and the sub-line construction below. Write
	$\bar\chi(s)=\chi(s)-1$, where $\chi$ is the one-based index function of
	Definition~\ref{def:ld-question-distribution}. Thus $\bar\chi$ takes
	values in $\{0,\dots,m-1\}$, and index $i$ here denotes coordinate $i+1$
	in the one-based notation of \cite{Ji2020MIPStar}.
	For any such $i$, the \emph{$i$-th restricted refined
	axis-parallel line distribution} $\widetilde D_{\mathrm{ALine},i}$ is the
	restriction of $\widetilde D_{\mathrm{ALine}}$ to seed-bearing pairs
	$(\ell,u)$ whose retained seed $s$ satisfies $\bar\chi(s)=i$, so the geometric
	direction of $\ell$ is $e_i$. Similarly, the \emph{$i$-th restricted
	refined diagonal line distribution} $\widetilde D_{\mathrm{DLine},i}$ is
	the restriction of $\widetilde D_{\mathrm{DLine}}$ to pairs $(\ell,u)$
	where $\ell=(u_0,s,v)$ satisfies $\bar\chi(s)=i$ and coordinates
	$0, \dots, i-1$ of $v$ are $0$, with no such coordinates when $i=0$.
	Their push-forwards under $\rho_{\mathrm{ALine}}$ and
	$\rho_{\mathrm{DLine}}$ from
	Definition~\ref{def:line-point-dist} are the source restrictions
	$D_{\mathrm{ALine},i}$ and $D_{\mathrm{DLine},i}$ of
	\cite{Ji2020MIPStar}. Both refined restricted distributions have total
	mass one.
\end{definition}

\begin{lemma}[Restricted line distributions and error inflation]
	\label{lem:restricted-line-mixture-bounds}
	\lean{MIPStarRE.QPBT.linePointDist_eq_mixture_restricted}
	\lean{MIPStarRE.QPBT.avg_restricted_le, MIPStarRE.QPBT.avg_restricted_prod_le}
	\leanok
	% Source: unlabelled prose in the paper --- the discussion following
	% def:ith-restricted-line at lines 1049-1061 of
	% references/qpbt-paper/14_analysis_of_the_pauli_basis_test.tex.
	\uses{def:ith-restricted-line, def:line-point-dist}
	The refined distribution $\widetilde D_{\mathrm{ALine}}$ is a uniform
	mixture of the distributions $\widetilde D_{\mathrm{ALine},i}$ for
	$i \in \{0, \dots, m-1\}$, and likewise
	$\widetilde D_{\mathrm{DLine}}$ is a uniform mixture of the
	$\widetilde D_{\mathrm{DLine},i}$. For nonnegative real-valued functions
	on seed-bearing line-point samples, the following bounds hold:
	\begin{enumerate}
		\item any average over $\widetilde D_{\mathrm{Line}}$ bounded by
		$\delta$ is bounded over each restricted distribution
		$\widetilde D_{\mathrm{ALine},i}$ and
		$\widetilde D_{\mathrm{DLine},i}$ by $2m\delta$;
		\item any average over
		$\widetilde D_{\mathrm{Line}}\times\widetilde D_{\mathrm{Line}}$
		bounded by $\delta$ is bounded over any product of refined restricted
		distributions by $4m^2\delta$.
	\end{enumerate}
	Pulling a source function back along the forgetful maps gives the same
	mixture identities and bounds for the untilded source laws.
\end{lemma}
\begin{proof}
	\leanok
	\uses{def:ith-restricted-line, def:line-point-dist}
	The refined line-point distribution places weight $\tfrac{1}{2}$ on each of
	the axis-parallel and diagonal components and, within each component,
	uniform weight on the index $i$, so each refined restriction carries mixture
	weight $\tfrac{1}{2m}$ in $\widetilde D_{\mathrm{Line}}$. A nonnegative
	average that is at most $\delta$ therefore bounds each mixture component by
	$2m\delta$, and products of two such mixtures give the factor
	$(2m)^2 = 4m^2$.
\end{proof}

\begin{lemma}[Restricted consistency of the source line measurements]
	\label{lem:qld-xz-lines-restricted}
	\notready
	% tracked in issue #414
	\uses{lem:qld-xz-lines, lem:restricted-line-mixture-bounds,
		def:consistency}
	For the joint point and line measurements furnished by
	Lemma~\ref{lem:qld-xz-lines}, for any types
	$\mathrm{v}_1, \mathrm{v}_2 \in \{\mathrm{ALine},\mathrm{DLine}\}$ and
	any $i, j \in \{0, \dots, m-1\}$,
	\begin{equation}
		\label{eq:qld-xz-lines-restricted}
		\E_{(\ell_X,x) \sim D_{\mathrm{v}_1,i}}\
		\E_{(\ell_Z,z) \sim D_{\mathrm{v}_2,j}}
		\sum_{f_X,f_Z}
		\bra{\hat\psi}
		\big(T^{\ell_X,\ell_Z}_{f_X,f_Z}\big)_{\mathsf{A}\mathsf{A}'}
		\ot
		\big(\Id - \hat{Q}^{x,z}_{f_X(x),f_Z(z)}
		\big)_{\mathsf{B}\mathsf{A}''}
		\ket{\hat\psi}
		= O(m^2 \delta_P)\,.
	\end{equation}
\end{lemma}
\begin{proof}
	\uses{lem:qld-xz-lines, lem:restricted-line-mixture-bounds,
		def:consistency}
	The consistency conclusion of Lemma~\ref{lem:qld-xz-lines} bounds the
	displayed nonnegative integrand on average over two independent copies of
	$D_{\mathrm{Line}}$. Apply the product restriction estimate in
	Lemma~\ref{lem:restricted-line-mixture-bounds} to obtain the factor
	$4m^2$.
\end{proof}

\begin{theorem}[Formalization support: conditional restricted consistency]
	\label{thm:conditional-restricted-lines-consistency}
	\lean{MIPStarRE.QPBT.restricted_lines_consistency_bound}
	\leanok
	\uses{def:ith-restricted-line, def:consistency,
		def:expanded-state}
	There is a universal constant $C>0$ with the following property. Fix
	admissible parameters, real errors $\eps,\delta_Q,\delta_P$, and a
	projective setting with expanded state $\hat\psi$. Suppose one is given a
	projective joint point family whose self-consistency and two ordered-product
	consistency bounds are at most $\delta_Q$, together with a paired-line
	measurement family having both axis-degree support properties and evaluated
	consistency defect at most $\delta_P$ with that point family on every
	directed opposite placement. Then, for every directed opposite placement
	pair, every $\mathrm{v}_1,\mathrm{v}_2\in\{\mathrm{ALine},\mathrm{DLine}\}$,
	and every $i,j\in\{0,\dots,m-1\}$, the same evaluated consistency defect,
	with both line-point variables restricted to
	$\widetilde D_{\mathrm{v}_1,i}\times
	\widetilde D_{\mathrm{v}_2,j}$, is at most
	$C m^2\delta_P$. Evaluation includes an additional undefined outcome when
	the affine parameter is not uniquely determined.

	This conditional auxiliary theorem assumes the point and line families,
	including the unrestricted consistency estimate for the latter. Its
	restrictions retain the line seeds. Passing this bound to the source laws
	would require the evaluated measurement families to factor through the
	forgetful maps of Definition~\ref{def:line-point-dist}. No such
	factorization is asserted here; the theorem neither constructs the source
	line measurement nor certifies
	Lemma~\ref{lem:qld-xz-lines-restricted}.
\end{theorem}
\begin{proof}
	\leanok
	\uses{lem:restricted-line-mixture-bounds, def:consistency}
	Apply the product restriction estimate to the nonnegative consistency-defect
	integrand. The supplied unrestricted bound gives the result with $C=4$.
\end{proof}

% This lemma is not a statement of the source; it records the finite
% mass bookkeeping of the proof-only conditioning step above.  It is a
% root of the dependency graph: it uses no earlier blueprint entry.
\begin{lemma}[Formalization support: conditioning and restoring a finite average]
	\label{lem:restricted-average-mass-split}
	\lean{MIPStarRE.QPBT.avgOver_restrict_mul_mass}
	\lean{MIPStarRE.QPBT.avgOver_le_restrict_add_discarded_mass}
	\leanok
	Let $\mu$ be a finitely supported nonnegative weighting of a sample type.
	Normalization is not assumed: neither linked declaration carries a
	probability hypothesis, and $\mu$ ranges over the arbitrary finite
	nonnegative weights of the formalized notion of distribution. Let $G$ be a
	decidable event whose retained mass
	$p = \mu(G) = \sum_{s \in G} \mu(s)$ is positive, and write
	$\mu \mid G$ for the normalized restriction of $\mu$ to $G$, with weights
	$\mu(s)/p$ on $G$ and $0$ off $G$; that restriction is a probability
	distribution even when $\mu$ is not. Then, for every real-valued $f$,
	\begin{equation}
		\label{eq:restricted-average-normalization}
		p \cdot \E_{\mu \mid G} f = \sum_{s \in G} \mu(s) f(s)\,,
	\end{equation}
	and if moreover $f(s) \le 1$ for every sample $s$,
	\begin{equation}
		\label{eq:restricted-average-restoration}
		\E_{\mu} f \le p \cdot \E_{\mu \mid G} f + \mu(G^{c})\,.
	\end{equation}
	Here $\mu(G^{c}) = \sum_{s \notin G} \mu(s)$ is the mass discarded by the
	conditioning; it is retained in full in
	\eqref{eq:restricted-average-restoration} and no term is dropped. Since
	$\mu$ need not be normalized, $p$ and $\mu(G^{c})$ are masses rather than
	probabilities, and their sum is the total mass of $\mu$ rather than one.
\end{lemma}
\begin{proof}
	\leanok
	The conditioned weights are $\mu(s)/p$ on $G$ and $0$ off $G$, which gives
	\eqref{eq:restricted-average-normalization} termwise. For
	\eqref{eq:restricted-average-restoration}, split the sum defining
	$\E_\mu f$ along $G$ and its complement: the first part is the right-hand
	side of \eqref{eq:restricted-average-normalization}, and the second is at
	most $\mu(G^{c})$ because the weights are nonnegative and $f \le 1$.
\end{proof}

\begin{lemma}[The sub-line distribution]
	\label{lem:qld-sublines}
% No declaration link: the available construction uses the directly indexed
% extended line-point law, not the seed-indexed law over F_q^(2m+2) of this
% statement.  That construction is Definition~\ref{def:direct-subline-law}
% and Lemma~\ref{lem:direct-subline-witness-exists}; transport to this statement
% is the open obligation of \cite{gap:qpbt_ld-dimension-divisibility}.
	\notready
	\uses{def:line, def:line-point-dist, def:ith-restricted-line}
	% Statement as in the source (references/qpbt-paper/
	% 14_analysis_of_the_pauli_basis_test.tex, lem:qld-sublines). The
	% line-point distribution over F_q^{2m+2} is built from the index
	% function chi of def:ld-question-distribution at dimension 2m+2, which
	% is defined only when 2m+2 divides q; the source instantiates it
	% without securing this condition. See rem:qld-4-7-divisibility.
	There exists a distribution $D$
	over tuples $(\ell, \ell_X, \ell_Z)$, where $\ell$ is a line in
	$\F_q^{2m+2}$ and $\ell_X, \ell_Z$ are lines in $\F_q^m$, with the
	following properties:
	\begin{enumerate}
		\item The marginal distribution over $\ell$ is the same as the
		marginal of the line-point distribution over $\F_q^{2m+2}$.
		% The source phrases the next property as: the marginal over
		% (ell_X, ell_Z) is "a mixture of the product distributions
		% D_{ALine,i} x D_{ALine,j} and D_{DLine,i} x D_{DLine,j}". Those
		% products are distributions over pairs of line-point pairs, not
		% over pairs of lines, so the phrase does not define the asserted
		% marginal distribution on pairs of lines as written; the
		% formulation below states the mixture at the level of the joint
		% law of (ell_X, ell_Z, u), which is what the proofs of
		% lem:claim-17-2 and lem:claim-17-3 use.
		\item For any point of the form $u = (u_X, u_Z, \alpha, \beta)$
		lying on $\ell$, it holds that $u_X \in \ell_X$ and
		$u_Z \in \ell_Z$. Moreover, the joint distribution of
		$(\ell_X, \ell_Z, u)$, for $(\ell, \ell_X, \ell_Z) \sim D$ and
		$u$ chosen uniformly at random on $\ell$, is a mixture of
		components indexed by a type
		$\mathrm{v} \in \{\mathrm{ALine}, \mathrm{DLine}\}$ and indices
		$i, j \in \{0, \dots, m-1\}$ with the following two properties
		within the component $(\mathrm{v}, i, j)$: the pair
		$(\ell_X, \ell_Z)$ is distributed according to the product of the
		marginal distributions on lines of $D_{\mathrm{v},i}$ and of
		$D_{\mathrm{v},j}$, and, conditioned on $(\ell_X, \ell_Z)$, the
		point $u_X$ is uniformly distributed on $\ell_X$ and likewise the
		point $u_Z$ is uniformly distributed on $\ell_Z$. Equivalently,
		since under each restricted distribution the point is uniformly
		distributed on the line conditionally on the line: within the
		component $(\mathrm{v}, i, j)$, the pair
		$((\ell_X, u_X), (\ell_Z, z'))$, where $z'$ is a uniformly random
		point of $\ell_Z$ sampled independently of $u$, is distributed
		according to $D_{\mathrm{v},i} \times D_{\mathrm{v},j}$, and
		symmetrically the pair $((\ell_X, x'), (\ell_Z, u_Z))$, with $x'$
		a uniformly random point of $\ell_X$ sampled independently of
		$u$, is distributed according to
		$D_{\mathrm{v},i} \times D_{\mathrm{v},j}$. (No assertion is made
		about the joint conditional law of the pair $(u_X, u_Z)$.)
		\item If $\ell$ is axis-parallel, then $\ell_X$ and $\ell_Z$ are
		axis-parallel.
	\end{enumerate}
	(Property~3 is not part of the numbered statement in the source; see
	Remark~\ref{rem:qld-sublines-property-three}.)
\end{lemma}
\begin{proof}
	\uses{def:line-representative, def:ld-question-distribution,
		def:line-point-dist, def:ith-restricted-line, def:line,
		prop:line-equiv}
	The distribution $D$ is described by a sampling procedure. First sample
	$(\ell, u)$ from the line-point distribution over $\F_q^{2m+2}$ and write
	$u = (u_X, u_Z, \alpha, \beta)$, so that Property~1 holds by
	construction; all auxiliary randomness introduced below is sampled
	independently of $(\ell, u)$. As in
	Definition~\ref{def:ith-restricted-line}, write $\bar\chi$ for the
	zero-based shift of the index function at dimension $m$. Its classes
	$\bar\chi^{-1}(i)$ for $i \in \{0, \dots, m-1\}$ partition $\F_q$
	into sets of equal size $q/m$. Write $\bar\chi_{2m+2}$ for the analogous
	zero-based shift at dimension $2m+2$. For the extended-coordinate index in
	the construction below, write $k=r-1$ when the paper writes $r$; the
	source-line indices $i_X,i_Z$ are shifted in the same way. The function
	$\bar\chi_{2m+2}$ is defined only when $2m+2$ divides $q$, and hence so is
	the line-point distribution over $\F_q^{2m+2}$ from which $(\ell, u)$ is
	drawn
	(Definition~\ref{def:ld-question-distribution}); admissibility of
	$(q,m,d)$ supplies only that $m$ divides $q$, and for admissible
	parameters with $m \ge 2$ the required divisibility fails, so this
	first step is an unresolved obligation of the source argument, recorded
	in Remark~\ref{rem:qld-4-7-divisibility} and analyzed in
	\cite{gap:qpbt_ld-dimension-divisibility}; the construction below is
	carried out for the parameters for which the distribution is defined.
	By a
	\emph{fresh diagonal direction pair} we mean a pair $(s', w')$ with
	$s'$ uniformly random in $\F_q$ and $w' \in \F_q^m$ having coordinates
	$0, \dots, \bar\chi(s')-1$ equal to $0$ and the remaining coordinates
	uniformly random; this is exactly the law of the coordinate--direction
	data of the diagonal component of the line-point distribution over
	$\F_q^m$ (Definition~\ref{def:line-point-dist}).

	If $\ell$ is axis-parallel, $\ell = (v, e_k)$ with
	$v = (v_X, v_Z, \alpha_0, \beta_0)$ and $k \in \{0, \dots, 2m+1\}$, set
	$\ell_X = (v'_X, e_{i_X})$ and $\ell_Z = (v'_Z, e_{i_Z})$ where: for
	$0 \le k < m$, $i_X = k$ and $i_Z$ is uniformly random in
	$\{0,\dots,m-1\}$; for $m \le k < 2m$, $i_X$ is uniformly random and
	$i_Z = k - m$; for $2m \le k < 2m+2$, both $i_X$ and $i_Z$ are
	uniformly random; and $v'_X, v'_Z$ are the canonical line
	representatives of $v_X, v_Z$ for the chosen directions
	(Definition~\ref{def:line-representative}), so that $v_X \in \ell_X$
	and $v_Z \in \ell_Z$.

	If $\ell$ is diagonal, $\ell = (v, s, w)$, then by the construction of
	the diagonal component of the line-point distribution
	(Definition~\ref{def:ld-question-distribution} instantiated at
	dimension $2m+2$) the direction $w = (w_X, w_Z, \gamma, \delta)$ has
	coordinates $0, \dots, k-1$ equal to $0$, where
	$k = \bar\chi_{2m+2}(s)$,
	and its remaining coordinates are uniformly random. These remaining
	coordinates may also vanish, so $w$ need not have a nonzero coordinate
	at index $k$ and may be the zero vector, in which case $\ell$ is a
	singleton (Definition~\ref{def:line}).
	% In its one-based notation, the source proof describes w as "a direction
	% vector whose first nonzero coordinate is j = chi(s)"
	% (references/qpbt-paper/14_analysis_of_the_pauli_basis_test.tex, proof
	% of lem:qld-sublines); the distribution delivers only the weaker
	% property recorded here, which is the one the argument uses. The
	% source also leaves the laws of s_X and s_Z unspecified ("chosen such
	% that chi(s_X) = j" in its one-based notation); the uniform choices
	% below are needed for the exact marginals claimed in Property 2.
	Sample $\ell_X = (v'_X, s_X, w'_X)$ and $\ell_Z = (v'_Z, s_Z, w'_Z)$
	where: for $0 \le k < m$, $w'_X = w_X$ and $w'_Z = w_Z$, with $s_X$
	uniformly random in the class $\bar\chi^{-1}(k)$ and $s_Z$ uniformly
	random in the class $\bar\chi^{-1}(0)$; for $m \le k < 2m$, the pair
	$(s_X, w'_X)$ is a fresh diagonal direction pair and $w'_Z = w_Z$,
	with $s_Z$ uniformly random in the class $\bar\chi^{-1}(k-m)$; for
	$2m \le k < 2m+2$, both $(s_X, w'_X)$ and $(s_Z, w'_Z)$ are fresh
	diagonal direction pairs, sampled independently; and $v'_X, v'_Z$ are
	the canonical representatives of $v_X, v_Z$ for the directions
	$w'_X, w'_Z$.

	Property~3 is immediate from the procedure: when $\ell$ is
	axis-parallel, both $\ell_X$ and $\ell_Z$ are constructed to be
	axis-parallel. The containments in Property~2 hold for every point
	$u^\dagger = (u^\dagger_X, u^\dagger_Z, \alpha^\dagger,
	\beta^\dagger)$ of $\ell$: the block $u^\dagger_X$ differs from $v_X$
	by a multiple of the direction of $\ell_X$ when that direction is
	inherited from $\ell$ (namely $e_k$ for axis-parallel $\ell$ with
	$0 \le k < m$, and $w_X$ for diagonal $\ell$ with $0 \le k < m$),
	and equals $v_X$ in all other cases, since then the corresponding
	block of the direction of $\ell$ vanishes; in either case
	$u^\dagger_X \in \ell_X$ because $v_X \in \ell_X$, and symmetrically
	$u^\dagger_Z \in \ell_Z$.

	For the distributional part of Property~2, observe first that
	$(\ell_X, \ell_Z)$ is a function of $\ell$ and of the auxiliary
	randomness alone: the directions used are either read off from the
	direction data of $\ell$ or freshly sampled, and by
	Proposition~\ref{prop:line-equiv} the canonical representatives
	$v'_X, v'_Z$ do not depend on which base point of $\ell$ is used to
	compute them. Under both components of the line-point distribution the
	point $u$ is uniformly distributed on $\F_q^{2m+2}$ and, conditioned
	on the line $\ell$, uniformly distributed on the points of $\ell$
	(each of the $q$ values of the line parameter, or the single point
	when the direction vanishes, being equally likely). Hence, conditioned
	on $(\ell, \ell_X, \ell_Z)$, the sampled point $u$ is uniformly
	distributed on $\ell$, and it suffices to establish the claimed
	mixture description for the tuple $(\ell_X, \ell_Z, u)$ produced by
	the procedure itself.

	Condition on the component of the line-point distribution
	(axis-parallel or diagonal) and on $k$; the mixture of Property~2 is
	over these choices together with the fresh index choices, and the
	blocks $u_X, u_Z$ of $u$ are independent and uniform on $\F_q^m$,
	independent of all direction data. In the axis-parallel case with
	$0 \le k < m$, the pair $(\ell_X, u_X)$ consists of the uniformly
	random point $u_X$ together with the canonical axis-parallel line
	through it in direction $e_k$, so
	$(\ell_X, u_X) \sim D_{\mathrm{ALine},k}$, while, conditioned
	additionally on $i_Z$, the pair $(\ell_Z, u_Z)$ consists of the
	uniformly random point $u_Z$ together with the canonical line through
	it in direction $e_{i_Z}$, so
	$(\ell_Z, u_Z) \sim D_{\mathrm{ALine},i_Z}$; the two pairs are
	functions of the disjoint independent data $u_X$ and $(u_Z, i_Z)$,
	hence independent. In the diagonal case with $0 \le k < m$, the
	triple $(u_X, w_X, s_X)$ has $u_X$ uniform, $w_X$ with coordinates
	$0, \dots, k-1$ equal to $0$ and the remaining coordinates uniform,
	and $s_X$ uniform in $\bar\chi^{-1}(k)$, which is exactly the law
	generating $D_{\mathrm{DLine},k}$; hence
	$(\ell_X, u_X) \sim D_{\mathrm{DLine},k}$, and similarly
	$(\ell_Z, u_Z) \sim D_{\mathrm{DLine},0}$ since $w_Z$ is uniform and
	$s_Z$ is uniform in $\bar\chi^{-1}(0)$; the two pairs are functions of
	the disjoint independent data $(u_X, w_X, s_X)$ and $(u_Z, w_Z, s_Z)$,
	hence independent. The remaining cases are analogous: each of
	$(\ell_X, u_X)$ and $(\ell_Z, u_Z)$ is either built from the inherited
	direction block as above, or from a fresh axis index or fresh diagonal
	direction pair, in which case it is distributed as the uniform mixture
	over $i \in \{0, \dots, m-1\}$ of $D_{\mathrm{v},i}$ for the
	corresponding type $\mathrm{v}$; and the two pairs are always functions of
	disjoint independent collections of the variables $u_X$, $u_Z$, the
	direction blocks, and the fresh auxiliary choices, hence independent
	conditioned on the component. Within each component this yields exactly the
	description of Property~2: the pair $(\ell_X, \ell_Z)$ is
	product-distributed according to the line marginals of
	$D_{\mathrm{v},i}$ and $D_{\mathrm{v},j}$, and, conditioned on
	$(\ell_X, \ell_Z)$, each of $u_X$ and $u_Z$ is uniformly distributed
	on its line, because under each restricted distribution the point is
	uniform on the line conditionally on the line. The equivalent product
	formulations follow from the same conditional uniformity.
\end{proof}

\begin{remark}[Property~3 of the sub-line distribution]
	\label{rem:qld-sublines-property-three}
	The source proof of Lemma~\ref{lem:qld-4-13} in \cite{Ji2020MIPStar}
	appeals to a ``Property~3'' of the sub-line lemma, while the statement
	given there lists only two numbered properties, so the reference is
	dangling. The intended assertion is the axis-parallel closure property
	recorded as Property~3 of Lemma~\ref{lem:qld-sublines} above, which is
	immediate from the sampling procedure in the proof.
\end{remark}

\begin{lemma}[Auxiliary sub-line distribution with directly sampled indices]
	\label{lem:qld-sublines-direct}
	\lean{MIPStarRE.QPBT.SubLineWitness, MIPStarRE.QPBT.exists_subLineWitness}
	\leanok
	\uses{def:line, def:line-point-dist, def:ith-restricted-line,
		def:combine-map}
	This formalization-only variant of Lemma~\ref{lem:qld-sublines} uses
	uniformly sampled coordinate indices on the extended space. Throughout
	this lemma, coordinate $i\in\{1,\ldots,n\}$ corresponds to Lean index
	$i-1$ in \texttt{Fin n}, for $n=m$ or $n=2m+2$. Thus the unit vector
	$e_i$ is represented by \texttt{coordinateDirection} at index $i-1$,
	and vanishing coordinates $1,\ldots,i-1$ means vanishing at Lean indices
	$k<i-1$. The restricted law $D_{\mathrm{v},i}$ likewise uses Lean
	index $i-1$ in \texttt{Fin m}. For every
	admissible $(q,m,d)$, there is a probability distribution $D_{\rm dir}$
	on triples $(\ell,\ell_X,\ell_Z)$ with the following properties:
	\begin{enumerate}
		\item Its extended-line marginal is obtained by choosing the line
		type with equal probabilities and its coordinate index uniformly in
		$\{1,\ldots,2m+2\}$. In the axis case the direction is the
		corresponding unit vector; in the diagonal case preceding coordinates
		vanish and the remaining coordinates are uniform. In both cases the
		base is the canonical representative of an independent uniform point.
		\item Every point $u=(u_X,u_Z,\alpha,\beta)$ on a supported extended
		line satisfies $u_X\in\ell_X$ and $u_Z\in\ell_Z$. The affine
		parameterizations are compatible: there are scalars $a_X,b_X,a_Z,b_Z$
		such that, writing $\ell=(v,w)$ and
		$\ell_W=(v_W,w_W)$ for $W\in\{X,Z\}$,
		$(v+tw)_W=v_W+(a_W+b_Wt)w_W$ for every $t\in\F_q$.
		The $\alpha$ and $\beta$ coordinates are affine functions of $t$,
		as required by Definition~\ref{def:combine-map}.
		\item For a uniform affine parameter on $\ell$, the two laws of
		$(\ell_X,\ell_Z,u_X)$ and $(\ell_X,\ell_Z,u_Z)$ are separate
		mixtures with the same probability weights over
		$(\mathrm{v},i,j)\in\{\mathrm{ALine},\mathrm{DLine}\}
		\times\{1,\ldots,m\}^2$. Their components are respectively
		$D_{\mathrm{v},i}$ times the line marginal of $D_{\mathrm{v},j}$,
		and the line marginal of $D_{\mathrm{v},i}$ times
		$D_{\mathrm{v},j}$. No joint conditional law of $(u_X,u_Z)$ is asserted.
		\item If $\ell$ is axis-parallel, both projected lines are axis-parallel.
	\end{enumerate}
	Identification with the source's seed-indexed extended-line law remains
	open; see \cite{gap:qpbt_subline-claims-line-marginal} and
	\cite{gap:qpbt_ld-dimension-divisibility}.
\end{lemma}
\begin{proof}
	\leanok
	\uses{def:line-representative, def:ith-restricted-line, def:combine-map}
	Apply the subline sampling construction with the extended coordinate
	index chosen directly. Its sampling law has mass one and the prescribed
	extended-line marginal. The inherited direction blocks and the fresh
	auxiliary directions give incidence, affine compatibility, and axis-line
	closure. Separating the independent X and Z sampling data within each
	index case gives the two marginal mixtures with common weights.
\end{proof}

The proof of Lemma~\ref{lem:qld-4-13} rests on three scalar estimates, stated
next. In all three, $D$ is the distribution
of Lemma~\ref{lem:qld-sublines},
$T^{\ell_X,\ell_Z}$ is the measurement of Lemma~\ref{lem:qld-xz-lines} placed
on registers $\mathsf{A}\mathsf{A}'$, the second tensor factor is placed on
registers $\mathsf{B}\mathsf{A}''$, and the sums range over
$(f_X, f_Z) \in \deg_{md}(\ell_X) \times \deg_{md}(\ell_Z)$.

\begin{lemma}[Formalization support: ordered products on one placement]
	\label{lem:qld-4-10-same-placement}
	\lean{MIPStarRE.QPBT.Placement.exists_isOpposite}
	\lean{MIPStarRE.QPBT.CombinedPointsWitness.orderedZX_dist_le}
	\lean{MIPStarRE.QPBT.CombinedPointsWitness.orderedXZ_dist_le}
	\leanok
	\uses{lem:qld-4-10, fact:triangle, def:povm-distance}
	On one placement, the joint point measurement has squared distance at most
	$4\delta_Q$ from each of the ordered $Z$-then-$X$ and $X$-then-$Z$
	products.
\end{lemma}
\begin{proof}
	\leanok
	\uses{lem:qld-4-10, fact:triangle, def:symmetric-equivalents}
	Insert the joint measurement on the opposite placement and apply the
	squared-distance triangle inequality to the corresponding two
	$\delta_Q$ bounds.
\end{proof}

\begin{lemma}[Formalization support: replacing one factor of a real overlap]
	\label{lem:overlap-gap-distance}
	\lean{MIPStarRE.QPBT.abs_overlap_gap_le_sqrt_of_opFamilyDistSq}
	\leanok
	\uses{def:povm-distance, def:expanded-state}
	For a probability distribution and a unit state, the difference of the
	real parts of overlaps weighted by a complete measurement is at most
	the square root of the state-dependent squared distance of the two
	operator families.
\end{lemma}
\begin{proof}
	\leanok
	\uses{def:povm-distance, def:expanded-state, def:povm-conventions}
	Apply Cauchy--Schwarz to the outcome sum and then to the question average;
	completeness bounds the squared norm of the left factor by one.
\end{proof}

\begin{theorem}[Formalization support: replacing one factor of a complex overlap]
	\label{thm:complex-overlap-gap-distance}
	\lean{MIPStarRE.QPBT.norm_overlap_gap_le_sqrt_of_opFamilyDistSq}
	\leanok
	\uses{def:povm-distance, def:povm-conventions}
	This formalization-only auxiliary isolates the complex Cauchy--Schwarz
	estimate used in Lemma~\ref{lem:claim-17-1} of \cite{Ji2020MIPStar}.
	Let $\mu$ be a finitely supported probability distribution on $\mathcal X$,
	let $\psi$ be a unit vector in a finite-dimensional complex Hilbert space,
	and let $\{E^x_a\}_{a\in\mathcal O}$ be a complete POVM for each $x$,
	with finite outcome set $\mathcal O$. Let $B^x_a,C^x_a$ be arbitrary
	operators on the same space. If
	\[
		\E_{x\sim\mu}\sum_a\|(B^x_a-C^x_a)\psi\|^2\le\zeta,
	\]
	then
	\[
		\left|\E_{x\sim\mu}\sum_a\langle\psi,E^x_aB^x_a\psi\rangle
		-\E_{x\sim\mu}\sum_a\langle\psi,E^x_aC^x_a\psi\rangle\right|
		\le\sqrt\zeta.
	\]
	The absolute value is the complex modulus. Neither operator family is
	assumed Hermitian or commuting with the POVM.
\end{theorem}
\begin{proof}
	\leanok
	\uses{def:povm-distance, def:povm-conventions}
	Apply weighted complex Cauchy--Schwarz to $\sqrt{E^x_a}\psi$ and
	$\sqrt{E^x_a}(B^x_a-C^x_a)\psi$. Completeness and normalization make
	the first squared factor one. Since $0\le E^x_a\le I$, the second
	is at most the displayed squared distance, hence at most $\zeta$.
	In particular, the squared-norm operator is
	$(B^x_a-C^x_a)^\dagger(B^x_a-C^x_a)$.
\end{proof}

\begin{lemma}[Formalization support: axis lines have nonzero direction]
	\label{lem:axis-line-zero-direction-mass}
	\lean{MIPStarRE.QPBT.aLinePointDist_zero_direction_mass}
	\leanok
	\uses{def:line-point-dist, def:line}
	Under the axis-parallel line-point distribution of
	Definition~\ref{def:line-point-dist}, the indicator of the event that the
	sampled line has direction $0$ has expectation $0$.
\end{lemma}
\begin{proof}
	\leanok
	\uses{def:line-point-dist, def:line}
	An axis-parallel sample has direction $e_i$ for the sampled coordinate
	$i$, and the $i$-th entry of $e_i$ is $1 \ne 0$. The indicator therefore
	vanishes at every sample, hence so does its expectation.
\end{proof}

\begin{lemma}[Formalization support: the pointwise consistency defect is at most one]
	\label{lem:consistency-defect-integrand-le-one}
	\lean{MIPStarRE.QPBT.consistencyDefect_integrand_le_one}
	\leanok
	\uses{def:expanded-state, def:consistency}
	Let $\{A_a\}_{a \in \mathcal{O}}$ and $\{B_a\}_{a \in \mathcal{O}}$ be
	complete measurements with the same finite outcome set, placed on two
	opposite registers of the expanded state $\ket{\hat\psi}$ of
	Definition~\ref{def:expanded-state}. Then
	\[
		\sum_{a \ne b} \bra{\hat\psi} A_a B_b \ket{\hat\psi} \,\le\, 1\,.
	\]
\end{lemma}
\begin{proof}
	\leanok
	\uses{lem:distance-point-defect, lem:distance-qform-nonnegative,
		def:expanded-state}
	By Lemma~\ref{lem:distance-point-defect} the left-hand side equals
	$\|\hat\psi\|^2-\sum_a \bra{\hat\psi} A_a B_a \ket{\hat\psi}$, and
	$\|\hat\psi\|=1$. On opposite placements the two effects commute and are
	positive semidefinite, so each diagonal term is nonnegative and the
	subtracted sum is nonnegative.
\end{proof}

\begin{definition}[Directly indexed sub-line laws]
	\label{def:direct-subline-law}
	\lean{MIPStarRE.QPBT.SubLineWitness}
	\leanok
	\uses{def:ith-restricted-line, def:line-representative, def:combine-map}
	Put $n=2m+2$. The directly indexed extended line-point law samples
	$i\in\{0,\dots,n-1\}$ and $u,v\in\F_q^n$ independently and uniformly,
	and chooses the axis or diagonal kind with equal probabilities.
	For the axis kind set $w=e_i$; for the diagonal kind set $w_j=0$ for
	$j<i$ and $w_j=v_j$ for $j\ge i$. The line has base $L_w^{\mathrm{lnf}}(u)$
	and direction $w$, retaining its kind and coordinate index; the sampled
	point is $u$. Zero diagonal directions are included.

	A directly indexed sub-line law $D_{\mathrm{dir}}$ is a probability law
	on triples $(\ell,\ell_X,\ell_Z)$ with the following properties.
	Its extended-line marginal is the line marginal just defined.
	Every supported triple admits affine maps $t\mapsto a_W+b_Wt$ such that
	\[
		\pi_W(\ell(t))=\ell_W(a_W+b_Wt),\qquad W\in\{X,Z\},
	\]
	where $\ell(t)$ denotes base plus parameter times direction and the
	field identifications are understood. Thus each extended point projects
	onto its two source lines, with affine scalar coordinates $\alpha(t)$
	and $\beta(t)$ as in Definition~\ref{def:combine-map}. An axis-parallel
	extended line has axis-parallel source lines.

	Finally there is one probability law on
	$(\mathrm v,i,j)\in\{\mathrm{ALine},\mathrm{DLine}\}
	\times\{0,\dots,m-1\}^2$ with the following two mixture properties.
	Sample a triple from $D_{\mathrm{dir}}$ and an independent uniform
	parameter $t$ on $\ell$. The law of $(\ell_X,\ell_Z,\pi_X(\ell(t)))$
	is the mixture obtained by sampling $(\ell_X,x)$ from
	$\widetilde D_{\mathrm v,i}$ and independently $\ell_Z$ from the line
	marginal of $\widetilde D_{\mathrm v,j}$, and retaining
	$(\ell_X,\ell_Z,x)$.
	The law of $(\ell_X,\ell_Z,\pi_Z(\ell(t)))$ is the mixture obtained
	by sampling $\ell_X$ from the line marginal of
	$\widetilde D_{\mathrm v,i}$ and independently $(\ell_Z,z)$ from
	$\widetilde D_{\mathrm v,j}$, and retaining $(\ell_X,\ell_Z,z)$.
	These are separate one-point marginal identities on the seed-bearing
	projected lines, with the same mixture weights; they impose no joint
	conditional independence of the two projected points. The corresponding
	source marginal identities hold only after applying the forgetful maps
	of Definition~\ref{def:line-point-dist} to both projected lines.
	Identification with the source law $D$ is a separate obligation;
	see \cite{gap:qpbt_ld-dimension-divisibility}.
\end{definition}

\begin{lemma}[Formalization support: construction of the directly indexed sub-line law]
	\label{lem:direct-subline-law-construction}
	\lean{MIPStarRE.QPBT.subLineDist}
	\lean{MIPStarRE.QPBT.subLineDist_isProbability}
	\lean{MIPStarRE.QPBT.subLineDist_map_fst}
	\lean{MIPStarRE.QPBT.subLineDist_map_subLineXProjection}
	\lean{MIPStarRE.QPBT.subLineDist_map_subLineZProjection}
	\lean{MIPStarRE.QPBT.subLineDist_source_mixture}
	\lean{MIPStarRE.QPBT.exists_samplingData_of_mem_subLineDist_support}
	\lean{MIPStarRE.QPBT.subLineTripleOf_incidence}
	\lean{MIPStarRE.QPBT.subLineTripleOf_compatibility}
	\lean{MIPStarRE.QPBT.subLineTripleOf_axis_closure}
	\leanok
	\uses{def:direct-subline-law}
	Fix an admissible parameter tuple. The concrete law $D_{\mathrm{dir}}$ is
	the equal mixture of the axis and
	diagonal sampling branches in Definition~\ref{def:direct-subline-law}. It is
	a probability law, and its extended-line marginal is the directly indexed
	extended-line marginal. Its projected $X$- and $Z$-point marginals are the
	two mixtures in that definition, with one common probability law on the line
	kind and the two source coordinate indices.

	The extended coordinate, point, and direction are uniform. Conditional on
	the extended coordinate, the two source seeds are independent and uniform
	on their prescribed coordinate-index fibers; these fibers have positive
	mass. Every triple in the support is produced by such data from one of the
	two branches. For every point on its extended line, such a triple satisfies
	incidence of the two projected points with their source lines, the affine
	compatibility relations for both projections, and closure of both source lines under an
	axis-parallel extended line.
	The projected laws use an independent uniform affine parameter and are
	separate one-point marginals, not a joint conditional independence claim.
	Zero diagonal directions are retained, and the scalar fields of the direct
	and source coordinates are related by the fixed field identifications.
\end{lemma}
\begin{proof}
	\leanok
	\uses{def:direct-subline-law}
	Normalization follows from the uniform branch law and the normalized seed
	laws. The representative--parameter identities give the extended and
	projected marginals. Recovering the auxiliary sample from a supported triple
	gives the incidence, compatibility, and axis-closure conclusions.
\end{proof}

\begin{lemma}[Existence of a directly indexed sub-line witness]
	\label{lem:direct-subline-witness-exists}
	\lean{MIPStarRE.QPBT.exists_subLineWitness}
	\leanok
	\uses{def:direct-subline-law}
	This formalization-only auxiliary is not the source sub-line lemma. For every
	admissible parameter tuple, there exists a probability law on
	triples $(\ell,\ell_X,\ell_Z)$ satisfying all the incidence, affine
	compatibility, separate seed-bearing projected-marginal, and axis-closure
	properties of Definition~\ref{def:direct-subline-law}. The source
	projected-marginal identities require the componentwise forgetful maps
	specified there.
\end{lemma}
\begin{proof}
	\leanok
	\uses{lem:direct-subline-law-construction}
	Apply Lemma~\ref{lem:direct-subline-law-construction} and collect its
	normalization, marginal, incidence, compatibility, and axis-closure
	properties into the directly indexed sub-line witness.
\end{proof}

\begin{lemma}[Formalization support: a uniform point under the directly indexed law]
	\label{lem:qld-sublines-uniform-point}
	\lean{MIPStarRE.QPBT.uniformDistribution_map_uncurry}
	\lean{MIPStarRE.QPBT.uniformDistribution_map_directLine_add_smul}
	\lean{MIPStarRE.QPBT.avgOver_uniform_directLdSpace_line_add_smul}
	\lean{MIPStarRE.QPBT.avgOver_directLinePointDist_line_add_smul}
	\lean{MIPStarRE.QPBT.SubLineWitness.avgOver_projX_projZ}
	\leanok
	\uses{def:direct-subline-law, def:line-representative,
		def:combine-map}
	For any $D_{\mathrm{dir}}$ as in Definition~\ref{def:direct-subline-law},
	a uniform affine parameter on the sampled extended line gives a uniform
	point of $\F_q^{2m+2}$, whose $x$ and $z$ blocks are independent uniform
	points of $\F_q^m$. This statement does not transport that law to the
	source distribution $D$; see \cite{gap:qpbt_ld-dimension-divisibility}.
\end{lemma}
\begin{proof}
	\leanok
	\uses{def:direct-subline-law, def:line-representative, def:combine-map}
	For fixed direction $v$, the map $(u,t)\mapsto L_v(u)+tv$ on
	ambient-point--parameter pairs has fibers of size $q$. If $v\ne0$,
	$u$ varies over a coset, while its canonical representative and $t$ are
	fixed; if $v=0$, $u$ is fixed and $t$ is free. The coordinate split and
	scalar identification are bijections.
\end{proof}

\begin{lemma}[Formalization support: real-part replacement under the directly indexed law]
	\label{lem:claim-17-1-re-direct}
	\lean{MIPStarRE.QPBT.subline_replace_by_ordered_product_re_direct}
	\leanok
	\uses{lem:overlap-gap-distance, lem:qld-4-10,
		lem:qld-4-10-same-placement, lem:qld-sublines-uniform-point,
		lem:qld-xz-lines, def:direct-subline-law}
	There is a universal $C>0$ such that the following holds.
	Let $D_{\mathrm{dir}}$ satisfy Definition~\ref{def:direct-subline-law},
	and let $Q,T$ satisfy the joint point and line measurement properties of
	Lemmas~\ref{lem:qld-4-10} and~\ref{lem:qld-xz-lines} with errors
	$\delta_Q,\delta_P$. Write
	$A_{\mathrm{dir}},B_{\mathrm{dir}}$ for the two expectations displayed in
	Lemma~\ref{lem:claim-17-1}, using $D_{\mathrm{dir}}$ and partial
	evaluation in $\F_q \cup \{\bot\}$, with zero point effects for undefined
	outcomes. Then
	$|\Re A_{\mathrm{dir}}-\Re B_{\mathrm{dir}}|\le C\sqrt{\delta_Q}$.
	This formalization-only estimate asserts neither a complex-modulus bound
	nor transport to the source distribution.
\end{lemma}
\begin{proof}
	\leanok
	\uses{lem:overlap-gap-distance, lem:qld-4-10-same-placement,
		lem:qld-sublines-uniform-point}
	The real overlap estimate and the same-placement squared-distance bound
	$4\delta_Q$ give $2\sqrt{\delta_Q}$. Use the uniform-point identity to
	rewrite the point average, and take $C=2$.
\end{proof}

\begin{theorem}[Formalization support: complex replacement under the directly indexed law]
	\label{thm:claim-17-1-direct}
	\lean{MIPStarRE.QPBT.subline_replace_by_ordered_product_direct}
	\leanok
	\uses{lem:qld-4-10, def:concrete-paired-line-measurement,
		def:direct-subline-law, def:expanded-point-measurement}
	This is an auxiliary version of Lemma~\ref{lem:claim-17-1} for completed
	evaluations and the directly indexed law. Fix admissible parameters and
	a projective setting with error $\eps$. Let $Q$ be a projective joint
	point family satisfying all three squared-distance bounds of
	Lemma~\ref{lem:qld-4-10}, each with upper bound $\delta_Q$, for every
	directed opposite placement. Use the concrete paired-line POVM
	\[
		T^{\ell_X,\ell_Z}_{f_X,f_Z}
		=\hat M^{(\mathrm{Line},X),\ell_X}_{f_X}
		 \hat M^{(\mathrm{Line},Z),\ell_Z}_{f_Z}
		 \hat M^{(\mathrm{Line},X),\ell_X}_{f_X}.
	\]
	Sample $(\ell,\ell_X,\ell_Z)$ from any $D_{\mathrm{dir}}$ of
	Definition~\ref{def:direct-subline-law}, then an independent uniform
	affine parameter $t\in\F_q$, and write
	$\ell(t)=(x,z,\alpha,\beta)$ using the fixed field identifications.
	Denote this expectation by $\E_{\mathrm{dir}}$. For each pair of
	polynomial representatives $f_X,f_Z$ of degree at most $md$, put
	$a=\mathrm{eval}_x(f_X)$ and $b=\mathrm{eval}_z(f_Z)$ in
	$\F_q\cup\{\bot\}$. Extend the joint and individual point measurements
	by zero effects at every answer containing $\bot$. Set
	\begin{align*}
		J_{\mathrm{dir}}&=\E_{\mathrm{dir}}\sum_{f_X,f_Z}
		\bra{\hat\psi}(T^{\ell_X,\ell_Z}_{f_X,f_Z})_{\mathsf{A}\mathsf{A}'}
		(\hat Q^{x,z}_{a,b})_{\mathsf{B}\mathsf{A}''}\ket{\hat\psi},\\
		A_{\mathrm{dir}}&=\E_{\mathrm{dir}}\sum_{f_X,f_Z}
		\bra{\hat\psi}(T^{\ell_X,\ell_Z}_{f_X,f_Z})_{\mathsf{A}\mathsf{A}'}
		(\hat M^{(\mathrm{Point},Z),z}_b
		 \hat M^{(\mathrm{Point},X),x}_a)_{\mathsf{B}\mathsf{A}''}
		\ket{\hat\psi}.
	\end{align*}
	Then $|J_{\mathrm{dir}}-A_{\mathrm{dir}}|\le2\sqrt{\delta_Q}$ in
	complex modulus. No paired-line consistency hypothesis is required.
	The source-law identification remains separate; see
	\cite{gap:qpbt_subline-claims-line-marginal}.
\end{theorem}
\begin{proof}
	\leanok
	\uses{thm:complex-overlap-gap-distance, lem:qld-4-10-same-placement,
		lem:qld-sublines-uniform-point}
	Regroup the complete line POVM by its two completed evaluations and apply
	Theorem~\ref{thm:complex-overlap-gap-distance}. Undefined point effects
	vanish. The uniform-point identity and the same-placement estimate bound
	the averaged squared distance by $4\delta_Q$, giving $2\sqrt{\delta_Q}$.
\end{proof}

\begin{lemma}[Formalization support: the concrete X marginal]
	\label{lem:concrete-line-X-marginal}
	\lean{MIPStarRE.QPBT.ProjectiveSetting.combinedLineMeasurement_sum_Z}
	\leanok
	\uses{def:concrete-paired-line-measurement, def:expanded-line-measurement}
	For the concrete X-Z-X measurement,
	$\sum_{f_Z}T^{\ell_X,\ell_Z}_{f_X,f_Z}
	=\hat M^{(\mathrm{Line},X),\ell_X}_{f_X}$.
\end{lemma}
\begin{proof}
	\leanok
	\uses{def:concrete-paired-line-measurement, def:expanded-line-measurement}
	Sum the middle factor to the identity and use projectivity of the X factor.
\end{proof}

\begin{lemma}[Formalization support: the concrete X deficit under the directly indexed law]
	\label{lem:concrete-X-deficit-direct}
	\lean{MIPStarRE.QPBT.concreteXPointOverlap}
	\lean{MIPStarRE.QPBT.exists_concreteXPointOverlap_deficit_le}
	\leanok
	\uses{def:concrete-paired-line-measurement, lem:qld-comm-line-cons,
		def:direct-subline-law}
	Let $D_{\mathrm{dir}}$ satisfy Definition~\ref{def:direct-subline-law}.
	For the concrete X-Z-X measurement, let
	$A_X(\ell_X,\ell_Z,x)$ be the real overlap obtained by summing
	$T_{f_X,f_Z}^{\ell_X,\ell_Z}\ot
	\hat M^{(\mathrm{Point},X),x}_{f_X(x)}$ over both outcomes, using
	partial evaluation in $\F_q \cup \{\bot\}$, with zero point effects for
	undefined outcomes. There is a universal $C>0$ such that
	$1-\E_{D_{\mathrm{dir}}}\E_{(x,z,\alpha,\beta)\in\ell} A_X
	\le C m^2\delta_{\mathrm{Line}}(\eps)$.
\end{lemma}
\begin{proof}
	\leanok
	\uses{lem:concrete-line-X-marginal, lem:restricted-line-mixture-bounds,
		lem:qld-comm-line-cons, def:direct-subline-law}
	The concrete marginal identity reduces the overlap to X-line/point
	consistency. The separate X marginal of the directly indexed subline law
	is a mixture of restricted laws. Error inflation and $m\ge1$ give the bound.
\end{proof}

\begin{lemma}[Formalization support: complex X-factor removal under the directly indexed law]
	\label{lem:claim-17-2-direct}
	\lean{MIPStarRE.QPBT.subline_remove_X_factor_direct}
	\leanok
	\uses{lem:concrete-X-deficit-direct}
	For the concrete X-Z-X measurement and the directly indexed law
	$D_{\mathrm{dir}}$, with partial evaluation in $\F_q \cup \{\bot\}$ and
	zero point effects for undefined outcomes, the complex
	modulus of the difference in Lemma~\ref{lem:claim-17-2} is at most
	$C m\sqrt{\delta_{\mathrm{Line}}(\eps)}$ for a universal $C>0$.
	Transport to the source law is separate; see
	\cite{gap:qpbt_subline-claims-line-marginal,gap:qpbt_ld-dimension-divisibility}.
\end{lemma}
\begin{proof}
	\leanok
	\uses{lem:concrete-X-deficit-direct}
	Complex weighted Cauchy--Schwarz bounds the difference by the square root
	of the X-overlap deficit. The deficit estimate gives $K m^2
	\delta_{\mathrm{Line}}(\eps)$, so $C=\sqrt K$ suffices.
\end{proof}

\begin{lemma}[Formalization support: the real Z overlap under the directly indexed law]
	\label{lem:claim-17-3-re-direct}
	\lean{MIPStarRE.QPBT.subline_Z_term_near_one_re_direct}
	\leanok
	\uses{lem:restricted-line-mixture-bounds, lem:qld-4-10,
		lem:qld-xz-lines, lem:qld-4-10-same-placement,
		lem:overlap-gap-distance, def:direct-subline-law}
	There is a universal $C>0$ such that the following holds.
	Let $Q,T$ satisfy the joint point and line measurement properties of
	Lemmas~\ref{lem:qld-4-10} and~\ref{lem:qld-xz-lines} with errors
	$\delta_Q,\delta_P$, and let $D_{\mathrm{dir}}$ be a directly indexed
	subline law as in Definition~\ref{def:direct-subline-law}. If $B_Z$ is the
	expectation in Lemma~\ref{lem:claim-17-3} under that law, with partial
	evaluation in $\F_q \cup \{\bot\}$ and zero point effects for undefined
	outcomes, then
	$|\Re B_Z-1|\le
	C\sqrt m(\delta_P^{1/4}+\delta_Q^{1/4}+\eps^{1/4})$.
\end{lemma}
\begin{proof}
	\leanok
	\uses{lem:restricted-line-mixture-bounds, lem:qld-4-10,
		lem:qld-4-10-same-placement, lem:overlap-gap-distance, def:direct-subline-law,
		def:expanded-point-measurement}
	The separate Z marginal and restricted consistency estimates bound the
	nonnegative overlap deficit; taking a square root gives the displayed
	bound with $C=2$.
\end{proof}

\begin{theorem}[Formalization support: reality of the directly indexed Z overlap]
	\label{thm:subline-z-overlap-real-direct}
	\lean{MIPStarRE.QPBT.subline_Z_overlap_eq_real_direct}
	\leanok
	\uses{def:direct-subline-law, def:expanded-point-measurement,
		def:expanded-state, def:povm-conventions}
	This formalization-only auxiliary supplies the reality identity for the
	directly indexed analogue of Lemma~\ref{lem:claim-17-3}. Fix admissible
	parameters, a projective setting with error $\eps$, and a law
	$D_{\mathrm{dir}}$ as in Definition~\ref{def:direct-subline-law}.
	Let $T^{\ell_X,\ell_Z}$ be any complete POVM on Alice's expanded local
	space with outcomes $(f_X,f_Z)$ consisting of polynomial representatives
	of degree at most $md$. Sample a subline triple and an independent uniform
	affine parameter $t$, and let $z=\pi_Z(\ell(t))$, with the scalar fields
	identified as above. Define
	\[
		B_{\mathrm{dir}}(T)=
		\E_{(\ell,\ell_X,\ell_Z)\sim D_{\mathrm{dir}}}\E_{t\in\F_q}
		\sum_{f_X,f_Z}\bra{\hat\psi}
		(T^{\ell_X,\ell_Z}_{f_X,f_Z})_{\mathsf{A}\mathsf{A}'}
		(\hat M^{(\mathrm{Point},Z),z}_{\mathrm{eval}_z(f_Z)})_{
		\mathsf{B}\mathsf{A}''}\ket{\hat\psi},
	\]
	where the evaluation takes values in $\F_q\cup\{\bot\}$ and the point
	effect at $\bot$ is zero. Then
	$B_{\mathrm{dir}}(T)=\operatorname{Re}B_{\mathrm{dir}}(T)$, with the
	real number on the right embedded in $\C$. No projectivity or consistency
	of $T$ is assumed.
\end{theorem}
\begin{proof}
	\leanok
	\uses{def:expanded-state, def:expanded-point-measurement,
		def:povm-conventions}
	Each line effect and the corresponding completed point effect are positive
	and act on opposite registers. Their product is positive semidefinite, so
	its expectation is real. Finite sums and averages with real weights
	preserve this identity.
\end{proof}

\begin{theorem}[Formalization support: the complex Z bound with line consistency]
	\label{thm:claim-17-3-direct}
	\lean{MIPStarRE.QPBT.subline_Z_term_near_one_direct}
	\leanok
	\uses{lem:qld-4-10, thm:subline-z-overlap-real-direct,
		def:line-point-dist, def:consistency, def:symmetric-equivalents}
	This auxiliary estimate assumes joint point and paired-line consistency
	data explicitly. There is a universal $C>0$ with the following property.
	Fix admissible parameters, a projective setting with error $\eps$, and
	any law $D_{\mathrm{dir}}$ of Definition~\ref{def:direct-subline-law}.
	Let $Q$ satisfy the projectivity and three squared-distance bounds of
	Lemma~\ref{lem:qld-4-10}, each at most $\delta_Q$ on every directed
	opposite placement. Let $T$ be a paired-line POVM on each player's
	expanded local space, with polynomial representatives of degree at most
	$md$ as outcomes. Assume that effects vanish when an axis-parallel
	$\ell_W$ has an outcome of degree greater than $d$, and that its completed
	evaluations have consistency defect at most $\delta_P$ with the completed joint point
	family, averaged over two independent samples from the refined line-point
	law of Definition~\ref{def:line-point-dist}, on every directed opposite
	placement. Then, for Alice's member of this family,
	\[
		|B_{\mathrm{dir}}(T)-1|\le
		C\sqrt m\bigl(\delta_P^{1/4}+\delta_Q^{1/4}+\eps^{1/4}\bigr).
	\]
	The overlap and its zero effects at undefined evaluations are those of
	Theorem~\ref{thm:subline-z-overlap-real-direct}. This is a direct-law
	auxiliary, not the source assertion of Lemma~\ref{lem:claim-17-3}.
\end{theorem}
\begin{proof}
	\leanok
	\uses{thm:subline-z-overlap-real-direct, lem:claim-17-3-re-direct}
	The reality identity gives
	$|B_{\mathrm{dir}}(T)-1|=|\operatorname{Re}B_{\mathrm{dir}}(T)-1|$.
	Apply Lemma~\ref{lem:claim-17-3-re-direct} with its unchanged constant;
	its proof permits $C=2$.
\end{proof}

\begin{theorem}[Formalization support: the complex Z bound for the constructed measurement]
	\label{thm:claim-17-3-concrete-direct}
	\lean{MIPStarRE.QPBT.subline_concrete_Z_term_near_one_direct}
	\leanok
	\uses{lem:qld-4-10, def:concrete-paired-line-measurement,
		thm:subline-z-overlap-real-direct}
	This auxiliary specializes the complex Z estimate to the constructed
	X-Z-X measurement without assuming paired-line consistency data.
	Fix $\delta_Q:\R\to\R$ such that, for some $K_Q\ge1$ and $r_Q>0$,
	$0\le\delta_Q(u)\le K_Q u^{r_Q}$ for all $u\ge0$.
	There exist a function $\delta_P:\R^2\to\R$, constants $C>0$ and
	$K_P\ge1$, and exponents $r_P,s_P>0$ such that
	\[
		0\le\delta_P(u,v)\le K_P(u^{r_P}+v^{s_P})\qquad(u,v\ge0),
	\]
	and the following holds for every admissible parameter tuple, projective
	setting with error $\eps$, projective joint point family satisfying the
	three bounds of Lemma~\ref{lem:qld-4-10} with upper bound
	$\delta_Q(\eps)$ on every directed opposite placement, and law
	$D_{\mathrm{dir}}$ of Definition~\ref{def:direct-subline-law}.
	The concrete X-Z-X POVM $T$ of
	Definition~\ref{def:concrete-paired-line-measurement} satisfies
	\[
		|B_{\mathrm{dir}}(T)-1|\le C\sqrt m
		\bigl(\delta_P(\eps,md/q)^{1/4}
		+\delta_Q(\eps)^{1/4}+\eps^{1/4}\bigr).
	\]
	The expectation uses the completed evaluations and uniform affine
	parameter of Theorem~\ref{thm:subline-z-overlap-real-direct}, including
	zero diagonal directions. The two-variable error is in the corrected sum
	sense of \cite{gap:qpbt_pasting-product-error}. Identification with the
	source law and its evaluation convention remains open.
\end{theorem}
\begin{proof}
	\leanok
	\uses{lem:combined-line-measurement-consistency,
		lem:concrete-paired-line-degree-support, thm:claim-17-3-direct}
	After fixing $\delta_Q$ and its polynomial bound, apply
	Lemma~\ref{lem:combined-line-measurement-consistency} to obtain
	$\delta_P$ for the concrete measurement. Its degree support and this
	consistency estimate supply all the line hypotheses of
	Theorem~\ref{thm:claim-17-3-direct}. Apply that theorem without changing
	its constant, so the same choice $C=2$ is available.
\end{proof}

\begin{definition}[Conditional extended-line construction]
	\label{def:qld-conditional-extended-line-measurement}
	\lean{MIPStarRE.QPBT.SubLineWitness.Fiber}
	\lean{MIPStarRE.QPBT.SubLineWitness.conditionalLaw}
	\lean{MIPStarRE.QPBT.SubLineWitness.affineData}
	\lean{MIPStarRE.QPBT.SubLineWitness.combinedPolynomial}
	\lean{MIPStarRE.QPBT.SubLineWitness.extendedMeasurement}
	\leanok
	\uses{def:direct-subline-law, def:combine-map,
		def:concrete-paired-line-measurement}
	Fix a directly indexed extended line $\ell$. Its fiber consists of the
	supported subline triples $(\ell,\ell_X,\ell_Z)$ above it. Normalize the
	original subline law on this fiber to obtain the conditional law
	$D_{\mid\ell}$.

	For each supported triple, choose the affine coordinate data supplied by
	Definition~\ref{def:direct-subline-law}. If $f_X$ and $f_Z$ are the two
	source-line polynomials, let $h$ be their affine combination
	\[
		h(t)=\alpha(t)f_X(a_X+b_Xt)+\beta(t)f_Z(a_Z+b_Zt).
	\]
	The extended-line measurement at $\ell$ is the conditional average, under
	$D_{\mid\ell}$, of the concrete paired-line measurements postprocessed by
	$(f_X,f_Z)\mapsto h$.
\end{definition}

\begin{example}[Formalization support: conditional extended-line identities]
	\label{lem:qld-conditional-extended-line-properties}
	\lean{MIPStarRE.QPBT.directLinePointDist_line_weight_pos}
	\lean{MIPStarRE.QPBT.SubLineWitness.fiber_mass_pos}
	\lean{MIPStarRE.QPBT.SubLineWitness.conditionalLaw_isProbability}
	\lean{MIPStarRE.QPBT.SubLineWitness.avgOver_conditionalLaw}
	\lean{MIPStarRE.QPBT.SubLineWitness.affineData_compatible}
	\lean{MIPStarRE.QPBT.SubLineWitness.combinedPolynomial_evaluatesTo}
	\lean{MIPStarRE.QPBT.SubLineWitness.combinedPolynomial_evalOpt}
	\leanok
	\uses{def:qld-conditional-extended-line-measurement,
		def:direct-subline-law, def:combine-map}
	Every directly indexed extended line has positive marginal mass, and hence
	each of its subline fibers has positive mass. Consequently, the conditional
	law $D_{\mid\ell}$ is a probability distribution, and finite disintegration
	gives
	\[
		\mathbb E_{\ell}\mathbb E_{D_{\mid\ell}}[F]=\mathbb E_D[F]
	\]
	for every real-valued function $F$ on supported subline triples. The selected
	affine data satisfy the projection identities in
	Definition~\ref{def:direct-subline-law}. If the two source polynomials
	evaluate to $a$ and $b$ at the projected point, their combined polynomial
	evaluates there to $\alpha a+\beta b$. The same identity holds for completed
	evaluation whenever both source evaluations are defined.
\end{example}
\begin{proof}
	\leanok
	\uses{def:qld-conditional-extended-line-measurement,
		def:direct-subline-law, def:combine-map}
	Positivity follows from the explicit direct-line sampler. Normalizing each
	fiber gives total mass one, and grouping the original finite sum by its
	extended-line coordinate proves the disintegration identity. The evaluation
	identities follow by substituting the compatible affine data in the polynomial
	combination formula.
\end{proof}

\begin{lemma}[Formalization support: axis degree of the combined polynomial]
	\label{lem:qld-combined-polynomial-axis-degree}
	\lean{MIPStarRE.QPBT.combineLinePoly_axis_degree}
	\lean{MIPStarRE.QPBT.SubLineWitness.affineData_axis}
	\lean{MIPStarRE.QPBT.SubLineWitness.combinedPolynomial_axis_degree}
	\lean{MIPStarRE.QPBT.SubLineWitness.extendedMeasurement_axis_degree}
	\leanok
	\uses{def:qld-conditional-extended-line-measurement,
		lem:qld-conditional-extended-line-properties,
		lem:concrete-paired-line-degree-support}
	Suppose two univariate coefficient families vanish above degree $d\geq1$.
	Their affine combination also vanishes above degree $d$ whenever either
	both coefficient functions $\alpha,\beta$ are constant or both projected
	line parameters are constant. The affine data of every supported
	axis-parallel extended line satisfy one of these alternatives. Consequently,
	the combined polynomial has degree at most $d$, and the conditional
	extended-line measurement has zero effect on every outcome of degree greater
	than $d$.
\end{lemma}
\begin{proof}
	\leanok
	\uses{def:qld-conditional-extended-line-measurement,
		lem:qld-conditional-extended-line-properties,
		lem:concrete-paired-line-degree-support}
	Affine substitution preserves the source degree bounds. On an extended
	coordinate line, either the $\alpha$ and $\beta$ slopes vanish or both
	projected direction parameters vanish. Average the resulting zero effects
	over the conditional law.
\end{proof}

\begin{definition}[Paired, extended, and ordered subline overlaps]
	\label{def:qld-subline-overlaps}
	\lean{MIPStarRE.QPBT.pairedSublineOverlap}
	\lean{MIPStarRE.QPBT.extendedPointOverlap}
	\lean{MIPStarRE.QPBT.orderedSublineOverlap}
	\leanok
	\uses{def:qld-conditional-extended-line-measurement,
		def:direct-subline-law, def:expanded-state}
	For a supported subline triple, the paired overlap is the uniform-parameter
	average of the overlap between the paired-line measurement and the completed
	joint point measurement evaluated at the two projected points. The extended
	point overlap compares the conditional extended-line measurement with the
	point measurement postprocessed by $(a,b)\mapsto\alpha a+\beta b$. The
	ordered subline overlap replaces the joint point effect in the paired overlap
	by the ordered $Z$-then-$X$ product.
\end{definition}

\begin{lemma}[Formalization support: comparison with the extended overlap]
	\label{lem:qld-extended-line-overlap-comparison}
	\lean{MIPStarRE.QPBT.SubLineWitness.completedPoint_effect_le}
	\lean{MIPStarRE.QPBT.avgOver_directLinePointDist_resample_parameter}
	\lean{MIPStarRE.QPBT.SubLineWitness.extendedMeasurement_overlap}
	\lean{MIPStarRE.QPBT.SubLineWitness.extendedPointOverlap_ge}
	\lean{MIPStarRE.QPBT.SubLineWitness.avgOver_extendedPointOverlap_ge}
	\lean{MIPStarRE.QPBT.SubLineWitness.extended_consistencyDefect_le}
	\leanok
	\uses{def:qld-subline-overlaps,
		lem:qld-conditional-extended-line-properties, def:consistency}
	At a point of a supported extended line, the completed joint-point effect at
	the two source evaluations is bounded above by the effect obtained by
	postprocessing with the combined value $\alpha a+\beta b$. Uniformly
	resampling the affine parameter preserves the full direct line-point law.
	Expanding the conditional extended-line measurement in a placed overlap
	therefore shows that the average extended point overlap is at least the
	average paired subline overlap.

	For every directed pair of opposite placements, the consistency defect
	between the evaluated extended-line measurement and the combined-value point
	measurement is consequently at most one minus the average paired overlap.
	Both measurements here have alphabet $\F_q\cup\{\bot\}$: the point
	measurement sends $(a,b)$ to the single value $\alpha a+\beta b$ and has
	zero effect at $\bot$, while the line measurement uses completed evaluation.
\end{lemma}
\begin{proof}
	\leanok
	\uses{def:qld-subline-overlaps, def:qld-conditional-extended-line-measurement,
		lem:qld-conditional-extended-line-properties, def:consistency}
	Expand the conditional measurement, group equal postprocessed answers, and
	use positivity on opposite registers to preserve the pointwise effect
	inequality. Disintegrate the subline law and use the identity expressing a
	consistency defect as one minus its matched-outcome overlap.
\end{proof}

\begin{definition}[The completed X-point overlap]
	\label{def:qld-x-point-overlap}
	\lean{MIPStarRE.QPBT.xPointOverlapAt}
	\leanok
	\uses{def:qld-subline-overlaps, def:expanded-point-measurement}
	For a pair of source lines and an $X$-point, the completed $X$-point overlap
	is the sum, over both line-polynomial outcomes, of the paired-line effect
	against the completed $X$-point effect indexed by the $X$-line evaluation.
\end{definition}

\begin{lemma}[Formalization support: the X-point overlap deficit]
	\label{lem:qld-x-point-overlap-deficit}
	\lean{MIPStarRE.QPBT.SubLineWitness.avgOver_regrouped_eq_at}
	\lean{MIPStarRE.QPBT.xPointOverlap_eq_at}
	\lean{MIPStarRE.QPBT.SubLineWitness.one_sub_avgOver_xPointOverlap_le_at}
	\lean{MIPStarRE.QPBT.ProjectiveSetting.completedPair_norm_sq_sum_XZ}
	\lean{MIPStarRE.QPBT.CombinedLinesWitness.one_sub_X_overlap_le_at}
	\lean{MIPStarRE.QPBT.CombinedLinesWitness.one_sub_X_overlap_restricted_le_at}
	\leanok
	\uses{def:qld-x-point-overlap, lem:qld-4-10-same-placement,
		lem:restricted-line-mixture-bounds, def:consistency}
	Regrouping the two line-answer sums identifies the subline average with the
	overlap of the evaluated paired-line measurement and a completed point
	family. Completing the joint point measurement and the ordered
	$X$-then-$Z$ product preserves their squared distance exactly.

	More generally, let $\nu$ be any probability law on pairs of line-point
	samples. If the evaluated paired-line measurement has consistency defect at
	most $\eta$ with the completed joint points, and those joint points have
	squared distance at most $\zeta$ from the completed $X$-then-$Z$ product,
	then the completed $X$-point overlap has deficit at most
	$2\sqrt\eta+2\sqrt\zeta$. On a product of two restricted line-point laws this
	gives
	\[
		2\sqrt{4m^2\delta_P}+2\sqrt{4\delta_Q}.
	\]
	The same bound holds for the directly indexed subline average by its
	one-point mixture identity.
\end{lemma}
\begin{proof}
	\leanok
	\uses{def:qld-x-point-overlap, lem:qld-4-10-same-placement,
		lem:restricted-line-mixture-bounds, def:consistency}
	Apply the consistency-to-overlap estimate and Cauchy--Schwarz to the two
	comparisons. Restricted line consistency inflates by $4m^2$, while the
	same-placement $X$-then-$Z$ point estimate contributes $4\delta_Q$.
	Average these bounds over the subline mixture.
\end{proof}

\begin{lemma}[Formalization support: components of the first overlap route]
	\label{lem:qld-first-route-overlap-components}
	\lean{MIPStarRE.QPBT.subline_remove_X_factor_at}
	\lean{MIPStarRE.QPBT.subline_Z_term_near_one_at}
	\lean{MIPStarRE.QPBT.SubLineWitness.paired_ordered_overlap_gap_le}
	\leanok
	\uses{def:qld-subline-overlaps, lem:qld-x-point-overlap-deficit,
		lem:qld-4-10-same-placement, lem:restricted-line-mixture-bounds}
	Let $I_{\mathrm{pair}}$ and $I_{ZX}$ denote the averages of the paired and
	ordered subline overlaps. For arbitrary placements,
	\[
		|I_{\mathrm{pair}}-I_{ZX}|\leq\sqrt{4\delta_Q}.
	\]
	For opposite placements, let $I_Z$ be the same ordered overlap after
	removing the $X$-point factor. There are universal constants $C_X,C_Z>0$
	such that
	\[
		|I_{ZX}-I_Z|
		\leq C_X\sqrt m\bigl(\delta_P^{1/4}+\delta_Q^{1/4}\bigr)
	\]
	and
	\[
		|I_Z-1|\leq C_Z\sqrt m
		\bigl(\delta_P^{1/4}+\delta_Q^{1/4}+\eps^{1/4}\bigr).
	\]
\end{lemma}
\begin{proof}
	\leanok
	\uses{def:qld-subline-overlaps, lem:qld-x-point-overlap-deficit,
		lem:qld-4-10-same-placement, lem:restricted-line-mixture-bounds}
	The first estimate replaces the completed joint points by the ordered
	product. For the second, apply Cauchy--Schwarz to the $X$-point deficit. The
	last estimate is the $Z$-line/point overlap bound under the separate
	one-point mixture law.
\end{proof}

\begin{lemma}[Replacing the combined point measurement by the ordered product]
	\label{lem:claim-17-1}
	\notready
	% claim:17-1 in the paper; the original label is recorded here as a
	% comment for cross-reference fidelity.
	\uses{lem:qld-sublines, lem:qld-xz-lines, lem:qld-4-10,
		def:expanded-point-measurement, def:expanded-state}
	\begin{align*}
		\E_{(\ell,\ell_X,\ell_Z) \sim D}\
		\E_{(x,z,\alpha,\beta) \in \ell}
		&\sum_{f_X,f_Z} \bra{\hat\psi}
		\big(T^{\ell_X,\ell_Z}_{f_X,f_Z}\big)_{\mathsf{A}\mathsf{A}'} \ot
		\big(\hat{Q}^{x,z}_{f_X(x),f_Z(z)}\big)_{\mathsf{B}\mathsf{A}''}
		\ket{\hat\psi} \\
		\approx_{\delta_Q^{1/2}}\
		\E_{(\ell,\ell_X,\ell_Z) \sim D}\
		\E_{(x,z,\alpha,\beta) \in \ell}
		&\sum_{f_X,f_Z} \bra{\hat\psi}
		\big(T^{\ell_X,\ell_Z}_{f_X,f_Z}\big)_{\mathsf{A}\mathsf{A}'} \ot
		\big(\hat{M}^{(\mathrm{Point},Z),z}_{f_Z(z)} \,
		\hat{M}^{(\mathrm{Point},X),x}_{f_X(x)}
		\big)_{\mathsf{B}\mathsf{A}''}
		\ket{\hat\psi}\,.
	\end{align*}
\end{lemma}
\begin{proof}
	\uses{lem:qld-4-10, lem:qld-4-10-same-placement, lem:qld-sublines,
		lem:qld-xz-lines, lem:overlap-gap-distance,
		lem:qld-sublines-uniform-point, def:line-point-dist}
	The scalar approximation here is a complex-modulus bound. Its formal
	proof over the source law remains open; the real-part auxiliary above
	does not establish it. See \cite{gap:qpbt_subline-claims-line-marginal}.
	Write
	$W^{x,z}_{a,b} = \hat{Q}^{x,z}_{a,b}
	- \hat{M}^{(\mathrm{Point},Z),z}_{b}
	\hat{M}^{(\mathrm{Point},X),x}_{a}$
	and apply the Cauchy--Schwarz inequality to the difference of the two
	sides. The first factor is at most $1$ because
	$\sum_{f_X,f_Z \,:\, f_X(x)=a,\ f_Z(z)=b}
	T^{\ell_X,\ell_Z}_{f_X,f_Z} \le \Id$
	for every $x, z, a, b$. For the second factor, note that when $\ell$ is
	drawn from the line-point distribution over $\F_q^{2m+2}$, a uniformly
	random point $(x,z,\alpha,\beta)$ on $\ell$ is a uniformly random point
	of $\F_q^{2m+2}$, so the coordinates $x, z$ are independent and uniform
	in $\F_q^m$; hence
	$\E_{x,z} \sum_{a,b}
	\bra{\hat\psi} \Id \ot (W^{x,z}_{a,b})^\dagger W^{x,z}_{a,b}
	\ket{\hat\psi}\le4\delta_Q$
	by Lemma~\ref{lem:qld-4-10-same-placement}. The total error is therefore
	at most $2\sqrt{\delta_Q}$, with the same asymptotic order as the
	source claim. The adjoint is necessary because the ordered product need
	not be Hermitian. The complex modulus estimate and source-distribution
	transport remain separate obligations; the proved auxiliary real-part
	estimate is stated below.
\end{proof}

\begin{lemma}[Removing the $X$-point factor]
	\label{lem:claim-17-2}
	\notready
	% claim:17-2 in the paper; the original label is recorded here as a
	% comment for cross-reference fidelity. The source right-hand side
	% carries a typographical superscript "Z" in place of the point "z";
	% the corrected form is stated here.
	\uses{lem:qld-sublines, lem:qld-xz-lines,
		def:expanded-point-measurement, lem:qld-comm-line-cons,
		def:expanded-state}
	Let $T$ be the X-Z-X measurement defined in the proof of
	Lemma~\ref{lem:qld-xz-lines}. The construction is part of this claim's
	domain: its degree and evaluated consistency properties alone do not
	imply the following estimate, as shown in
	\cite{gap:qpbt_subline-claims-line-marginal}.
	\begin{align*}
		\E_{(\ell,\ell_X,\ell_Z) \sim D}\
		\E_{(x,z,\alpha,\beta) \in \ell}
		&\sum_{f_X,f_Z} \bra{\hat\psi}
		\big(T^{\ell_X,\ell_Z}_{f_X,f_Z}\big)_{\mathsf{A}\mathsf{A}'} \ot
		\big(\hat{M}^{(\mathrm{Point},Z),z}_{f_Z(z)} \,
		\hat{M}^{(\mathrm{Point},X),x}_{f_X(x)}
		\big)_{\mathsf{B}\mathsf{A}''}
		\ket{\hat\psi} \\
		\approx_{m\,\delta_{\mathrm{Line}}^{1/2}}\
		\E_{(\ell,\ell_X,\ell_Z) \sim D}\
		\E_{(x,z,\alpha,\beta) \in \ell}
		&\sum_{f_X,f_Z} \bra{\hat\psi}
		\big(T^{\ell_X,\ell_Z}_{f_X,f_Z}\big)_{\mathsf{A}\mathsf{A}'} \ot
		\big(\hat{M}^{(\mathrm{Point},Z),z}_{f_Z(z)}
		\big)_{\mathsf{B}\mathsf{A}''}
		\ket{\hat\psi}\,,
	\end{align*}
	where $\delta_{\mathrm{Line}}$ is the error function of
	Lemma~\ref{lem:qld-comm-line-cons}.
\end{lemma}
\begin{proof}
	\uses{lem:qld-comm-line-cons, lem:qld-sublines, lem:qld-xz-lines,
		lem:restricted-line-mixture-bounds, def:ith-restricted-line,
		def:expanded-line-measurement}
	By Cauchy--Schwarz, the difference of the two sides is bounded by the
	square root of the deficit
	\[
		1 - \E_{(\ell,\ell_X,\ell_Z) \sim D}\
		\E_{(x,z,\alpha,\beta) \in \ell} \sum_{a}
		\bra{\hat\psi}
		T^{\ell_X,\ell_Z,x}_{a} \ot
		\hat{M}^{(\mathrm{Point},X),x}_{a}
		\ket{\hat\psi}\,,
		\qquad
		T^{\ell_X,\ell_Z,x}_{a}
		= \sum_{f_X,f_Z \,:\, f_X(x)=a}
		T^{\ell_X,\ell_Z}_{f_X,f_Z}\,.
	\]
	By the sandwich definition of $T$ and the projectivity of the lines
	measurement $\{\hat{M}^{(\mathrm{Line},X),\ell_X}_{f_X}\}$, the
	subtracted expectation equals
	$\E \sum_a \bra{\hat\psi}
	\hat{M}^{(\mathrm{Line},X),\ell_X}_{[\mathrm{eval}_x(\cdot)=a]} \ot
	\hat{M}^{(\mathrm{Point},X),x}_{a} \ket{\hat\psi}$,
	which by Property~2 of Lemma~\ref{lem:qld-sublines} is a mixture of the
	corresponding averages over the restricted distributions
	$D_{\mathrm{v},i}$. By the error inflation of
	Lemma~\ref{lem:restricted-line-mixture-bounds} applied to the
	line--point consistency in item~3 of Lemma~\ref{lem:qld-comm-line-cons},
	each restricted term has deficit $O(m\,\delta_{\mathrm{Line}})$, so the
	total deficit is $O(m^2\,\delta_{\mathrm{Line}})$ (the final line of the
	source proof asserts this bound for the consistency term itself rather
	than for its deficit; see Remark~\ref{rem:qld-4-13-source-defects}).
	Taking the square root gives the claimed error
	$m\,\delta_{\mathrm{Line}}^{1/2}$.
\end{proof}

\begin{lemma}[The remaining $Z$-point term is close to one]
	\label{lem:claim-17-3}
	\notready
	% claim:17-3 in the paper; the original label is recorded here as a
	% comment for cross-reference fidelity.
	\uses{lem:qld-sublines, lem:qld-xz-lines, lem:qld-4-10,
		def:expanded-point-measurement, def:expanded-state}
	\[
		\E_{(\ell,\ell_X,\ell_Z) \sim D}\
		\E_{(x,z,\alpha,\beta) \in \ell}
		\sum_{f_X,f_Z} \bra{\hat\psi}
		\big(T^{\ell_X,\ell_Z}_{f_X,f_Z}\big)_{\mathsf{A}\mathsf{A}'} \ot
		\big(\hat{M}^{(\mathrm{Point},Z),z}_{f_Z(z)}
		\big)_{\mathsf{B}\mathsf{A}''}
		\ket{\hat\psi}
		\ \approx_{\sqrt{m}\,(\delta_P^{1/4} + \delta_Q^{1/4} +
			\eps^{1/4})}\ 1\,.
	\]
\end{lemma}
\begin{proof}
	\uses{lem:restricted-line-mixture-bounds, lem:qld-sublines,
		lem:qld-xz-lines, lem:qld-4-10, lem:qld-comm-cons,
		def:ith-restricted-line}
	By Cauchy--Schwarz it suffices to bound the deficit
	$1 - \E \sum_b \bra{\hat\psi} T^{\ell_X,\ell_Z,z}_{b} \ot
	\hat{M}^{(\mathrm{Point},Z),z}_{b} \ket{\hat\psi}$,
	where
	$T^{\ell_X,\ell_Z,z}_{b}
	= \sum_{f_X} T^{\ell_X,\ell_Z}_{f_X,[\mathrm{eval}_z(\cdot)=b]}$.
	By Property~2 of Lemma~\ref{lem:qld-sublines}, the expectation can be
	rewritten over the mixture of product distributions
	$D_{\mathrm{v},i} \times D_{\mathrm{v},j}$ with $z$ uniformly random on
	$\ell_Z$, and, since
	$T^{\ell_X,\ell_Z,z}_{b}
	= \E_{x \in \ell_X} \sum_a
	T^{\ell_X,\ell_Z}
	_{[\mathrm{eval}_x(\cdot)=a,\ \mathrm{eval}_z(\cdot)=b]}$,
	one may additionally average over $x$ uniformly random on $\ell_X$. Up
	to error $m(\sqrt{\delta_P} + \sqrt{\delta_Q})$ --- obtained by
	applying the error inflation of
	Lemma~\ref{lem:restricted-line-mixture-bounds} to the conclusion of
	Lemma~\ref{lem:qld-xz-lines} (transported to the sub-line distribution
	via Lemma~\ref{lem:qld-sublines}) and to the consistency relations of
	Lemma~\ref{lem:qld-4-10} --- the resulting quantity is close to
	$\E \sum_{a,b} \bra{\hat\psi}
	\hat{M}^{(\mathrm{Point},X),x}_{a}
	\hat{M}^{(\mathrm{Point},Z),z}_{b} \ot
	\hat{M}^{(\mathrm{Point},Z),z}_{b} \ket{\hat\psi}$,
	which, since $\{\hat{M}^{(\mathrm{Point},X),x}_{a}\}$ is a measurement,
	equals
	$\sum_b \bra{\hat\psi} \hat{M}^{(\mathrm{Point},Z),z}_{b} \ot
	\hat{M}^{(\mathrm{Point},Z),z}_{b} \ket{\hat\psi}
	\ge 1 - \sqrt{\eps}$
	by the self-consistency of the points measurements
	(Lemma~\ref{lem:qld-comm-cons}). The deficit is therefore
	$O(m(\sqrt{\delta_P} + \sqrt{\delta_Q}) + \sqrt{\eps})$, and taking the
	square root gives the claim.
\end{proof}

The following auxiliary real-part estimate uses the directly indexed
extended-line law $D_{\rm dir}$ of Lemma~\ref{lem:qld-sublines-direct} and
Definition~\ref{def:direct-subline-law}. Fix admissible parameters and a
projective setting with error $\eps$.

For the concrete paired-line measurement $T$ of
Definition~\ref{def:concrete-paired-line-measurement}, write $A_{ZX}(T)$ and
$A_Z(T)$ for the ordered-product and Z-point overlap averages displayed in
Lemmas~\ref{lem:claim-17-1}--\ref{lem:claim-17-3}, with $D_{\rm dir}$ in place
of $D$ and a uniform affine parameter on each extended line. Line answers are
polynomial representatives of degree at most $md$, evaluations take values in
$\F_q\cup\{\perp\}$, and completed point effects vanish at undefined
evaluations. The placements remain $\mathsf{A}\mathsf{A}'$ and
$\mathsf{B}\mathsf{A}''$. No joint-point or paired-line consistency hypothesis
enters this estimate: it uses only the concrete sandwich.

\begin{lemma}[Formalization support: real-part X-factor removal under the
	directly indexed law]
	\label{lem:claim-17-2-direct-real}
	\lean{MIPStarRE.QPBT.subline_remove_X_factor_re_direct}
	\leanok
	\uses{def:concrete-paired-line-measurement, lem:qld-sublines-direct,
		def:expanded-point-measurement}
	There is a universal $C>0$ such that, for every projective setting and
	every directly indexed subline law above,
	\[
	 |\operatorname{Re}A_{ZX}(T)-\operatorname{Re}A_Z(T)|
	 \le C m\sqrt{\delta_{\mathrm{Line}}(\eps)}.
	\]
	This is the real-part analogue of Lemma~\ref{lem:claim-17-2-direct}; it
	asserts neither a complex-modulus bound nor transport to the source
	distribution.
\end{lemma}
\begin{proof}
	\leanok
	\uses{lem:concrete-line-X-marginal, lem:concrete-X-deficit-direct}
	The concrete X marginal identity and the concrete X-overlap deficit bound
	the deficit by $K m^2\delta_{\mathrm{Line}}(\eps)$, and projective
	Cauchy--Schwarz bounds the real-part difference by its square root, so
	$C=\sqrt K$ suffices. The setup and this deficit bound are shared with
	Lemma~\ref{lem:claim-17-2-direct}, which applies the complex
	Cauchy--Schwarz estimate to the same data. No paired-line consistency
	hypothesis is used in this calculation.
\end{proof}

The real-part estimates of Lemmas~\ref{lem:claim-17-1-re-direct},
\ref{lem:claim-17-2-direct-real} and~\ref{lem:claim-17-3-re-direct} remain
available. On the directly indexed law, the complex estimates for the
concrete X-Z-X measurement are Theorem~\ref{thm:claim-17-1-direct},
Lemma~\ref{lem:claim-17-2-direct}, and
Theorem~\ref{thm:claim-17-3-concrete-direct}. The last obtains its line
consistency data internally from a polynomially controlled joint point family.
The source-law transport, zero-direction evaluation convention, and
extended-line error bound remain separate obligations recorded in
\cite{gap:qpbt_subline-claims-line-marginal}. None of the three
source-labelled claims is certified by these auxiliary entries.

\begin{proof}[Proof of Lemma~\ref{lem:qld-4-13}]
	\proves{lem:qld-4-13}
	\uses{lem:qld-xz-lines, lem:qld-sublines, def:combine-map,
		lem:qld-4-12, lem:claim-17-1, lem:claim-17-2, lem:claim-17-3}
	Let $D$ be the distribution of Lemma~\ref{lem:qld-sublines} and let
	$D_{|\ell}$ denote $D$ conditioned on the first element being $\ell$.
	For every line $\ell$ in $\F_q^{2m+2}$ and $f \in \deg_{md+1}(\ell)$
	define
	\[
		\hat{Q}^{\ell}_{f}
		= \E_{(\ell_X,\ell_Z) \sim D_{|\ell}}
		\sum_{\substack{f_X \in \deg_{md}(\ell_X),\
			f_Z \in \deg_{md}(\ell_Z):\\
			f = \mathrm{combine}_{f_X,f_Z}}}
		T^{\ell_X,\ell_Z}_{f_X,f_Z}\,,
	\]
	a valid POVM; the polynomial $\mathrm{combine}_{f_X,f_Z}$ of
	Definition~\ref{def:combine-map} is well defined here because Property~2
	of Lemma~\ref{lem:qld-sublines} guarantees that $\ell_X$ and $\ell_Z$
	are compatible with $\ell$. If $\ell$ is axis-parallel, then by
	Property~3 of Lemma~\ref{lem:qld-sublines} so are $\ell_X$ and
	$\ell_Z$, hence $f_X, f_Z \in \deg_d$ by Lemma~\ref{lem:qld-xz-lines};
	along such a line either $(\alpha,\beta)$ is constant, in which case
	$f = \mathrm{combine}_{f_X,f_Z}$ is a fixed linear combination of
	polynomials of degree at most $d$, or $(x,z)$ is constant, in which case
	$f$ is affine in the line parameter, so in either case
	$f \in \deg_d(\ell)$. Chaining Lemmas~\ref{lem:claim-17-1},
	\ref{lem:claim-17-2} and~\ref{lem:claim-17-3} with the triangle
	inequality for scalars gives the consistency bound
	\begin{equation}
		\label{eq:qld-combined-lines-consistency}
		\E_{(\ell,\ell_X,\ell_Z) \sim D}\
		\E_{(x,z,\alpha,\beta) \in \ell}
		\sum_{f_X,f_Z} \bra{\hat\psi}
		\big(T^{\ell_X,\ell_Z}_{f_X,f_Z}\big)_{\mathsf{A}\mathsf{A}'} \ot
		\big(\Id - \hat{Q}^{x,z}_{f_X(x),f_Z(z)}
		\big)_{\mathsf{B}\mathsf{A}''}
		\ket{\hat\psi}
		= O\big(m\,(\eps^{1/4} + \delta_P^{1/4} + \delta_Q^{1/4}
		+ \delta_{\mathrm{Line}}^{1/2})\big)\,;
	\end{equation}
	since
	$\delta_P = \poly(\eps, md/q)$, $\delta_Q = \poly(\eps)$, and
	$\delta_{\mathrm{Line}} = \poly(\eps)$, the right-hand side is a
	function of the form $m \cdot \poly(\eps, md/q)$. This established
	bound is not of the stated form $\poly(m^2\eps, md/q)$, and no
	function of that form dominates it, so the error function claimed in
	the statement remains unproved along this route; the source's second
	route rests on a joint-law decomposition not provided by Property~2 of
	Lemma~\ref{lem:qld-sublines}, and the established bound is the one
	used in the sequel (Remark~\ref{rem:qld-4-13-source-defects}; the two
	routes are analyzed in \cite{gap:qpbt_combined-lines-error-term}). Now
	if the pair $(f_X, f_Z)$ evaluates to $(a, b)$ at $(x, z)$, then
	$\mathrm{combine}_{f_X,f_Z}$ evaluates to $\alpha a + \beta b$ at
	$(x, z, \alpha, \beta)$; hence the consistency of $T$ with
	$\hat{Q}^{x,z}$ along the sub-line pairs expressed by
	\eqref{eq:qld-combined-lines-consistency}, the definition of
	$\hat{Q}^{\ell}$ as the image of $T$ under the combining map, and the
	construction of $\hat{Q}^{x,z,\alpha,\beta}$ in
	Lemma~\ref{lem:qld-4-12} as the corresponding coarse-graining of
	$\hat{Q}^{x,z}$ together yield the first relation of
	\eqref{eq:qld-4-13}, with the established error
	$m \cdot \poly(\eps, md/q)$ in place of the stated
	$\delta_{\mathrm{combine}}$. The second relation follows from the
	analogous argument with the registers swapped.
\end{proof}

\begin{remark}[The error term in the source proof of
	Lemma~\ref{lem:qld-4-13}]
	\label{rem:qld-4-13-source-defects}
	The statement of Lemma~\ref{lem:qld-4-13} retains the source error
	form $\delta_{\mathrm{combine}} = \poly(m^2\eps, md/q)$, which
	remains unproved. The source proof in \cite{Ji2020MIPStar} asserts
	two different bounds for the quantity
	\eqref{eq:qld-combined-lines-consistency}: the bound
	$O(m(\eps^{1/4} + \delta_P^{1/4} + \delta_Q^{1/4} +
	\delta_{\mathrm{Line}}^{1/2}))$, obtained by chaining the three
	claims (Lemmas~\ref{lem:claim-17-1}--\ref{lem:claim-17-3}), and the
	bound $O(m^2 \delta_P)$, whose derivation decomposes
	\eqref{eq:qld-combined-lines-consistency} into
	restricted-distribution terms bounded by
	\eqref{eq:qld-xz-lines-restricted}. The second derivation requires a
	joint product law for the pair $((\ell_X, x), (\ell_Z, z))$ that
	Property~2 of Lemma~\ref{lem:qld-sublines} does not provide, and
	neither bound is dominated by a function of $m^2\eps$ and $md/q$
	alone that vanishes as its arguments vanish. The proof above
	therefore follows the first route only and establishes the bound
	$m \cdot \poly(\eps, md/q)$, which is the form consumed by
	Lemma~\ref{lem:qld-4-7}, whose error function absorbs the prefactor
	$m$ into $(md)^a$. The full analysis --- the two routes, the scalings
	showing that the claimed form follows from neither, the corrected
	deficit assertion in the proof of Lemma~\ref{lem:claim-17-2}, and the
	candidate repairs with their open obligations --- is given in
	\cite{gap:qpbt_combined-lines-error-term}.
	The printed error form is therefore an open source gap and not a Lean proof
	obligation of any consumer: every formalized consumer uses the established
	form of Lemma~\ref{lem:qld-4-13-established}, and the gap is tracked by
	GitHub issue 598 after the obstruction pull requests of the earlier issues
	closed those issues without proving the printed assertion.
	Since 2026-09-19 the printed claim is carried in Lean by the definition
	\path{MIPStarRE.QPBT.PrintedExtendedLinesWitnessClaim}, a proposition that is
	stated but not asserted, which replaced a theorem of the same statement whose
	proof was an open \path{sorry}. The statement the source prints is thus
	neither proved nor dropped, and the Lean development carries no proof debt at
	this site. What the source proof does establish is formalized by
	\path{MIPStarRE.QPBT.exists_extendedLinesWitness_established} of
	Lemma~\ref{lem:qld-4-13-established}; the scalar obstruction to the first
	route is \path{MIPStarRE.QPBT.not_exists_combining_quarter_power_bound}.
	No implication between the unasserted proposition and the established theorem
	is recorded in either direction, because none follows by a short faithful
	argument: the printed form applies an arbitrary polynomial error function to
	$m^2\eps$, whose exponent may exceed the single prefactor $m$ that the
	established form carries outside its error function.
\end{remark}

\begin{example}[The paired-line construction given a joint point family]
	\label{lem:combined-lines-given-points}
	\lean{MIPStarRE.QPBT.exists_combinedLinesWitness_ofPointsWitness}
	\leanok
	\uses{lem:qld-4-10, def:line-point-dist, def:consistency,
		def:ideg-deg-polynomials}
	Fix a function $\delta_Q(\eps)=\poly(\eps)$. There is a function
	$\delta_P(\eps,md/q)=\poly(\eps,md/q)$ with the following property. Let a
	projective setting with error $\eps$ be given, together with complete
	projective joint point measurements $Q^{x,z}_{a,b}$ on both expanded local
	spaces. Suppose that, for every directed opposite placement, these
	measurements are self-consistent with squared-distance error at most
	$\delta_Q(\eps)$ and are within the same error of both ordered products
	$\hat M^X_a(x)\hat M^Z_b(z)$ and
	$\hat M^Z_b(z)\hat M^X_a(x)$.

	Then there exists a complete paired-line POVM family
	$T^{\ell_X,\ell_Z}_{f_X,f_Z}$ on each expanded local space, with outcomes
	$(f_X,f_Z)$ of degree at most $md$ on their respective lines. Outcomes outside
	degree $d$ have zero effect in the $X$ coordinate when $\ell_X$ is
	axis-parallel, and likewise in the $Z$ coordinate when $\ell_Z$ is
	axis-parallel. Independently sample $(\ell_X,x)$ and $(\ell_Z,z)$ from the
	line-point law. For every directed opposite placement, postprocessing the line
	outcome to $(f_X(x),f_Z(z))$ and comparing it with the joint point outcome on
	the completed alphabet $(\F_q\cup\{\bot\})^2$ gives consistency defect at most
	$\delta_P(\eps,md/q)$. The point measurement is extended to this alphabet by
	assigning zero effect to every pair containing $\bot$.

	Because the joint point family is an explicit hypothesis, this auxiliary is
	separate from Lemma~\ref{lem:qld-xz-lines} and does not prove that source
	lemma.
\end{example}
\begin{proof}
	\leanok
	\uses{lem:combined-line-measurement-consistency,
		def:concrete-paired-line-measurement,
		lem:concrete-paired-line-degree-support}
	Choose the concrete sandwich of
	Definition~\ref{def:concrete-paired-line-measurement}. The three required
	properties are
	Lemma~\ref{lem:concrete-paired-line-degree-support} for the two axis-degree
	properties and Lemma~\ref{lem:combined-line-measurement-consistency} for the
	completed-answer consistency bound on the explicitly given point family.
\end{proof}

\begin{definition}[The paired subline overlap]
	\label{def:paired-subline-overlap}
	\lean{MIPStarRE.QPBT.pairedSublineOverlap}
	\leanok
	\uses{def:direct-subline-law, def:consistency}
	For a joint point family, a paired-line family, two placements, and a triple
	$(\ell,\ell_X,\ell_Z)$, define the paired subline overlap to be the
	average over a uniform affine parameter on $\ell$ of the real correlation
	between $T^{\ell_X,\ell_Z}_{f_X,f_Z}$ on the first placement and the completed
	joint point effect indexed by the evaluations of $(f_X,f_Z)$ at the projected
	point on the second placement, summed over $(f_X,f_Z)$.
\end{definition}

\begin{definition}[The extended-line measurement from a subline witness]
	\label{def:extended-line-measurement}
	\lean{MIPStarRE.QPBT.SubLineWitness.extendedMeasurement}
	\leanok
	\uses{def:direct-subline-law, def:combine-map, def:povm-conventions}
	For a fixed directly indexed extended line $\ell$, condition the sub-line law
	on the triples whose first coordinate is $\ell$. Push each paired-line POVM
	forward by the polynomial combination map and average over this conditional
	law. The resulting effects form a POVM on polynomials of degree at most
	$md+1$ on $\ell$.
\end{definition}

\begin{lemma}[Auxiliary estimates for the paired subline overlap]
	\label{lem:paired-subline-overlap-estimates}
	\lean{MIPStarRE.QPBT.SubLineWitness.paired_ordered_overlap_gap_le}
	\lean{MIPStarRE.QPBT.subline_remove_X_factor_at}
	\lean{MIPStarRE.QPBT.subline_Z_term_near_one_at}
	\leanok
	\uses{def:paired-subline-overlap, def:direct-subline-law}
	These formalization-only auxiliary estimates give the directly indexed real
	overlap steps used in the first route of Lemma~\ref{lem:qld-4-13}. Let the
	directly indexed sub-line law, a joint point family of error
	$\delta_Q$, and a paired-line family of error $\delta_P$ be fixed. Replacing
	the joint point effect in the paired overlap by the ordered $Z$-then-$X$
	product changes the average by at most $\sqrt{4\delta_Q}$ for arbitrary
	placements. For opposite placements, there are universal positive constants
	such that removing the remaining $X$ factor changes the average by at most
	\[
		C\sqrt m\bigl(\delta_P^{1/4}+\delta_Q^{1/4}\bigr),
	\]
	and the resulting $Z$ overlap differs from one by at most
	\[
		C'\sqrt m\bigl(\delta_P^{1/4}+\delta_Q^{1/4}+\eps^{1/4}\bigr).
	\]
\end{lemma}
\begin{proof}
	\leanok
	\uses{lem:qld-4-10-same-placement, lem:overlap-gap-distance,
		lem:qld-x-point-overlap-deficit, lem:qld-sublines-uniform-point,
		lem:restricted-line-mixture-bounds, thm:qld-quarter-power-deficit-bound}
	The first estimate uses the squared-distance comparison with the ordered point
	product. For the second, the separate $X$ marginal of the directly indexed law
	and Lemma~\ref{lem:qld-x-point-overlap-deficit} bound the $X$-overlap deficit by
	\[
		2\sqrt{4m^2\delta_P}+2\sqrt{4\delta_Q}.
	\]
	This bound holds for every paired-line family with consistency error
	$\delta_P$ relative to the joint point family of error $\delta_Q$.
	The source proof of Claim~\ref{lem:claim-17-2} uses the concrete
	$X$-$Z$-$X$ sandwich; no such identification of the paired-line family is
	assumed here. Cauchy--Schwarz bounds the change on removing the $X$ factor
	by the square root of this deficit bound. The consistency hypotheses imply
	$\delta_P,\delta_Q\geq0$, and admissibility gives $m\geq1$, so
	Theorem~\ref{thm:qld-quarter-power-deficit-bound} gives
	\[
		\sqrt{2\sqrt{4m^2\delta_P}+2\sqrt{4\delta_Q}}
		\leq 2\sqrt m\bigl(\delta_P^{1/4}+\delta_Q^{1/4}\bigr).
	\]
	Thus the quarter powers arise from the two successive square roots.
	The third estimate uses the separate $Z$ marginal, positivity, and point
	self-consistency.
\end{proof}

\begin{lemma}[Formalization support: the paired subline overlap]
	\label{lem:subline-joint-overlap}
	\lean{MIPStarRE.QPBT.subline_joint_overlap_near_one_at}
	\leanok
	\uses{def:paired-subline-overlap, def:direct-subline-law}
	There is a constant $c>0$ such that, for every projective strategy with
	error $\eps$, every joint point family with error $\delta_Q$, every paired
	line family with error $\delta_P$, every subline witness, and every directed
	opposite placement, the average paired subline overlap differs from $1$ by at
	most
	\[
		c\left(\delta_Q^{1/2}+\sqrt{m}\,
			\bigl(\delta_P^{1/4}+\delta_Q^{1/4}+\eps^{1/4}\bigr)\right).
	\]
	This is the first-route estimate of the proof of
	Lemma~\ref{lem:qld-4-13}, stated for both opposite placements. It compares
	real parts of overlaps; it does not assert the complex modulus comparisons of
	Lemmas~\ref{lem:claim-17-1}--\ref{lem:claim-17-3} over the source
	distribution, which remain open in their complex-modulus, source-law forms. See
	\cite{gap:qpbt_combined-lines-error-term}.
\end{lemma}
\begin{proof}
	\leanok
	\uses{lem:paired-subline-overlap-estimates}
	Chain the paired-to-ordered overlap gap, the removal of the $X$ factor, and
	the near-one estimate for the remaining $Z$ term, each taken over the
	directly indexed subline law, and collect the resulting quarter powers with
	the scalar estimate of that chaining.
\end{proof}

\begin{lemma}[Obstruction to absorbing the first combined-line estimate]
	\label{lem:qld-combining-quarter-power-obstruction}
	\lean{MIPStarRE.QPBT.not_exists_combining_quarter_power_bound}
	\leanok
	\uses{def:admissible}
	There is no function $f:\mathbb R^2\to\mathbb R$ satisfying
	$0\leq f(x,y)\leq C(x^r+y^s)$ on the closed nonnegative quadrant,
	for constants $C\geq 1$ and
	$r,s>0$, which satisfies
	\[
		\min\{1,m\eps^{1/4}\}\leq f(m^2\eps,md/q)
	\]
	for every admissible $(q,m,d)$ and every $\eps\geq0$.
	This is a scalar obstruction to the first proof route in
	Remark~\ref{rem:qld-4-13-source-defects}, not a counterexample to
	Lemma~\ref{lem:qld-4-13}.
\end{lemma}
\begin{proof}
	\leanok
	Set $m_n=2^{2n+1}$, $d_n=1$, $q_n=m_n^3$, and $\eps_n=m_n^{-4}$.
	These parameters are admissible: $q_n=2^{6n+3}$ has odd exponent and
	$m_n\mid q_n$. Both arguments of $f$ equal $m_n^{-2}$, so the assumed
	polynomial upper bound tends to zero. But
	$\min\{1,m_n\eps_n^{1/4}\}=1$ for every $n$, a contradiction.
\end{proof}

\begin{example}[The extended-line construction given a joint point family]
	\label{lem:qld-4-13-established-given-points}
	\lean{MIPStarRE.QPBT.exists_extendedLinesWitness_established_ofPointsWitness}
	\leanok
	\uses{lem:qld-4-10, lem:combined-lines-given-points, def:line,
		def:line-point-dist, def:consistency, def:bracket, def:expanded-state,
		def:ideg-deg-polynomials}
	This formalization-only auxiliary is not a named statement of the paper, and it
	does not claim the printed error form of Lemma~\ref{lem:qld-4-13}, which remains
	an open source gap tracked by GitHub issue 598.
	Fix a function $\delta_Q(\eps)=\poly(\eps)$. There are a constant $C>0$ and a
	function $\delta_{\mathrm{combine}}(\eps,md/q)=\poly(\eps,md/q)$ with the
	following property. Let a projective setting with error $\eps$ be given,
	together with a complete projective joint point family satisfying all the
	conclusions of Lemma~\ref{lem:qld-4-10} with error $\delta_Q(\eps)$. Then,
	relative to that same point family, there is an extended-line witness whose
	axis-degree property and both completed-answer consistency bounds hold with
	error
	\[
		C m\,\delta_{\mathrm{combine}}\!\left(\eps,\frac{md}{q}\right).
	\]
	The point family is an explicit hypothesis, so this auxiliary is separate from
	Lemma~\ref{lem:qld-4-13} and does not prove that source lemma. It shares the
	three replacements described in Lemma~\ref{lem:qld-4-13-established}: the
	directly indexed question carrier, the directly indexed line-point law, and the
	comparison on the completed answer alphabet $\F_q\cup\{\bot\}$.
	The hypothesis is the polynomially controlled family, not an arbitrary supplied
	point measurement at an arbitrary error: the unrestricted form is false, by the
	deterministic quadratic point answers exhibited in
	\cite{gap:qpbt_combined-lines-error-term}.
\end{example}
\begin{proof}
	\leanok
	\uses{lem:combined-lines-given-points,
		lem:direct-subline-witness-exists,
		def:qld-conditional-extended-line-measurement,
		lem:qld-combined-polynomial-axis-degree,
		lem:qld-extended-line-overlap-comparison,
		thm:qld-direct-line-real-overlap,
		lem:consistency-defect-integrand-le-one,
		thm:qld-combining-polynomial-bound}
	Apply Lemma~\ref{lem:combined-lines-given-points} to the supplied point family.
	The direct subline witness of Lemma~\ref{lem:direct-subline-witness-exists}
	supplies the conditional extended-line measurement. Its axis-degree theorem
	proves that no outcome above degree $d$ occurs on an axis-parallel extended
	line, while the extended-overlap comparison and
	Theorem~\ref{thm:qld-direct-line-real-overlap} give both opposite-placement
	consistency estimates. The independent unit bound handles large errors, and
	Theorem~\ref{thm:qld-combining-polynomial-bound} collects the remaining terms
	into $m\cdot\poly(\eps,md/q)$.
\end{proof}

\begin{lemma}[Joint projected law of the concrete direct sub-line sampler]
\label{lem:qld-subline-joint-mixture}
\lean{MIPStarRE.QPBT.subLineJointProjection, MIPStarRE.QPBT.subLineDist_map_joint}
\lean{MIPStarRE.QPBT.avgOver_subLineJointProjection_le}
\leanok
\uses{def:ith-restricted-line, lem:restricted-line-mixture-bounds}
This is a formalization-only assertion for the directly indexed law
$D_{\mathrm{dir}}$ defined by \texttt{subLineDist}, with its concrete choices
of source indices. Choose a line kind $v$ uniformly from the two kinds and
$k$ uniformly from $\{0,\ldots,2m+1\}$. Put $i_X(k)=k$ for $k<m$ and
$i_X(k)=0$ otherwise; put $i_Z(k)=k-m$ for $m\leq k<2m$ and $i_Z(k)=0$
otherwise. Let $\mu$ be the law of $(v,i_X(k),i_Z(k))$.
Sample $(\ell,\ell_X,\ell_Z)$ from $D_{\mathrm{dir}}$ and a uniform affine
parameter $t$, and set $u=\ell(t)$. Then
\[
 \operatorname{Law}((\ell_X,u_X),(\ell_Z,u_Z))
 =\sum_{v,i,j}\mu(v,i,j)\bigl(D_{v,i}\times D_{v,j}\bigr).
\]
In particular, for every nonnegative function $F$ of two source line-point
pairs,
\[
 \mathbb{E}_{D_{\mathrm{dir}},t}F((\ell_X,u_X),(\ell_Z,u_Z))
 \leq 4m^2\mathbb{E}_{D_{\mathrm{Line}}\times D_{\mathrm{Line}}}F.
\]
Zero directions and singleton lines are included. The assertion concerns
this concrete direct law, not every distribution with the separate marginal
properties of Lemma~\ref{lem:qld-sublines}.
\end{lemma}
\begin{proof}
\leanok
For fixed $v,k$, replacing the original uniform extended point by a uniform
point of its line preserves its uniform law and both source-line
representatives. Its two coordinate blocks, the two direction blocks, and
the two source seeds then form independent restricted samples. Averaging
this joint identity over $v,k$ gives the displayed law. Apply the
nonnegative restricted-product estimate of
Lemma~\ref{lem:restricted-line-mixture-bounds} in each component.
\end{proof}

\begin{remark}[Conditional direct sub-line consistency]
\label{rem:qld-subline-joint-consistency}
In a projective setting, suppose a point witness $Q$ and a joint line
witness $T$ are supplied, with completed-answer consistency defect at most
$\delta_P$ over two independent source line-point samples. For any directed
pair of opposite expanded placements, their completed evaluations at
$(u_X,u_Z)$, sampled from the concrete law $D_{\mathrm{dir}}$ above, have
consistency defect at most $4m^2\delta_P$.
This conditional auxiliary proves the distribution comparison used in the
second route at source lines 1241--1245. Construction of $T$ with a suitable
error, and the error form of Lemma~\ref{lem:qld-4-13}, remain separate
obligations.

The integrand defining the completed consistency defect is nonnegative
because the two measurements act on opposite registers. Apply
Lemma~\ref{lem:qld-subline-joint-mixture} to that integrand and use the
assumed consistency bound on $T$.
\end{remark}

\begin{lemma}[Formalization support for Lemma~\ref{lem:qld-4-13}]
	\label{lem:qld-4-13-established}
	\lean{MIPStarRE.QPBT.exists_extendedLinesWitness_established}
	\leanok
	\uses{lem:qld-4-10, def:line, def:line-point-dist, def:consistency,
		def:bracket, def:expanded-state, def:ideg-deg-polynomials}
	This formalization-only auxiliary is not a named statement of the paper.
	It records a directly indexed auxiliary with the numerical error estimate
	$C \cdot m \cdot \poly(\eps, md/q)$ obtained from the first route in the
	printed proof, for a universal constant $C > 0$. It does not claim either
	the printed error form $\poly(m^2\eps, md/q)$ or the source question carrier,
	line-point law, and field-valued answer sum of Lemma~\ref{lem:qld-4-13}; see
	Remark~\ref{rem:qld-4-13-source-defects}.
	More precisely, there exist a polynomially controlled function
	$\delta_Q:\mathbb R\to\mathbb R$, a constant $C>0$, and a polynomially
	controlled function
	$\delta_{\mathrm{combine}}:\mathbb R^2\to\mathbb R$ such that, for every admissible
	$(q,m,d)$ and every projective strategy passing with error $\eps$, there is a
	complete projective joint point family satisfying all the conclusions of
	Lemma~\ref{lem:qld-4-10} with error $\delta_Q(\eps)$ and, relative to that
	same point family, an extended-line witness whose axis-degree property and
	both completed-answer consistency bounds hold with error
	\[
		C m\,\delta_{\mathrm{combine}}\!\left(\eps,\frac{md}{q}\right).
	\]
	Within the directly indexed, completed-answer setting described below, the
	quantitative change from the printed estimate is the error form. Compared
	with Lemma~\ref{lem:qld-4-13}, the auxiliary construction also differs in
	three respects. Its extended questions are
	drawn from the directly indexed line space of
	\cite{gap:qpbt_ld-dimension-divisibility} rather than from the lines of
	$\F_q^{2m+2}$ used above; its average is taken over the corresponding
	directly indexed line-point law rather than over the line-point
	distribution of Definition~\ref{def:line-point-dist} instantiated at
	dimension $2m+2$, whose definition is blocked by the divisibility
	obstruction of Remark~\ref{rem:qld-4-7-divisibility}; and both sides of
	its two consistency displays are compared on the completed answer
	alphabet $\F_q \cup \{\bot\}$, whose extra outcome collects the line
	polynomials that admit no evaluation at the sampled point, rather than
	under the answer summation over $a \in \F_q$ used above. Thus the auxiliary
	construction does not establish the source formulation until the question carrier
	and law are related by the transport obligation of
	\cite{gap:qpbt_ld-dimension-divisibility} and the completed answer comparison
	is related to the source sum by
	\cite{gap:qpbt_combined-lines-error-term}.
\end{lemma}
\begin{proof}
\leanok
\uses{lem:qld-4-10, lem:qld-4-13-established-given-points}
Construct the polynomially controlled joint point family of
Lemma~\ref{lem:qld-4-10}, and apply
Lemma~\ref{lem:qld-4-13-established-given-points} to that same family. This
proves the stated directly indexed auxiliary. Its comparison with the printed source formulation remains subject
to the separate question-law and completed-answer issues in
\cite{gap:qpbt_ld-dimension-divisibility} and
\cite{gap:qpbt_combined-lines-error-term}.
\end{proof}

% This lemma is not a statement of the source; it records the resampling
% identity behind the coefficient-consistency estimate of the auxiliary above.
\begin{lemma}[Formalization support: coefficient consistency under the directly indexed line laws]
	\label{lem:direct-coefficient-resampling}
	\lean{MIPStarRE.QPBT.directCoefficientEvaluationDefect}
	\lean{MIPStarRE.QPBT.directJointCoefficientEvaluationDefect}
	\lean{MIPStarRE.QPBT.consistencyDefect_directALine_coefficients_le_joint_add}
	\lean{MIPStarRE.QPBT.consistencyDefect_directDLine_coefficients_le_joint_add}
	\leanok
	\uses{def:consistency, def:ld-game, def:admissible-size,
		def:ideg-deg-polynomials}
	Fix directly indexed parameters $D = (q,m,d,k)$ as in
	Lemma~\ref{lem:qld-4-13-established}, a degree bound $\mathrm{deg}$, a
	state $\psi$ with $\|\psi\| = 1$, and two measurement families $A$ and $B$
	indexed by directly indexed lines whose outcomes are the coefficient
	vectors of univariate polynomials of degree at most $\mathrm{deg}$.

	For a line $\ell$ and a scalar $t$, the \emph{evaluated coefficient
	defect} is the off-diagonal consistency integrand of $A^{\ell}$ against
	$B^{\ell}$ on $\psi$, after postcomposing each with the evaluation of its
	coefficient vector at $t$. Its \emph{joint} form at a line-point sample
	$(\ell, u)$ is that quantity at the affine parameter of $u$ on $\ell$ when
	the direction of $\ell$ is nonzero, and the full uniform average over
	$t \in \F_q$ when the direction is zero; in the degenerate case the point
	does not determine a parameter, so no distinguished parameter is selected
	and the whole parameter average is retained.

	Then, for the axis-parallel line marginal and for the diagonal line
	marginal of the directly indexed line-point law alike, the consistency
	defect of $A$ against $B$ over that marginal is at most the average of the
	joint evaluated coefficient defect over the original line-point law, plus
	$\mathrm{deg}/q$. The diagonal statement carries no nonzero-direction
	restriction.
\end{lemma}
\begin{proof}
	\leanok
	\uses{lem:sandwich-codeword-defect, lem:coefficient-polynomial-agreement,
		def:consistency, def:ld-game, def:admissible-size,
		def:ideg-deg-polynomials}
	Comparing the coefficient outcomes with their evaluations at an independent
	uniform parameter is the codeword step of
	Lemma~\ref{lem:sandwich-codeword-defect}. Its collision hypothesis holds
	with $\eps = \mathrm{deg}/q$: two distinct coefficient vectors of degree at
	most $\mathrm{deg}$ agree at a uniformly random parameter with probability
	at most $\mathrm{deg}/q$, which is the uniform-parameter form of the root
	count of Lemma~\ref{lem:coefficient-polynomial-agreement}.
	The resulting product law of a line marginal with an independent uniform
	parameter is then reexpressed over the original line-point law: for a
	nonzero direction the representative-parameter map identifies the two
	averages termwise, and for a zero direction both sides are the uniform
	parameter average of the same integrand. That reexpression is carried in
	Lean by the axis and diagonal resampling identities of
	\texttt{DirectLowDegree/Transport/LineResampling.lean} and the evaluated
	defect identities of
	\texttt{DirectLowDegree/ResampledCoefficientConsistency.lean}. These are
	the directly indexed analogues of
	Corollary~\ref{cor:line-point-parameter-resampling} and carry no blueprint
	node of their own.
\end{proof}

The measurements of Lemma~\ref{lem:qld-4-12} and Lemma~\ref{lem:qld-4-13}
constitute a strategy for the classical low individual degree game in $2m+2$
variables. Applying the quantum soundness of that game yields a single
projective measurement whose outcomes are pairs of global polynomials.

\begin{definition}[Restriction and completion of combined polynomial outcomes]
	\label{def:qld-pair-restriction-completion}
	\lean{MIPStarRE.QPBT.recoverCombinedPair}
	\lean{MIPStarRE.QPBT.combinedPolyEmbedding}
	\lean{MIPStarRE.QPBT.combinedPairSubmeasurement}
	\lean{MIPStarRE.QPBT.completedPairLabel}
	\lean{MIPStarRE.QPBT.completedPairMeasurement}
	\leanok
	\uses{def:combine-map, def:polyfunc, def:povm-conventions}
	For admissible $(q,m,d)$, write
	$j(g_X,g_Z)=\mathrm{combine}_{g_X,g_Z}$. Given a complete measurement
	$T$ on bounded polynomials in $2m+2$ variables, define the restricted
	sub-measurement $R_{g_X,g_Z}=T_{j(g_X,g_Z)}$. For a specified pair $p_0$,
	complete it by relabeling every outcome outside the range of $j$ by $p_0$.
	The resulting measurement is denoted by $C$.
\end{definition}

\begin{lemma}[Properties of polynomial-pair restriction and completion]
	\label{lem:qld-combined-pair-completion}
	\lean{MIPStarRE.QPBT.recoverCombinedPair_combinePoly}
	\lean{MIPStarRE.QPBT.combinePoly_injective}
	\lean{MIPStarRE.QPBT.combinedPolyEmbedding_eval}
	\lean{MIPStarRE.QPBT.combinedPairSubmeasurement_projective}
	\lean{MIPStarRE.QPBT.combinedPairSubmeasurement_missing}
	\lean{MIPStarRE.QPBT.completedPairLabel_combinedPolyEmbedding}
	\lean{MIPStarRE.QPBT.completedPairLabel_of_not_mem_range}
	\lean{MIPStarRE.QPBT.completedPairMeasurement_effect}
	\lean{MIPStarRE.QPBT.completedPairMeasurement_projective}
	\leanok
	\uses{def:qld-pair-restriction-completion}
	The combining map $j$ is injective and satisfies
	$j(g_X,g_Z)(x,z,\alpha,\beta)=\alpha g_X(x)+\beta g_Z(z)$.
	If $T$ is projective, then the restricted sub-measurement $R$ is projective,
	and its missing operator is
	\[
		\Id-\sum_{g_X,g_Z}R_{g_X,g_Z}=\sum_{g\notin\operatorname{range}j}T_g.
	\]
	The completion satisfies $C_p=R_p+1_{p=p_0}(\Id-\sum R)$ and is projective.
	No bound on the state mass of the missing operator, or on consistency with
	point measurements, is asserted here.
\end{lemma}
\begin{proof}
	\leanok
	\uses{def:qld-pair-restriction-completion,
		lem:sandwich-postprocess-projective}
	Substitute $(\alpha,\beta)=(1,0)$ and set the $Z$ variables to zero
	to recover $g_X$ as a formal polynomial; the opposite substitution
	recovers $g_Z$. Thus distinct pairs label distinct effects, proving
	the sub-measurement assertion. Partitioning the outcomes into
	$\mathcal G$ and its complement gives the operator identity.
	Relabel each outcome in $\mathcal G$ by its unique pair, and all
	remaining outcomes by $p_0$. This is precisely the stated completion,
	and projectivity follows from projective postprocessing.
\end{proof}

\begin{definition}[Completion of the actual singleton outcomes]
	\label{def:qld-direct-pair-completion}
	\lean{MIPStarRE.QPBT.ExtendedLineGame.directPolynomialEquiv}
	\lean{MIPStarRE.QPBT.ExtendedLineGame.directCombinedEmbedding}
	\lean{MIPStarRE.QPBT.ExtendedLineGame.canonicalPolynomialMeasurement}
	\lean{MIPStarRE.QPBT.ExtendedLineGame.directPairMeasurement}
	\leanok
	\uses{def:qld-pair-restriction-completion}
	Taking the unique component of a singleton polynomial tuple and applying the
	canonical field isomorphism gives an equivalence $e$ with bounded extended
	polynomials over $\F_q$. Put $j_{\rm dir}=e^{-1}\circ j$, relabel the actual
	rounded measurement by $e$, and apply
	Definition~\ref{def:qld-pair-restriction-completion}.
\end{definition}

\begin{theorem}[Properties of the actual polynomial-pair completion]
	\label{thm:qld-direct-pair-completion-properties}
	\lean{MIPStarRE.QPBT.ExtendedLineGame.canonicalPolynomialMeasurement_effect}
	\lean{MIPStarRE.QPBT.ExtendedLineGame.directPairMeasurement_projective}
	\lean{MIPStarRE.QPBT.ExtendedLineGame.mem_range_directCombinedEmbedding}
	\lean{MIPStarRE.QPBT.ExtendedLineGame.directCombinedEmbedding_read}
	\leanok
	\uses{def:qld-direct-pair-completion, lem:qld-combined-pair-completion}
	The canonical relabeling has effect $R_{e^{-1}(h)}$ at $h$, and the completed
	measurement is projective whenever the rounded measurement is projective.
	Membership in the range of $j_{\rm dir}$ is precisely equality, after field
	transport, to $\mathrm{combine}_{g_X,g_Z}$ for a bounded pair. Moreover, the
	direct readout of $j_{\rm dir}(g_X,g_Z)$ is
	$\alpha g_X(x)+\beta g_Z(z)$ under the joint question equivalence.
\end{theorem}
\begin{proof}
	\leanok
	Support inclusion transports the individual-degree certificate in both
	directions. The field equivalence and the combining-map evaluation identity
	give the effect, range, and readout formulas. Projectivity follows from
	projective postprocessing and completion.
\end{proof}

\begin{lemma}[Agreement after completing the polynomial-pair measurement]
	\label{lem:qld-completed-pair-agreement}
	\lean{MIPStarRE.QPBT.completedPairMeasurement_evaluated_defect_le}
	\lean{MIPStarRE.QPBT.completedPairMeasurement_consistencyDefect_le}
	\leanok
	\uses{lem:qld-combined-pair-completion, def:consistency}
	This is a formalization-only estimate for the last paragraph of the proof
	of Lemma~\ref{lem:qld-4-7}. Let $T$ be any complete measurement on extended
	polynomials, $R_p=T_{\mathrm{combine}_p}$ its restriction to polynomial
	pairs, and $C$ its completion at a fixed pair $p_0$. Projectivity is not
	required for this estimate. Let $U$ be a nonempty finite question set,
	$e_u$ a map from polynomial pairs to a finite answer alphabet, and
	$\{B^u_a\}_a$ a complete measurement on the opposite tensor factor.
	For every vector $\psi$, the consistency defect satisfies
	\[
		\E_{u\in U}\sum_{a\ne b}
		\langle\psi|C_{[e_u=a]}\otimes B^u_b|\psi\rangle
		\le \|\psi\|^2-
		\E_{u\in U}\sum_p
		\langle\psi|R_p\otimes B^u_{e_u(p)}|\psi\rangle.
	\]
	The same inequality holds pointwise for each $u$. It does not assert a
	lower bound for the retained overlap on the right-hand side.
\end{lemma}
\begin{proof}
	\leanok
	\uses{lem:qld-combined-pair-completion}
	Each difference $C_p-R_p$ is positive semidefinite. Tensoring it with
	$B^u_{e_u(p)}$ therefore gives a nonnegative quadratic form, so
	completion increases the diagonal overlap. Since $C$ and $B^u$ are
	complete, their off-diagonal overlap equals $\|\psi\|^2$ minus their
	diagonal overlap. Postprocessing groups the latter by the fibers of $e_u$.
	Apply the resulting inequality and average over $u$.
\end{proof}

\begin{lemma}[Formalization support for uniform specialization in a block of variables]
	\label{lem:qpbt-block-specialization}
	\lean{MIPStarRE.QPBT.totalDegree_combinedCoef_le_totalDegree}
	\lean{MIPStarRE.QPBT.totalDegree_combinedRestrict_le_totalDegree}
	\lean{MIPStarRE.QPBT.exists_nonzero_combinedCoef_of_depends_on_block}
	\lean{MIPStarRE.QPBT.exists_block_specialization_exceptional_coefficient}
	\lean{MIPStarRE.QPBT.block_specialization_weighted_avg_le}
	\leanok
	This is not a named statement of \cite{Ji2020MIPStar}; it isolates the
	specialization estimate between the paper equations labelled
	\texttt{eq:qld-g-2} and \texttt{eq:qld-g-prime-xpt-bound}, which the source
	carries out inside its proof of Lemma~\ref{lem:qld-4-7} with
	$D = (2m+2)d$ and with the weights given by the expectations of a complete
	measurement.

	Let $p$ be a polynomial on $\F_q^m \times \F_q^k$, written $p(z,x)$, of
	total degree at most $D$, and assume that $p$ is not a polynomial in $z$
	alone.  Write $p = \sum_\mu c_\mu(z)\,x^\mu$.  Then there is an exponent
	$\mu \neq 0$ with $c_\mu \neq 0$ such that $c_\mu$ has total degree at most
	$D$,
	\[
		\Pr_{z \in \F_q^m}[c_\mu(z) = 0] \le \frac{D}{q}\,,
	\]
	and, for every $z$ with $c_\mu(z) \neq 0$, simultaneously for every
	$b \in \F_q$,
	\[
		\Pr_{x \in \F_q^k}[p(z,x) = b] \le \frac{D}{q}\,.
	\]
	Consequently, for any weights $w(z,b) \ge 0$ with
	$\sum_{b \in \F_q} w(z,b) = M$ for every $z$,
	\[
		\E_{z \in \F_q^m} \sum_{b \in \F_q}
		\Pr_{x \in \F_q^k}[p(z,x) = b] \, w(z,b)
		\le \frac{2D}{q} \, M\,.
	\]
\end{lemma}
\begin{proof}
	\leanok
	\uses{lem:schwartz-zippel-total-degree}
	If $c_\mu = 0$ for every $\mu \neq 0$ then $p = c_0(z)$, contrary to the
	assumption; fix $\mu \neq 0$ with $c_\mu \neq 0$.  Splitting the exponent
	of each monomial of $p$ into its $z$-part and its $x$-part shows that
	neither extracting the coefficient of an $x$-monomial nor substituting a
	value for $z$ raises the total degree, so $c_\mu$ and every
	$p(z,\cdot)$ have total degree at most $D$.  Applying
	Lemma~\ref{lem:schwartz-zippel-total-degree} to $c_\mu$ and the zero
	polynomial bounds $\Pr_z[c_\mu(z) = 0]$ by $D/q$.  If $c_\mu(z) \neq 0$
	then $p(z,\cdot)$ differs from the constant polynomial $b$ for every
	$b$, because its coefficient at $\mu \neq 0$ equals $c_\mu(z)$, so the
	same lemma bounds $\Pr_x[p(z,x) = b]$ by $D/q$, uniformly in $b$.
	Hence $\Pr_x[p(z,x) = b] \le \mathbf{1}[c_\mu(z) = 0] + D/q$ for all $z$
	and $b$.  Multiplying by $w(z,b) \ge 0$, summing over $b$, using
	$\sum_b w(z,b) = M$ and averaging over $z$ gives
	$(\Pr_z[c_\mu(z) = 0] + D/q)\,M \le (2D/q)\,M$.  One exceptional set
	serves every answer, so the answer alphabet contributes no factor.
	Nonconstancy is used as a property of the polynomial and not of the
	function it induces, so no bound on $D$ in terms of $q$ is needed.
\end{proof}

\begin{lemma}[The global polynomial-pair measurement]
	\label{lem:qld-4-7}
	\lean{MIPStarRE.QPBT.Poly, MIPStarRE.QPBT.PolyPair, MIPStarRE.QPBT.evalAt}
	\lean{MIPStarRE.QPBT.GlobalPairWitness}
	\lean{MIPStarRE.QPBT.exists_globalPairWitness}
	\leanok
	\uses{def:polyfunc, def:expanded-state, def:expanded-point-measurement,
		def:consistency, def:bracket, def:ideg-deg-polynomials}
	% Statement as in the source (references/qpbt-paper/
	% 14_analysis_of_the_pauli_basis_test.tex, lem:qld-4-7). The proof
	% instantiates the classical low individual degree game, and Theorem
	% lem:ld-soundness, at dimension 2m+2, which def:ld-game admits only
	% when the dimension divides q; the source performs the instantiation
	% without securing this condition. See rem:qld-4-7-divisibility.
	There exists a function
	$\delta_S(\eps,m,d,q) = a(md)^a(\eps^b + q^{-b} + 2^{-bmd})$, for some
	universal constants $a > 1$ and $0 < b < 1$, such that the following
	holds. For
	$\hilb \in \{\hilb_{\mathsf{A}}, \hilb_{\mathsf{B}}\}$ there
	exists a projective measurement $\{\hat{S}_{g_X,g_Z}\}$ acting on
	$\hilb \ot (\C^q)^{\ot M}$, with outcomes consisting of pairs
	$(g_X, g_Z)$ of polynomials, each in $\mathrm{ideg}_{d,m}(\F_q)$, such
	that for $W \in \{X, Z\}$, on average over uniformly random
	$u \in \F_q^m$,
	\begin{align}
		\big(\hat{S}_{[\mathrm{eval}_u(\cdot_W) = a]}
		\big)_{\mathsf{A}\mathsf{A}'}
		&\simeq_{\delta_S}
		\big(\hat{M}^{(\mathrm{Point},W),u}_{a}
		\big)_{\mathsf{B}\mathsf{A}''}\,,
		\label{eq:qld-s-point-con-alice} \\
		\big(\hat{S}_{[\mathrm{eval}_u(\cdot_W) = a]}
		\big)_{\mathsf{B}\mathsf{B}'}
		&\simeq_{\delta_S}
		\big(\hat{M}^{(\mathrm{Point},W),u}_{a}
		\big)_{\mathsf{A}\mathsf{B}''}\,,
		\label{eq:qld-s-point-con-bob}
	\end{align}
	where $\mathrm{eval}_u(\cdot_W)$ denotes the evaluation at the point $u$
	of the first outcome $g_X$ if $W = X$, and of the second outcome $g_Z$
	if $W = Z$.
\end{lemma}
\begin{proof}
\leanok
\uses{lem:qld-4-13-established, thm:qld-completed-pair-actual-error,
  thm:qld-actual-rounded-error-bound}
Lemma~\ref{lem:qld-4-13-established} supplies the actual point and directly
indexed line witnesses with errors $\delta_Q=p(\eps)$ and
$\delta_L=Cm\,g(\eps,md/q)$, for polynomially small $p,g$ and universal $C>0$.
Theorem~\ref{thm:qld-completed-pair-actual-error} constructs both projective
polynomial-pair measurements, with all four point-consistency defects at most
\[
 D=8(4\eta+8\delta_Q)+\frac{12md+4d+14}{q},\qquad
 \eta=\delta+\sqrt{220\delta^{1/4}}+2\sqrt{2\delta},
\]
where
\[
 \delta=\delta_{\rm ld}\left(a,b,
 3\left(\sqrt{\delta_Q+\delta_L}+md/q\right),q,2m+2,d,1\right).
\]
Each defect is also at most one, since both measurements are complete and
the state is normalized. Theorem~\ref{thm:qld-actual-rounded-error-bound}
bounds $\min\{1,D\}$ by the required universal error function for every
nonnegative error. The same two measurements therefore satisfy both
displayed relations for both bases.

This proof uses soundness for the directly indexed game. The separate
comparison with the source's seed-bearing auxiliary game, and the
discrepancies documented in \cite{gap:qpbt_ld-dimension-divisibility} and
\cite{gap:qpbt_combined-lines-error-term}, remain unchanged.
\end{proof}

\begin{lemma}[Projective refinement of an operator comparison]
\label{lem:qld-projective-refinement-ordered}
\lean{MIPStarRE.QPBT.projective_refinement_ordered_dist_le}
\leanok
This is an auxiliary quantitative form of the insertion calculation in the
paper equations labelled \texttt{eq:qld-g-42} and \texttt{eq:qld-g-43}
in \cite{Ji2020MIPStar}. Let $R=\{R_g\}$ be a
complete projective measurement, let $e_x(g)$ be a finite outcome readout,
and let $A^x,B^x$ be complete measurements on the same finite-dimensional
space. Let $T^x_a$ be arbitrary operators. For a finite distribution $\mu$
and a state $\psi$, suppose that
\[
 \operatorname{Cons}_\mu(R^{e_x},B^x)\leq\eta,\qquad
 \mathbb E_\mu\sum_a\|(B^x_a-A^x_a)\psi\|^2\leq\delta_0,\qquad
 \mathbb E_\mu\sum_a\|(A^x_a-T^x_a)\psi\|^2\leq\delta_1.
\]
Then
\[
 \mathbb E_\mu\sum_g\|R_g(I-T^x_{e_x(g)})\psi\|^2
 \leq 4\eta+4\delta_0+4\delta_1.
\]
\end{lemma}
\begin{proof}
\leanok
\uses{fact:agreement, fact:add-a-proj2}
Projectivity gives $R_gR^{e_x}_{e_x(g)}=R_g$. Agreement and projective
refinement bound the insertion of $B^x_{e_x(g)}$ by $2\eta$.
The squared-distance triangle inequality bounds the comparison from $B$
to $T$ by $2\delta_0+2\delta_1$. Left multiplication by $R_g$, with
the full outcome sum retained, contracts this bound. A second application
of the triangle inequality gives the displayed constants.
\end{proof}

\begin{lemma}[Both ordered products on opposite expanded blocks]
\label{lem:qld-opposite-ordered-products}
\lean{MIPStarRE.QPBT.CombinedPointsWitness.ordered_dist_le}
\leanok
Fix a projective Pauli strategy and a combined-point witness with error
$\delta_Q$. Let $p_1,p_2$ be any directed opposite pair of expanded register
placements. Suppose a complete projective measurement $R$ on $p_1$, read
by $e_{x,z,\alpha,\beta}$, has consistency defect at most $\eta$ against
$\hat Q^{x,z,\alpha,\beta}$ on $p_2$, on the actual expanded state.
For either order $XZ$ or $ZX$, define
\[
 T^{XZ}_{x,z,\alpha,\beta,c}
 =\sum_{\alpha r+\beta s=c}\hat M^{(\mathrm{Point},X),x}_r
                              \hat M^{(\mathrm{Point},Z),z}_s,
 \qquad
 T^{ZX}_{x,z,\alpha,\beta,c}
 =\sum_{\alpha r+\beta s=c}\hat M^{(\mathrm{Point},Z),z}_s
                              \hat M^{(\mathrm{Point},X),x}_r.
\]
Then, for each order $o\in\{XZ,ZX\}$,
\[
 \mathbb E_{x,z,\alpha,\beta}\sum_g
 \|(R_g)_{p_1}(I-(T^o_{x,z,\alpha,\beta,e_{x,z,\alpha,\beta}(g)})_{p_2})
       \hat\psi\|^2\leq4\eta+8\delta_Q.
\]
This auxiliary uses the actual point-comparison error $\eta$, including any
rounding loss. It does not identify that error with the unrounded
low-degree soundness error.
\end{lemma}
\begin{proof}
\leanok
\uses{lem:qld-projective-refinement-ordered, lem:qld-4-12}
Use the opposite-block $Q$ as $B$ and the same-block $Q$ as $A$.
The two comparison errors in Lemma~\ref{lem:qld-projective-refinement-ordered}
are both bounded by $\delta_Q$, separately for the two orders.
\end{proof}

\begin{definition}[Extended-polynomial ordered error]
\label{def:qld-extended-polynomial-ordered-error}
\lean{MIPStarRE.QPBT.ExtendedLineGame.extendedPolynomialRead}
\lean{MIPStarRE.QPBT.ExtendedLineGame.extendedPolynomialOrderedError}
\leanok
For the directly indexed extended game at $k=1$, its polynomial outcome is
a singleton tuple $g$. Identify its scalar field with the Pauli scalar
field and evaluate its single polynomial at $(x,z,\alpha,\beta)$; denote
the resulting readout by $e_{x,z,\alpha,\beta}(g)$.
For a measurement $R$, let $E^o_{p_1,p_2}(R)$ be the full averaged
squared-norm sum in Lemma~\ref{lem:qld-opposite-ordered-products} with
this readout, for an order parameter $o\in\{XZ,ZX\}$. The outcome sum is over
every singleton polynomial tuple, without a support restriction.
\end{definition}

\begin{theorem}[Formalization support: direct soundness for arbitrary strategies]
\label{thm:qld-direct-soundness-any-strategy}
\lean{MIPStarRE.QPBT.exists_direct_ld_soundness_of_k_eq_one_any_strategy}
\leanok
\uses{lem:direct-ld-soundness-k-one, def:consistency, def:ld-game}
This is not a named statement of \cite{Ji2020MIPStar}. There are universal
constants $a\geq1$ and $0<b\leq1$ such that the following holds. Let
$D=(q,n,d,k)$ be directly indexed low-degree parameters with $k=1$, let
$\xi>0$, and let $S$ be an arbitrary, not necessarily projective, strategy for
the directly indexed game with value at least $1-\xi$. There are complete
polynomial-tuple POVMs $G^{\mathrm A},G^{\mathrm B}$ on the original player
spaces for which both point--polynomial consistency defects and the mutual
polynomial consistency defect are at most
$\delta_{\rm ld}(a,b,\xi,q,n,d,1)$. The point readout sends every malformed
answer to the zero tuple, and no projectivity conclusion is asserted for
$G^{\mathrm A},G^{\mathrm B}$.
\end{theorem}
\begin{proof}
\leanok
\uses{lem:direct-ld-soundness-k-one}
Dilate both question-indexed POVM families by Naimark's theorem, apply
Theorem~\ref{lem:direct-ld-soundness-k-one} to the resulting projective
strategy, and compress the two polynomial POVMs at the distinguished ancillary
coordinates. The game value and all three consistency defects are preserved.
\end{proof}

\begin{theorem}[Formalization-only supplied-witness polynomial-tuple consistency]
\label{thm:qld-supplied-direct-polynomial-consistency}
\lean{MIPStarRE.QPBT.ExtendedLineGame.exists_direct_polynomial_measurements_of_supplied_witnesses}
\leanok
\uses{thm:qld-direct-soundness-any-strategy, lem:qld-4-12, lem:qld-4-13,
	def:consistency, def:ld-meas}
This is not a named statement of \cite{Ji2020MIPStar}. Fix admissible parameters
$(q,m,d)$, a projective Pauli strategy, a supplied combined-point witness with
error $\delta_Q$, and a supplied extended-line witness over it with error
$\delta_L$. There are universal constants $a\geq1$ and $0<b\leq1$ and complete
polynomial-tuple POVMs $G^{\mathrm A},G^{\mathrm B}$ on the two original
expanded local spaces such that both point--polynomial consistency defects,
averaged over uniform $u\in\F_q^{2m+2}$, and the mutual polynomial consistency
defect are at most
\[
 \delta_{\rm ld}\!\left(a,b,
  3\left(\sqrt{\delta_Q+\delta_L}+\frac{md}{q}\right),
  q,2m+2,d,1\right).
\]
The point families are the completed-answer readouts of the supplied strategy,
with malformed answers sent to the zero singleton tuple. The witnesses are
assumptions: this theorem neither constructs them nor asserts projectivity,
polynomial separation, or a paired-polynomial measurement. It is only support
for the first paragraph of the proof of Lemma~\ref{lem:qld-4-7}; in particular,
it does not discharge the correspondence or supplied-line construction gaps.
\end{theorem}
\begin{proof}
\leanok
\uses{thm:qld-direct-soundness-any-strategy, lem:qld-4-12, lem:qld-4-13}
The supplied witnesses define the directly indexed strategy at parameters
$(q,2m+2,d,1)$. Its passing error is bounded by
$3(\sqrt{\delta_Q+\delta_L}+md/q)$. Admissibility gives $m,d\geq1$ and $q>0$,
so this error is positive. Apply
Theorem~\ref{thm:qld-direct-soundness-any-strategy}; its three conclusions have
exactly the displayed bound and already act on the original expanded spaces.
\end{proof}

\begin{theorem}[Formalization-only supplied-witness scalar polynomial consistency]
\label{thm:qld-supplied-scalar-polynomial-consistency}
\lean{MIPStarRE.QPBT.ExtendedLineGame.exists_scalar_polynomial_measurements_of_supplied_witnesses}
\leanok
\uses{thm:qld-supplied-direct-polynomial-consistency,
	lem:qld-supplied-scalar-point-measurement, fact:data-processing,
	def:consistency, def:ld-meas}
Under the assumptions of
Theorem~\ref{thm:qld-supplied-direct-polynomial-consistency}, there are ordinary
polynomial POVMs $G^{\mathrm A},G^{\mathrm B}$ on the same original expanded
local spaces. For uniform $u\in\F_q^{2m+2}$, the supplied scalar point POVM for
Alice is consistent with evaluation of $G^{\mathrm B}$, evaluation of
$G^{\mathrm A}$ is consistent with the supplied scalar point POVM for Bob, and
the two polynomial POVMs are mutually consistent. Each of these three defects
is at most
\[
 \delta_{\rm ld}\!\left(a,b,
  3\left(\sqrt{\delta_Q+\delta_L}+\frac{md}{q}\right),
  q,2m+2,d,1\right).
\]
Here the scalar point POVM is the supplied joint point measurement
$\hat Q^{x,z}$ coarse-grained by $(r,s)\mapsto\alpha r+\beta s$, and polynomial
evaluation is transported through the canonical scalar-field identification.
This is a formalization-only auxiliary conditional on the supplied witnesses.
It asserts no projectivity, polynomial separation, paired measurement, or
completion of Lemma~\ref{lem:qld-4-7}.
\end{theorem}
\begin{proof}
\leanok
\uses{thm:qld-supplied-direct-polynomial-consistency,
	lem:qld-supplied-scalar-point-measurement, fact:data-processing}
Take the unique coordinate marginal of each polynomial-tuple POVM. Polynomial
evaluation commutes with this marginal. Apply data processing to the first two
tuple consistency defects and use the exact supplied-point readout identity;
apply data processing to both polynomial families in the third defect. These
postprocessings preserve the displayed direct-soundness bound.
\end{proof}

\begin{theorem}[Formalization support: simultaneous rounding under postprocessing]
\label{thm:qld-simultaneous-postprocessed-rounding}
\lean{MIPStarRE.QPBT.projective_rounding_preserves_postprocessed_consistency}
\leanok
\uses{lem:ortho, def:consistency, def:povm-distance}
Let $\psi$ be a unit vector on a finite-dimensional bipartite space, and let
$A,B$ be complete POVMs with a common finite outcome set whose cross-player
consistency defect is at most $\delta\geq0$. There are complete projective
measurements $P_A,P_B$, chosen before any subsequent question distribution or
postprocessing, whose squared state-dependent distances from $A,B$ are each at
most $220\delta^{1/4}$. Moreover, for every finite probability distribution
$\mu$, every question-dependent outcome map $f$, and every comparison family
$N$, a consistency bound $\eta$ between $f_*A$ and $N$ implies
\[
 \operatorname{Cons}_\mu(f_*P_A,N)
 \leq\delta+\sqrt{220\delta^{1/4}}+2\sqrt{\delta+\eta},
\]
and the analogous implication holds between $N$ and $f_*P_B$.
\end{theorem}
\begin{proof}
\leanok
\uses{lem:ortho, def:consistency, def:povm-distance}
Apply the one-sided projective rounding and consistency transport on each
tensor factor. The two choices are independent of the later distribution,
postprocessing map, comparison family, and comparison error, so they form one
simultaneous pair with the stated bounds.
\end{proof}

\begin{theorem}[Rounded extended-polynomial ordered estimates]
\label{thm:qld-rounded-polynomial-ordered}
\lean{MIPStarRE.QPBT.ExtendedLineGame.exists_rounded_polynomial_ordered_estimates}
\leanok
There are universal direct-soundness constants $a\geq1$ and $0<b\leq1$
with the following property. Fix admissible Pauli parameters, a projective
Pauli strategy, a combined-point witness with error $\delta_Q$, and an
extended-line witness with error $\delta_L$ on the directly indexed
completed-answer domain. Set
\[
 e=3\bigl(\sqrt{\delta_Q+\delta_L}+md/q\bigr),\qquad
 \delta=\delta_{\mathrm{ld}}(a,b,e,q,2m+2,d,1),\qquad
 \eta=\delta+\sqrt{220\delta^{1/4}}+2\sqrt{2\delta}.
\]
There exist complete projective extended-polynomial measurements $R_A,R_B$
on the original expanded local spaces such that, for both $o=XZ,ZX$,
\[
 E^o_{AA',BA''}(R_A)\leq4\eta+8\delta_Q,\qquad
 E^o_{BB',AB''}(R_B)\leq4\eta+8\delta_Q.
\]
The line witness is an explicit input to this auxiliary theorem. Its
domain restrictions remain those of
Remark~\ref{rem:qld-4-13-source-defects}; this theorem does not certify
Lemma~\ref{lem:qld-4-7} or assert concentration on separated polynomials.
\end{theorem}
\begin{proof}
\leanok
\uses{thm:qld-direct-soundness-any-strategy,
  thm:qld-simultaneous-postprocessed-rounding,
  lem:qld-opposite-ordered-products, def:qld-extended-polynomial-ordered-error}
Apply the established direct soundness theorem for arbitrary strategies
at $k=1$. Compression preserves all three conclusions, so the polynomial
consistency error and both evaluated errors are at most $\delta$.
Simultaneous projective rounding therefore gives the displayed $\eta$
for both evaluated comparisons. Read the singleton scalar coordinate,
transport the joint uniform question law by its coordinate equivalence,
and transport the bipartite correlations to the two expanded placements.
Lemma~\ref{lem:qld-opposite-ordered-products} gives both conclusions.
\end{proof}

\begin{definition}[Scalar coefficient expansion]
\label{def:qld-scalar-coefficient-expansion}
\lean{MIPStarRE.QPBT.ExtendedLineGame.baseCoordinateEquiv}
\lean{MIPStarRE.QPBT.ExtendedLineGame.scalarBaseCoordinateEquiv}
\lean{MIPStarRE.QPBT.ExtendedLineGame.scalarBaseQuestionEquiv}
\lean{MIPStarRE.QPBT.ExtendedLineGame.scalarPolynomial}
\lean{MIPStarRE.QPBT.ExtendedLineGame.scalarNonlinearMass}
\leanok
Identify the first $2m$ coordinates with the ordered blocks $(x,z)$ and
the last two with $(\alpha,\beta)$. Identify the direct scalar field with
the Pauli scalar field $F=\mathbb F_q$. For a singleton polynomial outcome
$g$, write its transported polynomial as
\[
 p_g(\alpha,\beta)=\sum_{i,j}c_{g,i,j}(x,z)\alpha^i\beta^j
 \in F[x,z][\alpha,\beta].
\]
For an expanded placement $p$ and a complete measurement $R$, define
\[
 N_p(R)=\sum_{g:\,p_g\ne\alpha r+\beta t\ \text{for all }r,t\in F[x,z]}
   \|(R_g)_p\hat\psi\|^2.
\]
The coefficient polynomials $r,t$ in this definition may depend on both
base blocks. These are auxiliary definitions for the scalar-linearity
argument in the proof of Lemma~\ref{lem:qld-4-7}.
\end{definition}

\begin{theorem}[Faithful scalar expansion and degree bounds]
\label{thm:qld-scalar-expansion-degrees}
\lean{MIPStarRE.QPBT.ExtendedLineGame.scalarPolynomial_injective}
\lean{MIPStarRE.QPBT.ExtendedLineGame.scalarPolynomial_eval}
\lean{MIPStarRE.QPBT.ExtendedLineGame.scalarPolynomial_degrees}
\lean{MIPStarRE.QPBT.ExtendedLineGame.scalarPolynomial_degree_term_le}
\leanok
For admissible parameters, the map $g\mapsto p_g$ is injective and
evaluating $p_g$ after specializing its coefficients at $(x,z)$ gives
exactly $e_{x,z,\alpha,\beta}(g)$. The coordinate splitting is an
equivalence of the full joint sample spaces. Moreover,
\[
 \deg c_{g,i,j}\leq2md,\qquad \deg_{\alpha,\beta}p_g\leq2d.
\]
Writing
$D(p)=\max_{e\in\operatorname{supp}p}\deg(\operatorname{coeff}_e p)
+\max\{1,\deg p\}+1$, with the empty maximum equal to zero, gives
$D(p_g)\leq(2m+2)d+1$.
\end{theorem}
\begin{proof}
\leanok
\uses{def:qld-scalar-coefficient-expansion}
Field transport, variable renaming, and the coefficient expansion are
injective ring maps. The singleton tuple adds no information. Evaluation
commutes with each map. A nonzero coefficient of a coefficient polynomial
is the coefficient of the corresponding combined monomial of $g$.
Every exponent is at most $d$; summing over the $2m$ base coordinates
or the two scalar coordinates gives the degree bounds. Use $d\geq1$
for the maximum in $D$.
\end{proof}

\begin{definition}[Ordered and diagonal indicators]
\label{def:qld-polynomial-indicators}
\lean{MIPStarRE.QPBT.PolynomialImageBounds.orderedIndicator}
\lean{MIPStarRE.QPBT.PolynomialImageBounds.diagonalIndicator}
\leanok
For complete measurements $X,Z$ with outcomes in a finite field $F$, set
\[
 T_{\alpha,\beta,a}=\sum_{\alpha b+\beta c=a}X_bZ_c,
 \qquad H_{\alpha,\beta,a}=\sum_{\alpha b+\beta c=a}Z_cX_bZ_c.
\]
These operators are the ordered indicator and diagonal sandwich in the
scalar-linearity argument of Lemma~\ref{lem:qld-4-7}.
\end{definition}

\begin{theorem}[Projective outcome mass]
\label{thm:qld-projective-outcome-mass}
\lean{MIPStarRE.QPBT.PolynomialImageBounds.sum_projective_state_norm_sq}
\leanok
For a complete projective measurement $R$ on any finite-dimensional
space and any vector $\psi$, one has $\sum_g\|R_g\psi\|^2=\|\psi\|^2$.
\end{theorem}
\begin{proof}
\leanok
Expand each squared norm, use $R_g^*R_g=R_g$, and sum the effects.
\end{proof}

\begin{theorem}[Scalar-nonlinear mass and ordered error]
\label{thm:qld-nonlinear-mass-ordered-error}
\lean{MIPStarRE.QPBT.PolynomialImageBounds.nonlinear_mass_le_ordered_error}
\leanok
Let $F$ be a finite field, let $s$ be a finite set of labels, and let
$p_g\in F[u_1,\ldots,u_n][\alpha,\beta]$ for $g\in s$, with
$p_g\ne\alpha r+\beta t$ for all coefficient polynomials $r,t$.
Let $X(u),Z(u)$ be complete projective measurements on a common finite
space, and let $\psi_g$ be arbitrary vectors on that space. With
$T_{u,v,g}=T_{v_0,v_1,p_g(u,v)}$ and $D$ as above,
\[
 \sum_{g\in s}\|\psi_g\|^2\leq
 2\mathbb E_{u\in F^n}\mathbb E_{v\in F^2}
   \sum_{g\in s}\|\psi_g-T_{u,v,g}\psi_g\|^2
 +2\sum_{g\in s}\frac{D(p_g)}{|F|}\|\psi_g\|^2.
\]
Both expectations are uniform. The vectors need not be normalized, and
no commutation assumption between $X$ and $Z$ is imposed.
\end{theorem}
\begin{proof}
\leanok
\uses{def:qld-polynomial-indicators}
For each $g$, choose a nonzero coefficient outside the two scalar-linear
monomials. Its zero set is a common exceptional set of base points for
all measurement outcomes. Apply Schwartz--Zippel to that coefficient,
then to the specialized scalar polynomial minus each linear answer.
Positivity and completeness of the diagonal sandwich preserve the weight
$\|\psi_g\|^2$. The event $\beta=0$ contributes at most
$|F|^{-1}\|\psi_g\|^2$. Finally apply the squared-norm triangle inequality
and sum over the labels.
\end{proof}

\begin{theorem}[Actual scalar-nonlinear mass]
\label{thm:qld-actual-scalar-nonlinear-mass}
\lean{MIPStarRE.QPBT.ExtendedLineGame.scalarNonlinearMass_le_ordered_error}
\leanok
For a projective Pauli setting with admissible parameters, any opposite
placements $p_1,p_2$, and any complete projective extended-polynomial
measurement $R$ on $p_1$,
\[
 N_{p_1}(R)\leq2E^{XZ}_{p_1,p_2}(R)+\frac{2((2m+2)d+1)}q.
\]
\end{theorem}
\begin{proof}
\leanok
\uses{thm:qld-nonlinear-mass-ordered-error, thm:qld-scalar-expansion-degrees,
thm:qld-projective-outcome-mass, def:qld-extended-polynomial-ordered-error}
Use $\psi_g=(R_g)_{p_1}\hat\psi$ and the point projectors on $p_2$.
Opposite placements commute, so the residual in the preceding theorem
equals the residual defining the ordered error. Transport the full joint
uniform distribution through the coordinate equivalence. Enlarge the
nonnegative restricted sums to all polynomial outcomes and use
$\sum_g\|\psi_g\|^2=1$, together with the proved degree bound.
\end{proof}

\begin{theorem}[Rounded scalar-linearity concentration]
\label{thm:qld-rounded-scalar-linearity}
\lean{MIPStarRE.QPBT.ExtendedLineGame.exists_rounded_polynomial_scalar_mass}
\leanok
There are universal constants $a\geq1$ and $0<b\leq1$ with the following
property. Under the explicit point and directly indexed completed-answer
line-witness hypotheses of Theorem~\ref{thm:qld-rounded-polynomial-ordered},
with exactly its definitions of $\delta$ and $\eta$, there exist complete
projective measurements $R_A,R_B$ satisfying both ordered estimates of
that theorem and, simultaneously,
\[
 N_{AA'}(R_A),\ N_{BB'}(R_B)
 \leq 2(4\eta+8\delta_Q)+\frac{2((2m+2)d+1)}q.
\]
The line witness remains an input on its stated restricted domain.
Neither an ordered-error bound nor a mass bound is an input.
Dependence of the coefficient of $\alpha$ on $z$, dependence of the
coefficient of $\beta$ on $x$, and the retained point overlaps are not
conclusions of this auxiliary theorem.
\end{theorem}
\begin{proof}
\leanok
\uses{thm:qld-rounded-polynomial-ordered, thm:qld-actual-scalar-nonlinear-mass}
Choose both measurements by the rounded ordered-estimate theorem.
Apply the actual mass bound to $AA'\mid BA''$ and $BB'\mid AB''$.
\end{proof}

\begin{lemma}[Uniform polynomial and scalar collisions]
\label{lem:qld-fiber-collisions}
\lean{MIPStarRE.QPBT.PolynomialImageBounds.avg_eval_eq_le}
\lean{MIPStarRE.QPBT.affine_collision_average}
\leanok
Over a finite field $F$, distinct polynomials of total degree at most $D$
agree at a uniform point with probability at most $D/|F|$.
For two pairs $u,u'\in F^2$, the probability that
$\alpha u_1+\beta u_2=\alpha u'_1+\beta u'_2$ is one if $u=u'$ and
$1/|F|$ otherwise. These are auxiliary forms of Schwartz--Zippel and
the scalar collision calculation in \texttt{eq:qld-g-separable}.
\end{lemma}
\begin{proof}
\leanok
Apply Schwartz--Zippel in the first assertion. In the second, a nonzero
coefficient difference determines exactly one value of one scalar for
each value of the other. These proofs reuse the existing polynomial-image
and affine coarse-graining calculations.
\end{proof}

\begin{lemma}[Ordered norm and residual comparisons]
\label{lem:qld-fiber-norm-comparisons}
\lean{MIPStarRE.QPBT.PolynomialImageBounds.abs_orderedIndicator_norm_sq_sub_diagonal_le}
\lean{MIPStarRE.QPBT.PolynomialImageBounds.sum_mass_le_residual_add_image}
\lean{MIPStarRE.QPBT.ExtendedLineGame.placed_residual}
\leanok
For complete projective measurements $X,Z$ and any vector $\psi$, the
ordered and diagonal indicators satisfy
\[
 \bigl|\|T_{\alpha,\beta,a}\psi\|^2
       -\langle\psi,H_{\alpha,\beta,a}\psi\rangle\bigr|
 \leq 1_{\beta=0}\|\psi\|^2.
\]
For any finite set $s$ and vectors $\psi_g,A_g$,
$\sum_{g\in s}\|\psi_g\|^2\leq
2\sum_{g\in s}\|\psi_g-A_g\|^2+2\sum_{g\in s}\|A_g\|^2$.
For opposite expanded placements, the residual
$(R_g)_{p_1}\hat\psi-T_{p_2}(R_g)_{p_1}\hat\psi$
equals $(R_g)_{p_1}(I-T_{p_2})\hat\psi$.
\end{lemma}
\begin{proof}
\leanok
\uses{def:qld-polynomial-indicators}
For nonzero $\beta$, orthogonality gives the diagonal Gram identity.
For zero $\beta$, both quadratic forms are between zero and $\|\psi\|^2$.
The second assertion is the squared-norm triangle inequality.
The final identity follows from commutation of opposite placements.
The first two calculations are also used in the polynomial-image estimate.
\end{proof}

The next five assertions are formalization-only auxiliary calculations for
the ordered products in \texttt{eq:qld-g-prime} and
\texttt{eq:qld-g-43} of \cite{Ji2020MIPStar}. They do not assert the
concentration or polynomial-pair conclusion of Lemma~\ref{lem:qld-4-7}.

\begin{theorem}[Uniform exceptional-coefficient error]
\label{thm:aux-ordered-exceptional-average}
\lean{MIPStarRE.QPBT.PolynomialImageBounds.avg_abs_orderedIndicator_norm_sq_sub_diagonal_le}
\leanok
\uses{def:qld-polynomial-indicators}
Let $F$ be a finite field, let $X,Z$ be complete projective measurements
on a common finite-dimensional space, and let $\psi$ be any vector. For
each fixed $\alpha\in F$ and arbitrary function $a:F\to F$,
\[
 \mathbb E_{\beta\in F}
 \left|\|T_{\alpha,\beta,a(\beta)}\psi\|^2
 -\langle\psi,H_{\alpha,\beta,a(\beta)}\psi\rangle\right|
 \leq |F|^{-1}\|\psi\|^2.
\]
The expectation is uniform and $\psi$ need not be normalized.
\end{theorem}
\begin{proof}
\leanok
\uses{lem:qld-fiber-norm-comparisons}
Apply the pointwise estimate and average the indicator of $\beta=0$;
its uniform probability is $|F|^{-1}$.
\end{proof}

\begin{theorem}[Summed exceptional-coefficient error]
\label{thm:aux-ordered-exceptional-sum}
\lean{MIPStarRE.QPBT.PolynomialImageBounds.avg_sum_abs_orderedIndicator_norm_sq_sub_diagonal_le}
\leanok
\uses{def:qld-polynomial-indicators}
Let $F$ be a finite field, $X,Z$ complete projective measurements,
$\alpha\in F$, $\Gamma$ a finite outcome set, $\psi_g$ arbitrary vectors
on their common space, and $a_g:F\to F$ arbitrary functions. Then
\[
 \mathbb E_{\beta\in F}\sum_{g\in\Gamma}
 \left|\|T_{\alpha,\beta,a_g(\beta)}\psi_g\|^2
 -\langle\psi_g,H_{\alpha,\beta,a_g(\beta)}\psi_g\rangle\right|
 \leq |F|^{-1}\sum_{g\in\Gamma}\|\psi_g\|^2.
\]
\end{theorem}
\begin{proof}
\leanok
\uses{thm:aux-ordered-exceptional-average}
Exchange the finite sum with the uniform average and apply the preceding
inequality separately to each $\psi_g$.
\end{proof}

\begin{theorem}[Projective exceptional-coefficient error]
\label{thm:aux-projective-exceptional-average}
\lean{MIPStarRE.QPBT.PolynomialImageBounds.avg_sum_abs_projective_orderedIndicator_le}
\leanok
\uses{def:qld-polynomial-indicators}
Let $F$ be a finite field and $S=\{S_g\}_{g\in\Gamma}$ a complete projective measurement on
the same space as complete projective measurements $X,Z$. For any vector
$\psi$, fixed $\alpha\in F$, and arbitrary functions $a_g:F\to F$,
\[
 \mathbb E_{\beta\in F}\sum_{g\in\Gamma}
 \left|\|T_{\alpha,\beta,a_g(\beta)}S_g\psi\|^2
 -\langle S_g\psi,H_{\alpha,\beta,a_g(\beta)}S_g\psi\rangle\right|
 \leq |F|^{-1}\|\psi\|^2.
\]
No commutation of the three measurements is required.
\end{theorem}
\begin{proof}
\leanok
\uses{thm:aux-ordered-exceptional-sum, thm:qld-projective-outcome-mass}
Apply the summed bound with $\psi_g=S_g\psi$ and use
$\sum_g\|S_g\psi\|^2=\|\psi\|^2$.
\end{proof}

\begin{theorem}[Uniform base-point and scalar exceptional error]
\label{thm:aux-projective-exceptional-uniform}
\lean{MIPStarRE.QPBT.PolynomialImageBounds.avg_uniform_sum_abs_projective_orderedIndicator_le}
\leanok
\uses{def:qld-polynomial-indicators}
Over a finite field $F$, let $U$ be a finite nonempty base set and
$S=\{S_g\}_{g\in\Gamma}$ a
complete projective measurement, and $X(t),Z(t)$ complete projective
measurements for every $t\in U$, all on the same finite-dimensional
space. For arbitrary field-valued answers $a_{g,t}(\alpha,\beta)$ and
any vector $\psi$,
\[
 \mathbb E_{(t,\alpha,\beta)\in U\times(F\times F)}\sum_{g\in\Gamma}
 \left|\|T^{(t)}_{\alpha,\beta,a_{g,t}(\alpha,\beta)}S_g\psi\|^2
 -\langle S_g\psi,
 H^{(t)}_{\alpha,\beta,a_{g,t}(\alpha,\beta)}S_g\psi\rangle\right|
 \leq |F|^{-1}\|\psi\|^2,
\]
where the superscript denotes the indicators formed from $X(t),Z(t)$
and the expectation is uniform on the entire product.
\end{theorem}
\begin{proof}
\leanok
\uses{thm:aux-projective-exceptional-average}
Decompose the product uniform law into successive averages over $t$,
$\alpha$, and $\beta$. Apply the preceding inequality at each fixed
$t,\alpha$, then average the constant bound.
\end{proof}

\begin{theorem}[Reversal of the ordered indicator]
\label{thm:aux-ordered-indicator-reverse}
\lean{MIPStarRE.QPBT.PolynomialImageBounds.orderedIndicator_reverse}
\leanok
\uses{def:qld-polynomial-indicators}
For complete measurements $X,Z$ over a finite field $F$ and any
$\alpha,\beta,a\in F$,
\[
 \sum_{\alpha b+\beta c=a}Z_cX_b
 =T_{\beta,\alpha,a}(Z,X).
\]
Projectivity and commutation are not assumed.
\end{theorem}
\begin{proof}
\leanok
\uses{def:qld-polynomial-indicators}
Exchange the pair of summation indices and commute the field addition;
the order of the operators remains unchanged.
\end{proof}

\begin{definition}[Partial coefficient polynomials]
\label{def:qld-block-coefficient-polynomial}
\lean{MIPStarRE.QPBT.ExtendedLineGame.blockCoefficientPolynomial}
\lean{MIPStarRE.QPBT.ExtendedLineGame.fiberQuestionEquiv}
\leanok
For a scalar coefficient $c_{g,a}(x,z)$, view it as a polynomial in $x$
with coefficients in $F[z]$, or in $z$ with coefficients in $F[x]$.
An order parameter $o\in\{x\mid z,z\mid x\}$ selects the varying block.
The associated joint question equivalence takes a fixed block, a varying
block, and two scalars to the actual extended question. Reversal exchanges
both the point blocks and the scalars.
\end{definition}

\begin{theorem}[Partial degrees and the separated image]
\label{thm:qld-block-coefficient-degrees}
\lean{MIPStarRE.QPBT.ExtendedLineGame.blockCoefficientPolynomial_degrees}
\lean{MIPStarRE.QPBT.ExtendedLineGame.blockCoefficientPolynomial_eval}
\lean{MIPStarRE.QPBT.ExtendedLineGame.scalarPolynomial_linear_read_fiber}
\lean{MIPStarRE.QPBT.ExtendedLineGame.exists_polyPair_of_scalar_separated}
\leanok
For an actual singleton outcome, every coefficient of either partial
polynomial and the partial polynomial itself have total degree at most $md$.
Partial evaluation is exactly evaluation of the original coefficient at the
corresponding two blocks. For a scalar-linear outcome, this recovers the
actual readout in either block and scalar order.
If the coefficient of $\alpha$ is independent of $z$ and the coefficient
of $\beta$ is independent of $x$, there is a pair of individual-degree-$d$
polynomials $(g_X,g_Z)$ such that the field-transported outcome equals
$\mathrm{combine}_{g_X,g_Z}$ as a formal polynomial.
\end{theorem}
\begin{proof}
\leanok
\uses{def:qld-block-coefficient-polynomial, thm:qld-scalar-expansion-degrees,
def:combine-map}
Nonzero partial coefficients are coefficients of the original polynomial;
each exponent is at most $d$. The ring equivalences commute with evaluation.
A constant partial polynomial is the renamed polynomial of the fixed block.
Its coefficients inherit the individual-degree certificate. Apply the
injective scalar expansion to identify the resulting pair with the existing
combining map.
\end{proof}

\begin{theorem}[Common exceptional coefficient and weighted fibers]
\label{thm:qld-common-exceptional-fibers}
\lean{MIPStarRE.QPBT.PolynomialImageBounds.exists_common_exceptional_coefficient}
\lean{MIPStarRE.QPBT.PolynomialImageBounds.avg_weighted_fiber_le}
\leanok
Let $p\in F[z_1,\ldots,z_k][x_1,\ldots,x_n]$ be nonconstant in $x$,
with coefficient degrees at most $C$ and degree in $x$ at most $D$.
There is a nonzero coefficient $c(z)$ of a nonconstant $x$-monomial whose
zero set has probability at most $C/|F|$. For every $z$ outside this one set and every
$b\in F$, $\Pr_x[p(x,z)=b]\leq D/|F|$.
Consequently, if $w(z,b)\geq0$ and $\sum_b w(z,b)=M$ for every $z$, then
\[
 \mathbb E_z\sum_b\Pr_x[p(x,z)=b]w(z,b)
 \leq \frac{C+D}{|F|}M.
\]
\end{theorem}
\begin{proof}
\leanok
\uses{lem:qld-fiber-collisions}
Choose one nonzero coefficient at a nonconstant monomial. Specialization
outside its zero set leaves a polynomial distinct from every constant.
Apply Schwartz--Zippel twice. Sum the weights on the common exceptional
set before averaging it; no outcome-cardinality factor appears.
\end{proof}

\begin{theorem}[Wrong-variable ordered norm]
\label{thm:qld-wrong-variable-norm}
\lean{MIPStarRE.QPBT.PolynomialImageBounds.avg_orderedIndicator_norm_sq_le_of_nonconstant}
\leanok
With $p,C,D$ as above, let $X(x),Z(z)$ be complete projective measurements,
let $r(x,z)$ be any field-valued function, and let $\psi$ be any vector.
The ordered indicator for the readout $\alpha r(x,z)+\beta p(x,z)$ satisfies
\[
 \mathbb E_z\mathbb E_x\mathbb E_{\alpha,\beta}
   \|T_{\alpha,\beta,\alpha r(x,z)+\beta p(x,z)}\psi\|^2
 \leq \frac{C+D+2}{|F|}\|\psi\|^2.
\]
\end{theorem}
\begin{proof}
\leanok
\uses{thm:qld-common-exceptional-fibers, lem:qld-fiber-norm-comparisons,
lem:qld-fiber-collisions}
The zero second scalar and the unequal scalar-answer collision each cost
$|F|^{-1}\|\psi\|^2$. Remove the middle $X$ contraction from $ZXZ$.
Its remaining $Z$ weight is independent of $x$, so the weighted fiber
theorem applies. Interchanging both blocks and measurements gives $XZX$.
\end{proof}

\begin{definition}[Wrong-variable and nonseparated masses]
\label{def:qld-wrong-variable-masses}
\lean{MIPStarRE.QPBT.ExtendedLineGame.scalarWrongVariableMass}
\lean{MIPStarRE.QPBT.ExtendedLineGame.separatedPolynomialMass}
\leanok
For a complete extended-polynomial measurement $R$ on placement $p$, let
$W_p^{XZ}(R)$ be the sum of $\|(R_g)_p\hat\psi\|^2$ over scalar-linear
outcomes whose beta coefficient depends on $x$. Define $W_p^{ZX}(R)$
using the alpha coefficient depending on $z$.
Let $G_p(R)$ be the same mass sum outside the field-transported image of
the existing bounded combining map $(g_X,g_Z)\mapsto\mathrm{combine}_{g_X,g_Z}$.
\end{definition}

\begin{theorem}[Actual wrong-variable and separated masses]
\label{thm:qld-actual-separated-mass}
\lean{MIPStarRE.QPBT.ExtendedLineGame.scalarWrongVariableMass_le_ordered_error}
\lean{MIPStarRE.QPBT.ExtendedLineGame.separatedPolynomialMass_le_ordered_errors}
\leanok
For every opposite pair $p_1,p_2$ and complete projective measurement $R$
on $p_1$, for both $o=XZ,ZX$,
\[
 W_{p_1}^{o}(R)\leq2E^o_{p_1,p_2}(R)+\frac{2(2md+2)}q,
\qquad
 G_{p_1}(R)\leq4E^{XZ}_{p_1,p_2}(R)+2E^{ZX}_{p_1,p_2}(R)
                  +\frac{12md+4d+10}q.
\]
\end{theorem}
\begin{proof}
\leanok
\uses{thm:qld-wrong-variable-norm, thm:qld-block-coefficient-degrees,
thm:qld-projective-outcome-mass, thm:qld-actual-scalar-nonlinear-mass,
def:qld-wrong-variable-masses}
Use the actual vectors $(R_g)_{p_1}\hat\psi$ and opposite placed point
measurements. Their squared masses sum to one. The norm-to-residual
comparison retains the full outcome sum. The complement of the combining
image lies in the union of the three excluded classes; sum their bounds.
\end{proof}

\begin{theorem}[Rounded wrong-variable concentration]
\label{thm:qld-rounded-wrong-variable-mass}
\lean{MIPStarRE.QPBT.ExtendedLineGame.exists_rounded_polynomial_wrong_variable_mass}
\leanok
Under exactly the hypotheses and definitions of $\delta,\eta$ in
Theorem~\ref{thm:qld-rounded-polynomial-ordered}, there exist actual rounded
measurements satisfying both ordered estimates with $E=4\eta+8\delta_Q$,
both wrong-variable bounds $2E+2(2md+2)/q$, and the scalar-nonlinear bound
of Theorem~\ref{thm:qld-rounded-scalar-linearity}.
\end{theorem}
\begin{proof}
\leanok
\uses{thm:qld-rounded-scalar-linearity, thm:qld-actual-separated-mass}
Apply the wrong-variable mass inequality to each of the two measurements
supplied by rounded scalar-linearity concentration.
\end{proof}

\begin{theorem}[Rounded concentration on polynomial pairs]
\label{thm:qld-rounded-separated-mass}
\lean{MIPStarRE.QPBT.ExtendedLineGame.exists_rounded_polynomial_separated_mass}
\leanok
Under the same hypotheses, there exist actual rounded measurements satisfying
both ordered estimates with $E=4\eta+8\delta_Q$ and, on both heterogeneous
placements,
\[
 G_{AA'}(R_A),\ G_{BB'}(R_B)\leq6E+\frac{12md+4d+10}q.
\]
The directly indexed completed-answer line-witness restriction remains
explicit. This proves concentration, not the subsequent retained overlaps,
completion consistency, or the source global-pair theorem.
This existential choice is separate from the measurements in the preceding theorem.
\end{theorem}
\begin{proof}
\leanok
\uses{thm:qld-rounded-polynomial-ordered, thm:qld-actual-separated-mass}
Apply the actual mass inequalities to $AA'\mid BA''$ and $BB'\mid AB''$.
\end{proof}

\begin{theorem}[Polynomial consistency from point comparisons]
\label{thm:qld-polynomial-consistency-from-points}
\lean{MIPStarRE.QPBT.direct_polynomial_consistency_le_point_bounds}
\leanok
This is an auxiliary estimate for the polynomial measurements in
Lemma~\ref{lem:qld-4-7}. Let $(q,n,d,k)$ be directly indexed low-degree
parameters, and let $A$ and $B$ be complete POVMs on two finite-dimensional
spaces whose outcomes are $k$-tuples of individual-degree-$d$ polynomials in
$n$ variables over $\F_q$. Let $Q_A^u,Q_B^u$ be complete POVMs with outcomes
in $\F_q^k$, and let the bipartite state be normalized. Write $A^u,B^u$ for
the evaluation of every component at $u$, and write
$\operatorname{Cons}$ for the off-diagonal consistency defect, averaged over
uniform $u$ when the measurements depend on $u$. If
\[
\operatorname{Cons}(A^u,Q_B^u)\leq\eta_A,\qquad
\operatorname{Cons}(Q_A^u,B^u)\leq\eta_B,\qquad
\operatorname{Cons}(Q_A^u,Q_B^u)\leq\delta_Q,
\]
then
\[
\operatorname{Cons}(A,B)
\leq \eta_A+2\sqrt{\delta_Q+\eta_B}+\frac{nd}{q}.
\]
No projectivity assumption is required.
\end{theorem}
\begin{proof}
\lean{MIPStarRE.QPBT.direct_poly_tuple_collision_le}
\leanok
The consistency comparison through $Q_B^u,Q_A^u$ bounds
$\operatorname{Cons}(A^u,B^u)$ by
$\eta_A+2\sqrt{\delta_Q+\eta_B}$. Distinct polynomial tuples differ in
at least one component. The Schwartz--Zippel bound for that component gives
collision probability at most $nd/q$ for the entire tuple. Summing the
nonnegative bipartite outcome weights therefore bounds the full polynomial
defect by the evaluated defect plus $nd/q$.
\end{proof}

\begin{theorem}[Formalization support: postprocessed consistency under ground compression]
\label{thm:qld-ground-compression-postprocessing}
\lean{MIPStarRE.QPBT.consistency_defect_ground_compress_measurement_postprocess}
\lean{MIPStarRE.QPBT.ground_compress_dilated_measurement}
\leanok
\uses{def:consistency}
Let $\psi$ be a bipartite vector, let $k_0,l_0$ be distinguished coordinates
of two ancillary spaces, and let $A^x,B^x$ be finite families of complete POVMs
on the corresponding padded spaces. For question-dependent maps $f_x,g_x$ to
a common finite outcome set, the consistency defect of the postprocessed
families on the state padded at $(k_0,l_0)$ is exactly the consistency defect of
their ground-coordinate compressions on $\psi$:
\[
 \operatorname{Cons}_\mu(f_*A,g_*B;\operatorname{pad}_{k_0,l_0}\psi)
 =\operatorname{Cons}_\mu(f_*A^0,g_*B^0;\psi).
\]
In particular, compressing the completed Naimark dilation of a POVM at its
distinguished residual coordinate recovers the original POVM.
\end{theorem}
\begin{proof}
\leanok
\uses{def:consistency}
Expand the off-diagonal defect and use the equality of each bipartite quadratic
form before and after ground-coordinate compression. Compression commutes with
postprocessing. The final assertion is the defining compression identity for
the completed Naimark dilation.
\end{proof}

\begin{theorem}[Consistency after ground compression]
\label{thm:qld-compressed-polynomial-consistency}
\lean{MIPStarRE.QPBT.direct_point_consistency_le_point_rejection}
\lean{MIPStarRE.QPBT.ExtendedLineGame.point_consistency_le}
\lean{MIPStarRE.QPBT.ExtendedLineGame.compressed_polynomial_consistency_le}
\leanok
This is an auxiliary estimate for the directly indexed strategy constructed
from combined point and extended line measurements with errors $\delta_Q$
and $\delta_L$. Its point readout maps wrong-format answers to the zero tuple.
For any directly indexed strategy, the consistency defect of these readouts
is at most the point/point rejection probability. For the constructed
extended strategy it is therefore at most $\delta_Q$.

Let $A,B$ be polynomial POVMs on the Naimark dilation of this strategy,
with the two opposite point-consistency defects at most $\eta_A,\eta_B$.
If $J_A,J_B$ insert the distinguished ancillary coordinates and
$A_g^0=J_A^*A_gJ_A$, $B_g^0=J_B^*B_gJ_B$, then on the original expanded
bipartite state,
\[
\operatorname{Cons}(A^0,B^0)
\leq\eta_A+2\sqrt{\delta_Q+\eta_B}+\frac{(2m+2)d}{q}.
\]
In applying direct soundness, both point-comparison errors retain the
passing error $3(\sqrt{\delta_Q+\delta_L}+md/q)$.
\end{theorem}
\begin{proof}
\leanok
\uses{thm:qld-ground-compression-postprocessing,
  thm:qld-polynomial-consistency-from-points}
Different point readouts force rejection; wrong-format answers also reject.
The established point/point rejection estimate gives the bound $\delta_Q$.
Ground compression preserves every bipartite correlation on the padded state
and commutes with outcome postprocessing. Apply
Theorem~\ref{thm:qld-polynomial-consistency-from-points} to the compressed
POVMs with $n=2m+2$ and $k=1$.
\end{proof}

% This lemma is not a statement of the source; it records the branch
% accounting and the Naimark transport for a strategy assembled from supplied
% measurements.  It does not construct those measurements.
\begin{lemma}[Formalization support: passing value of the supplied extended-line strategy]
	\label{lem:direct-passing-value}
	\lean{MIPStarRE.QPBT.directPassingErrorEnvelope}
	\lean{MIPStarRE.QPBT.ExtendedLineGame.seven_branch_rejections_le}
	\lean{MIPStarRE.QPBT.ExtendedLineGame.strategy_value_ge_directPassingErrorEnvelope}
	\lean{MIPStarRE.QPBT.ExtendedLineGame.projectiveStrategy_value_ge_directPassingErrorEnvelope}
	\leanok
	\uses{lem:qld-4-10, lem:qld-4-13-established, def:ld-combined-strategy,
		def:tensor-product-value, thm:naimark}
	Write $\mathrm{env}(x,y) = 3(\sqrt{x} + y)$ for the direct passing-error
	function. Let $\mathcal S$ be a projective setting carrying combined point
	measurements with error $\delta_Q$ as in Lemma~\ref{lem:qld-4-10} and
	supplied extended-line measurements with error $\delta_{\mathrm{Line}}$ as
	in Lemma~\ref{lem:qld-4-13-established}, and let $T$ be the strategy for
	the directly indexed low individual degree game at parameters
	$(q,2m+2,d,1)$ assembled from them. Then the seven potentially rejecting
	question branches of $T$ reject with total probability at most
	\[
		\delta_Q + 16\delta_{\mathrm{Line}}
		+ 4\sqrt{\delta_Q + 2\delta_{\mathrm{Line}}}
		+ 5\,\frac{md}{q}\,,
	\]
	and consequently the value of $T$, which is one minus one ninth of that
	total, satisfies
	\[
		\mathrm{val}(T) \ge 1 - \mathrm{env}\Big(
			\delta_Q + \delta_{\mathrm{Line}},\ \frac{md}{q}\Big)\,.
	\]
	The strategy obtained from $T$ by Naimark dilation
	(Theorem~\ref{thm:naimark}) satisfies the same lower bound.

	Both measurement families are hypotheses here, not constructions, so this
	lemma does not establish Lemma~\ref{lem:qld-4-7}: the questions are the
	directly indexed ones, and transport to the seed-indexed question
	distribution of Definition~\ref{def:ld-game} remains open; see
	\cite{gap:qpbt_ld-dimension-divisibility} and
	\cite{gap:qpbt_combined-lines-error-term}.
\end{lemma}
\begin{proof}
	\leanok
	\uses{lem:qld-4-10, lem:qld-4-13-established, def:ld-combined-strategy,
		def:tensor-product-value, thm:naimark, lem:sandwich-codeword-defect,
		lem:coefficient-polynomial-agreement}
	The seven branches are estimated separately: the point/point branch by the
	combined point consistency, the four mixed line/point branches by four
	times the shared line/point estimate, and the two same-kind line branches
	by the axis and the diagonal parameter-evaluation defects together with
	the coefficient-collision loss $d/q + (2m+2)d/q \le 5md/q$. The two
	same-kind branches reach those defects through the codeword step of
	Lemma~\ref{lem:sandwich-codeword-defect}, whose collision hypothesis is
	supplied by Lemma~\ref{lem:coefficient-polynomial-agreement} at degree $d$
	for the axis branch and at degree $(2m+2)d$ for the diagonal branch. The
	value of the strategy is one minus one ninth of the branch total; capping
	the rejection probability by one and applying the capped scalar estimate
	for $\mathrm{env}$ gives the stated bound. Naimark dilation preserves the
	value of the strategy, hence the same bound.

	In Lean the branch accounting is assembled from
	\texttt{point\_point\_rejection\_le},
	\texttt{mixed\_line\_point\_rejections\_le\_four\_mul} and the two
	same-line bounds of
	\texttt{ExtendedLineGame/SameLineCoefficientBound.lean}; the latter are
	transported to the parameter-evaluation defects by the two
	\texttt{parameter\_evaluated\_coefficient\_defect\_eq} identities and
	bounded by \texttt{axis\_parameter\_evaluation\_defect\_le} and
	\texttt{diagonal\_parameter\_evaluation\_defect\_le}. The collision
	arithmetic $d/q + (2m+2)d/q \le 5md/q$ is
	\texttt{extendedDirectLd\_coefficientCollisionLoss\_le\_five\_mul}, the
	cap on the rejection probability is
	\texttt{directLdRejectionProbability\_le\_one}, and the capped scalar
	envelope step is
	\texttt{seven\_branch\_capped\_error\_le\_directPassingErrorEnvelope}.
	None of these auxiliaries carries a blueprint node of its own.
\end{proof}

% This remark records a defect of the source argument; it is not a statement
% from the paper.
\begin{remark}[Divisibility obstruction in the application of the classical
	soundness theorem]
	\label{rem:qld-4-7-divisibility}
	Definition~\ref{def:ld-game} defines the classical low individual
	degree game only when its dimension parameter divides $q$: the index
	function $\chi$ of the question distribution
	(Definition~\ref{def:ld-question-distribution}) partitions $\F_q$ into
	classes of equal size, one for each coordinate. The same divisibility
	underlies the line-point distribution over $\F_q^{2m+2}$
	(Definition~\ref{def:line-point-dist} instantiated at dimension
	$2m+2$), and with it the statements of Lemmas~\ref{lem:qld-sublines}
	and~\ref{lem:qld-4-13} and the proofs of all three of
	Lemmas~\ref{lem:qld-sublines}, \ref{lem:qld-4-13}
	and~\ref{lem:qld-4-7}. The condition that $2m+2$ divide $q$ does not
	follow from admissibility of $(q,m,d)$
	(Definition~\ref{def:admissible}), which supplies only that $m$
	divides $q$; indeed the two are jointly satisfiable only in a
	degenerate range. An admissible field size
	(Definition~\ref{def:admissible-size}) is a power of two, so $m$
	dividing $q$ forces $m$ to be a power of two; for $m \ge 2$ the
	integer $m+1$ is odd and exceeds $1$, so $2m+2 = 2(m+1)$ divides no
	power of two, while for $m = 1$ the condition $4 \mid q$ holds for
	every admissible $q$ except $q = 2$. The source proof in
	\cite{Ji2020MIPStar} performs the instantiation at dimension $2m+2$
	without addressing the divisibility. The statements above are
	therefore retained in their source form, and for admissible
	parameters with $m \ge 2$ their proofs, like the source's, are
	blocked at the marked steps, so the conclusions are not established
	for such parameters, here or in the source. A proof covering every
	admissible parameter tuple requires a formulation of the game, of its
	question distribution, and of Theorem~\ref{lem:ld-soundness} that
	applies to dimensions not dividing $q$; the missing step and the
	candidate repairs are analyzed in
	\cite{gap:qpbt_ld-dimension-divisibility}.
\end{remark}

\begin{lemma}[Formalization support for the supplied scalar point measurement]
	\label{lem:qld-supplied-scalar-point-measurement}
	\lean{MIPStarRE.QPBT.ExtendedLineGame.point_values_measurement_eq_suppliedQ}
	\leanok
	\uses{lem:qld-4-10, lem:qld-4-12, lem:qld-4-13, def:ld-game,
		def:bracket, def:expanded-state}
	This is not a named statement of \cite{Ji2020MIPStar}; it identifies
	the point measurement used in the first paragraph of the proof of
	Lemma~\ref{lem:qld-4-7} when the point and line measurements are supplied.
	Fix an admissible tuple $(q,m,d)$ and a projective setting as above.
	Suppose joint point measurements $\hat{Q}^{x,z}$ and extended line
	measurements $\hat{Q}^{\ell}$ satisfying the respective conclusions of
	Lemmas~\ref{lem:qld-4-10} and~\ref{lem:qld-4-13} are given, with error
	parameters $\delta_Q,\delta_L$.
	Use them to form the directly indexed game strategy at parameters
	$(q,2m+2,d,1)$. Let $K$ be its scalar field, and let
	$\iota:K\simeq\F_q$ be the canonical field identification with the
	Pauli scalar field. For $u\in K^{2m+2}$, write
	$\iota(u)=(x,z,\alpha,\beta)$ coordinatewise, and let $A^u_s$ be
	the answer measurement at the point question $u$ for player
	$s\in\{\mathsf{A},\mathsf{B}\}$. Let $r$ read the singleton point
	value tuple from an answer, returning the zero tuple on other answer
	forms. Then the following is an exact equality of measurements on
	$\hilb_s\ot(\C^q)^{\ot M}$, for both choices of $s$:
	\[
		\bigl(A^u_s\bigr)_{[\iota(r(\cdot)_0)=c]}
		=\bigl(\hat{Q}^{x,z}_s\bigr)_{[\alpha a+\beta b=c]},
		\qquad c\in\F_q.
	\]
	This identity assumes the supplied measurements and does not assert
	the existence of the polynomial-pair measurement in
	Lemma~\ref{lem:qld-4-7}.
\end{lemma}
\begin{proof}
	\leanok
	\uses{lem:qld-4-12, def:bracket}
	At a point question, the answer measurement first coarse-grains
	$\hat{Q}^{x,z}_s$ by $(a,b)\mapsto\alpha a+\beta b$ and then records
	the result $c$ as the singleton point tuple $(\iota^{-1}(c))$.
	Compose these outcome maps with $r$ and with the map reading coordinate
	$0$ and applying $\iota$. Since $\iota(\iota^{-1}(c))=c$, the composite
	is exactly $(a,b)\mapsto\alpha a+\beta b$. Equality of the finite fiber
	sums proves the measurement identity for either player.
\end{proof}

% The two lemmas below are not statements of the source; they record the
% coordinate change for the LDT consistency conclusions and the compression of
% the correlated residue register discussed in the preceding remark.
\begin{lemma}[Formalization support for the consistency conclusions in directly indexed coordinates]
	\label{lem:direct-ld-consistency-coordinates}
	\lean{MIPStarRE.QPBT.directPointPolynomial_consistencyDefect_le}
	\lean{MIPStarRE.QPBT.directPolynomialPoint_consistencyDefect_le}
	\lean{MIPStarRE.QPBT.directPolynomialPolynomial_consistencyDefect_le}
	\leanok
	\uses{def:consistency, def:ld-game, def:admissible-size,
		lem:direct-ld-consistency-defect-equivalence}
	This is not a named statement of \cite{Ji2020MIPStar}. It records three
	consistency implications in the directly indexed coordinates, in which the
	sampled coordinate is carried by the question itself rather than by a
	residue class of $\F_q$; see Remark~\ref{rem:qld-4-7-divisibility}. The
	hypothesis of each implication is an arbitrary consistency relation
	between measurement families of the stated kinds; no implication produces
	such a relation or constructs the measurements it relates. They may in
	particular be applied to the three consistency relations in the conclusion
	of Theorem~\ref{thm:main-formal}.

	Fix a directly indexed parameter tuple $D=(q,m,d,k)$, where $m,d,k \ge 1$
	are integers and $q$ is an admissible field size
	(Definition~\ref{def:admissible-size}). In contrast with the parameters of
	the classical game (Definition~\ref{def:ld-game}), no divisibility
	condition $m \mid q$ is imposed; dropping that condition is the only
	difference between the two parameter domains. Let $\delta\in\R$.
	The following three implications hold:
	\begin{enumerate}
		\item let $S=(\psi,A,B)$ be a projective strategy for the
		      directly indexed game, fix
		      $r\in\operatorname{Fin}(k)=\{0,\ldots,k-1\}$, and let
		      $G^{\mathrm B}$ be a projective measurement on
		      $\hilb_{\mathrm B}$ indexed by single $m$-variate
		      polynomials of individual degree at most $d$, consistency with
		      error $\delta$ between the coordinate-$r$ LDT point family of
		      Alice and the evaluation family of $G^{\mathrm B}$ implies
		      that the corresponding directly indexed point/polynomial
		      off-diagonal defect, evaluated on $\psi$ and averaged uniformly
		      over $u\in\F_q^m$, is at most $\delta$;
		\item let $S=(\psi,A,B)$ be a projective strategy for the
		      directly indexed game, fix
		      $r\in\operatorname{Fin}(k)=\{0,\ldots,k-1\}$, and let
		      $G^{\mathrm A}$ be a projective measurement on
		      $\hilb_{\mathrm A}$ indexed by the same single polynomials,
		      consistency with error $\delta$ between its evaluation family
		      and the coordinate-$r$ LDT point family of Bob implies the
		      corresponding directly indexed polynomial/point defect bound
		      $\delta$, again averaged uniformly over $u\in\F_q^m$;
		\item let $S=(\psi,A,B)$ be an arbitrary strategy for the
		      directly indexed game, with no projectivity assumption on $S$,
		      and let $G^{\mathrm A}$ and $G^{\mathrm B}$ be projective
		      measurements on the two players' spaces indexed by single
		      polynomials. Consistency with error $\delta$ between these two
		      constant LDT measurement families over the one-element question
		      set implies that the corresponding directly indexed
		      polynomial/polynomial defect is at most $\delta$.
	\end{enumerate}
	Thus the first two implications concern one fixed coordinate and the
	third concerns no coordinate. All three use single-polynomial projective
	measurements; they do not construct a joint measurement of a
	$k$-tuple of polynomials.
\end{lemma}
\begin{proof}
	\leanok
	For each implication, the density-state consistency relation is changed
	to the identical off-diagonal defect on the vector $\psi$. The first two
	defects are then reindexed by the point and scalar bijections and by
	polynomial evaluation; the third is reindexed by the polynomial-outcome
	bijection. These changes preserve the bound $\delta$ exactly.
\end{proof}

\begin{lemma}[Formalization support for the compression of a correlated residue register]
	\label{lem:direct-ld-residue-compression}
	\lean{MIPStarRE.QPBT.ldStrategyToDirect_pointPolynomial_compression}
	\lean{MIPStarRE.QPBT.ldStrategyToDirect_polynomialPoint_compression}
	\leanok
	\uses{def:consistency, def:ld-game, def:ld-question-distribution, def:admissible-size}
	This is not a named statement of \cite{Ji2020MIPStar}. It records two
	compression identities between the classical low individual degree game
	and its directly indexed presentation.

	Fix parameters $\mathsf{ldparams}=(q,m,d,k)$ of the classical game as in
	Definition~\ref{def:ld-game}: integers $m,d,k \ge 1$, an admissible field
	size $q$ (Definition~\ref{def:admissible-size}), and $m \mid q$. Unlike
	in Lemma~\ref{lem:direct-ld-consistency-coordinates}, the divisibility
	belongs to the domain here, since it is what gives the residue register
	below its dimension; the directly indexed parameters appearing on the
	other side of the identities are this tuple with the divisibility
	condition dropped. Let $(\psi, A, B)$ be a strategy for the classical
	game, and extend it to the directly indexed game by adjoining to each
	player a register of dimension $q/m$, the common size of the classes of
	the index function $\chi$ of
	Definition~\ref{def:ld-question-distribution}; the two copies of that
	register are perfectly correlated in the extended state. For a
	polynomial-tuple POVM $G^{\mathrm B}$ on
	$\hilb_{\mathrm B}\ot\C^{q/m}$, let $\bar G^{\mathrm B}$ be the
	normalized average of its diagonal blocks over the residue register.
	Then, over a uniformly random point $u\in\F_q^m$, the off-diagonal
	defect between Alice's point measurement and the evaluation of
	$G^{\mathrm B}$ on the extended strategy equals the defect between
	Alice's original point measurement and the evaluation of
	$\bar G^{\mathrm B}$ on $(\psi,A,B)$.

	Similarly, for every polynomial-tuple POVM $G^{\mathrm A}$ on
	$\hilb_{\mathrm A}\ot\C^{q/m}$, compression gives equality between
	the extended polynomial/point defect and the defect between the
	evaluation of $\bar G^{\mathrm A}$ and Bob's original point
	measurement. No projectivity hypothesis is imposed on the strategy or
	on either polynomial-tuple POVM.
\end{lemma}
\begin{proof}
	\leanok
	The point measurement of the extended strategy is the block-diagonal
	amplification of the original point measurement, while polynomial
	evaluation commutes with averaging the diagonal blocks. The quadratic
	form of the correlated state against each mixed pair of operators is
	therefore the quadratic form of the original state against the
	compressed pair, and summing the off-diagonal terms gives both identities.
\end{proof}

\begin{remark}[Compression of two polynomial measurements]
	The preceding lemma does not compare defects in which both families are
	polynomial measurements on the extended carriers. The extended state
	pairs equal residues, whereas separate compression averages independently
	over two residues; in general neither resulting quantity bounds the other.
	This obstruction is analyzed in
	\cite{gap:qpbt_ld-dimension-divisibility}.
\end{remark}

% The results below are not statements of the source. They record the
% coordinatewise application of Theorem~\ref{thm:main-formal} to the directly
% indexed game and the one-coordinate polynomial-tuple conclusion discussed in
% Remark~\ref{rem:ld-soundness-simultaneity}.
\begin{lemma}[Formalization support: consistency relations as bipartite consistency defects]
	\label{lem:direct-ld-consistency-defect-equivalence}
	\lean{MIPStarRE.QPBT.consRel_iff_consistencyDefect}
	\leanok
	\uses{def:consistency, def:simeq, def:povm-conventions}
	This is not a named statement of \cite{Ji2020MIPStar}. Let $\psi$ be a unit
	vector on a bipartite space with nonempty index set, let $\mu$ be a
	probability distribution on a finite question set, and let $A$ and $B$ be
	families of complete projective measurements with a common finite outcome set
	on the two factors. The consistency relation of Definition~\ref{def:simeq},
	formed for the density state of $\psi$, holds with error $\delta$ if and only
	if the off-diagonal defect of Definition~\ref{def:consistency} between the
	left placement of $A$ and the right placement of $B$, evaluated on the
	coordinate vector of $\psi$, is at most $\delta$. Passing between the two
	formulations therefore loses nothing in the error parameter.
\end{lemma}
\begin{proof}
	\leanok
	Both quantities are the same average over questions of the same finite sum of
	quadratic forms: the expectation of an operator in the density state of
	$\psi$ is the Euclidean quadratic form of the coordinate vector of $\psi$ at
	that operator, and the two placements realize the tensor factors of the
	operators appearing in either formulation.
\end{proof}

\begin{lemma}[Formalization support for the collision estimate in directly indexed coordinates]
	\label{lem:direct-ld-polynomial-agreement}
	\lean{MIPStarRE.QPBT.directPolynomialAgreement_avg_le_mdq}
	\leanok
	\uses{lem:polynomial-agreement-mdq, def:ld-game}
	This is not a named statement of \cite{Ji2020MIPStar}. Fix directly indexed
	parameters $D=(q,m,d,k)$. Two distinct $m$-variate polynomial representatives
	of individual degree at most $d$ over the scalar field of $D$ take the same
	value at a uniformly random point of $\F_q^m$ with probability at most
	$md/q$. This is the collision hypothesis of Lemma~\ref{lem:ld-sandwich},
	stated for the directly indexed point and polynomial carriers.
\end{lemma}
\begin{proof}
	\leanok
	\uses{lem:polynomial-agreement-mdq}
	The point and polynomial identifications of the directly indexed presentation
	are bijections, so the agreement event is carried to the corresponding event
	for the seed-indexed carriers, where
	Lemma~\ref{lem:polynomial-agreement-mdq} applies.
\end{proof}

\begin{lemma}[Formalization support for the coordinatewise soundness conclusion]
	\label{lem:direct-ld-coordinate-soundness}
	\lean{MIPStarRE.QPBT.directCoordinateMainFormal}
	\leanok
	\uses{thm:main-formal, def:ld-game, def:consistency, def:ld-meas}
	This is not a named statement of \cite{Ji2020MIPStar}. Fix directly indexed
	parameters $D=(q,m,d,k)$, let $\eps\in\R$, let $S$ be a projective strategy
	for the directly indexed game with value at least $1-\eps$, and fix a
	coordinate $r\in\{0,\ldots,k-1\}$. Reading coordinate $r$ of $S$ as a
	projective strategy for the $(m,q,d)$ low individual degree test of
	Chapter~\ref{chap:test} gives a strategy that passes that test with
	probability at least $1-3\eps$. Theorem~\ref{thm:main-formal}, applied to it
	at the sampling parameter $2560000\,m^3d$, then produces projective
	measurements on the two players' spaces indexed by single $m$-variate
	polynomials of individual degree at most $d$ satisfying the three consistency
	relations of that theorem, with error the error function of
	Theorem~\ref{thm:main-formal} at that sampling parameter and at pass bound
	$3\eps$. Both numerical side conditions of Theorem~\ref{thm:main-formal},
	positivity of the sampling parameter and the lower bound $400md$, hold for
	this choice because the dimension and the degree are positive.
\end{lemma}
\begin{proof}
	\leanok
	\uses{thm:main-formal}
	Instantiate Theorem~\ref{thm:main-formal} at the coordinate reading of $S$
	and at the stated sampling parameter, supplying the pass bound of that
	reading and the two numerical side conditions.
\end{proof}

\begin{lemma}[Formalization support: the polynomial-tuple conclusion for a single coordinate]
	\label{lem:direct-ld-simultaneous-single}
	\lean{MIPStarRE.QPBT.exists_directSimultaneousPolynomialMeasurements_of_k_eq_one}
	\leanok
	\uses{lem:direct-ld-coordinate-soundness, lem:direct-ld-consistency-coordinates,
		def:ld-meas, def:consistency, def:ld-game}
	This is not a named statement of \cite{Ji2020MIPStar}. It is the
	polynomial-tuple conclusion of Theorem~\ref{lem:ld-soundness} for
	simultaneity parameter $k=1$, obtained for the directly indexed game rather
	than for the seed-indexed game of Definition~\ref{def:ld-game}; see
	Remark~\ref{rem:qld-4-7-divisibility}.

	Fix directly indexed parameters $D=(q,m,d,k)$ with $k=1$, let $\eps\in\R$,
	and let $S=(\psi,A,B)$ be a projective strategy for the directly indexed game
	with value at least $1-\eps$. Then there are POVMs on the two players' spaces
	whose outcomes are $k$-tuples of $m$-variate polynomial representatives of
	individual degree at most $d$ for which the three consistency defects of the
	conclusion of Theorem~\ref{lem:ld-soundness} --- the two point/polynomial
	defects averaged uniformly over $u\in\F_q^m$ and the polynomial/polynomial
	defect over the one-element question set --- are each at most the error
	function of Theorem~\ref{thm:main-formal} at the sampling parameter
	$2560000\,m^3d$ and pass bound $3\eps$. For $k\ge2$ the coordinate
	conclusions do not yield such POVMs;
	Remark~\ref{rem:ld-soundness-simultaneity} records the obstruction.
\end{lemma}
\begin{proof}
	\leanok
	\uses{lem:direct-ld-coordinate-soundness, lem:direct-ld-consistency-coordinates}
	Apply Lemma~\ref{lem:direct-ld-coordinate-soundness} at the unique coordinate
	and transport its three consistency relations to directly indexed consistency
	defects by Lemma~\ref{lem:direct-ld-consistency-coordinates}. Since the index
	set of coordinates is a singleton, sending a value to the constant tuple
	carrying it is a bijection of outcome sets, under which each defect is
	unchanged; relabeling a single-polynomial measurement by one-coordinate
	tuples likewise leaves its effects unchanged.
\end{proof}

\begin{remark}[Constant regimes in the quotation of the classical soundness
	theorem]
	\label{rem:qld-4-7-constants}
	Theorem~\ref{lem:ld-soundness} states its error function with universal
	constants $a \ge 1$ and $0 < b \le 1$, while the source proof of
	Lemma~\ref{lem:qld-4-7} in \cite{Ji2020MIPStar} quotes it, and the
	statement of Lemma~\ref{lem:qld-4-7} is phrased, with $a > 1$ and
	$0 < b < 1$. The discrepancy is harmless: enlarging $a$ and shrinking
	$b$ only weakens a bound of this form, so any admissible pair for
	Theorem~\ref{lem:ld-soundness} yields an admissible pair for
	Lemma~\ref{lem:qld-4-7}. The closing computation of the source proof
	also mixes the terms $d/q$ and $md/q$ in the bound for $\delta_S$; the
	weaker $md/q$ is the one recorded above.
\end{remark}

\section{Scalar error estimates}

The following formalization-only auxiliary statements concern real
functions and numerical parameters, independently of the standing strategy
assumptions of this chapter. They isolate calculations supporting
Lemmas~\ref{lem:qld-xz-lines}, \ref{lem:qld-4-13}, and \ref{lem:qld-4-7}
in \cite{Ji2020MIPStar}; they assert no measurement construction.
Here $p(x)=\poly(x)$ means that $p:\R\to\R$ satisfies
$0\le p(x)\le A x^r$ for all $x\ge0$, for some $A\ge1$ and $r>0$.
Similarly, $f(x,y)=\poly(x,y)$ means
$0\le f(x,y)\le B(x^s+y^t)$ for all $x,y\ge0$, for some $B\ge1$ and
$s,t>0$. These are the corrected error conventions of
\cite{gap:qpbt_polynomial-error-square-root,gap:qpbt_pasting-product-error}.
The functions are defined on all real arguments; these bounds concern their
nonnegative domains.

\begin{theorem}[Quarter powers and square-root deficits]
	\label{thm:qld-quarter-power-deficit-bound}
	\lean{MIPStarRE.QPBT.rpow_quarter_nonneg}
	\lean{MIPStarRE.QPBT.sqrt_deficit_bound_le}
	\leanok
	For $x\in\R$, write
	\[
	 \operatorname{RePow}_{1/4}(x)=
	 \begin{cases}
	  \exp\bigl(\frac14\log x\bigr),&x>0,\\
	  0,&x=0,\\
	  \exp\bigl(\frac14\log(-x)\bigr)\cos(\pi/4),&x<0.
	 \end{cases}
	\]
	This is the real part of the principal complex quarter power. In particular,
	the negative-base branch is part of the operation rather than the ordinary
	real fourth root. It is nonnegative for every real $x$.
	If $m\geq1$ and $\delta_P,\delta_Q\geq0$, then
	\[
	 \sqrt{2\sqrt{4m^2\delta_P}+2\sqrt{4\delta_Q}}
	 \leq2\sqrt m\bigl(
	  \operatorname{RePow}_{1/4}(\delta_P)
	  +\operatorname{RePow}_{1/4}(\delta_Q)\bigr).
	\]
	These are formalization-only scalar estimates used in the real-part bounds
	of Claims 17-2 and 17-3.
\end{theorem}
\begin{proof}
	\leanok
	For a negative base, the displayed convention is a product of a positive
	exponential and $\cos(\pi/4)>0$; the nonnegative-base case is immediate.
	For the inequality, separate the two nonnegative summands under the outer
	square root, simplify the square roots of $4m^2\delta_P$ and $4\delta_Q$,
	and compare the remaining factor $2$ with $2\sqrt m$.
\end{proof}

\begin{theorem}[Polynomial error after conditioning]
	\label{thm:qld-conditioned-polynomial-bound}
	\lean{MIPStarRE.QPBT.exists_conditioned_polynomial_bound}
	\leanok
	Let $p(x)=\poly(x)$, let $f(x,y)=\poly(x,y)$, and fix $c\ge1$.
	There exists $g(x,y)=\poly(x,y)$ such that for every
	$\eps,\rho\ge0$ and every $\mu\in[1/2,1]$,
	\[
		\min\left\{1,\,
			\mu f\left(\rho,
				\frac{8p(\eps)+c(\eps+\sqrt{\eps})}{\mu}\right)+\rho
		\right\}
		\le g(\eps,\rho).
	\]
	This scalar estimate supports the pasting substitution in the proof of
	Lemma~\ref{lem:qld-xz-lines}. The interval for $\mu$ is an explicit
	hypothesis of this auxiliary estimate.
\end{theorem}
\begin{proof}
	\leanok
	Choose $A,r,B,s,t$ as in the polynomial bounds above, and set
	$\alpha=\min\{r,1/2\}$, $\beta=\min\{s,1\}$, and $L=8A+2c$.
	For $0\le\eps,\rho\le1$, the numerator in the argument of $f$ is at
	most $L\eps^\alpha$, and division by $\mu\ge1/2$ costs at most a
	factor of two. The polynomial bound on $f$, together with $\mu\le1$,
	therefore gives the required inequality with
	\[
		g(\eps,\rho)=\bigl(B(1+(2L)^t)+1\bigr)
			\bigl(\eps^{\alpha t}+\rho^\beta\bigr).
	\]
	If either $\eps>1$ or $\rho>1$, this function is at least one, which
	bounds the minimum. All its exponents are positive.
\end{proof}

\begin{theorem}[Polynomial composition with the dimension factor]
	\label{thm:qld-combining-polynomial-bound}
	\lean{MIPStarRE.QPBT.exists_combining_polynomial_bound}
	\leanok
	Let $p(x)=\poly(x)$, let $f(x,y)=\poly(x,y)$, and fix $c\ge0$.
	There exists $g(x,y)=\poly(x,y)$ such that for all
	$\eps,\rho\ge0$ and all real $D\ge1$,
	\[
		\min\left\{1,\,
			c\left(p(\eps)^{1/2}+\sqrt D
				\left(f(\eps,\rho)^{1/4}
					+p(\eps)^{1/4}+\eps^{1/4}\right)\right)
		\right\}
		\le D\,g(\eps,\rho).
	\]
	The factor $D$ is retained in this numerical consequence of the first
	route in the proof of Lemma~\ref{lem:qld-4-13}, as explained in
	Remark~\ref{rem:qld-4-13-source-defects}.
\end{theorem}
\begin{proof}
	\leanok
	Choose $A,r,B,s,t$ as in the polynomial bounds and put
	$\alpha=\min\{r/4,s/4,1/4\}$ and
	$K=1+c(A^{1/2}+B^{1/4}+A^{1/4}+1)$.
	For $0\le\eps\le1$, use
	$(u+v)^{1/4}\le u^{1/4}+v^{1/4}$ for nonnegative $u,v$ to bound
	$f(\eps,\rho)^{1/4}$, compare the powers of $\eps$ with
	$\eps^\alpha$, and use $\sqrt D\le D$. This gives the inequality
	with $g(\eps,\rho)=K(\eps^\alpha+\rho^{t/4})$. For $\eps>1$,
	$Dg(\eps,\rho)\ge1$, so the cap proves the same bound.
\end{proof}

\begin{theorem}[Incorrect outer projection of the ordered product]
\label{thm:qld-retained-local-mismatch}
\lean{MIPStarRE.QPBT.PolynomialImageBounds.orderedIndicator_gram_le_one}
\lean{MIPStarRE.QPBT.PolynomialImageBounds.avg_wrong_ordered_norm_sq_le}
\lean{MIPStarRE.QPBT.PolynomialImageBounds.point_mismatch_le_ordered_error}
\leanok
Let $X,Z$ be complete projective measurements with outcomes in a finite
field $F$, let $r,t\in F$, and let $v$ be any vector. Set
$T_{\alpha,\beta}=\sum_{\alpha b+\beta c=\alpha r+\beta t}X_bZ_c$.
Each $T_{\alpha,\beta}$ is a contraction, and
\[
 \mathbb E_{\alpha,\beta}\|(I-X_r)T_{\alpha,\beta}v\|^2
 \leq \frac{2}{|F|}\|v\|^2,
 \qquad
 \|(I-X_r)v\|^2\leq
 2\mathbb E_{\alpha,\beta}\|v-T_{\alpha,\beta}v\|^2
 +\frac{4}{|F|}\|v\|^2.
\]
No normalization or commutation of $X$ with $Z$ is required.
\end{theorem}
\begin{proof}
\leanok
When $\beta\ne0$, orthogonality identifies the squared norm with the
sum of $Z_cX_bZ_c$ selected by the scalar check and $b\ne r$.
The affine collision probability is $1/|F|$ for each such pair. Sum its
nonnegative sandwich weights before using completeness. The exception
$\beta=0$ contributes at most another $\|v\|^2/|F|$ by contraction.
Apply the squared triangle inequality to
$(I-X_r)v=(I-X_r)(v-Tv)+(I-X_r)Tv$.
\end{proof}

\begin{theorem}[Retained point overlap on both register placements]
\label{thm:qld-retained-point-overlap}
\lean{MIPStarRE.QPBT.ExtendedLineGame.placed_fiber_orderedIndicator}
\lean{MIPStarRE.QPBT.ExtendedLineGame.retainedPointMismatch}
\lean{MIPStarRE.QPBT.ExtendedLineGame.retainedPointMismatch_le_ordered_error}
\lean{MIPStarRE.QPBT.ExtendedLineGame.retainedPointOverlap}
\lean{MIPStarRE.QPBT.ExtendedLineGame.retainedPointMismatch_eq}
\lean{MIPStarRE.QPBT.ExtendedLineGame.retainedPointOverlap_ge_ordered_error}
\leanok
\uses{def:qld-direct-pair-completion, thm:qld-retained-local-mismatch}
Let $R$ be a complete projective singleton-polynomial measurement on a
placement $p_1$, and let $p_2$ be opposite. For a bounded pair $g$, put
$v_g=(R_{j_{\rm dir}(g)})_{p_1}\hat\psi$. Define
\[
 K_W=\mathbb E_u\sum_g\|(I-\hat M^{W,u}_{g_W(u)})_{p_2}v_g\|^2,
 \qquad
 H_W=\mathbb E_u\sum_g
 \langle\hat\psi|(R_{j_{\rm dir}(g)})_{p_1}
   (\hat M^{W,u}_{g_W(u)})_{p_2}|\hat\psi\rangle.
\]
Write $G$ for the actual mass outside the bounded combining image, and
$E^{XZ},E^{ZX}$ for the full ordered errors. Then
\[
 K_W=1-G-H_W,\qquad
 K_X\leq2E^{XZ}+4/q,\qquad K_Z\leq2E^{ZX}+4/q.
\]
Thus $H_X\geq1-G-2E^{XZ}-4/q$ and
$H_Z\geq1-G-2E^{ZX}-4/q$. These statements hold for every opposite
placement, including both $AA'\mid BA''$ and $BB'\mid AB''$.
\end{theorem}
\begin{proof}
\leanok
\uses{thm:qld-retained-local-mismatch, thm:qld-projective-outcome-mass,
  thm:qld-direct-pair-completion-properties}
Apply Theorem~\ref{thm:qld-retained-local-mismatch} on each actual vector
$v_g$, reversing both base and scalar blocks for Z. The joint coordinate
equivalence preserves the uniform law and the placed ordered product.
Injectivity bounds the retained sums by the full outcome sums.
Projectivity and normalization give total squared mass one. Opposite
placements commute, yielding the exact identity for $K_W$.
\end{proof}

\begin{theorem}[Consistency of the direct polynomial-pair completion]
\label{thm:qld-direct-pair-consistency}
\lean{MIPStarRE.QPBT.ExtendedLineGame.directPairMeasurement_consistency_le}
\leanok
\uses{def:qld-direct-pair-completion, thm:qld-retained-point-overlap,
  def:consistency}
Let $R$ be a complete projective singleton-polynomial measurement on a
placement $p_1$, let $p_2$ be opposite, and complete the direct restriction of
$R$ at a fixed polynomial pair. For either basis $W$, the consistency defect of
the completed pair measurement against the point measurement on $p_2$ is at
most
\[
 G_{p_1}(R)+K_W,
\]
where $G_{p_1}(R)$ is the mass outside the bounded combining image and $K_W$ is
the retained mismatch of Theorem~\ref{thm:qld-retained-point-overlap}.
\end{theorem}
\begin{proof}
\leanok
\uses{thm:qld-retained-point-overlap,
  thm:qld-direct-pair-completion-properties,
  lem:qld-combined-pair-completion}
The completion adds a positive operator at the distinguished pair. On opposite
placements this operator commutes with each point effect, so the diagonal
overlap cannot decrease. The defect is therefore bounded by one minus the
retained overlap, which equals the excluded mass plus the retained mismatch.
\end{proof}

\begin{theorem}[Ordered-error bound for the direct polynomial-pair completion]
\label{thm:qld-direct-pair-ordered-consistency}
\lean{MIPStarRE.QPBT.ExtendedLineGame.directPairMeasurement_consistency_le_ordered_error}
\leanok
\uses{thm:qld-direct-pair-consistency, thm:qld-retained-point-overlap,
  def:qld-extended-polynomial-ordered-error}
In the setting of Theorem~\ref{thm:qld-direct-pair-consistency}, let
$o\in\{XZ,ZX\}$ select the corresponding point basis $W$. Then the completed
measurement has point-consistency defect at most
\[
 G_{p_1}(R)+2E^o_{p_1,p_2}(R)+\frac4q.
\]
\end{theorem}
\begin{proof}
\leanok
\uses{thm:qld-direct-pair-consistency, thm:qld-retained-point-overlap}
Apply Theorem~\ref{thm:qld-direct-pair-consistency} and substitute the retained
mismatch estimate from Theorem~\ref{thm:qld-retained-point-overlap}.
\end{proof}

\begin{theorem}[Completed pair witness with the actual rounded error]
\label{thm:qld-completed-pair-actual-error}
\lean{MIPStarRE.QPBT.ExtendedLineGame.exists_pairWitness_of_points_lines}
\leanok
\uses{def:qld-direct-pair-completion, thm:qld-rounded-separated-mass,
  thm:qld-direct-pair-ordered-consistency,
  thm:qld-direct-pair-completion-properties}
There are universal $a\geq1$ and $0<b\leq1$ with the following property.
Given an admissible projective setting, a combined-point witness of error
$\delta_Q$, and a directly indexed completed-answer extended-line witness
of error $\delta_L$ for that same point witness, put
\[
 \delta=\delta_{\rm ld}\left(a,b,
  3\left(\sqrt{\delta_Q+\delta_L}+md/q\right),q,2m+2,d,1\right),
 \quad \eta=\delta+\sqrt{220\delta^{1/4}}+2\sqrt{2\delta},
 \quad E=4\eta+8\delta_Q.
\]
The actual rounded measurements give complete projective polynomial-pair
measurements whose four point-consistency defects, for X and Z on each
player placement, are all at most
\[
 8E+\frac{12md+4d+14}{q}.
\]
This auxiliary explicitly assumes the two displayed witnesses.
Lemma~\ref{lem:qld-4-13-established} supplies them from the original strategy,
and Theorem~\ref{thm:qld-actual-rounded-error-bound} absorbs their actual error.
\end{theorem}
\begin{proof}
\leanok
\uses{def:qld-direct-pair-completion, thm:qld-rounded-separated-mass,
  thm:qld-direct-pair-ordered-consistency,
  thm:qld-direct-pair-completion-properties}
Apply Theorem~\ref{thm:qld-rounded-separated-mass} to obtain one rounded
measurement on each player space with both ordered estimates and the stated
excluded-mass bound. Complete each measurement at the zero polynomial pair.
The completion is projective by
Theorem~\ref{thm:qld-direct-pair-completion-properties}, and
Theorem~\ref{thm:qld-direct-pair-ordered-consistency} bounds each point defect.
Combine these bounds with $G\leq6E+(12md+4d+10)/q$.
The two rounded measurements are chosen before either basis estimate, so
one completed measurement on each player side supplies both conclusions.
\end{proof}

\begin{theorem}[Absorption into the global polynomial-pair error]
	\label{thm:qld-global-pair-error-bound}
	\lean{MIPStarRE.QPBT.exists_global_pair_error_bound}
	\leanok
	\uses{def:admissible}
	Let $p(x)=\poly(x)$ and $g(x,y)=\poly(x,y)$.
	Fix $c,C\ge0$, $a\ge1$, and $0<b\le1$, and write the numerical
	low-degree error with its constants explicit as
	\[
		\delta_{\mathrm{ld}}(a,b,\eta,q,n,d,k)
		=a(dnk)^a\bigl(\eta^b+q^{-b}+2^{-bnd}\bigr).
	\]
	There exist $a'>1$ and $0<b'<1$, depending only on these fixed
	functions and constants, such that for every admissible tuple
	$(q,m,d)$ and every $\eps\ge0$, setting $\rho=md/q$ gives
	\[
		\begin{split}
		\min\bigl\{1,\,
			C\bigl(&\delta_{\mathrm{ld}}
				(a,b,cm\,g(\eps,\rho)+\eps,q,2m+2,d,1)\\
				&+\sqrt{p(\eps)}+\rho\bigr)\bigr\}
		\le a'(md)^{a'}
			\bigl(\eps^{b'}+q^{-b'}+2^{-b'md}\bigr).
		\end{split}
	\]
	This is the scalar substitution in the proof of
	Lemma~\ref{lem:qld-4-7}, including the slack $\eps$ in the
	low-degree input. It supplies no projective measurement or comparison
	between games.
\end{theorem}
\begin{proof}
	\leanok
	\uses{def:admissible}
	Choose bounds $p(\eps)\le A\eps^r$ and
	$g(\eps,\rho)\le B(\eps^s+\rho^t)$, and take
	$b'=\min\{1/2,b,sb,tb,r/2\}$. Admissibility gives $m,d,q\ge1$.
	For $0\le\eps\le1$, subadditivity of the $b$-th power separates the
	low-degree input into its polynomial error and slack terms. Use
	$\rho^{tb}=(md)^{tb}q^{-tb}$, $m\le md$, and
	$md\le(2m+2)d\le4md$ to bound all terms by a fixed constant times
	a fixed power of $md$ times
	$\eps^{b'}+q^{-b'}+2^{-b'md}$. Enlarging $a'$ absorbs both the
	constant and the power of $md$. For $\eps>1$, this final expression
	is at least one and hence bounds the capped error.
\end{proof}

\begin{theorem}[Direct-line real-part overlap]
\label{thm:qld-direct-line-real-overlap}
\lean{MIPStarRE.QPBT.subline_joint_overlap_near_one_at}
\leanok
\uses{def:qld-subline-overlaps, lem:qld-first-route-overlap-components}
This auxiliary concerns the actual directly indexed subline law. Let the
combined point and paired-line witnesses have errors $Q$ and $L$, and let
$I$ be their paired overlap, averaged over the subline law and a uniform
affine parameter, using completed evaluations of both line polynomials.
For every directed opposite placement there is a universal $C>0$ such that
\[
 |I-1|\leq C\left(Q^{1/2}+\sqrt m
   (L^{1/4}+Q^{1/4}+\eps^{1/4})\right).
\]
\end{theorem}
\begin{proof}
\leanok
\uses{def:qld-subline-overlaps, lem:qld-first-route-overlap-components}
Replace the joint points by the ordered product, remove the X factor by
Cauchy--Schwarz, and bound the remaining Z deficit. The one-point mixture
laws suffice; no joint conditional product law is assumed.
\end{proof}

\begin{theorem}[Direct passing-error absorption]
\label{thm:qld-direct-passing-error-bound}
\lean{MIPStarRE.QPBT.exists_direct_passing_polynomial_bound}
\lean{MIPStarRE.QPBT.exists_direct_global_pair_error_bound}
\leanok
Let $p(x)=\poly(x)$ and $g(x,y)=\poly(x,y)$, with the corrected sum
convention, and fix $C,s\geq0$, $a\geq1$, $0<b\leq1$.
For $x,y\geq0$ and $n\geq1$, the capped quantity
$\min\{1,3(\sqrt{p(x)+Cn g(x,y)}+y)\}$ is at most $n h(x,y)$ for a
polynomially small function $h$. For admissible parameters, put
$\rho=md/q$ and $v=3(\sqrt{p(\eps)+Cm g(\eps,\rho)}+\rho)$.
There exist universal $A_0>1$ and $0<B_0<1$ such that, for every $\eps\geq0$,
\[
 \min\{1,s(\delta_{\rm ld}(a,b,v+\eps,q,2m+2,d,1)
     +\sqrt{p(\eps)}+\rho)\}
 \leq\delta_{\rm qld}(A_0,B_0,\eps,m,d,q).
\]
\end{theorem}
\begin{proof}
\leanok
\uses{thm:qld-global-pair-error-bound}
Bound the capped passing error by $m h(\eps,\rho)$ and apply the existing
polynomial absorption. When $v\geq1$, the capped low-degree error is at
least one, which justifies replacing the capped passing error by the full v.
\end{proof}

\begin{theorem}[Square roots and coefficients of the global error]
\label{thm:pauli-error-square-root-support}
\lean{MIPStarRE.QPBT.sqrt_deltaQld_le}
\lean{MIPStarRE.QPBT.scale_deltaQld_le}
\leanok
For admissible parameters, $a\geq1$, $\eps\geq0$, $b\in\mathbb R$ and
$K\geq1$, the global error satisfies
\[
 \sqrt{\delta_{\rm qld}(a,b,\eps,m,d,q)}
 \leq\delta_{\rm qld}(a,b/2,\eps,m,d,q),\qquad
 K\delta_{\rm qld}(a,b,\eps,m,d,q)
 \leq\delta_{\rm qld}(Ka,b,\eps,m,d,q).
\]
\end{theorem}
\begin{proof}
\leanok
The square root is subadditive on the three nonnegative error terms,
and the prefactor is at least one. Increasing the prefactor exponent
preserves order because $md\geq1$. Neither argument restricts $\eps$ to one.
\end{proof}

\begin{theorem}[Absorption of the actual rounded error]
\label{thm:qld-actual-rounded-error-bound}
\lean{MIPStarRE.QPBT.actual_rounding_error_le_root}
\lean{MIPStarRE.QPBT.exists_actual_rounded_global_pair_error_bound}
\leanok
Fix polynomially small $p,g$, $C\geq0$, $a\geq1$ and $0<b\leq1$.
For admissible parameters and $\eps\geq0$, put $Q=p(\eps)$,
$L=Cm g(\eps,md/q)$ and define $\delta,\eta,D$ as in the proof of
Lemma~\ref{lem:qld-4-7}. There are universal $A>1$, $0<B<1$, chosen
before these parameters, such that
\[
 \min\{1,D\}\leq\delta_{\rm qld}(A,B,\eps,m,d,q).
\]
More generally, if $\delta,Q\geq0$, $\delta\leq t$, $\sqrt Q\leq t$ and
$\rho\leq t$, then
\[
 \min\{1,32\delta+32\sqrt{220\delta^{1/4}}+
 64\sqrt{2\delta}+64Q+30\rho\}
 \leq1024\sqrt{\sqrt{\sqrt{\min\{1,t\}}}}.
\]
\end{theorem}
\begin{proof}
\leanok
\uses{thm:qld-direct-passing-error-bound, thm:pauli-error-square-root-support}
Take $t=\delta_{\rm ld}(a,b,v+\eps,q,2m+2,d,1)+\sqrt Q+md/q$,
with the actual passing error $v=3(\sqrt{Q+L}+md/q)$.
Monotonicity gives $\delta\leq t$. If $t\leq1$, all rounding terms
are bounded by their coefficients times $t^{1/8}$, and $Q\leq t$.
Also $(12md+4d+14)/q\leq30md/q$. Their coefficient sum is at most 1024.
If $t\geq1$, use the unit cap. Apply the square-root bound three times
to the direct absorption constants $A_0,B_0$ and choose
$A=1024A_0$, $B=B_0/8$.
\end{proof}

\section{Closed quantitative baseline}
\label{sec:qld-combining-explicit}

The results in this section expose the constants used by the current formal
proof.  They are Lean-only quantitative specializations of the constructions
above.  All ratios below are real numbers, and for admissible parameters we
write
\[
  \rho=\frac{md}{q}.
\]

\begin{definition}[Explicit pasting error]
\label{def:qld-heterogeneous-pasting-error-explicit}
\lean{MIPStarRE.QPBT.heterogeneousPastingError}
\leanok
The heterogeneous one-sided pasting error is
\[
  H(\eta,\delta)=115\bigl(\eta^{1/4}+\delta^{1/8}\bigr).
\]
\end{definition}

\begin{theorem}[Explicit commutator and heterogeneous pasting bounds]
\label{thm:qld-pasting-explicit}
\lean{MIPStarRE.QPBT.op_dist_sq_commutator_le_explicit}
\lean{MIPStarRE.QPBT.op_dist_sq_commutator_right_le_explicit}
\lean{MIPStarRE.QPBT.coarse_commutator_bound_explicit}
\lean{MIPStarRE.QPBT.pasting_error_of_marginal_consistency_explicit}
\lean{MIPStarRE.QPBT.pasting_error_heterogeneous_explicit}
\lean{MIPStarRE.QPBT.heterogeneous_pasting_error_is_poly_err₂}
\leanok
\uses{def:qld-heterogeneous-pasting-error-explicit, lem:pasting}
If two marginals are each within squared operator-family distance $\delta$ of
the corresponding marginal of one projective joint measurement, then their
commutator family has squared distance at most $16\delta$.  This holds on
either tensor factor.  In the codeword-pasting application, the two forward
marginal defects therefore give the coarse commutator bound $32\delta$.

Let the sampling law be a probability distribution, let the state be
normalized, and let $\eta,\delta\geq0$.  If the second codeword family has
conditional collision at most $\eta$, the joint and second-codeword
measurements are projective, and the two forward marginal consistency defects
are at most $\delta$, then the pasted joint family has consistency defect at
most $H(\eta,\delta)$.  The same statement holds for unequal
finite-dimensional local spaces.  Moreover, $H$ is polynomially small on the
nonnegative quadrant, with coefficient $115$ and exponents $1/4$ and $1/8$.
\end{theorem}
\begin{proof}
\leanok
The commutator estimate compares both products to the projective refinement.
The pasting estimate then applies Cauchy--Schwarz to the collision term and
uses the coarse commutator bound; adjoining fixed unit vectors transports the
same calculation to unequal local spaces.
\end{proof}

\begin{theorem}[Explicit combined-point construction]
\label{thm:qld-combined-points-explicit}
\lean{MIPStarRE.QPBT.ProjectiveSetting.avg_sandwich_defect_bound_le_explicit}
\lean{MIPStarRE.QPBT.ProjectiveSetting.sandwich_point_ordered_dist_le_explicit}
\lean{MIPStarRE.QPBT.ProjectiveSetting.ordered_cross_dist_le_explicit}
\lean{MIPStarRE.QPBT.ProjectiveSetting.exp_point_comm_explicit}
\lean{MIPStarRE.QPBT.exists_combined_points_witness_explicit}
\leanok
\uses{thm:qld-expanded-comparisons-explicit,
thm:qld-twisted-commutation-explicit, lem:qld-4-10}
For opposite placements, the average sandwich defect is at most
\[
  (344+3C_{\mathrm{tw}})(\eps+\sqrt\eps).
\]
On one placement the sandwich measurement is within
$C_{\mathrm{tw}}\sqrt\eps$ of the ordered product, the ordered products on
opposite placements are within $688\eps$, and reversing the two factors on one
placement costs at most $C_{\mathrm{tw}}\sqrt\eps$.

Consequently there is a projective combined-point witness whose error is
\[
  E_{\mathrm{pt}}(\eps)=C_{\mathrm{pt}}\eps^{1/8},
\]
where $C_{\mathrm{pt}}$ is defined in
Definition~\ref{def:qld-baseline-scalar-functions} below.  The witness contains
all four directed opposite-placement comparisons and both ordered marginals.
\end{theorem}
\begin{proof}
\leanok
Orthonormalize the sandwich measurement pointwise and chain the four displayed
distance estimates.  The eighth root absorbs the square roots introduced by
orthonormalization and by the triangle inequalities.
\end{proof}

\begin{theorem}[Explicit combined-line construction]
\label{thm:qld-combined-lines-explicit}
\lean{MIPStarRE.QPBT.exp_line_point_same_placement_distance_le_explicit}
\lean{MIPStarRE.QPBT.combined_points_line_marginal_distance_le_explicit}
\lean{MIPStarRE.QPBT.combined_points_line_marginal_defect_le_explicit}
\lean{MIPStarRE.QPBT.combined_points_conditioned_line_marginal_defect_le_explicit}
\lean{MIPStarRE.QPBT.combined_line_conditioned_defect_le_explicit}
\lean{MIPStarRE.QPBT.combined_line_restored_defect_le_explicit}
\lean{MIPStarRE.QPBT.combined_line_measurement_consistency_explicit}
\lean{MIPStarRE.QPBT.exists_combined_lines_witness_of_points_witness_explicit}
\leanok
\uses{thm:qld-pasting-explicit, thm:qld-combined-points-explicit,
lem:qld-xz-lines}
Let a supplied combined-point witness have error $\delta_Q$, and let $r$ be
the mass of the nondegenerate line-pasting event.  The same-placement
line--point distance is at most $1040(\eps+\sqrt\eps)$.  Each of the two
completed point-to-line marginal distances, and hence each corresponding
consistency defect, is at most
\[
  8\delta_Q+2080(\eps+\sqrt\eps).
\]
After conditioning, both defects are at most this quantity divided by $r$.
The conditioned pasted line defect is therefore bounded by
\[
 H\left(\rho,\frac{8\delta_Q+2080(\eps+\sqrt\eps)}{r}\right),
\]
and restoring the discarded zero direction gives
\[
 rH\left(\rho,\frac{8\delta_Q+2080(\eps+\sqrt\eps)}{r}\right)
   +\frac1{2q}.
\]
The consistency defect is independently at most one.  Consequently, for
$\delta_Q=E_{\mathrm{pt}}(\eps)$,
\[
 \min\!\left\{1,
 rH\!\left(\rho,\frac{8E_{\mathrm{pt}}(\eps)
   +2080(\eps+\sqrt\eps)}{r}\right)+\frac1{2q}\right\}
 \leq E_{\mathrm{line}}(\eps,\rho),
\]
and hence there is a combined-line witness with precisely that error.  All
measurement outcomes in these comparisons belong to the completed field
alphabet.
\end{theorem}
\begin{proof}
\leanok
\uses{thm:qld-baseline-scalar-bounds}
Use the expanded line--point comparison twice, condition on nonzero direction,
apply Theorem~\ref{thm:qld-pasting-explicit}, and restore the discarded mass.
The final scalar inequality is included in
Theorem~\ref{thm:qld-baseline-scalar-bounds}.
\end{proof}

\begin{theorem}[Explicit subline overlap estimates]
\label{thm:qld-subline-overlaps-explicit}
\lean{MIPStarRE.QPBT.subline_remove_x_factor_at_explicit}
\lean{MIPStarRE.QPBT.subline_z_term_near_one_at_explicit}
\lean{MIPStarRE.QPBT.subline_joint_overlap_near_one_at_explicit}
\leanok
\uses{thm:qld-combined-lines-explicit, lem:qld-sublines}
For point error $\delta_Q$, line error $\delta_P$, and any directed opposite
placements, removing the $X$ factor from the ordered subline overlap costs at
most
\[
  2\sqrt m\bigl(\delta_P^{1/4}+\delta_Q^{1/4}\bigr).
\]
The remaining $Z$ overlap differs from one by at most
\[
  2\sqrt m\bigl(\delta_P^{1/4}+\delta_Q^{1/4}+\eps^{1/4}\bigr).
\]
Thus the paired subline overlap $I$ satisfies the closed estimate
\[
 |I-1|\leq6\left(\delta_Q^{1/2}+\sqrt m
   (\delta_P^{1/4}+\delta_Q^{1/4}+\eps^{1/4})\right).
\]
\end{theorem}
\begin{proof}
\leanok
Apply Cauchy--Schwarz to remove the first point factor, use the line and point
consistency defects for the remaining factor, and chain the three overlap
comparisons.
\end{proof}

\begin{theorem}[Explicit extended-line witness]
\label{thm:qld-extended-lines-explicit}
\lean{MIPStarRE.QPBT.exists_extended_lines_witness_established_of_points_witness_explicit}
\leanok
\uses{thm:qld-subline-overlaps-explicit}
Given the witness of Theorem~\ref{thm:qld-combined-points-explicit}, there is a
directly indexed extended-line witness with completed answer alphabet and
error
\[
  mE_{\mathrm{ext}}(\eps,\rho).
\]
The conclusion includes both player orientations and the direct axis-degree
bound.
\end{theorem}
\begin{proof}
\leanok
\uses{thm:qld-baseline-scalar-bounds}
Use the subline measurement and the paired-overlap estimate, then absorb the
result with the scalar combining inequality of
Theorem~\ref{thm:qld-baseline-scalar-bounds}.
\end{proof}

\begin{theorem}[Explicit one-coordinate direct low-degree soundness]
\label{thm:qld-direct-soundness-explicit}
\lean{MIPStarRE.QPBT.direct_ld_transport_constants_explicit}
\lean{MIPStarRE.QPBT.direct_ld_soundness_of_k_eq_one_explicit}
\lean{MIPStarRE.QPBT.direct_ld_soundness_of_k_eq_one_any_strategy_explicit}
\leanok
\uses{lem:ld-soundness}
For every $C_0\geq1$, the transport estimate is absorbed by the direct
low-degree error with constants
\[
  a=2500000000C_0,\qquad b=\frac1{80000}.
\]
More precisely, for $0<\eta\leq1$ it bounds
\[
 C_0k\sqrt{\delta_{\mathrm{LDT}}(3\eta)+md/q+\eta}
 \leq\delta_{\mathrm{ld}}(a,b,\eta,q,m,d,k).
\]
When $k=1$, every projective strategy of value at least $1-\eta$ has two
global polynomial-tuple measurements satisfying the two point-consistency
bounds and their mutual consistency bound, each at
\[
 \delta_{\mathrm{ld}}(2500000000,1/80000,\eta,q,m,d,1).
\]
The same three bounds hold for an arbitrary strategy after Naimark dilation
and compression; projectivity of the compressed polynomial measurements is
not asserted.
\end{theorem}
\begin{proof}
\leanok
Specialize the quantitative low-degree transport to one simultaneous
coordinate.  For an arbitrary strategy, apply the projective result to its
questionwise Naimark dilation and compress the resulting measurements.
\end{proof}

\begin{theorem}[Explicit rounded polynomial estimates]
\label{thm:qld-rounded-polynomials-explicit}
\lean{MIPStarRE.QPBT.ExtendedLineGame.rounded_polynomial_ordered_estimates_explicit}
\lean{MIPStarRE.QPBT.ExtendedLineGame.rounded_polynomial_scalar_mass_explicit}
\lean{MIPStarRE.QPBT.ExtendedLineGame.rounded_polynomial_wrong_variable_mass_explicit}
\lean{MIPStarRE.QPBT.ExtendedLineGame.rounded_polynomial_separated_mass_explicit}
\lean{MIPStarRE.QPBT.ExtendedLineGame.pair_witness_of_points_lines_explicit}
\leanok
\uses{thm:qld-direct-soundness-explicit, thm:qld-extended-lines-explicit,
lem:qld-4-7}
For supplied point and extended-line errors $\delta_Q,\delta_L$, put
\[
 \begin{split}
 \delta&=\delta_{\mathrm{ld}}\left(2500000000,\frac1{80000},
   3(\sqrt{\delta_Q+\delta_L}+\rho),q,2m+2,d,1\right),\\
 \eta&=\delta+\sqrt{220\delta^{1/4}}+2\sqrt{2\delta},
 \qquad E=4\eta+8\delta_Q.
 \end{split}
\]
There are projective rounded polynomial measurements for Alice and Bob whose
ordered errors, in both orderings and on both player placements, are at most
$E$.  There also exist such measurements satisfying, in addition, the two
scalar-nonlinearity mass bounds
\[
 2E+2\frac{(2m+2)d+1}{q}.
\]
There exist such measurements satisfying those ordered and scalar bounds and,
for both orderings, the two wrong-variable mass bounds
\[
 2E+2\frac{2md+2}{q}.
\]
Separately, there exist projective rounded measurements satisfying the ordered
bounds and the two non-separated-mass bounds
\[
 6E+\frac{12md+4d+10}{q}.
\]
Finally, the supplied point and extended-line witnesses determine a projective
global polynomial-pair witness with error
\[
 8E+\frac{12md+4d+14}{q}.
\]
Each existential assertion concerns the measurements constructed in its own
linked declaration; no conjunction between independently quantified witnesses
is asserted.  All four estimates use the directly indexed, completed-answer
extended-line witness; no unproved source-domain conversion is hidden in these
statements.
\end{theorem}
\begin{proof}
\leanok
Apply Theorem~\ref{thm:qld-direct-soundness-explicit}, round both polynomial
measurements on their own spaces, and combine the ordered estimates with the
three finite-field root-counting bounds.
\end{proof}

\begin{definition}[Scalar functions of the closed baseline]
\label{def:qld-baseline-scalar-functions}
\lean{MIPStarRE.QPBT.pauliBaselinePointConstant}
\lean{MIPStarRE.QPBT.pauliBaselinePointError}
\lean{MIPStarRE.QPBT.pauliBaselineLineConstant}
\lean{MIPStarRE.QPBT.pauliBaselineLineError}
\lean{MIPStarRE.QPBT.pauliBaselineExtendedLineConstant}
\lean{MIPStarRE.QPBT.pauliBaselineExtendedLineError}
\lean{MIPStarRE.QPBT.pauliBaselinePassingConstant}
\lean{MIPStarRE.QPBT.pauliBaselinePassingError}
\lean{MIPStarRE.QPBT.pauliBaselineLowDegreeConstant}
\lean{MIPStarRE.QPBT.pauliBaselineLowDegreePower}
\lean{MIPStarRE.QPBT.pauliBaselineGlobalAbsorptionConstant}
\lean{MIPStarRE.QPBT.pauliBaselineGlobalAbsorptionPower}
\lean{MIPStarRE.QPBT.pauliBaselineGlobalPairConstant}
\lean{MIPStarRE.QPBT.pauliBaselineGlobalPairPower}
\leanok
\uses{def:pauli-observable-baseline-constants}
Define
\[
\begin{aligned}
C_{\mathrm{pt}}&=10560(344+3C_{\mathrm{tw}})+24C_{\mathrm{tw}}+11008,
&E_{\mathrm{pt}}(x)&=C_{\mathrm{pt}}x^{1/8},\\
C_{\mathrm{line}}&=115\left(1+(16C_{\mathrm{pt}}+8320)^{1/8}\right)+1,
&E_{\mathrm{line}}(x,y)&=C_{\mathrm{line}}(x^{1/64}+y^{1/4}),\\
C_{\mathrm{ext}}&=1+6\left(\sqrt{C_{\mathrm{pt}}}
 +C_{\mathrm{line}}^{1/4}+C_{\mathrm{pt}}^{1/4}+1\right),
&E_{\mathrm{ext}}(x,y)&=C_{\mathrm{ext}}(x^{1/256}+y^{1/16}),\\
C_{\mathrm{pass}}&=4+3\sqrt{C_{\mathrm{pt}}+C_{\mathrm{ext}}},
&E_{\mathrm{pass}}(x,y)&=C_{\mathrm{pass}}(x^{1/512}+y^{1/32}).
\end{aligned}
\]
The direct low-degree constants are
\[
 C_{\mathrm{ld}}=2500000000,
 \qquad b_{\mathrm{ld}}=\frac1{80000}.
\]
Finally, set
\[
\begin{aligned}
C_{\mathrm{abs}}={}&2\left(C_{\mathrm{ld}}4^{C_{\mathrm{ld}}}
  (C_{\mathrm{pass}}^{b_{\mathrm{ld}}}+2)+\sqrt{C_{\mathrm{pt}}}+1\right)
  +C_{\mathrm{ld}}+\frac{33}{32}b_{\mathrm{ld}}+3,\\
b_{\mathrm{abs}}={}&\frac1{40960000},\qquad
C_{\mathrm{pair}}=1024C_{\mathrm{abs}},\qquad
b_{\mathrm{pair}}=\frac1{327680000}.
\end{aligned}
\]
\end{definition}

\begin{theorem}[Polynomiality and scalar absorption of the baseline]
\label{thm:qld-baseline-scalar-bounds}
\lean{MIPStarRE.QPBT.one_le_pauli_baseline_point_constant}
\lean{MIPStarRE.QPBT.one_le_pauli_baseline_line_constant}
\lean{MIPStarRE.QPBT.one_le_pauli_baseline_extended_line_constant}
\lean{MIPStarRE.QPBT.one_le_pauli_baseline_passing_constant}
\lean{MIPStarRE.QPBT.pauli_baseline_point_error_is_poly_err}
\lean{MIPStarRE.QPBT.pauli_baseline_line_error_is_poly_err₂}
\lean{MIPStarRE.QPBT.pauli_baseline_extended_line_error_is_poly_err₂}
\lean{MIPStarRE.QPBT.pauli_baseline_passing_error_is_poly_err₂}
\lean{MIPStarRE.QPBT.pauli_baseline_conditioned_line_bound}
\lean{MIPStarRE.QPBT.pauli_baseline_combining_bound}
\lean{MIPStarRE.QPBT.pauli_baseline_direct_passing_bound}
\lean{MIPStarRE.QPBT.pauli_baseline_global_absorption_bound}
\lean{MIPStarRE.QPBT.pauli_baseline_direct_global_pair_error_bound}
\lean{MIPStarRE.QPBT.pauli_baseline_actual_rounded_global_pair_error_bound}
\leanok
\uses{def:qld-baseline-scalar-functions,
thm:qld-global-pair-bound-explicit,
thm:qld-actual-rounded-error-bound,
thm:pauli-error-square-root-support}
The four coefficients $C_{\mathrm{pt}},C_{\mathrm{line}},C_{\mathrm{ext}}$
and $C_{\mathrm{pass}}$ are at least one.  Their error functions are
polynomially small on the nonnegative domain with the coefficients and powers
displayed in Definition~\ref{def:qld-baseline-scalar-functions}.

If $x,y\geq0$, $1/2\leq r\leq1$, and $n\geq1$, then
\[
 \min\left\{1,rH\left(y,
   \frac{8E_{\mathrm{pt}}(x)+2080(x+\sqrt x)}r\right)+y\right\}
 \leq E_{\mathrm{line}}(x,y),
\]
\[
 \min\left\{1,6\left(E_{\mathrm{pt}}(x)^{1/2}+\sqrt n
   \left(E_{\mathrm{line}}(x,y)^{1/4}
    +E_{\mathrm{pt}}(x)^{1/4}+x^{1/4}\right)\right)\right\}
 \leq nE_{\mathrm{ext}}(x,y),
\]
and
\[
 \min\{1,3(\sqrt{E_{\mathrm{pt}}(x)+nE_{\mathrm{ext}}(x,y)}+y)\}
 \leq nE_{\mathrm{pass}}(x,y).
\]
Moreover $C_{\mathrm{abs}}>1$, $0<b_{\mathrm{abs}}<1$, and for every
admissible parameter tuple and $x\geq0$,
\[
\begin{split}
 \min\{1,2(&\delta_{\mathrm{ld}}(C_{\mathrm{ld}},b_{\mathrm{ld}},
    mE_{\mathrm{pass}}(x,\rho)+x,q,2m+2,d,1)\\
   &+\sqrt{E_{\mathrm{pt}}(x)}+\rho)\}
 \leq\delta_{\mathrm{qld}}(C_{\mathrm{abs}},b_{\mathrm{abs}},x,m,d,q).
\end{split}
\]
The direct-envelope version is
\[
\begin{split}
 \min\{1,{}&\delta_{\mathrm{ld}}(C_{\mathrm{ld}},b_{\mathrm{ld}},
  3(\sqrt{E_{\mathrm{pt}}(x)+mE_{\mathrm{ext}}(x,\rho)}+\rho)+x,
  q,2m+2,d,1)\\
  &+\sqrt{E_{\mathrm{pt}}(x)}+\rho\}
 \leq\delta_{\mathrm{qld}}(C_{\mathrm{abs}},b_{\mathrm{abs}},x,m,d,q).
\end{split}
\]

Finally, put
\[
\begin{split}
 \delta_0={}&\delta_{\mathrm{ld}}\left(C_{\mathrm{ld}},b_{\mathrm{ld}},
  3\left(\sqrt{E_{\mathrm{pt}}(x)+mE_{\mathrm{ext}}(x,\rho)}+\rho\right),
  q,2m+2,d,1\right),\\
 \eta_0={}&\delta_0+\sqrt{220\delta_0^{1/4}}+2\sqrt{2\delta_0}.
\end{split}
\]
The final rounded error satisfies
\[
 \min\left\{1,8(4\eta_0+8E_{\mathrm{pt}}(x))
   +\frac{12md+4d+14}{q}\right\}
 \leq\delta_{\mathrm{qld}}(C_{\mathrm{pair}},b_{\mathrm{pair}},x,m,d,q).
\]
\end{theorem}
\begin{proof}
\leanok
Each assertion is a direct comparison of nonnegative powers.  Split into the
unit range and its complement, use the cap by one in the latter case, and
apply the eighth-root absorption to the rounded error.
\end{proof}

\begin{definition}[Generic global-pair absorption witnesses]
\label{def:qld-global-pair-generic-constants}
\lean{MIPStarRE.QPBT.globalPairBoundPower}
\lean{MIPStarRE.QPBT.globalPairBoundConstant}
\leanok
For real powers $p,e,r,s$, define
\[
 B=\min\left\{\frac12,s,es,rs,\frac p2\right\}.
\]
For coefficients $P,L,K,A$ and a final scale $F$, put
\[
 \begin{split}
 \lambda&=s+rs,\qquad G=A+\lambda+1,\\
 L_0&=A4^A\bigl((KL)^s+2\bigr),\\
 C&=F(L_0+\sqrt P+1)+G+2.
 \end{split}
\]
These are respectively the generic global-pair power and coefficient.
\end{definition}

\begin{theorem}[Generic global-pair absorption]
\label{thm:qld-global-pair-bound-explicit}
\lean{MIPStarRE.QPBT.global_pair_error_bound_explicit}
\leanok
\uses{def:qld-global-pair-generic-constants}
Suppose $0\leq f(x)\leq Px^p$ and
$0\leq g(x,y)\leq L(x^e+y^r)$ on the nonnegative domain, with $P,L\geq1$
and $p,e,r>0$.  Suppose also $A\geq1$, $0<s\leq1$, and $K,F\geq0$.
Then the constants $C,B$ of
Definition~\ref{def:qld-global-pair-generic-constants} satisfy
$C>1$ and $0<B<1$, and
\[
 \min\{1,F(\delta_{\mathrm{ld}}(A,s,Km g(x,\rho)+x,q,2m+2,d,1)
   +\sqrt{f(x)}+\rho)\}
 \leq\delta_{\mathrm{qld}}(C,B,x,m,d,q).
\]
\end{theorem}
\begin{proof}
\leanok
Bound each summand by the smallest displayed exponent and absorb all
dimension growth into the coefficient $C$.
\end{proof}

\begin{theorem}[Explicit global polynomial-pair witness]
\label{thm:qld-global-pair-witness-explicit}
\lean{MIPStarRE.QPBT.exists_global_pair_witness_explicit}
\leanok
Every projective Pauli-test setting has a projective global polynomial-pair
witness with error
\[
 \delta_{\mathrm{qld}}(C_{\mathrm{pair}},b_{\mathrm{pair}},
   \eps,m,d,q).
\]
The point and extended-line witnesses are constructed internally and are not
hypotheses of this statement.
\end{theorem}
\begin{proof}
\leanok
\uses{thm:qld-combined-points-explicit, thm:qld-extended-lines-explicit,
thm:qld-rounded-polynomials-explicit, thm:qld-baseline-scalar-bounds}
Construct the point and extended-line witnesses, apply the rounded pair
construction, and use the final inequality of
Theorem~\ref{thm:qld-baseline-scalar-bounds}.
\end{proof}

\section{Coefficient-30 quantitative global-pair construction}
\label{sec:qld-coefficient-thirty-global-pair}

The results in this section are Lean-only quantitative refinements of the
argument through Lemma~\ref{lem:qld-4-7}. They use the complete-measurement
linear-triangle estimate of Theorem~\ref{thm:main-formal-linear-triangle} and
do not alter the statement or status of any source-labelled result. The final
construction is restricted to the small regime $0\leq e\leq1$ and
$md/q\leq1$; it is not the unrestricted Pauli soundness theorem.

\begin{definition}[Coefficient-30 error scales and envelopes]
\label{def:qld-coefficient-thirty-scales}
\lean{MIPStarRE.QPBT.quantitativeLowDegreePower}
\lean{MIPStarRE.QPBT.quantitativeGlobalPairPower}
\lean{MIPStarRE.QPBT.quantitativePreRoundingPower}
\lean{MIPStarRE.QPBT.directQuantitativeEnvelope}
\lean{MIPStarRE.QPBT.quantitativePreRoundingEnvelope}
\lean{MIPStarRE.QPBT.quantitativeGlobalPairEnvelope}
\leanok
Put
\[
 \tau=\frac1{8192},\qquad
 u=\frac1{4194304},\qquad
 b=\frac1{33554432}=\frac u8.
\]
For direct low-degree parameters $D=(q,m,d,k)$ and $\eps\in\mathbb R$, define
\[
 E_D(\eps)=\eps^\tau+q^{-\tau}+2^{-\tau md}.
\]
For admissible Pauli-test parameters $P=(q,m,d)$ and $e\in\mathbb R$, define
the pre-rounding and post-rounding envelopes
\[
 E_{\mathrm{pre}}(P,e)=e^u+q^{-u}+2^{-umd},
 \qquad
 E_{\mathrm{pair}}(P,e)=e^b+q^{-b}+2^{-bmd}.
\]
All powers in these definitions are real powers.
\end{definition}

\begin{definition}[Native direct and rounded-pair errors]
\label{def:qld-native-error-scales}
\lean{MIPStarRE.QPBT.directNativeError}
\lean{MIPStarRE.QPBT.nativeRoundingError}
\lean{MIPStarRE.QPBT.nativeOrderedPolynomialError}
\lean{MIPStarRE.QPBT.nativeGlobalPairRawError}
\lean{MIPStarRE.QPBT.nativeGlobalPairError}
\leanok
\uses{def:qld-coefficient-thirty-scales}
For direct low-degree parameters $D=(q,m,d,k)$ and $t\in\mathbb R$, define
\[
 \Lambda_D(t)=\min\!\left\{1,
 400000m^{5/4}d^{1/4}
 \left((3t)^\tau+(d/q)^\tau+\exp(-4md)\right)\right\}.
\]
For admissible Pauli-test parameters $P=(q,m,d)$ and real numbers
$\delta,\delta_Q$, define
\[
 \eta(\delta)=\delta+\sqrt{220}\,\delta^{1/8}
   +2\sqrt2\,\delta^{1/2},
 \qquad
 J(\delta,\delta_Q)=4\eta(\delta)+8\delta_Q,
\]
\[
 R_P(\delta,\delta_Q)=32\eta(\delta)+64\delta_Q
   +\frac{12md+4d+14}{q},
 \qquad
 G_P(\delta,\delta_Q)=\min\{1,R_P(\delta,\delta_Q)\}.
\]
All powers are real powers.
\end{definition}

\begin{theorem}[Scalar bounds for the coefficient-30 direct LDT]
\label{thm:qld-coefficient-thirty-ldt-scalars}
\lean{MIPStarRE.QPBT.three_rpow_quantitative_low_degree_power_le_two}
\lean{MIPStarRE.QPBT.direct_ld_aux_parameter_linear_triangle_exp_arg}
\lean{MIPStarRE.QPBT.direct_ld_aux_parameter_quarter_le}
\lean{MIPStarRE.QPBT.direct_ld_aux_parameter_quarter_le_sharp}
\lean{MIPStarRE.QPBT.direct_ld_field_term_quantitative_eq}
\lean{MIPStarRE.QPBT.direct_ld_exponential_term_quantitative_le}
\lean{MIPStarRE.QPBT.direct_ld_test_term_quantitative_le}
\lean{MIPStarRE.QPBT.direct_quantitative_envelope_nonneg}
\lean{MIPStarRE.QPBT.direct_ld_linear_triangle_raw_error_le}
\lean{MIPStarRE.QPBT.main_formal_linear_triangle_error_le_delta_ld_quantitative}
\leanok
\uses{def:qld-coefficient-thirty-scales,
  thm:main-formal-linear-triangle}
Let $D=(q,m,d,k)$ be direct low-degree parameters and let $N_D$ be its
auxiliary sampling parameter. Then
\[
 3^\tau\leq2,
 \qquad \frac{N_D}{640000m^2}=4md,
\]
\[
 N_D^{1/4}\leq40md,
 \qquad N_D^{1/4}\leq40m d^{1/4},
\]
\[
 (d/q)^\tau=d^\tau q^{-\tau},
 \qquad
 \exp\!\left(-\frac{N_D}{640000m^2}\right)\leq2^{-\tau md}.
\]
For every $\eps\geq0$ one also has
\[
 (3\eps)^\tau\leq2\eps^\tau,
 \qquad 0\leq E_D(\eps),
\]
and the uncapped complete-measurement linear-triangle error satisfies
\[
 T_D(\eps)\leq10^6(md)^2E_D(\eps).
\]
If $k=1$, its capped value is bounded without any additional upper bound on
$\eps$ by
\[
 \min\{1,T_D(\eps)\}
 \leq\delta_{\mathrm{ld}}(30,\tau,\eps,q,m,d,k)
 =30(md)^{30}E_D(\eps).
\]
\end{theorem}
\begin{proof}
\leanok
\uses{def:qld-coefficient-thirty-scales,
  thm:main-formal-linear-triangle}
Substitute $N_D=2560000m^3d$ in the linear-triangle estimate, separate the
field and exponential terms, and absorb the resulting quadratic factor into
the coefficient-$30$ error. The cases $md=1$ and $md\geq2$ respectively use
the unit cap and the power $(md)^{30}$.
\end{proof}

\begin{theorem}[Elementary properties of the native errors]
\label{thm:qld-native-error-scalars}
\lean{MIPStarRE.QPBT.direct_native_error_nonneg}
\lean{MIPStarRE.QPBT.native_rounding_error_eq}
\lean{MIPStarRE.QPBT.native_rounding_error_nonneg}
\lean{MIPStarRE.QPBT.native_rounding_error_mono}
\lean{MIPStarRE.QPBT.native_ordered_polynomial_error_nonneg}
\lean{MIPStarRE.QPBT.native_global_pair_raw_error_nonneg}
\lean{MIPStarRE.QPBT.native_global_pair_error_nonneg}
\lean{MIPStarRE.QPBT.direct_native_error_le_delta_ld_quantitative}
\leanok
\uses{def:qld-native-error-scales,
  thm:qld-coefficient-thirty-ldt-scalars}
If $t\geq0$, then $\Lambda_D(t)\geq0$.  If
$0\leq\delta\leq\delta'$, then
\[
 \delta+\sqrt{220\delta^{1/4}}+2\sqrt{2\delta}=\eta(\delta),
 \qquad
 0\leq\eta(\delta)\leq\eta(\delta').
\]
If also $\delta_Q\geq0$, then
\[
 0\leq J(\delta,\delta_Q),\qquad
 0\leq R_P(\delta,\delta_Q),\qquad
 0\leq G_P(\delta,\delta_Q).
\]
Finally, if $k=1$ and $t\geq0$, then
\[
 \Lambda_D(t)\leq
 \delta_{\mathrm{ld}}(30,\tau,t,q,m,d,k).
\]
\end{theorem}
\begin{proof}
\leanok
\uses{def:qld-native-error-scales,
  thm:qld-coefficient-thirty-ldt-scalars}
The square-root identity follows by separating the positive constants and
halving the real-power exponents.  Nonnegativity and monotonicity then follow
term by term.  The last inequality identifies $\Lambda_D(t)$ with the capped
linear-triangle error and applies the coefficient-$30$ estimate.
\end{proof}

\begin{theorem}[Fixed numerical certificates and the uncapped passing bound]
\label{thm:qld-coefficient-thirty-passing-scalars}
\lean{MIPStarRE.QPBT.pauli_baseline_point_constant_le}
\lean{MIPStarRE.QPBT.pauli_baseline_line_constant_le}
\lean{MIPStarRE.QPBT.pauli_baseline_extended_line_constant_le}
\lean{MIPStarRE.QPBT.pauli_baseline_passing_constant_le}
\lean{MIPStarRE.QPBT.pauli_baseline_direct_passing_bound_uncapped}
\lean{MIPStarRE.QPBT.direct_passing_error_envelope_quantitative_bound}
\leanok
\uses{def:qld-baseline-scalar-functions}
The fixed coefficients of
Definition~\ref{def:qld-baseline-scalar-functions} satisfy
\[
 C_{\mathrm{pt}}\leq10^{15},\qquad
 C_{\mathrm{line}}\leq10^5,\qquad
 C_{\mathrm{ext}}\leq10^9,\qquad
 C_{\mathrm{pass}}\leq10^8.
\]
Write $V(x,r)=3(\sqrt{x}+r)$ for the direct passing envelope. For all real
$e,r,n$ with $0\leq e\leq1$, $0\leq r\leq1$, and $n\geq1$,
\[
 V\bigl(E_{\mathrm{pt}}(e)+nE_{\mathrm{ext}}(e,r),r\bigr)
 \leq nE_{\mathrm{pass}}(e,r).
\]
In particular, if $P=(q,m,d)$ is admissible, then for the same domain of
$e$ and $r$,
\[
 V\bigl(E_{\mathrm{pt}}(e)+mE_{\mathrm{ext}}(e,r),r\bigr)
 \leq10^8m\bigl(e^{1/512}+r^{1/32}\bigr).
\]
These are uncapped inequalities.
\end{theorem}
\begin{proof}
\leanok
\uses{def:qld-baseline-scalar-functions}
Bound the square root of the point and extended-line sum by the sum of the
$1/512$ and $1/32$ powers, use $\sqrt n\leq n$, and substitute the four
fixed coefficient estimates.
\end{proof}

\begin{theorem}[Pre-rounding and rounded-pair scalar absorption]
\label{thm:qld-coefficient-thirty-rounding-scalars}
\lean{MIPStarRE.QPBT.quantitative_low_degree_error_bound}
\lean{MIPStarRE.QPBT.quantitative_pre_rounding_envelope_nonneg}
\lean{MIPStarRE.QPBT.quantitative_point_error_root_bound}
\lean{MIPStarRE.QPBT.quantitative_ratio_bound_of_nonneg}
\lean{MIPStarRE.QPBT.quantitative_pre_rounding_envelope_eighth_root_le}
\lean{MIPStarRE.QPBT.quantitative_actual_rounded_global_pair_error_bound}
\leanok
\uses{def:qld-coefficient-thirty-scales,
  thm:qld-coefficient-thirty-passing-scalars}
Let $P=(q,m,d)$ be admissible, put $n=md$ and $\rho=n/q$, and define
\[
 v=V\bigl(E_{\mathrm{pt}}(e)+mE_{\mathrm{ext}}(e,\rho),\rho\bigr),
 \qquad
 \delta=\delta_{\mathrm{ld}}(30,\tau,v,q,2m+2,d,1).
\]
If $0\leq e\leq1$ and $\rho\leq1$, then
\[
 \delta\leq10^{30}n^{32}E_{\mathrm{pre}}(P,e).
\]
For every $e\geq0$,
\[
 0\leq E_{\mathrm{pre}}(P,e),
 \qquad
 \rho\leq nE_{\mathrm{pre}}(P,e),
 \qquad
 \sqrt{\sqrt{\sqrt{E_{\mathrm{pre}}(P,e)}}}
 \leq E_{\mathrm{pair}}(P,e).
\]
If in addition $e\leq1$, then
\[
 \sqrt{E_{\mathrm{pt}}(e)}\leq4\cdot10^7E_{\mathrm{pre}}(P,e).
\]
Finally, under $0\leq e\leq1$ and $\rho\leq1$, set
\[
 \eta=\delta+\sqrt{220\delta^{1/4}}+2\sqrt{2\delta},
 \qquad
 R=8\bigl(4\eta+8E_{\mathrm{pt}}(e)\bigr)
   +\frac{12md+4d+14}{q}.
\]
Then the exact rounded-pair error obeys
\[
 \min\{1,R\}\leq10^7n^4E_{\mathrm{pair}}(P,e).
\]
\end{theorem}
\begin{proof}
\leanok
\uses{def:qld-coefficient-thirty-scales,
  thm:qld-coefficient-thirty-passing-scalars}
The uncapped passing estimate bounds the low-degree input. Each field, point,
and tail term is absorbed into $E_{\mathrm{pre}}$. Three successive square
roots convert that envelope to $E_{\mathrm{pair}}$, while the unit cap handles
the complementary range.
\end{proof}

\begin{theorem}[Native one-coordinate direct soundness]
\label{thm:qld-native-direct-soundness}
\lean{MIPStarRE.QPBT.direct_ld_aux_parameter_quarter_eq}
\lean{MIPStarRE.QPBT.main_formal_linear_triangle_error_eq_direct_native_error}
\lean{MIPStarRE.QPBT.direct_coordinate_main_formal_at_native_error}
\lean{MIPStarRE.QPBT.exists_direct_simultaneous_polynomial_measurements_at_native_error_of_k_eq_one}
\lean{MIPStarRE.QPBT.direct_ld_soundness_of_k_eq_one_at_native_error}
\lean{MIPStarRE.QPBT.direct_ld_soundness_of_k_eq_one_any_strategy_at_native_error}
\leanok
\uses{def:qld-native-error-scales, thm:main-formal-linear-triangle}
This is a Lean-only sharp sibling of the common-error presentation in
Lemma~\ref{lem:ld-soundness}.  For direct parameters $D=(q,m,d,k)$ put
$N_D=2560000m^3d$.  Then
$N_D^{1/4}=40m^{3/4}d^{1/4}$, and the complete-measurement linear-triangle
error at $N_D$ is exactly $\Lambda_D(t)$ from
Definition~\ref{def:qld-native-error-scales}.

For every real $t$, every projective direct-game strategy of value at least
$1-t$, and every coordinate, there are projective polynomial measurements
whose two point--polynomial defects and polynomial self-consistency defect are
at most $\Lambda_D(t)$.  If $k=1$, relabelling the same measurements as
singleton tuples preserves all three bounds; neither conclusion requires
$t>0$.

The separately stated projective direct-game soundness theorem assumes
$k=1$, $t>0$, projectivity, and value at least $1-t$, and returns the same
three tuple bounds.  Under $k=1$, $t>0$, and the same passing premise, but
without projectivity of the input strategy, Naimark dilation and compression
give tuple measurements with the same bounds; projectivity of the compressed
measurements is not asserted.
\end{theorem}
\begin{proof}
\leanok
\uses{def:qld-native-error-scales, thm:main-formal-linear-triangle}
Substitute the exact fourth root of $N_D$ and the exact exponential argument
into the linear-triangle estimate.  Apply the singleton-tuple relabelling and,
for arbitrary strategies, the existing two-sided Naimark compression.
\end{proof}

\begin{theorem}[Coefficient-30 one-coordinate direct soundness]
\label{thm:qld-coefficient-thirty-direct-soundness}
\lean{MIPStarRE.QPBT.direct_coordinate_main_formal_quantitative}
\lean{MIPStarRE.QPBT.exists_direct_simultaneous_polynomial_measurements_quantitative_of_k_eq_one}
\lean{MIPStarRE.QPBT.direct_ld_soundness_of_k_eq_one_quantitative}
\lean{MIPStarRE.QPBT.direct_ld_soundness_of_k_eq_one_any_strategy_quantitative}
\leanok
\uses{def:qld-coefficient-thirty-scales,
  thm:qld-coefficient-thirty-ldt-scalars,
  thm:main-formal-linear-triangle}
Let $D=(q,m,d,k)$ be direct low-degree parameters with $k=1$, let
$\eps\geq0$, and put
\[
 \Delta_D(\eps)=\delta_{\mathrm{ld}}(30,\tau,\eps,q,m,d,k).
\]
For every projective direct-game strategy of value at least $1-\eps$ and its
unique coordinate, there are projective polynomial measurements
$G^{\mathrm A},G^{\mathrm B}$ such that the point measurement of Alice is
consistent with the evaluation of $G^{\mathrm B}$, the evaluation of
$G^{\mathrm A}$ is consistent with Bob's point measurement, and
$G^{\mathrm A}$ is consistent with $G^{\mathrm B}$, each with defect at most
$\Delta_D(\eps)$. Relabelling these measurements as one-coordinate tuples
gives tuple measurements with the same three bounds.

If moreover $\eps>0$, the projective direct-game soundness theorem returns
such tuple measurements directly under the hypotheses $k=1$, projectivity,
and value at least $1-\eps$. Under the same strict positivity and passing
hypotheses but without projectivity of the input strategy, Naimark dilation
and compression produce tuple measurements satisfying the same three defect
bounds; projectivity of these compressed measurements is not asserted.
\end{theorem}
\begin{proof}
\leanok
\uses{thm:qld-coefficient-thirty-ldt-scalars,
  thm:main-formal-linear-triangle}
Apply the complete-measurement linear-triangle theorem to the unique
coordinate and use the scalar absorption above. Constant-tuple relabelling
preserves all three defects. For an arbitrary strategy, dilate every answer
measurement, apply the projective result, and compress at the distinguished
ancillary outcome.
\end{proof}

\begin{theorem}[Native rounded polynomial and pair witnesses]
\label{thm:qld-native-rounded-polynomials}
\lean{MIPStarRE.QPBT.ExtendedLineGame.rounded_polynomial_ordered_estimates_at_native_error}
\lean{MIPStarRE.QPBT.ExtendedLineGame.rounded_polynomial_scalar_mass_at_native_error}
\lean{MIPStarRE.QPBT.ExtendedLineGame.rounded_polynomial_wrong_variable_mass_at_native_error}
\lean{MIPStarRE.QPBT.ExtendedLineGame.rounded_polynomial_separated_mass_at_native_error}
\lean{MIPStarRE.QPBT.ExtendedLineGame.pair_witness_of_points_lines_at_native_error}
\leanok
\uses{thm:qld-native-direct-soundness, thm:qld-rounded-polynomials-explicit}
Let $P=(q,m,d)$ be admissible, put $n=md$ and $\rho=n/q$, and suppose that
combined-point and extended-line witnesses have errors $Q$ and $L_{\rm ext}$.
Set
\[
 T=3\bigl(\sqrt{Q+L_{\rm ext}}+\rho\bigr),\qquad
 \lambda=\Lambda_{P_{\rm ext}}(T),
\]
where the extended direct dimension is $2m+2$, and define
\[
 \eta(\lambda)=\lambda+\sqrt{220}\,\lambda^{1/8}
   +2\sqrt2\,\lambda^{1/2},\qquad
 J=4\eta(\lambda)+8Q.
\]
There are rounded projective polynomial measurements with both ordered errors
at most $J$.  There are such measurements retaining the scalar-nonlinear and
wrong-variable mass bounds.  Separately, there are rounded measurements with
the two ordered bounds and non-separated mass at most
\[
 6J+\frac{12md+4d+10}{q}.
\]
Their completion gives projective polynomial-pair measurements satisfying all
four point comparisons at the uncapped error
\[
 32\eta(\lambda)+64Q+\frac{12md+4d+14}{q}.
\]
Each declaration constructs its measurements through the same ordered-measurement
construction; no conjunction between independently quantified witnesses is
asserted.
\end{theorem}
\begin{proof}
\leanok
\uses{thm:qld-native-direct-soundness, thm:qld-rounded-polynomials-explicit}
Use native arbitrary-strategy soundness in the existing projective-rounding
construction.  The scalar concentration, variable separation, and completion
arguments preserve the displayed ordered error without a common-exponent
replacement.
\end{proof}

\begin{theorem}[Coefficient-30 rounded polynomial witnesses]
\label{thm:qld-coefficient-thirty-rounded-polynomials}
\lean{MIPStarRE.QPBT.ExtendedLineGame.rounded_polynomial_ordered_estimates_quantitative}
\lean{MIPStarRE.QPBT.ExtendedLineGame.rounded_polynomial_scalar_mass_quantitative}
\lean{MIPStarRE.QPBT.ExtendedLineGame.rounded_polynomial_wrong_variable_mass_quantitative}
\lean{MIPStarRE.QPBT.ExtendedLineGame.rounded_polynomial_separated_mass_quantitative}
\lean{MIPStarRE.QPBT.ExtendedLineGame.pair_witness_of_points_lines_quantitative}
\leanok
\uses{thm:qld-coefficient-thirty-direct-soundness,
  thm:qld-rounded-polynomials-explicit}
Let $P=(q,m,d)$ be admissible, let $\eps,\delta_Q,\delta_L\in\mathbb R$,
and let $S$ be a projective Pauli-test setting with a supplied combined-point
witness of error $\delta_Q$ and a supplied directly indexed extended-line
witness of error $\delta_L$. Put $\rho=md/q$ and
\[
 \delta=\delta_{\mathrm{ld}}
   \bigl(30,\tau,V(\delta_Q+\delta_L,\rho),q,2m+2,d,1\bigr),
 \qquad
 \eta=\delta+\sqrt{220\delta^{1/4}}+2\sqrt{2\delta},
 \qquad E=4\eta+8\delta_Q.
\]
No sign or small-regime hypothesis is imposed on
$\eps,\delta_Q,\delta_L$ in these conditional statements.

There exist projective rounded polynomial measurements on the two expanded
local spaces whose ordered errors, in both orderings and on both player
placements, are at most $E$. There also exist such measurements satisfying,
in addition, the two scalar nonlinear-mass bounds
\[
 2E+2\frac{(2m+2)d+1}{q}.
\]
There exist such measurements satisfying those ordered and scalar bounds and,
for both orderings, the two wrong-variable mass bounds
\[
 2E+2\frac{2md+2}{q}.
\]
Separately, there exist projective rounded measurements satisfying the ordered
bounds and the two non-separated-mass bounds
\[
 6E+\frac{12md+4d+10}{q}.
\]
Finally, the supplied point and line witnesses determine complete projective
polynomial-pair measurements on both player spaces. All four point defects,
one for each player placement and Pauli basis, are bounded by the uncapped
quantity
\[
 8E+\frac{12md+4d+14}{q}.
\]
Each existential assertion above concerns the measurements constructed in its
own declaration; no conjunction between independently quantified witnesses is
being asserted.
\end{theorem}
\begin{proof}
\leanok
\uses{thm:qld-coefficient-thirty-direct-soundness,
  thm:qld-rounded-polynomials-explicit}
Apply arbitrary-strategy one-coordinate soundness to the direct combined game,
round the two resulting polynomial measurements projectively, and retain both
ordered estimates. The scalar, wrong-variable, and separated-mass bounds are
the corresponding finite-field root-counting estimates. Completion at the
zero polynomial pair gives the two projective pair measurements and the four
displayed defects.
\end{proof}

\begin{definition}[Separated native scalar errors]
\label{def:qld-native-separated-errors}
\lean{MIPStarRE.QPBT.quantitativeNativeSeparatedRaw}
\lean{MIPStarRE.QPBT.quantitativeNativeSeparatedError}
\lean{MIPStarRE.QPBT.quantitativeNativeGlobalPairSeparatedError}
\leanok
\uses{def:qld-coefficient-thirty-scales}
Let $h=2m+2$, $r=md/q$, and let $a$ denote the native direct low-degree
exponent.  For $s\in\mathbb R$, define
\[
\begin{split}
 S_s(P,e)={}&h^{5s/4}d^{s/4}\Bigl[800000^s\bigl(
 e^{as/16}+m^{as/2}e^{as/512}+m^{as/2}r^{as/32}+r^{as}\bigr)\\
 &\hspace{42mm}+400000^s\bigl((d/q)^{as}+\exp(-4hds)\bigr)\Bigr],
 \qquad \bar S_s(P,e)=\min\{1,S_s(P,e)\}.
\end{split}
\]
The separated global-pair bound is
\[
 B_G(P,e)=\min\left\{1,
 32\bar S_1+32\sqrt{220}\,\bar S_{1/8}
 +64\sqrt2\,\bar S_{1/2}+64E_{\rm pt}(e)
 +\frac{12md+4d+14}{q}\right\}.
\]
These quantities retain the six native contributions and do not replace a
numerical coefficient by a parameter exponent.
\end{definition}

\begin{theorem}[Separated native scalar certificates]
\label{thm:qld-native-separated-scalar-support}
\lean{MIPStarRE.QPBT.quantitative_native_error_rpow_le_separated}
\lean{MIPStarRE.QPBT.quantitative_native_global_pair_error_le_separated}
\leanok
\uses{def:qld-native-separated-errors,
  def:qld-native-error-scales,
  thm:qld-native-error-scalars,
  thm:qld-coefficient-thirty-passing-scalars}
For $e\geq0$ and $0\leq s\leq1$, the $s$-th real power of the exact native
direct low-degree error is at most $\bar S_s(P,e)$.  Consequently the exact
native global-pair error, with its point and collision terms unchanged, is at
most $B_G(P,e)$.
\end{theorem}
\begin{proof}
\leanok
\uses{def:qld-native-separated-errors,
  def:qld-native-error-scales,
  thm:qld-native-error-scalars,
  thm:qld-coefficient-thirty-passing-scalars}
Apply subadditivity of real powers separately to the four passing-error terms
and the two direct low-degree terms.  Monotonicity through the unit cap then
bounds the three powers occurring in the native rounding error.  Substitution
in the exact global-pair expression proves the second inequality.
\end{proof}

\begin{theorem}[Native global polynomial pair and small-regime bounds]
\label{thm:qld-native-small-regime-global-pair}
\lean{MIPStarRE.QPBT.quantitative_native_global_pair_error_le_fractional}
\lean{MIPStarRE.QPBT.quantitative_native_global_pair_error_bound}
\lean{MIPStarRE.QPBT.exists_quantitative_global_pair_witness_native}
\lean{MIPStarRE.QPBT.exists_quantitative_global_pair_witness_at_native_error}
\leanok
\uses{thm:qld-native-rounded-polynomials,
  thm:pauli-final-fractional-scalar-support}
Let $P=(q,m,d)$ be admissible, $n=md$, and $\rho=n/q$.  For $e\geq0$, take
\[
 Q=E_{\rm pt}(e),\qquad
 L_{\rm ext}=mE_{\rm ext}(e,\rho),\qquad
 \lambda=\Lambda_{P_{\rm ext}}
   \bigl(3(\sqrt{Q+L_{\rm ext}}+\rho)\bigr),
\]
and set
\[
 G=\min\!\left\{1,
 32\eta(\lambda)+64Q+\frac{12n+4d+14}{q}\right\}.
\]
For every projective Pauli-test setting $S$ of value at least $1-e$, there is
one global-pair witness at the exact error $G$.  Thus one common pair of
complete projective polynomial-pair measurements realizes all four point
comparisons, with defect at most $G$.  This assertion requires neither
$e\leq1$ nor $\rho\leq1$.

If in addition $e\leq1$, then, without any restriction on $\rho$, the exact
native error satisfies
\[
 G\leq277248\,m^{20481/131072}d^{1/32}E_{\rm pair}(P,e).
\]
If also $\rho\leq1$, the restricted existential theorem records the same exact
error, its nonnegativity, the global-pair witness, and the common-envelope
consequence
\[
 G\leq10^7n^4E_{\rm pair}(P,e).
\]
\end{theorem}
\begin{proof}
\leanok
\uses{thm:qld-native-rounded-polynomials,
  thm:pauli-final-fractional-scalar-support}
For the unrestricted existence assertion, construct the point and extended-line
witnesses and apply the native pair construction.  Under $e\leq1$, the
fractional scalar certificate gives
$\sqrt G\leq304m^{20481/262144}d^{1/64}E_b$, where
$b=1/67108864$.  Square this estimate, use the exact dimension-power identity,
and apply $(x+y+z)^2\leq3(x^2+y^2+z^2)$ to obtain the displayed bound with
$E_{2b}=E_{\rm pair}$ and coefficient $304^2\cdot3=277248$.  When also
$\rho\leq1$, the restricted existential theorem packages the same exact error
and witness with the common-envelope consequence, which follows from
$m^{20481/131072}d^{1/32}\leq(md)^4$ and $277248\leq10^7$; it does not use the
coefficient-$30$ low-degree comparison.
\end{proof}

\begin{theorem}[Small-regime quantitative global polynomial pair]
\label{thm:qld-small-regime-global-pair}
\lean{MIPStarRE.QPBT.exists_quantitative_global_pair_witness}
\leanok
\uses{thm:qld-native-small-regime-global-pair}
Let $P=(q,m,d)$ be admissible, let $0\leq e\leq1$, suppose
\[
 \frac{md}{q}\leq1,
\]
and let $S$ be a projective Pauli-test setting with error $e$. Then there is
a real number $g$ such that
\[
 0\leq g\leq10^7(md)^4
 \left(e^{1/33554432}+q^{-1/33554432}
       +2^{-md/33554432}\right),
\]
together with complete projective polynomial-pair measurements on both
expanded player spaces. For each of the two Pauli bases, Alice's evaluated
global measurement is consistent with Bob's expanded point measurement with
defect at most $g$, and Bob's evaluated global measurement is consistent with
Alice's expanded point measurement with defect at most $g$. Thus all four
point defects are bounded by the same $g$.

The combined-point witness, the extended-line witness, and the global pair are
constructed inside the theorem; no global-pair witness is a hypothesis. The
two displayed small-regime assumptions are additional to the source statement
of Lemma~\ref{lem:qld-4-7}, so this Lean-only theorem is not presented as that
unrestricted source result or as the final Pauli soundness theorem.
\end{theorem}
\begin{proof}
\leanok
\uses{thm:qld-native-small-regime-global-pair}
Take the real number and polynomial-pair measurements supplied by
Theorem~\ref{thm:qld-native-small-regime-global-pair}.  Discard their
identifying equality with the native capped error, and retain the
nonnegativity, common-envelope bound, and four consistency estimates.
\end{proof}
